%% 本文件由 工程文件/scripts/assemble.py 生成，勿手改。

%% begin units/ch07_u01.tex
% ch07_u01 —— 原书 p.236--p.245（7.1.1 -- 7.3.3，第 7 章开头）
% 本单元自第 7 章标题起，至 7.3.3 末句止；尾部接 p.246。

\bdchap{ch:appstack}{附录：代数叠上的 D-模理论与 Hecke 模式}

\SourcePage{236}

\bdsec{7.1. 引言}

\bdp{7.1.1} 本节的中心目标是陈述一个一般的 Hecke 模式，它用于主定理的证明。
其（非扭）有限维版本如下。设 $G$ 为代数群，$K\subset G$ 为代数子群，
$\mathfrak{g}$ 为 $G$ 的李代数，$Y$ 为带 $G$-作用的光滑簇。记
$\mathcal{H}:=D(K\backslash G/K)$ 为叠 $K\backslash G/K$ 的 $\mathcal{D}$-模导出范畴。
类似地有导出范畴 $D(K\backslash Y)$，以及 $(\mathfrak{g},K)$-模的范畴
$\mathcal{M}(\mathfrak{g},K)$ 的导出范畴 $D(\mathfrak{g},K)$。于是有如下
「Hecke 模式」：
\begin{enumerate}[label=(\alph*),leftmargin=2.2em]
\item $\mathcal{H}$ 是幺半三角范畴；
\item $D(K\backslash Y)$ 是 $\mathcal{H}$-模；
\item $D(\mathfrak{g},K)$ 是 $\mathcal{H}$-模；
\item 标准函子
\[
{}^{L}\Delta:D(\mathfrak{g},K)\longrightarrow D(K\backslash Y),
\qquad
R\Gamma:D(K\backslash Y)\longrightarrow D(\mathfrak{g},K)
\]
是 $\mathcal{H}$-模的态射。
\end{enumerate}

这里 ${}^{L}\Delta$ 与 $R\Gamma$ 是 1.2.4 中函子 $\Delta$ 与 $\Gamma$ 的导出版本。
$\mathcal{H}$ 上的张量积，以及 (b)、(c) 中的 $\mathcal{H}$-作用，
都是适当的「卷积」函子 $\circledast$。例如考察 $K=\{1\}$ 的情形。
记 $\delta_g$ 为 $g\in G$ 处的 $\delta$-函数对应的 $\mathcal{D}$-模，则有
$\delta_{g_1}\circledast\delta_{g_2}=\delta_{g_1g_2}$。对 $Y$ 上的
$\mathcal{D}$-模 $M$，$\delta_g\circledast M$ 是 $M$ 的 $g$-平移；对
$\mathfrak{g}$-模 $V$，$\delta_g\circledast V$ 就是 $V$ 配上经
$\operatorname{Ad}_g$ 扭过的 $\mathfrak{g}$-作用。$M$ 上的 $\mathcal{D}$-模结构
把无限接近的 $g$ 所对应的 $\delta_g\circledast M$ 典范地等同起来；类似地，
$V$ 上的 $\mathfrak{g}$-作用把这些 $\delta_g\circledast V$ 等同起来。
这就允许对 $G$ 上任意的 $\mathcal{D}$-模定义卷积函子。

\bdp{7.1.2} Hecke 函子的精确构造需要某种叠上的 $\mathcal{D}$-模形式化。例如，
人们需要光滑叠 $\mathcal{Y}$ 的 $\mathcal{D}$-模%
\SourcePage{237}%
导出范畴 $D(\mathcal{Y})$ 的定义（它可能与 $\mathcal{Y}$ 上 $\mathcal{D}$-模范畴的
导出范畴并不相同！）。似乎没有现成的参考文献。例外是 $\mathcal{Y}$ 为轨道叠的
情形，即光滑簇对仿射群作用的商；[BL]、[Gi87] 处理过它，但方式不太便于
Hecke 函子的应用。因此我们只好补上一些一般废话，以免沉没。

我们仿照 Kapranov [Kap91] 与 Saito [Sa89]，在 7.2 中从如下典范等价开始：
$\mathcal{D}$-模的导出范畴与 $\Omega$-模的导出范畴之间的等价
（这里 $\Omega$ 是微分形式的 DG 代数），它把 $\mathcal{D}$-模等同于其 de Rham 复形。
处理叠时，$\Omega$-模更易驾驭：原因在于 $\Omega$ 是光滑拓扑上的环层，
而 $\mathcal{D}$ 不是。在重要的特殊情形——叠的对角态射为仿射——这一「超」%
\footnote{这是某位数学家对 Mary Poppins 生造的词
「supercalifragelistic-expialidocious」的缩写。
% [Tutor Note: 原书 p.237 脚注把该生造词印作「supercalifragelistic-」，通行拼法为
% 「supercalifragilisticexpialidocious」；底本如此，译文照录。]
}%
格式尤为方便。此时可以（见 7.3）直接用「整体」$\Omega$-复形来定义 $\mathcal{D}$-模
导出范畴。在 7.5 中，经过 7.4 的一般同调代数离题之后，我们给出 $\mathcal{D}$-模
导出范畴的一个「局部」定义，它对任意的光滑叠都适用。

7.6 解释 Hecke 模式的 (a)、(b) 两部分；我们还证明，轨道叠的
$\mathcal{D}$-模导出范畴等价于 [BL]、[Gi87] 中的等变导出范畴。7.7 描述
Harish-Chandra 模的类似「超」格式；作为额外收获，我们在 7.7.12 得到 [BL]
主要结果的一个简单证明。Hecke 模式的 Harish-Chandra 部分 (c)、(d) 在 7.8 中处理。
正文所需的带额外对称与参数的版本在 7.9 中给出。在转入无穷维设定之前，
我们在 7.10 讨论 $\mathcal{D}$-模的晶体方法，处理奇异空间时它尤为方便
（本节要归功于 1980 年与 J. Bernstein 的讨论）。7.11 与 7.12 两节包含关于
归纳概形、Mittag-Leffler 模的一些基本材料，%
\SourcePage{238}%
以及形式光滑归纳概形上的 $\mathcal{D}$-模。7.13 节综述 BRST 约化。Hecke 模式
(c)、(d) 两部分的无穷维版本在 7.14 中。最后，7.15 证明仿射旗簇上的正扭
$\mathcal{D}$-模本质上等同于小于临界水平的仿射 Kac–Moody 李代数的表示。
在 $\mathcal{D}$-模沿 Schubert 分层光滑的特殊情形，Kashiwara 与 Tanisaki [KT95]
得到了类似的结果（[KT95] 的作者不使用归纳概形上 $\mathcal{D}$-模的语言）。
我们还把相应的 de Rham 上同调群与 BRST 上同调群等同起来。

我们对 $\mathcal{D}$-模理论的叙述相当不完整；基本上只处理本文正文用到的题材。
例外是 7.4、7.5 两节（叠 $\operatorname{Bun}_G$ 落入 7.3 的形式化）、7.10 节
（我们遇到的奇异空间都是仿射 Grassmann 流形上的分层，故可用 7.11）以及 7.15 节
（它收进来只为博读者一乐）。

回忆：$\mathcal{M}^{\ell}(X)$（相应地 $\mathcal{M}^{r}(X)$）表示光滑簇 $X$ 上
左（相应地右）$\mathcal{D}$-模的范畴；我们常把这些范畴等同起来，并记作
$\mathcal{M}(X)$。若 $F$ 是复形，则记 $F^{\cdot}$ 为相应的分次对象
（忘掉微分）。

\bdsec{7.2. $\mathcal{D}$- 与 $\Omega$-模}

\bdp{7.2.1} 设 $X$ 为光滑代数簇%
\footnote{或者更一般地，$\mathbb{C}$ 上光滑拟紧的代数空间，其对角态射
$X\to X\times X$ 为仿射。本节（7.2.10 除外）的构造与陈述都是局部的，
故对任意光滑代数空间都有意义。对 $X$ 的这一要求是为了保证我们定义的导出范畴
满足适当的局部-整体（下降）性质。我们在 7.5 中于栈的更一般设定下讨论此事。}%
。记 $\Omega_X$ 为 $X$ 上微分形式的 DG 代数。则 $(X,\Omega_X)$ 是 DG 环化空间，
于是有 $\Omega_X$-复形的范畴（$:=$ DG $\Omega_X$-模）。$\Omega_X$-复形
$F=(F^{\cdot},d)$ 称为拟凝聚的，如果各 $F^i$ 都是拟凝聚的 $\mathcal{O}_X$-模；
拟凝聚的 $\Omega_X$-复形通常称为 $X$ 上的 \emph{$\Omega$-复形}。记%
\SourcePage{239}%
为 $X$ 上 $\Omega$-复形的 DG 范畴 $C(X,\Omega)$。这是张量 DG 范畴。

\emph{注.} 对 $\Omega_X$-复形 $F$，微分 $d:F^{\cdot}\to F^{\cdot+1}$
是阶 $\le1$ 的微分算子，其象征等于乘积映射
$\Omega_X^1\otimes F^{\cdot}\to F^{\cdot+1}$。由此可见，$F^{\cdot}$ 上的
$\Omega_X^{\cdot}$-模结构可以由 $\mathcal{O}_X$-模结构与 $d$ 重构出来。事实上，
忘掉 $\Omega_X^{\ge1}$-作用，就把 $C(X,\Omega)$ 与如下复形 $(F^{\cdot},d)$
的范畴等同起来：各 $F^{\cdot}$ 是拟凝聚 $\mathcal{O}_X$-模，$d$ 是阶 $\le1$
的微分算子。

\bdp{7.2.2} 设 $C(X,\mathcal{D}):=C(\mathcal{M}^{r}(X))$ 为 $X$ 上右
$\mathcal{D}$-模的复形（\emph{右 $\mathcal{D}$-复形}，或简称
\emph{$\mathcal{D}$-复形}）的 DG 范畴，$K(X,\mathcal{D})$ 为相应的同伦范畴。
我们有一对伴随 DG 函子
\begin{equation}\tag{306}
\mathcal{D}:C(X,\Omega)\longrightarrow C(X,\mathcal{D}),
\qquad
\Omega:C(X,\mathcal{D})\longrightarrow C(X,\Omega)
\end{equation}
其定义如下。记 $DR_X$ 为 $\mathcal{D}_X$ 视为左 $\mathcal{D}$-模时的 de Rham
复形，故 $DR_X=\Omega_X^{\cdot}\otimes_{\mathcal{O}_X}\mathcal{D}_X$。
这是一个带 $\mathcal{D}_X$ 右作用的 $\Omega$-复形。现在，对 $\Omega$-复形 $F$
与右 $\mathcal{D}$-复形 $M$，有
\begin{equation}\tag{307}
\mathcal{D}F=F\otimes_{\Omega_X}DR_X,
\qquad
\Omega M:=\operatorname{Hom}_{\mathcal{D}_X}(DR_X,M).
\end{equation}
伴随性质是显然的。

\bdp{7.2.3} \emph{注.} (i) 有
$\mathcal{D}F^{\cdot}=F^{\cdot}\otimes_{\mathcal{O}_X}\mathcal{D}_X
=\operatorname{Diff}(\mathcal{O},F^{\cdot})$；微分
$d_{\mathcal{D}F}:\mathcal{D}F^{\cdot}\to\mathcal{D}F^{\cdot+1}$ 把微分算子
$a:\mathcal{O}_X\to F^{\cdot}$ 送到复合 $d\cdot a$。$\Omega$-复形 $\Omega M$
就是 $M$ 的 de Rham 复形，其中
$(\Omega M)^i=\bigoplus_{a-b=i}M^a\otimes\Lambda^b\Theta_X$。

(ii) $X$ 上左 $\mathcal{D}$-模的范畴 $\mathcal{M}^{\ell}(X)$ 按通常方式
（张量积取在 $\mathcal{O}_X$ 上）是张量范畴，故左 $\mathcal{D}$-复形的范畴
$C(\mathcal{M}^{\ell}(X))$ 是张量 DG 范畴。DG 函子 $\Omega:$%
\SourcePage{240}%
$C(\mathcal{M}^{\ell}(X))\to C(X,\Omega)$，它把左 $\mathcal{D}$-复形 $N$ 送到其
de Rham 复形 $(\Omega N)^{\cdot}=\Omega_X^{\cdot}\otimes_{\mathcal{O}_X}N$，
是张量函子。

(iii) DG 范畴 $C(X,\Omega)$ 与 $C(X,\mathcal{D})$ 都是张量 DG 范畴
$C(\mathcal{M}^{\ell}(X))$ 上的模。函子 $\mathcal{D}$ 与 $\Omega$ 都是
$C(\mathcal{M}^{\ell}(X))$-模的态射。

\bdp{7.2.4} \emph{引理.} 对任意 $\mathcal{D}$-复形 $M$，典范态射
$\mathcal{D}\Omega M\to M$ 是拟同构。

\emph{证明.} 令
\[
V_j^{i}:=\bigoplus_{\substack{a-b=i\\ b+c=j}}
M^a\otimes\Lambda^b\Theta_X\otimes\mathcal{D}_X^{\leq c}
\subset(\mathcal{D}\Omega M)^i.
\]
则 $V_{*}$ 是 $\mathcal{D}\Omega M$ 的递增滤链，各项为 $\mathcal{O}$-子复形，
且 $V_0\xrightarrow{\sim}M$，而对 $i\ge1$，$V_i/V_{i-1}$ 无环
（因为 $V_i/V_{i-1}$ 是 $M$ 与 $\Theta_X$ 的第 $i$ 个 Koszul 复形之张量积）。
\hfill$\square$

\bdp{7.2.5} 对 $\Omega$-复形 $F$，置 $H_{\mathcal{D}}^{\cdot}F=H^{\cdot}\mathcal{D}F$。
于是 $H_{\mathcal{D}}$ 是 $K(X,\Omega)$ 上的上同调函子，取值于阿贝尔范畴
$\mathcal{M}^{r}(X)$。$\Omega$-复形态射 $\phi:F_1\to F_2$ 称为
\emph{$\mathcal{D}$-拟同构}，如果 $\mathcal{D}$-复形态射
$\mathcal{D}\phi:\mathcal{D}F_1\to\mathcal{D}F_2$ 是拟同构，亦即
$H_{\mathcal{D}}^{\cdot}F_1\to H_{\mathcal{D}}^{\cdot}F_2$ 是同构。我们有如下
简单性质（用 7.2.4 证明 (ii)、(iii)）：
\begin{enumerate}[label=(\roman*),leftmargin=2.2em]
\item 若 $\phi$ 是 $\mathcal{D}$-拟同构，$N$ 是作为 $\mathcal{O}$-模平坦的左
$\mathcal{D}$-模，则 $\phi\otimes\operatorname{id}_N:F_1\otimes N\to F_2\otimes N$
是 $\mathcal{D}$-拟同构。
\item 典范态射 $\alpha_F:F\to\Omega\mathcal{D}F$ 是 $\mathcal{D}$-拟同构。
\item $\Omega$ 把拟同构送到 $\mathcal{D}$-拟同构。
\end{enumerate}

下面的引理在后文中不会用到；读者可以跳过。若 $\Omega$-复形态射
$\phi:F_1\to F_2$ 是层复形（取值于向量空间）的拟同构，则称之为
\emph{朴素的拟同构}。

\SourcePage{241}

\bdp{7.2.6} \emph{引理.} (i) 任何 $\mathcal{D}$-拟同构都是朴素的拟同构。

(ii) 如上态射 $\phi$ 是 $\mathcal{D}$-拟同构，当且仅当对任意下有界的局部自由
$\Omega$-模复形 $A$，态射 $\phi\otimes\operatorname{id}_A:F_1\otimes A\to F_2\otimes A$
都是朴素的拟同构。

(iii) 假设或者 $\Omega^{\ge1}F_i^{\cdot}=0$（即微分是 $\mathcal{O}$-线性的），
或者各 $F_i^{\cdot}$ 有界且 $\mathcal{O}$-凝聚。则任何朴素的拟同构 $\phi$ 都是
$\mathcal{D}$-拟同构。对一般的 $\Omega$-复形，这未必成立。

\emph{证明.} (i) 对任意 $\Omega$-复形 $F$，典范态射
$\alpha_F:F\to\Omega\mathcal{D}F$ 是朴素的拟同构。由于 $\Omega$ 把
$\mathcal{D}$-复形的拟同构送到朴素的拟同构，$\Omega(\mathcal{D}\phi)$ 是朴素的
拟同构。现在由 $\alpha_{F_2}\phi=\Omega((\mathcal{D}\phi)\alpha_{F_1}$ 即得结论。
% [Tutor Note: 原书 p.241 此处 $\Omega((\mathcal{D}\phi)\alpha_{F_1}$ 缺一个右括号，
% 底本如此，译文照录。]

(ii) 为证「若」的陈述，只需取 $A=DR_X$。反之，设 $\phi$ 是
$\mathcal{D}$-拟同构。$A$ 上有下有界的递增滤链 $A_i$，使得 $\cup A_i=A$，
且每个 $\operatorname{gr}_iA$ 是生成元位于 $i$ 次处的局部自由 $\Omega_X$-模
（取 $A_i:=\Omega_X\cdot A_{\le i}$）。于是，若所有
$\phi\otimes\operatorname{id}_{\operatorname{gr}_iA}$ 都是朴素的拟同构，
则 $\phi\otimes\operatorname{id}_A$ 也是。这样便可假设 $A$ 是生成元位于固定次数
（比如 0 次）的局部自由 $\Omega_X$-模，即 $A=\Omega N$，其中 $N$ 是作为
$\mathcal{O}$-模局部自由的左 $\mathcal{D}$-模。此时
$\phi\otimes\operatorname{id}_A=\phi\otimes\operatorname{id}_N$，
由 7.2.5 的 (i) 即得结论。

(iii) $\mathcal{O}$-线性的情形显然（因为此时
$\mathcal{D}F=F\otimes_{\mathcal{O}_X}\mathcal{D}_X$）。$\mathcal{O}$-凝聚的
情形由下面应用于 $\mathcal{D}\phi$ 的子引理推出（注意由 7.2.5 的性质 (ii)，
$\mathcal{D}F$ 在 $x$ 处的纤维与 $R\Gamma_x(X,F)$ 一致）。

\emph{子引理.} 设 $\psi:M_1\to M_2$ 是 $X$ 上凝聚 $\mathcal{D}$-模的有限复形
之间的态射。假设对任意 $x\in X(\mathbb{C})$，相应的纤维%
\footnote{这里当然是在通常的导出范畴意义下考虑 $\mathcal{O}$-模纤维。
% [Tutor Note: 原书 p.241 脚注印作「the $\mathcal{O}$-moduli fibers」，其中
% 「moduli」当为「module」；底本如此，译文按原意译出。]
}%
态射 $M_{1x}\to M_{2x}$ 都是拟同构。则 $\psi$ 是拟同构。

\SourcePage{242}

\emph{子引理的证明.} 置 $C=\operatorname{Cone}(\psi)$；记 $Y$ 为 $H^{\cdot}(C)$
的支撑。假设 $\psi$ 不是拟同构，即 $Y$ 非空。必要时限制 $X$，可设 $Y$ 是 $X$
的光滑子簇，且凝聚 $\mathcal{D}_Y$-模
$P^{\cdot}:=i_Y^{!}H^{\cdot}(C)=H^{\cdot}i_Y^{!}(C)$ 作为 $\mathcal{O}_Y$-模
自由。由于对 $x\in Y$ 有 $H^{\cdot}(C_x)=P_x^{\cdot+n}$，其中 $n$ 是 $Y$ 在
$X$ 中的余维，故 $P^{\cdot}=0$，矛盾。

为得到朴素拟同构而非 $\mathcal{D}$-拟同构的例子，只需找到非零 $\mathcal{D}$-模
$M$，使 $\Omega M$ 是层的无环复形。取 $M$ 为常数层，其值为 $\mathcal{D}_X$-模，
等于微分算子环（在 $X$ 的一般点处）的分式域。\hfill$\square$

\bdp{7.2.7} 由于 $H_{\mathcal{D}}$ 是上同调函子，$\mathcal{D}$-拟同构在
$C(X,\Omega)$ 的同伦范畴中构成局部化族。因此相应的局部化 $D(X,\Omega)$ 是
三角范畴（见 [Ve]）；我们称之为 $\Omega$-复形的 \emph{$\mathcal{D}$-导出范畴}。
函子 $\mathcal{D}$、$\Omega$ 给出三角范畴的互逆等价
\begin{equation}\tag{308}
\mathcal{D}:D(X,\Omega)\longrightarrow D(X,\mathcal{D}),
\qquad
\Omega:D(X,\mathcal{D})\longrightarrow D(X,\Omega).
\end{equation}
这里 $D(X,\mathcal{D})=D\mathcal{M}^{r}(X)$。我们常把这样等同起来的这些三角范畴
记为 $D(X)$。也可以考虑有界导出范畴。

\emph{注.} 对下有界的单射 $\mathcal{D}$-模复形 $M$，相应的 $\Omega$-复形
$\Omega M$ 是单射的。因此同伦范畴 $K^{+}(X,\Omega)$ 有许多单射对象。

\bdp{7.2.8} 设 $f:Y\to Z$ 是光滑簇之间的态射。它给出 DG 环化空间的态射
$f_{\Omega}:(Y,\Omega_Y)\to(Z,\Omega_Z)$。于是有相应的 DG 函子
$f_{\Omega}^{\cdot}:C(Z,\Omega)\to C(Y,\Omega)$，
$f_{\cdot}=f_{\Omega\cdot}:C(Y,\Omega)\to C(Z,\Omega)$。
我们先考察拉回函子。

\SourcePage{243}

对左 $\mathcal{D}$-模有通常的拉回函子
$f^{\dagger}:\mathcal{M}^{\ell}(Z)\to\mathcal{M}^{\ell}(Y)$，
$f^{\dagger}(N)=\mathcal{O}_Y\otimes_{f^{-1}\mathcal{O}_Z}f^{-1}N$。有
$\Omega f^{\dagger}(N)=f_{\Omega}^{\cdot}(\Omega N)$。可以把左 $\mathcal{D}$-模
换成右 $\mathcal{D}$-模%
\footnote{用标准等价
$\mathcal{M}^{\ell}(Z)\xrightarrow{\sim}\mathcal{M}(Z)$，
$N\mapsto N\otimes\omega_Z$。}%
，并考虑相应的函子
$f^{\dagger}:\mathcal{M}^{r}(Z)\to\mathcal{M}^{r}(Y)$；此时
$f_{\Omega}^{\cdot}(\Omega M)=\Omega f^{\dagger}M[-\dim Y/Z]$。

若 $f$ 光滑，则对任意 $F\in C(Z,\Omega)$ 有
$H_{\mathcal{D}}^{\cdot}f_{\Omega}^{\cdot}F
=f^{\dagger}H_{\mathcal{D}}^{-\dim Y/Z}F$。于是 $f_{\Omega}^{\cdot}$ 保持
$\mathcal{D}$-拟同构，从而有函子
$f_{\Omega}^{\cdot}:D(Z,\Omega)\to D(Y,\Omega)$。伴随态射
$\mathcal{D}f_{\Omega}^{\cdot}(\Omega M)\to f^{\dagger}M[-\dim U/X]$ 是拟同构。
% [Tutor Note: 原书 p.243 此处写作 $[-\dim U/X]$，而上一式用的是 $\dim Y/Z$；
% 底本如此，译文照录。]

\bdp{7.2.9} \emph{引理.} $\Omega$-复形是光滑拓扑下的局部对象，即拉回函子使
$C(U,\Omega)$（$U\in X_{sm}$）成为 $X$ 的光滑拓扑上的 DG 范畴层。
$\mathcal{D}$-拟同构的概念在 $X_{sm}$ 上是局部的。\hfill$\square$

\bdp{7.2.10} 让我们回到 7.2.8 的情形，考察 DG 函子
$f_{\cdot}:C(Y,\Omega)\to C(Z,\Omega)$。右导出函子
$Rf_{\cdot}:D(Y,\Omega)\to D(Z,\Omega)$ 的定义是正确的。事实上，设
$U_{\cdot}$ 是 $Y$ 的（有限）仿射覆盖（平展的或 Zariski 的均可）。对
$F\in C(Y,\Omega)$，设 $F\to C(F)$ 为 $F$ 相应的 Čech 消解。则
\footnote{例如可由 7.3.9 之后的注推出。}
$f_{\cdot}C(F)\xrightarrow{\sim}Rf_{\cdot}F$。

我们把相应的函子 $D(Y)\to D(Z)$ 记为 $f_{*}$。它与通常的 $\mathcal{D}$-模前推
函子一致。事实上，对 $Y$ 上的 $\mathcal{D}$-复形 $M$ 有
\[
\mathcal{D}f_{\cdot}\Omega M=f_{\cdot}(\Omega M\otimes f^{\dagger}\mathcal{D}_Z)
=f_{\cdot}(\mathcal{D}(\Omega M)\otimes_{\mathcal{D}_Y}f^{\dagger}\mathcal{D}_Z).
\]
由于 $f^{\dagger}\mathcal{D}_Z$ 是平坦的 $\mathcal{O}_Y$-模，而
$\mathcal{D}(\Omega M)$ 是 $M$ 的消解，可见
$\mathcal{D}(\Omega M)\otimes_{\mathcal{D}_Y}f^{\dagger}\mathcal{D}_Z
=M\otimes^{L}_{\mathcal{D}_Y}f^{\dagger}\mathcal{D}_Z$。于是
$f_{*}M=Rf_{\cdot}(M\otimes^{L}_{\mathcal{D}_Y}f^{\dagger}\mathcal{D}_Z)$，证毕。

我们留给读者验证：$Rf_{\cdot}$ 与 $f$ 的复合相容，即典范态射
$R(fg)_{\cdot}\to Rf_{\cdot}Rg_{\cdot}$ 是同构%
\footnote{见 7.3.10(ii) 中该陈述在更一般情形下的证明。}%
，并且这一等同 $(fg)_{*}=f_{*}g_{*}$ 与 $\mathcal{D}$-模理论中的标准等同一致。

\SourcePage{244}

\bdp{7.2.11} 对 $Y$ 上的 $\mathcal{D}$-复形 $M$，记
$M_{\mathcal{O}}\in D(Y,\mathcal{O})$ 为把同一个 $M$ 视为
$\mathcal{O}^{!}$-模的复形。有一个典范的\emph{积分}态射
\begin{equation}\tag{309}
i_f:Rf_{\cdot}(M_{\mathcal{O}})\longrightarrow(f_{*}M)_{\mathcal{O}}
\end{equation}
它位于 $D(Y,\mathcal{O})$ 中，定义如下。只需定义态射
$i_f:f_{\cdot}(M_{\mathcal{O}})\to\mathcal{D}(f_{\cdot}\Omega M)$。现在 $i_f$
是如下复合：
\[
f_{\cdot}(M_{\mathcal{O}})\longrightarrow[\mathcal{D}(f_{\cdot}(M_{\mathcal{O}}))]_{\mathcal{O}}
\longrightarrow[\mathcal{D}(f_{\cdot}\Omega M)]_{\mathcal{O}},
\]
其中箭头来自典范态射 $N\to(\mathcal{D}N)_{\mathcal{O}}$
（取 $N=f_{\cdot}(M_{\mathcal{O}})$）与 $M_{\mathcal{O}}\to\Omega M$。换言之，
$i_f$ 由对显然的态射
$M_{\mathcal{O}}\to(M\otimes^{L}_{\mathcal{D}_Y}f^{\dagger}\mathcal{D}_Z)_{\mathcal{O}}$
施用 $Rf_{\cdot}$ 而得。
% [Tutor Note: 原书 p.244 此处说 $i_f$ 位于 $D(Y,\mathcal{O})$ 中，但按定义
% $i_f$ 取值于 $D(Z,\mathcal{O})$；底本如此，译文照录。]

我们留给读者验证 $i_f$ 与 $f$ 的复合相容。

\bdsec{7.3. 光滑叠上的 $\mathcal{D}$-模理论 I}

% [Tutor Note: 原书 p.244 的 7.3 节标题是行内（run-in）粗斜体，
% 与正文同一行开始；译文按全书的节标题体例单独成行，编号与标题文字不变。]

我们为满足下文条件（310）的光滑叠建立基本的 $\mathcal{D}$-模形式化。
在 7.3.12 中我们修改定义，使得可以去掉拟紧性假设。任意的光滑叠将在 7.5 中处理。

\bdp{7.3.1} 设 $\mathcal{Y}$ 为光滑拟紧代数叠。假设它满足如下条件%
\footnote{这一条件是为了保证我们定义的范畴 $D(\mathcal{Y})$ 具有右局部-整体性质，
见 7.5.3。7.3.1--7.3.3 的构造对任意光滑代数栈都有意义。}%
：
\begin{equation}\tag{310}
\text{对角态射 }\mathcal{Y}\to\mathcal{Y}\times\mathcal{Y}\text{ 是仿射的。}
\end{equation}
等价地说，这意味着存在光滑仿射满射 $U\to\mathcal{Y}$，使 $U$ 是仿射概形。
换言之，把光滑代数簇 $X$ 模去光滑群胚 $Q$%
\footnote{$Q=X\times_{\mathcal{Y}}X$。}%
的作用，所得商即 $\mathcal{Y}$，且结构态射 $Q\to X\times X$ 是仿射的。

\SourcePage{245}

注意 $\Omega(U)$（$U\in\mathcal{Y}_{sm}$）构成 $\mathcal{Y}_{sm}$ 上的 DG 代数层
$\Omega_{\mathcal{Y}}$。$\mathcal{Y}$ 上的 \emph{$\Omega$-复形}是指这样的
DG $\Omega_{\mathcal{Y}}$-模：它作为 $\mathcal{O}_{\mathcal{Y}}$-模是拟凝聚的。
记 $\mathcal{Y}$ 上 $\Omega$-复形的 DG 范畴为 $C(\mathcal{Y},\Omega)$。

\emph{注.} 范畴 $C(U,\Omega)$（$U\in\mathcal{Y}_{sm}$）构成 $\mathcal{Y}_{sm}$
上的 DG 范畴层 $C(\mathcal{Y}_{sm},\Omega)$（见 7.2.9），而 $\mathcal{Y}$ 上的
$\Omega$-复形就是 $C(\mathcal{Y}_{sm},\Omega)$ 的笛卡尔截面。等价地，
$\mathcal{Y}$ 上的 $\Omega$-复形就是 $X$ 上 $Q$-等变的 $\Omega$-复形。

\bdp{7.3.2} 回忆：$\mathcal{D}$-模的范畴 $\mathcal{M}(U)$（$U\in\mathcal{Y}_{sm}$）
构成 $\mathcal{Y}_{sm}$ 上的阿贝尔范畴层，而 $\mathcal{Y}$ 上的
$\mathcal{D}$-模范畴 $\mathcal{M}(\mathcal{Y})$ 即其笛卡尔截面范畴。由 7.2.8 有
典范的上同调函子
$H_{\mathcal{D}}^{\cdot}:C(\mathcal{Y},\Omega)\to\mathcal{M}(\mathcal{Y})$，
$H_{\mathcal{D}}^{\cdot}(F)_U:=H_{\mathcal{D}}^{\cdot+\dim U/\mathcal{Y}}(F_U)$。
$\Omega$-复形的态射若诱导各 $H_{\mathcal{D}}^{\cdot}$ 的同构，则称为
$\mathcal{D}$-拟同构。把 $\Omega$-复形的同伦范畴对 $\mathcal{D}$-拟同构作局部化，
得到三角范畴 $\mathcal{D}(\mathcal{Y})=\mathcal{D}(\mathcal{Y},\Omega)$。
相应地也有有界导出范畴。

有一个全忠实嵌入 $\mathcal{M}(\mathcal{Y})\hookrightarrow\mathcal{D}(\mathcal{Y})$，
它把 $\mathcal{Y}$ 上的 $\mathcal{D}$-模 $M$ 送到其 de Rham 复形 $\Omega M$，
$(\Omega M)_U:=\Omega M_U[-\dim U/\mathcal{Y}]$。有
$H_{\mathcal{D}}^0\Omega M=M$，且对 $a\neq0$ 有 $H_{\mathcal{D}}^a\Omega M=0$。
容易看出，$\Omega$ 把 $\mathcal{M}(\mathcal{Y})$ 与 $\mathcal{D}(\mathcal{Y})$ 中
由那些满足「对 $a\neq0$ 有 $H_{\mathcal{D}}^a(F)=0$」的 $\Omega$-复形 $F$ 组成的
全子范畴等同起来。

\bdp{7.3.3} \emph{例.} 记 $\Omega\mathcal{D}_{\mathcal{Y}}$ 为 $\mathcal{Y}$ 上
由 $\Omega\mathcal{D}_{\mathcal{Y}U}:=\Omega_{U/\mathcal{Y}}[\dim\mathcal{Y}]$
定义的 $\Omega$-复形。注意对 $a>0$ 有
$H_{\mathcal{D}}^a(\Omega\mathcal{D}_{\mathcal{Y}})=0$。若 $\mathcal{Y}$ 是良叠，
则我们的 $\Omega$-复形属于 $\mathcal{M}(\mathcal{Y})$ 的本质像；相应的
$\mathcal{D}$-模 $\mathcal{D}_{\mathcal{Y}}
=H_{\mathcal{D}}^0(\Omega\mathcal{D}_{\mathcal{Y}})$ 与 1.1.3 中的左
$\mathcal{D}$-模 $\mathcal{D}_{\mathcal{Y}}$ 一致。更一般地，对 $\mathcal{Y}$ 上
任意 $\mathcal{O}$-模 $P$，有 $\Omega$-复形
$\Omega(\mathcal{D}_{\mathcal{Y}}\otimes P)$，其中
\[
\Omega(\mathcal{D}_{\mathcal{Y}}\otimes P)_U:=
\Omega_{U/\mathcal{Y}}\otimes_{\mathcal{O}_{\mathcal{Y}}}P_U[\dim\mathcal{Y}].
\]
若 $\mathcal{Y}$ 是良叠且 $P$ 局部自由，则我们的 $\Omega$-复形落在
$\mathcal{M}(\mathcal{Y})$ 中，且等于左 $\mathcal{D}$-模
$\mathcal{D}_{\mathcal{Y}}\otimes P
=\mathcal{D}_{\mathcal{Y}}\otimes_{\mathcal{O}_{\mathcal{Y}}}P$。

% 尾部接 p.246（本单元末句完结于 D_𝒴⊗P = D_𝒴⊗_{O_𝒴}P，7.3.3 至此结束）。
%% end units/ch07_u01.tex
%% begin units/ch07_u02.tex
% ch07_u02 —— 原书 p.246–p.255（第 7 章 7.3.4–7.4.10）
% 底本：A. Beilinson, V. Drinfeld, Quantization of Hitchin's Integrable System
% and Hecke Eigensheaves，附录「代数叠上的 D-模理论与 Hecke 模式」

\SourcePage{246}

以 $D(\mathcal{Y})^{\ge0}\subset D(\mathcal{Y})$ 记由下述 Ω-复形 $F$ 构成的满子范畴：当 $a<0$ 时 $H^{a}_{D}F=0$。类似地定义 $D(\mathcal{Y})^{\le0}$。

\bdp{7.3.4} \emph{命题.} 这给出 $D(\mathcal{Y})$ 上的一个 t-结构，其核心为 $\mathcal{M}(\mathcal{Y})$，上同调函子为 $H_{\mathcal{D}}$。

此命题由下文的引理 7.5.3 立即推出。当 $\mathcal{Y}$ 是轨道叠时，另有一个证明，见 7.6.11。

\bdp{7.3.5} \emph{注.} 考虑函子 $\Omega:C(\mathcal{M}(\mathcal{Y}))\to C(\mathcal{Y},\Omega)$。对 $M\in C(\mathcal{M}(\mathcal{Y}))$ 有 $H^{\cdot}M=H_{\mathcal{D}}(\Omega M)$，故 Ω 给出 t-正合函子 $\Omega:D(\mathcal{M}(\mathcal{Y}))\to D(\mathcal{Y})$，它延拓两个核心之间的「恒等」等价。若 $\mathcal{Y}$ 是 Deligne--Mumford 叠\footnote{即存在一个簇，它平展地覆盖 $\mathcal{Y}$。此时函子 $D:C(\mathcal{Y},\Omega)\to C(\mathcal{M}(\mathcal{Y}))$ 显然有意义（它给出逆等价 $D(\mathcal{M}(\mathcal{Y}))\to D(\mathcal{Y})$，见 7.2.7）。}，则此函子是范畴等价；但一般情形下不然。

\bdp{7.3.6} 设 $f:\mathcal{Y}\to\mathcal{Z}$ 是满足（310）的光滑叠的态射。它给出 DG 赋环拓扑的态射 $(\mathcal{Y}_{sm},\Omega_{\mathcal{Y}})\to(\mathcal{Z}_{sm},\Omega_{\mathcal{Z}})$，从而给出一对伴随 DG 函子
\begin{equation}\tag{311}
f_{\Omega}^{\cdot}:C(\mathcal{Z},\Omega)\to C(\mathcal{Y},\Omega),\qquad
f_{\cdot}:C(\mathcal{Y},\Omega)\to C(\mathcal{Z},\Omega),
\end{equation}
以及同伦范畴之间相应的伴随三角函子（因 $\mathcal{Y}$ 拟紧，$f_{\cdot}$ 保持拟凝聚）。

若 $f$ 光滑，则 $f_{\Omega}^{\cdot}$ 保持 $\mathcal{D}$-拟同构，故它定义 t-正合函子 $f^{\cdot}:D(\mathcal{Z})\to D(\mathcal{Y})$。显然它与 $f$ 的复合相容。
% [Tutor Note: 原书 p.246 此句印作 "f_Ω preservres D-quasi-isomorphisms"，"preservres" 系 "preserves" 之误，底本如此，保留原样。]

设 $f$ 任意。把推动函子 $f_{*}:D^{+}(\mathcal{Y})\to D^{+}(\mathcal{Z})$ 定义为右导出函子 $Rf_{\cdot}$。我们将在 7.3.10 证明 $f_{*}$ 定义正确。为此需要足够多的「脆弱」对象。

\SourcePage{247}

\bdp{7.3.7} \emph{定义.} 称 $\mathcal{Y}$ 上的 $\mathcal{O}$-模 $F$ 为松的，若对 $\mathcal{Y}$ 上任意平坦 $\mathcal{O}$-模 $P$ 都有 $H^{a}(\mathcal{Y},P\otimes F)=0$（$a>0$）。$\mathcal{O}$-复形或 Ω-复形 $F$ 为松的，若每个 $F^{i}$ 都松。

\bdp{7.3.8} \emph{引理.} (i) 对 $\mathcal{Y}$ 上任意 Ω-复形 $F'$，存在 $\mathcal{D}$-拟同构 $F'\to F$，且 $F$ 松。若 $F'$ 下有界，则可取 $F$ 下有界。(ii) 设 $f$（见 7.3.6）光滑且仿射。则 $f_{\Omega}^{\cdot}$ 与 $f_{\cdot}$ 把松 Ω-复形送到松 Ω-复形。(iii) 若 $F_{1},F_{2}$ 分别是叠 $\mathcal{Y}_{1},\mathcal{Y}_{2}$ 上的松 Ω-复形，则 $F_{1}\boxtimes F_{2}$ 是 $\mathcal{Y}_{1}\times\mathcal{Y}_{2}$ 上的松 Ω-复形。

\emph{证明.} (i) 因 $\mathcal{Y}$ 拟紧，存在 $\mathcal{Y}$ 的超覆叠 $U_{\cdot}$，其中各 $U_{a}$ 是仿射概形。因 $\mathcal{Y}$ 的对角态射仿射，投影 $\pi_{a}:U_{a}\to\mathcal{Y}$ 也仿射。取 $F$ 为 $F'$ 关于此超覆叠的 Čech 复形，即
\[
F^{i}=\bigoplus_{a\ge0}\pi_{a\cdot}\bigl(F'^{i-a}_{U_{a}}\bigr).
\]
(ii) 显然。(iii) 可设 $F_{i}$ 是松 $\mathcal{O}_{\mathcal{Y}_{i}}$-模。设 $P$ 是 $\mathcal{Y}_{1}\times\mathcal{Y}_{2}$ 上的平坦 $\mathcal{O}$-模。因 $F_{1}$ 松，对 $a>0$ 有 $R^{a}p_{2\cdot}(P\otimes p_{1}^{*}F_{1})=0$，且 $p_{2\cdot}(P\otimes p_{1}^{*}F_{1})$ 是 $\mathcal{Y}_{2}$ 上的平坦 $\mathcal{O}$-模（这里 $p_{i}:\mathcal{Y}_{1}\times\mathcal{Y}_{2}\to\mathcal{Y}_{i}$ 是投影）。于是
\[
H^{a}\bigl(\mathcal{Y}_{1}\times\mathcal{Y}_{2},P\otimes(F_{1}\boxtimes F_{2})\bigr)
=H^{a}\bigl(\mathcal{Y}_{2},\bigl(p_{2\cdot}(P\otimes p_{1}^{*}F_{1})\bigr)\otimes F_{2}\bigr),
\]
因 $F_{2}$ 松，此式当 $a>0$ 时为零。\hfill$\Box$

回到 7.3.6 末尾的情形。

\bdp{7.3.9} \emph{引理.} 若 $F$ 是 $\mathcal{Y}$ 上下有界的松 Ω-复形，则 $f_{\cdot}F=Rf_{\cdot}F$。

\emph{证明.} 只需验证：若 $F$ 还是 $\mathcal{D}$-无环的（即满足 $H^{\cdot}_{D}F=0$），则 $f_{\cdot}F$ 也 $\mathcal{D}$-无环（用 7.3.8(i)）。

a. 可设 $\mathcal{Z}$ 是光滑仿射概形 $Z$。事实上，待验证的陈述关于 $\mathcal{Z}$ 是局部的。把 $\mathcal{Z}$ 换成 $\mathcal{Z}_{sm}$ 中的仿射对象 $Z$，把 $\mathcal{Y}$ 换成 $\mathcal{Y}\times_{\mathcal{Z}}Z$，
\SourcePage{248}%
并把 $F$ 换成它到 $\mathcal{Y}\times_{\mathcal{Z}}Z$ 的拉回。新数据满足引理的全部条件。

b. 可设 $\mathcal{Y}$ 是光滑仿射概形 $Y$。事实上，取 (i) 中的 $U_{\cdot}$，记 $A$ 为项是 $A^{i}=\bigoplus_{a\ge0}(f\pi_{a})_{\cdot}(F'^{i-a}_{U_{a}})$ 的 Čech 复形。$A$ 是 $\mathcal{Z}$ 上的 Ω-复形。因 $F$ 松，显然的态射 $f_{\cdot}F\to A$ 是 $\mathcal{D}$-拟同构（用（310））。注意 $A$ 带显然的滤链，其相继商为 $(f\pi_{a})_{\cdot}(F_{U_{a}})[-a]$。若这些复形都 $\mathcal{D}$-无环，则 $A$ 也 $\mathcal{D}$-无环（此处用 $F$ 下有界），故 $f_{\cdot}F$ 是 $\mathcal{D}$-无环的。

c. 设 $i:Y\to Y\times Z$ 是 $f$ 的图像嵌入。则 $G:=i_{\cdot}F$ 是 $\mathcal{D}$-无环的。因 $f_{\cdot}F=p_{\cdot}G$（$p$ 是投影 $Y\times Z\to Z$），只需证 $p_{\cdot}G$ 是 $\mathcal{D}$-无环的。设 $T$ 是 $\mathcal{D}G$ 沿 $p$ 的纤维的相对 de Rham 复形。处于直积情形，故 $p_{\cdot}T$ 是 $Z$ 上的 $\mathcal{D}$-复形。有显然的 $\mathcal{D}$-复形态射 $Dp_{\cdot}G\to p_{\cdot}T$，它是拟同构。因 $p_{\cdot}T$ 无环（$T$ 带滤链，其相继商为 $DG\otimes\Lambda\Theta_{Y}$，而 $\mathcal{D}G$ 无环），证毕。\hfill$\Box$

\emph{注.} 设 $f$ 是仿射态射。则对任意 $F\in C(\mathcal{Y},\Omega)$ 有 $f_{\cdot}F=Rf_{\cdot}F$。事实上，该陈述关于 $\mathcal{Z}$ 是局部的，故可设 $\mathcal{Z}$ 是仿射概形。此时 $\mathcal{Y}$ 是仿射概形，于是 $\mathcal{Y}$ 上任意复形都松；再用 7.3.9 即可。

\bdp{7.3.10} \emph{推论.} (i) 函子 $f_{*}:=Rf_{\cdot}:D^{+}(\mathcal{Y})\to D^{+}(\mathcal{Z})$ 定义正确。(ii) $f_{*}$ 与 $f$ 的复合相容，即典范态射 $(f_{1}f_{2})_{*}\to f_{1*}f_{2*}$ 是同构。

\emph{证明.} (i) 用 7.3.8(i) 与 7.3.9。(ii) $f_{\cdot}$ 把松 Ω-复形送到松 Ω-复形。\hfill$\Box$

\bdp{7.3.11} \emph{注.} (i) 上述引理在 $\mathcal{O}$-复形的框架中也成立。

\SourcePage{249}
(ii) 设 $\mathcal{Y}$ 的 $\mathcal{O}$-模范畴上的函子 $f_{\cdot}$ 有有限的上同调维数（例如 $f$ 可表时如此）。则对任意边界条件的 Ω-复形的导出范畴，$f_{*}:=Rf_{\cdot}$ 都定义正确。事实上，7.3.9（连同其证明）对无界的松 Ω-复形仍成立。

(iii) 若我们的叠是光滑簇，则上述函子 $f_{*}$ 就是 $\mathcal{D}$-模论中标准的推动函子（见 7.2.10）。此时引理 7.3.9（连同其证明）仍成立，只需假定：对 $Y$ 的任意使 $U\to Y$ 为仿射态射的 Zariski 开集 $U$，上同调 $H^{a}(U,F^{i})$（$a>0$）为零。

\bdp{7.3.12} 现在设 $\mathcal{Y}$ 是任一光滑叠，其对对角态射 $\mathcal{Y}\to\mathcal{Y}\times\mathcal{Y}$ 仿射（即放弃拟紧性假设）。此时 $\mathcal{Y}$ 上 Ω-复形的范畴可能太小，不足以定义正确的 $\mathcal{D}$-模导出范畴。我们如下扩充上述形式体系。

为简化记号，设 $\mathcal{Y}$ 有可数个拟紧开集构成的覆盖。换言之，$\mathcal{Y}$ 是递增的开拟紧子叠序列 $\mathcal{Y}_{1}\subset\mathcal{Y}_{2}\subset\ldots$ 之并。
% [Tutor Note: 原书 p.249 此处印作 "open quasi-comact substacks"，"quasi-comact" 系 "quasi-compact" 之误，底本如此，保留原样。]
$\mathcal{Y}_{\cdot}$ 上的 Ω-复形是集合 $F=(F_{i},a_{i})$，其中 $F_{i}$ 是 $\mathcal{Y}_{i}$ 上的 Ω-复形，而 $a_{i}:F_{i+1}|_{\mathcal{Y}_{i}}\to F_{i}$ 是 $\mathcal{D}$-拟同构。这样的 Ω-复形构成 DG 范畴 $C(\mathcal{Y}_{\cdot},\Omega)$，于是有相应的同伦范畴 $K(\mathcal{Y}_{\cdot},\Omega)$。它带有上同调函子 $H_{\mathcal{D}}$，取值于 $\mathcal{Y}$ 上 $\mathcal{D}$-模的阿贝尔范畴 $\mathcal{M}(\mathcal{Y})$：
\[
H_{\mathcal{D}}(F)|_{\mathcal{Y}_{i}}=H_{\mathcal{D}}(F_{i}).
\]

把 $D(\mathcal{Y}_{\cdot},\Omega)$ 定义为 $K(\mathcal{Y}_{\cdot},\Omega)$ 关于 $\mathcal{D}$-拟同构的局部化。不同 $\mathcal{Y}_{\cdot}$ 给出的三角范畴 $D(\mathcal{Y}_{\cdot},\Omega)$ 被典范地等同。事实上，设 $\mathcal{Y}'_{j}$ 是 $\mathcal{Y}$ 的另一个如上所述的开子叠序列。选取递增函数 $j=j(i)$，使 $\mathcal{Y}_{i}\subset\mathcal{Y}'_{j(i)}$。令 $\mathcal{Y}'_{\cdot}$ 上的 Ω-复形 $F'$ 对应 $\mathcal{Y}_{\cdot}$ 上的 Ω-复形 $F$，其中 $F^{i}=F'^{j(i)}|_{\mathcal{Y}_{i}}$。此函子与 $H_{\mathcal{D}}$ 交换。\SourcePage{250}%
导出范畴之间的相应函子在显然意义下不依赖辅助选取 $j(i)$，且是范畴等价。

我们看到，范畴 $D(\mathcal{Y}_{\cdot},\Omega)$ 只依赖于 $\mathcal{Y}$，故把它记作 $D(\mathcal{Y},\Omega)$ 或简记为 $D(\mathcal{Y})$。我们的三角范畴带有上同调函子 $H_{\mathcal{D}}:D(\mathcal{Y})\to\mathcal{M}(\mathcal{Y})$，并有典范的全忠实嵌入 $\Omega:\mathcal{M}(\mathcal{Y})\to D(\mathcal{Y})$（见 7.3.2）。命题 7.3.4 仍然成立；其证明由 7.5.4 推出。

设 $f:\mathcal{Y}\to\mathcal{Z}$ 是满足我们假设的光滑叠的态射。若 $f$ 光滑，则按显然方式定义 t-正合拉回函子 $f^{\cdot}:D(\mathcal{Z})\to D(\mathcal{Y})$。若 $f$ 是任意拟紧态射，则有典范推动函子 $f_{*}:D(\mathcal{Y})^{+}\to D(\mathcal{Z})^{+}$。我们在关于松 Ω-复形的简短插叙之后定义它。

按定义，$F\in C(\mathcal{Y}_{\cdot},\Omega)$ 是松的，若所有 $F_{i}\in C(\mathcal{Y}_{i},\Omega)$ 都松。引理 7.3.8(i),(iii) 在我们的情形仍成立。这意味着，可以只用松复形来定义 $\mathcal{D}$-导出范畴。为证 7.3.8(i)，取覆盖 $\pi_{i}:V_{i}\to\mathcal{Y}_{i}$，其中 $V_{i}$ 是仿射概形。以 $U_{i}$ 记诸 $V_{j}$（$1\le j\le i$）的不交并，以 $U_{i\cdot}$ 记 $\mathcal{Y}_{i}$ 的相应超覆叠，$U_{ia}$ 是 $U_{i}$ 在 $\mathcal{Y}_{i}$ 上的 $a$ 重纤维积。现取任意 $F'\in C(\mathcal{Y}_{\cdot},\Omega)$。设 $F_{i}$ 是 $F'_{i}$ 关于超覆叠 $U_{i\cdot}$ 的 Čech 复形（见 7.3.8(i) 的证明）。则诸 $F_{i}$ 按显然方式构成 $\mathcal{Y}_{\cdot}$ 上的 Ω-复形 $F$。此 $F$ 是松的，且显然的态射 $F'\to F$ 是 $\mathcal{D}$-拟同构，证毕。

现定义 $f_{*}$。设 $\mathcal{Z}_{i}$ 是 $\mathcal{Z}$ 的如上所述的开拟紧子叠序列。则 $\mathcal{Y}_{i}:=f^{-1}\mathcal{Z}_{i}$ 是 $\mathcal{Y}$ 的相应序列。设 $F$ 是 $\mathcal{Y}_{\cdot}$ 上下有界的松 Ω-复形。则 $(f_{\cdot}F)_{i}:=f_{\cdot}(F_{i})$ 构成 $\mathcal{Z}_{\cdot}$ 上的 Ω-复形 $f_{\cdot}F$（用 7.3.9）。函子 $f_{\cdot}$ 保持 $\mathcal{D}$-拟同构（由 7.3.9）。我们的 $f_{*}$ 就是 $\mathcal{D}$-导出范畴之间的相应函子。推论 7.3.10(ii) 连同其证明仍成立。

\SourcePage{251}

再设所有函子 $f_{i\cdot}:\mathcal{M}(\mathcal{Y}_{i},\mathcal{O})\to\mathcal{M}(\mathcal{Z}_{i},\mathcal{O})$ 都有有限的上同调维数（例如 $f$ 可表时如此）。则函子 $f_{*}$ 在整个 $D(\mathcal{Y})$ 上定义正确。事实上，设 $F$ 是 $\mathcal{Y}_{\cdot}$ 上任意松 Ω-复形。则 $(f_{\cdot}F)_{i}:=f_{\cdot}(F_{i})$ 构成 $\mathcal{Z}_{\cdot}$ 上的 Ω-复形 $f_{\cdot}F$（用 7.3.11(ii)）。函子 $f_{\cdot}$ 保持 $\mathcal{D}$-拟同构，故把 $f_{*}:D(\mathcal{Y})\to D(\mathcal{Z})$ 定义为 $\mathcal{D}$-导出范畴之间的相应函子。

\bdp{7.3.13} \emph{注.} 设 $A$ 是交换代数。以 $\mathcal{M}(\mathcal{Y},A)$ 记 $\mathcal{Y}$ 上带 $A$-作用的 $\mathcal{D}$-模构成的阿贝尔范畴。如 7.3.12 那样，用带 $A$-作用的 Ω-复形定义以 $\mathcal{M}(\mathcal{Y},A)$ 为核心的 t-范畴 $D(\mathcal{Y},A)$。标准的函子毫无困难地迁移到 $A$-线性的情形。更一般地，设 $A_{\mathcal{Y}}$ 是 $\mathcal{Y}$ 上的交换 $\mathcal{D}$-代数（即张量范畴 $\mathcal{M}^{\ell}(\mathcal{Y})$ 中的交换代数）。有 $A_{\mathcal{Y}}$-模的阿贝尔范畴 $\mathcal{M}(\mathcal{Y},A_{\mathcal{Y}})$ 及其导出版本 $D(\mathcal{Y},A_{\mathcal{Y}})$，后者如 7.3.12 那样用带 $A_{\mathcal{Y}}$-作用的 Ω-复形定义。

\bdsec{7.4. 导出范畴的下降}

我们解释导出范畴的一般同伦逆向极限构造。为能给出 $\mathcal{D}$-模导出范畴的「局部」定义，需要这一构造。

\bdp{7.4.1} 以 $(\Delta)$ 记非空有限全序集 $\Delta_{n}=[0,n]$ 与递增单射构成的范畴。设 $\mathcal{M}_{\cdot}$ 是余纤维化于 $(\Delta)$ 上的一族阿贝尔范畴，使得对任意态射 $\alpha:\Delta_{n}\hookrightarrow\Delta_{m}$，相应的函子 $\alpha_{\cdot}:\mathcal{M}_{n}\to\mathcal{M}_{m}$ 正合。

以 $\mathcal{M}_{\mathrm{tot}}$ 记 $\mathcal{M}_{\cdot}$ 的余笛卡儿截面构成的范畴：$\mathcal{M}_{\mathrm{tot}}$ 的对象是集合 $M=\{M_{n},\alpha^{*}\}$，其中 $M_{n}\in\mathcal{M}_{n}$，$\alpha^{*}=\alpha^{*}_{M}:\alpha_{\cdot}M_{n}\xrightarrow{\sim}M_{m}$ 是同构，且满足 $(\alpha\beta)^{*}=\alpha^{*}\alpha_{\cdot}(\beta^{*})$（这里 $\beta:\Delta_{l}\hookrightarrow\Delta_{n}$）。这是阿贝尔范畴。注意 $\mathcal{M}_{\mathrm{tot}}$ 与对偶相容：有 $(\mathcal{M}_{\mathrm{tot}})^{\circ}=(\mathcal{M}^{\circ})_{\mathrm{tot}}$。

\SourcePage{252}

我们的目标是定义以 $\mathcal{M}_{\mathrm{tot}}$ 为核心的 t-范畴 $D_{\mathrm{tot}}(\mathcal{M}_{\cdot})$，使其具有下述关键性质：
\begin{equation}\tag{312}
\begin{gathered}
\text{对任意 }M,N\in\mathcal{M}_{\mathrm{tot}}\text{，存在收敛到 }\operatorname{Ext}^{p+q}_{D_{\mathrm{tot}}(\mathcal{M}_{\cdot})}(N,M)\text{ 的典范谱序列 }E^{p,q}_{r}\text{，且}\\
E^{p,q}_{1}=\operatorname{Ext}^{q}_{\mathcal{M}_{p}}(N_{p},M_{p}).
\end{gathered}
\end{equation}
$D_{\mathrm{tot}}(\mathcal{M}_{\cdot})$ 的构造与对偶相容。

\bdp{7.4.2} 考虑范畴 $\sec_{+}=\sec_{+}(\mathcal{M}_{\cdot})$，其对象是集合 $M=(M_{n},\alpha^{*})$，其中 $M_{n}\in\mathcal{M}_{n}$，$\alpha^{*}=\alpha^{*}_{M}:\alpha_{\cdot}M_{n}\to M_{m}$ 是使 $(\alpha\beta)^{*}=\alpha^{*}\alpha_{\cdot}(\beta^{*})$、$\mathrm{id}^{*}_{\Delta_{n}}=\mathrm{id}_{M_{n}}$ 的态射。这是阿贝尔范畴，它包含 $\mathcal{M}_{\mathrm{tot}}$ 作为对扩张封闭的满子范畴。用对偶定义 $\sec_{-}=\sec_{-}(\mathcal{M}_{\cdot})$：$\sec_{-}(\mathcal{M}_{\cdot}):=(\sec_{+}(\mathcal{M}^{\circ}_{\cdot}))^{\circ}$；于是 $\sec_{-}$ 的对象是集合 $N=(N_{n},\alpha_{*})$，$N_{n}\in\mathcal{M}_{n}$，$\alpha_{*}=\alpha^{N}_{*}:N_{m}\to\alpha_{\cdot}N_{n}$。

考虑 $\sec_{\pm}$ 中复形构成的 DG 范畴 $C\sec_{\pm}$ 及相应的同伦范畴 $K\sec_{\pm}$。有伴随 DG 函子
\begin{equation}\tag{313}
c_{+}:C\sec_{-}\longrightarrow C\sec_{+},\qquad c_{-}:C\sec_{+}\longrightarrow C\sec_{-},
\end{equation}
它们定义如下。取 $M\in C\sec_{+}$。则对任意 $m\ge0$，有单纯形 $\Delta_{m}$ 上取值于 $C\mathcal{M}_{m}$ 的「上同调型」系数系统 $\widetilde{M}_{m}$。确切地说，$\widetilde{M}_{m}$ 把面 $\alpha:\Delta_{n}\hookrightarrow\Delta_{m}$ 对应到复形 $\alpha_{\cdot}M_{n}$；若 $\alpha':\Delta_{l}\hookrightarrow\Delta_{m}$ 是 $\alpha$ 的面，即 $\alpha'=\alpha\beta$，则相应的连接态射 $\alpha'_{\cdot}M_{l}\to\alpha_{\cdot}M_{n}$ 是 $\alpha_{\cdot}(\beta^{*})$。此时 $(c_{-}M)_{m}$ 是全上链复形 $C^{\cdot}(\Delta_{m},\widetilde{M}_{m})$（即 $c_{-}(M)^{\cdot}_{m}=\bigoplus_{\alpha:\Delta_{n}\to\Delta_{m}}\alpha_{\cdot}M_{n}^{\cdot-n}$），而 $\alpha^{c_{-}(M)}_{*}$ 是显然的投影。用对偶定义 $c_{+}$。

为看出 $c_{\pm}$ 伴随，对如上 $N,M$ 考虑阿贝尔群复形 $Hom(N,M)$，其项为
\[
Hom(N,M)^{i}=\prod_{a,n}\operatorname{Hom}\bigl(N^{a+n}_{n},M^{a+i}_{n}\bigr),
\]
而微分为 $df$，它把 $f=(f_{a,n})\in Hom(N,M)^{i}$ 送到
% [Tutor Note: 原书 p.252 此处印作 "the differential wich sends"，"wich" 系 "which" 之误，底本如此，保留原样。]
\[
(df)_{a,n}=df_{a,n}-(-1)^{i+n}f_{a+1,n}d+\sum_{j=0,\ldots,n}(-1)^{j}\alpha^{*}_{j}\alpha_{j\cdot}(f_{a+1,n-1})\alpha_{j*}.
\]

\SourcePage{253}

这里 $\alpha_{j}:\Delta_{n-1}\to\Delta_{n}$ 是第 $j$ 个面嵌入。于是伴随性由同态复形的显然等同
\begin{equation}\tag{314}
Hom(c_{+}N,M)\xrightarrow{\sim}Hom(N,M)\xleftarrow{\sim}Hom(N,c_{-}M)
\end{equation}
推出。

\bdp{7.4.3} \emph{注.} 固定某个 $m\ge0$。对 $i=0,\ldots,m$，设 $\nu_{i}:c_{-}(M)_{m}\to M_{m}$ 是投影 $c_{-}(M)_{m}\to\alpha_{i\cdot}M_{0}$ 与 $\alpha^{*}_{i}:\alpha_{i\cdot}M_{0}\to M_{m}$ 的复合；这里 $\alpha_{i}:\Delta_{0}\to\Delta_{m}$ 是第 $i$ 个顶点。此时诸态射 $\nu_{i}$ 彼此同伦（伴有典范同伦与「高阶同伦」）。

\bdp{7.4.4} \emph{引理.} 函子 $c_{\pm}$ 保持拟同构。伴随态射 $c_{+}c_{-}M\to M$、$N\to c_{-}c_{+}N$ 都是拟同构。\hfill$\Box$

由此可见，$c_{\pm}$ 给出导出范畴 $D\sec_{\pm}$ 之间互逆的等价。把如此等同的这两个范畴记作 $D\sec$。于是 $D\sec$ 带两个 t-结构，其核心分别为 $\sec_{\pm}$，上同调函子为 $H_{\pm}:D\sec\to\sec_{\pm}$。

\bdp{7.4.5} 设 $C_{\mathrm{tot}+}\subset C\sec_{+}$ 是由下述复形 $M$ 构成的满子范畴：对任意 $i$ 有 $H^{i}M\in\mathcal{M}_{\mathrm{tot}}\subset\sec_{+}$。换言之，$M\in C\sec_{+}$ 属于 $C_{\mathrm{tot}+}$，当且仅当所有态射 $\alpha^{*}_{M}$ 都是拟同构。类似地定义 $C_{\mathrm{tot}-}\subset C\sec_{-}$。以 $K_{\mathrm{tot}\pm}\subset K\sec_{\pm}$、$D_{\mathrm{tot}\pm}\subset D\sec_{\pm}$ 记相应的同伦范畴与导出范畴；它们是三角范畴。

导出范畴 $D(\mathcal{M}_{n})$ 构成 $(\Delta)$ 上的余纤维化范畴。以 $D^{\mathrm{fake}}_{\mathrm{tot}}$ 记其余笛卡儿截面构成的范畴（这\emph{不是}三角范畴！）。$\mathcal{M}_{\cdot}$ 的上同调函子给出函子 $H:D^{\mathrm{fake}}_{\mathrm{tot}}\to\mathcal{M}_{\mathrm{tot}}$。有显然的函子 $\epsilon_{+}:D_{\mathrm{tot}+}\to D^{\mathrm{fake}}_{\mathrm{tot}}$，它把 $M$ 对应到视为 $D^{\mathrm{fake}}_{\mathrm{tot}}$ 对象的资料 $(M_{n},\alpha_{*})$。类似地有函子 $\epsilon_{-}:D_{\mathrm{tot}}\to D^{\mathrm{fake}}_{\mathrm{tot}}$。

\bdp{7.4.6} \emph{引理.} 对任意 $M\in D_{\mathrm{tot}+}$ 有 $c_{-}M\in D_{\mathrm{tot}-}$，且存在唯一同构 $\epsilon_{-}(c_{-}M)\xrightarrow{\sim}\epsilon_{+}(M)$，其第 $0$ 个分量是 $\mathrm{id}_{M_{0}}$。把 $+$ 与 $-$ 互换，即得对偶的陈述。

\emph{证明.} 用 7.4.3。\hfill$\Box$

\SourcePage{254}

\bdp{7.4.7} 我们看到，函子 $c_{\pm}$ 等同三角范畴 $D_{\mathrm{tot}\pm}$。换言之，子范畴 $D_{\mathrm{tot}\pm}\subset D\sec$ 重合；这就是 7.4.1 中许诺的范畴 $D_{\mathrm{tot}}=D_{\mathrm{tot}}(\mathcal{M}_{\cdot})$。函子 $\epsilon_{\pm}$ 被典范地等同，故有函子 $\epsilon:D_{\mathrm{tot}}\to D^{\mathrm{fake}}_{\mathrm{tot}}$。注意 $H_{\pm}=H\epsilon$，故有典范上同调函子 $H:D_{\mathrm{tot}}\to\mathcal{M}_{\mathrm{tot}}$。这是 $D_{\mathrm{tot}}$ 上以 $\mathcal{M}_{\mathrm{tot}}$ 为核心的非退化 t-结构的上同调函子。注意嵌入 $D_{\mathrm{tot}}\to D\sec$ 关于 $D\sec$ 上两个 $\pm$ t-结构中的任一个都 t-正合；它把核心 $\mathcal{M}_{\mathrm{tot}}$ 与核心 $\sec_{+}$、$\sec_{-}$ 之交等同。

\bdp{7.4.8} 我们从 7.4.1 导出谱序列（312）。更一般地，考虑对象 $N\in D^{-}\sec_{-}\subset D\sec$、$M\in D^{+}\sec_{+}\subset D\sec$。用复形 $N\in K^{-}\sec_{-}$、$M\in K^{+}\sec_{+}$ 表示它们。考虑复形 $Hom(N,M)$（见 7.4.2）。它带显然的递减滤过 $F^{\cdot}$，且 $gr^{n}_{F}=Hom(N_{n},M_{n})[-n]$。注意 $Hom(N,M)$ 是下有界复形，且滤过 $F^{\cdot}$ 在每个项 $Hom(N,M)^{i}$ 上诱导有限滤过。我们把 $Hom(N,M)$ 视为此类复形的滤过导出范畴 DF 的对象。设 $RHom(N,\cdot)$ 是函子 $K^{+}\sec_{+}\to DF$、$M\mapsto Hom(N,M)$ 的右导出函子。此函子定义正确；且对任意 $n$，显然的态射 $gr^{n}_{F}RHom(N,M)\to RHom(N_{n},M_{n})[-n]$ 是拟同构。这由下述事实推出：对 $\mathcal{M}_{n}$ 中任意拟同构 $f:M_{n}\to I$，存在 $K^{+}\sec_{+}$ 中的拟同构 $g:M\to J$ 与态射 $h:I\to J_{n}$，使 $g_{n}=hf$。考虑滤过复形 $RHom(N,M)$ 的谱序列 $E^{p,q}_{r}$。它收敛到 $H^{\cdot}RHom(N,M)$，且
\[
E^{p,q}_{1}=H^{q}RHom_{\mathcal{M}_{p}}(N_{p},M_{p}).
\]

\bdp{7.4.9} \emph{注.} 设范畴 $\mathcal{M}_{n}$ 有许多内射对象。则范畴 $K^{+}_{\mathrm{tot}-}$ 有许多内射对象（即函子 $K^{+}_{\mathrm{tot}-}\to D^{+}_{\mathrm{tot}}$ 有右伴随函子）。事实上，若 $I\in K^{+}_{\mathrm{tot}+}$ 是使每个 $I^{a}_{n}$ 都是 $\mathcal{M}_{n}$ 的内射对象的复形，则 $c_{-}I$
% [Tutor Note: 原书 p.254 此处印作 "eacb"，系 "each" 之误，底本如此，保留原样。]
\SourcePage{255}%
是 $K^{+}_{\mathrm{tot}-}$ 的内射对象，而 $K^{+}_{\mathrm{tot}-}$ 中任意对象都拟同构于这样的 $I$。对偶地，若 $\mathcal{M}_{n}$ 有许多投射对象，则 $K^{-}_{\mathrm{tot}+}$ 有许多投射对象。

\bdp{7.4.10} 本小节后文不再使用；读者可以跳过。可以用稍许不同的方式定义 $D\sec$，从而定义 $D_{\mathrm{tot}}$，这在某些应用中较方便\footnote{此构造可追溯到 Toledo 与 Tong 的工作。}。如下定义范畴 $\mathrm{hot}_{+}=\mathrm{hot}_{+}(\mathcal{M}_{\cdot})$：其对象是族 $A=(A_{m})$，$A_{m}\in\mathcal{M}_{m}$。态射 $f:A\to B$ 是集合 $(f_{\alpha})$，其中对箭 $\alpha:\Delta_{n}\to\Delta_{m}$，相应的 $f_{\alpha}$ 是态射 $\alpha_{\cdot}A_{n}\to B_{m}$。态射的复合为 $(fg)_{\alpha}=\sum_{\alpha=\beta\gamma}f_{\beta}\beta_{\cdot}(g_{\gamma})$。这是加法范畴。置 $\mathrm{hot}_{-}(\mathcal{M}_{\cdot})=(\mathrm{hot}_{+}(\mathcal{M}^{\circ}_{\cdot}))^{\circ}$。有相应的复形 DG 范畴 $C\mathrm{hot}_{\pm}$。

有 DG 函子 $t_{+}:C\sec_{+}\to C\mathrm{hot}_{+}$，它把 $M\in C\sec_{+}$ 送到复形 $t_{+}M\in C\mathrm{hot}_{+}$，其分量为 $(t_{+}M)^{a}_{m}=M^{a-m}_{m}$，而微分 $d=d_{t_{+}M}$ 满足
\[
d_{\mathrm{id}_{\Delta_{m}}}=(-1)^{m}d^{a-m}_{M_{m}}:M^{a-m}_{m}\to M^{a-m+1}_{m},
\]
且对第 $i$ 个边缘映射 $\alpha_{i}:\Delta_{m}\hookrightarrow\Delta_{m+1}$ 有
\[
d_{\alpha_{i}}=(-1)^{i}\alpha^{*}_{i}:\alpha_{i\cdot}M^{a-m}_{m}\to M^{a-m}_{m+1},
\]
$d$ 的其余分量全为零。对 $l\in Hom(M_{1},M_{2})$ 有 $t_{+}(l)_{\mathrm{id}_{\Delta_{m}}}=l_{m}$，其余分量全为零。

\emph{注.} 函子 $t_{+}$ 是忠实的。可以把 $C\mathrm{hot}_{+}$ 的对象视为 $\sec_{+}$ 中带额外高阶同伦的「广义复形」。

另有 DG 函子 $s_{-}:C\mathrm{hot}_{+}\to C\sec_{-}$，定义如下。对 $A\in C\mathrm{hot}_{+}$，复形 $s_{-}A$ 的分量为
\[
(s_{-}A)^{a}_{m}=\sum_{\beta:\Delta_{n}\to\Delta_{m}}\beta_{\cdot}A^{a}_{n}.
\]
对 $\alpha:\Delta_{m}\to\Delta_{l}$，相容态射 $\alpha_{*}:(s_{-}A)^{a}_{l}\to\alpha_{\cdot}(s_{-}A)^{a}_{m}$ 的分量 $\gamma_{\cdot}A^{a}_{k}\to\alpha_{\cdot}\beta_{\cdot}A^{a}_{n}$ 当 $k=n$、$\gamma=\alpha\beta$ 时等于 $\mathrm{id}_{\gamma_{\cdot}A^{a}_{k}}$，否则为零。微分 $d_{s_{-}A}:(s_{-}A)^{a}_{m}\to(s_{-}A)^{a+1}_{m}$ 的分量 $\gamma_{\cdot}A^{a}_{k}\to\alpha_{\cdot}\beta_{\cdot}A^{a}_{n}$ 当 $\beta=\gamma\delta$ 时等于 $\gamma_{\cdot}(d_{A\delta})$，否则为零。

\emph{注.} DG 函子 $s_{-}$ 是全忠实的。
%% end units/ch07_u02.tex
%% begin units/ch07_u03.tex
% ch07_u03 —— 原书 p.256--p.265（7.4.11 -- 7.6.10）
% 本单元起于 p.256 首段（承接 p.255 的 7.4.10），末尾停在 7.6.10 的中途；尾部接 p.266。

\SourcePage{256}

我们用对偶来定义 DG 函子 $t_{-}:C\sec_{-}\to C\mathrm{hot}_{-}$ 与 $s_{+}:C\mathrm{hot}_{-}\to C\sec_{+}$。注意，复合 $s_{+}t_{-}:C\sec_{-}\to C\sec_{+}$ 与 7.4.2 中的函子 $c_{+}$ 一致；类似地，$s_{-}t_{+}=c_{-}$。函子 $t_{-}s_{-}:C\mathrm{hot}_{+}\to C\mathrm{hot}_{-}$ 与 $t_{+}s_{+}:C\mathrm{hot}_{-}\to C\mathrm{hot}_{+}$ 互为伴随（正如 7.4.2 中的函子 $c_{\pm}$）。

设 $f:A\to B$ 是 $C\mathrm{hot}_{\pm}$ 的同伦范畴 $K\mathrm{hot}_{\pm}$ 中的态射。若所有态射 $f_{m}:=f_{\mathrm{id}_{\Delta_{m}}}:A_{m}\to B_{m}$ 都是拟同构，则称 $f$ 为\emph{拟同构}。拟同构构成局部化族。相应的局部化三角范畴记为 $D\mathrm{hot}_{\pm}$。

函子 $s_{\pm}$、$t_{\pm}$ 保持拟同构，因此给出导出范畴之间的函子。这些函子复合的伴随态射都是拟同构。于是导出范畴 $D\sec_{\pm}$ 与 $D\mathrm{hot}_{\pm}$ 典范地等同。

% [Tutor Note: 原书 p.256 此处拼作「canonicaly」（应为 canonically），底本如此。]

\emph{注.} (i) 复形 $A\in D\mathrm{hot}_{+}$ 属于 $D_{\mathrm{tot}}$，当且仅当对任意 $\alpha:\Delta_{m}\to\Delta_{m+1}$，其 $\alpha$-分量 $d_{A\alpha}:\alpha_{\cdot}A_{m}\to A_{m+1}$ 是复形的拟同构（$A_{m}$ 上的微分为 $d_{A\,\mathrm{id}_{\Delta_{m}}}$，$A_{m+1}$ 亦然）。

(ii) 若范畴 $\mathcal{M}_{n}$ 有许多内射对象，则 $K^{+}\mathrm{hot}_{+}$ 有许多内射对象。对偶地，若 $\mathcal{M}_{n}$ 有许多投射对象，则 $K^{-}\mathrm{hot}_{-}$ 有许多投射对象（参看 7.4.9）。

\bdp{7.4.11} 上述某些构造在下面稍一般的设定中仍有意义。考虑任一在 $(\Delta)$ 上余纤维化的 DG 范畴族 $\mathcal{C}_{\cdot}$。有 DG 范畴 $\mathcal{C}\sec_{\pm}=\sec_{\pm}(\mathcal{C}_{\cdot})$（其定义与 7.4.2 中的范畴 $\sec_{\pm}(\mathcal{M}_{\cdot})$ 完全相同），以及相应的同伦范畴。按 7.4.2 的方式，在 $\pm$ 两类范畴之间定义伴随函子 $c_{\pm}$。

再设我们有 7.4.1 中的 $\mathcal{M}$，以及一族与纤维范畴结构相容的上同调函子 $H:\mathcal{C}_{\cdot}\to\mathcal{M}_{\cdot}$。于是得到相应的上同调函子 $H_{\pm}:\mathcal{C}\sec_{\pm}\to\sec_{\pm}$。把同伦范畴对 $H$-拟同构局部化，得到导出范畴 $D\sec_{\pm}$。如引理 7.4.4 那样，函子 $c_{\pm}$ 把范畴 $D\sec_{\pm}$ 等同起来，故可简记为 $D\sec$。接着定义\SourcePage{257}范畴 $\mathcal{C}_{\mathrm{tot}\pm}$ 等，如 7.4.5 所述。引理 7.4.6 仍然成立，故有全三角子范畴 $D_{\mathrm{tot}}\subset D\sec$，以及上同调函子 $H:D_{\mathrm{tot}}\to\mathcal{M}_{\mathrm{tot}}$。

\bdsec{7.5. 光滑叠上的 $\mathcal{D}$-模理论 II}

\bdp{7.5.1} 设 $\mathcal{Y}$ 是任意的光滑代数叠。设 $U_{\cdot}$ 是 $\mathcal{Y}$ 的超覆叠，使得每个 $U_{n}$ 都是（光滑）拟紧分离代数空间的无交并（例如仿射概形）。称这样的 $U_{\cdot}$ 为\emph{容许}超覆叠。把 $U_{\cdot}$ 看作 $(\Delta)^{\circ}$-代数空间。范畴 $\mathcal{M}(U_{\cdot})$ 构成 7.4.1 中的 $(\Delta)$-族阿贝尔范畴；相应的范畴 $\mathcal{M}_{\mathrm{tot}}$ 就是 $\mathcal{M}(\mathcal{Y})$。依 7.4.7，得到相应的 t-范畴 $D_{\mathrm{tot}}=D_{\mathrm{tot}}(U_{\cdot},\mathcal{D})$，其核心为 $\mathcal{M}(\mathcal{Y})$。

我们也可以考虑 DG 范畴 $C(U_{\cdot},\Omega)$，连同上同调函子 $H_{\mathcal{D}\cdot}:C(U_{\cdot},\Omega)\to\mathcal{M}(U_{\cdot})$，其中 $H_{\mathcal{D}n}F_{n}=H_{\mathcal{D}}F_{n}[\dim U_{n}/\mathcal{Y}]$，$F_{n}\in C(U_{n},\Omega)$；然后应用 7.4.11。得到三角范畴 $D_{\mathrm{tot}}(U_{\cdot},\Omega)$，以及上同调函子 $H_{\mathcal{D}}:D_{\mathrm{tot}}(U_{\cdot},\Omega)\to\mathcal{M}(\mathcal{Y})$。

范畴 $D_{\mathrm{tot}}(U_{\cdot},\mathcal{D})$ 与 $D_{\mathrm{tot}}(U_{\cdot},\Omega)$ 典范地等同。确切地说，有函子 $\Omega_{\cdot}:C(U_{\cdot},\mathcal{D})\to C(U_{\cdot},\Omega)$，$\Omega_{n}(M_{n}):=\Omega M_{n}[-\dim U_{n}/\mathcal{Y}]$。该函子与 DG 结构、纤维范畴结构以及与上同调函子相容（即 $H=H_{\mathcal{D}\cdot}\Omega_{\cdot}$）。因此它给出正合函子
\begin{equation}\tag{315}
\Omega:D_{\mathrm{tot}}(U_{\cdot},\mathcal{D})\to D_{\mathrm{tot}}(U_{\cdot},\Omega).
\end{equation}

该函子是范畴等价。事实上，尽管 $C(U_{\cdot},\Omega)$ 与 $C(U_{\cdot},\mathcal{D})$ 之间的函子 $\mathcal{D}$ \emph{不}与纤维范畴结构相容，它仍给出函子 $\mathcal{D}:C\sec_{-}(U_{\cdot},\Omega)\to C\sec_{-}(U_{\cdot},\mathcal{D})$，$(\mathcal{D}F)_{n}=\mathcal{D}F_{n}[\dim U_{n}/\mathcal{Y}]$（用 7.2.8 来定义各个 $\alpha^{*}$）。这个 $\mathcal{D}$ 是相应 $\Omega$ 函子的左伴随，且与上同调函子相容。$\mathcal{D}$-$\Omega$ 伴随态射都是拟同构（见 7.2.4、7.2.5），故 $\mathcal{D}$ 给出 (315) 的逆函子。

\SourcePage{258}

我们把这样等同起来的范畴 $D_{\mathrm{tot}}(U_{\cdot},\mathcal{D})$ 与 $D_{\mathrm{tot}}(U_{\cdot},\Omega)$ 简记为 $D_{\mathrm{tot}}(U_{\cdot})$。

\bdp{7.5.2} \emph{命题.} 对 $\mathcal{Y}$ 的不同容许覆盖，t-范畴 $D_{\mathrm{tot}}(U_{\cdot})$ 之间存在典范等同。

证明见下文 7.5.5。把这样等同起来的范畴记为 $D(\mathcal{Y})$；它是核心为 $\mathcal{M}(\mathcal{Y})$ 的 t-范畴。

在证明 7.5.2 之前先说明：若 $\mathcal{Y}$ 满足条件 (310)，则确实得到与 7.3.2 相同的范畴 $D(\mathcal{Y})$。顺带可见 7.3.4。

选取 $\mathcal{Y}$ 的超覆叠 $U_{\cdot}$，使各 $U_{n}$ 为仿射概形。有显然的正合函子（限制到 $U_{\cdot}$）
\begin{equation}\tag{316}
r:D(\mathcal{Y},\Omega)\to D_{\mathrm{tot}}(U_{\cdot},\Omega).
\end{equation}

\bdp{7.5.3} \emph{引理.} 函子 $r$ 是范畴等价。

\emph{证明.} 我们来构造逆函子。对 $F\in K_{\mathrm{tot}+}(\Omega)$，用各项为 $\pi_{\cdot}F^{ab}:=\pi_{b}(F^{a})$ 的 Čech 双复形的全复形，定义 $\mathcal{Y}$ 上的 $\Omega$-复形 $\pi_{\cdot}F$；于是 $(\pi_{\cdot}F)^{n}=\bigoplus_{a+b=n}F^{ab}$，这里 $\pi_{b}$ 是投影 $U_{b}\to\mathcal{Y}$。这样得到正合函子 $\pi_{\cdot}:K_{\mathrm{tot}+}(\Omega)\to K(\mathcal{Y},\Omega)$。该函子保持 $\mathcal{D}$-拟同构（因为由 (310)，投影 $\pi_{b}$ 是仿射的），故给出函子 $D_{\mathrm{tot}}(U_{\cdot},\Omega)\to D(\mathcal{Y},\Omega)$。

验证这个函子是 $r$ 的逆，留给读者（提示：对如上 $F$，伴随拟同构 $\pi'_{\Omega}\pi_{\cdot}F\to F$ 来自 $C\sec_{-}(U_{\cdot},\Omega)$ 中的典范态射 $\pi'_{\Omega}\pi_{\cdot}F\to c_{-}F$）。\hfill$\Box$

\bdp{7.5.4} \emph{注.} 上述引理可按如下方式移植到 7.3.12 的设定中。设 $\mathcal{Y}$ 是任意的光滑叠，其对角态射 $\mathcal{Y}\to\mathcal{Y}\times\mathcal{Y}$ 为仿射。则 7.3.12 与 7.5.1 中定义的范畴 $D(\mathcal{Y})$ 典范等价。事实上，设 $\mathcal{Y}_{i}$ 是 7.3.12 中 $\mathcal{Y}$ 的一列开子叠，$V_{i}\to\mathcal{Y}$ 是覆盖，各 $V_{i}$ 为仿射概形。\SourcePage{259}于是各 $V_{i}$ 构成 $\mathcal{Y}$ 的覆盖。设 $U_{\cdot}$ 是相应的 Čech 超覆叠。因此 $U_{a}$ 是分支 $U_{\alpha}$ 的无交并，这些分支由序列 $\alpha=(\alpha_{1},\alpha_{2},\ldots)$ 标号，$\alpha_{i}\ge0$，$\Sigma\alpha_{i}=a+1$；其中 $U_{\alpha}$ 是 $\alpha_{1}$ 个 $V_{1}$、$\alpha_{2}$ 个 $V_{2}$……在 $\mathcal{Y}$ 上的纤维积。对 $F\in C(\mathcal{Y}_{\cdot},\Omega)$，置 $F_{U\alpha}:=F_{i_{\alpha}U_{\alpha}}$，其中 $i_{\alpha}$ 是使 $\alpha_{i}$ 非零的最小 $i$（注意 $U_{\alpha}\in\mathcal{Y}_{sm}$。这些 $F_{U_{\alpha}}$ 按显然的方式构成 $U_{\cdot}$ 上的 $\Omega$-复形 $F$，它属于 $C_{\mathrm{tot}}(U_{\cdot},\Omega)$。

% [Tutor Note: 原书 p.259「(note that U_α ∈ Y_sm.」左括号未闭合，底本如此，译文依原样。]

函子 $C(\mathcal{Y}_{\cdot},\Omega)\to C_{\mathrm{tot}}(U_{\cdot},\Omega)$ 与函子 $H_{\mathcal{D}}$ 交换，故给出三角函子
\begin{equation}\tag{317}
r:D(\mathcal{Y}_{\cdot},\Omega)\to D_{\mathrm{tot}}(U_{\cdot},\Omega).
\end{equation}

验证该函子是范畴等价，以及 7.3.12 与 7.5.2 意义下 $D(\mathcal{Y})$ 的相应等同不依赖 $\mathcal{Y}_{\cdot}$ 与 $V_{\cdot}$ 这些辅助数据，均留给读者。

\bdp{7.5.5} \emph{7.5.2 的证明.} 我们需要对不同的 $U_{\cdot}$ 典范地等同 t-范畴 $D_{\mathrm{tot}}(U_{\cdot})$。设 $U'_{\cdot}$ 是另一个容许超覆叠。首先用某些辅助数据 $V$ 来定义 t-正合函子 $\Phi=\Phi_{V}:D_{\mathrm{tot}}(U_{\cdot})\to D_{\mathrm{tot}}(U'_{\cdot})$。然后证明 $\Phi$ 实际上不依赖 $V$，且是范畴等价。

我们的 $V$ 是 $\mathcal{Y}$ 上的 $(\Delta)^{\circ}\times(\Delta)^{\circ}$-代数空间 $V_{\cdot\cdot}$，连同光滑态射 $\pi:V_{mn}\to U_{m}$、$\pi':V_{mn}\to U'_{n}$。假定 $\pi$、$\pi'$ 按显然的方式与 $(\Delta)$ 投影相容，$\pi'_{\cdot n}:V_{\cdot n}\to U'_{n}$ 是超覆叠，$\pi'_{mn}:V_{mn}\to U'_{n}$ 是仿射态射。

对 $F\in K_{\mathrm{tot}+}(U_{\cdot},\Omega)$，有 $\Omega$-复形 $F_{V_{n}}\in K_{\mathrm{tot}+}(V_{\cdot n},\Omega)$，即 $F$ 到 $V_{\cdot n}$ 的拉回。置 $\Phi_{V_{n}}F:=\pi'_{\cdot}F_{V_{n}}$（记号见 7.5.3 的证明）。这是 $U'_{n}$ 上的 $\Omega$-复形。这些 $\Omega$-复形 $\Phi_{V_{n}}F$ 按显然的方式构成 $\Omega$-复形 $\Phi_{V}F\in K_{\mathrm{tot}+}(U'_{\cdot},\Omega)$，并满足 $H_{\mathcal{D}}F=H_{\mathcal{D}}\Phi_{V}F$。因此有 t-正合函子 $\Phi_{V}:D_{\mathrm{tot}}(U_{\cdot},\Omega)\to D_{\mathrm{tot}}(U'_{\cdot},\Omega)$，它在核心 $\mathcal{M}(\mathcal{Y})$ 之间诱导恒等函子。

设有如上 $V_{1}$ 与 $V_{2}$。为等同函子 $\Phi_{V_{i}}$，选取另一个如上 $V$，连同嵌入 $V_{1},V_{2}\subset V$\SourcePage{260}，且这些嵌入与所有投影相容，这些投影把 $(V_{1})_{mn}$、$(V_{2})_{mn}$ 与 $V_{mn}$ 的若干连通分支之并等同。这些嵌入诱导投影 $\Phi_{V}F\to\Phi_{V_{1}}$、$\Phi_{V}F\to\Phi_{V_{2}}F$，它们显然是拟同构。

% [Tutor Note: 原书 p.260 第一个投影右端作「Φ_V F → Φ_{V_1}」，末端缺 F，底本如此，译文依原样。]

于是我们等同了导出范畴之间的函子 $\Phi_{V_{i}}$。验证该等同不依赖辅助数据 $V$，留给读者。

于是有典范函子 $\Phi=\Phi_{UU'}:D_{\mathrm{tot}}(U_{\cdot},\Omega)\to D_{\mathrm{tot}}(U'_{\cdot},\Omega)$。若 $U''_{\cdot}$ 是第三个超覆叠，则有函子的典范同构 $\Phi_{UU''}=\Phi_{U'U''}\Phi_{UU'}$；其定义以及通常相容性的验证，都留给读者。由于 $\Phi_{UU}$ 是恒等函子，可见这些 $\Phi$ 同时等同了所有范畴 $D_{\mathrm{tot}}(U_{\cdot})$。\hfill$\Box$

\bdp{7.5.6} 设 $f:\mathcal{Y}\to\mathcal{Z}$ 是光滑叠的拟紧态射。我们来定义前推函子 $f_{*}:D(\mathcal{Y})^{+}\to D(\mathcal{Z})^{+}$。为此，考虑 $\mathcal{Y}$ 的任一容许超覆叠 $U_{\cdot}$ 与 $\mathcal{Z}$ 的任一容许超覆叠 $W_{\cdot}$。得到 $(\Delta)^{\circ}\times(\Delta)^{\circ}$-代数空间 $U_{\cdot}\times_{\mathcal{Z}}W_{\cdot}$。可以找到 $(\Delta)^{\circ}\times(\Delta)^{\circ}$-代数空间 $V_{\cdot\cdot}$，连同态射 $\phi=(\phi_{1},\phi_{2}):V_{\cdot\cdot}\to U_{\cdot}\times_{\mathcal{Z}}W_{\cdot}$，使得投影 $V_{mn}\to U_{m}$ 光滑，$V_{mn}\to W_{n}$ 仿射，且 $V_{\cdot n}\to\mathcal{Y}\times_{\mathcal{Z}}W_{n}$ 是超覆叠。现在对 $F\in K^{+}_{\mathrm{tot}+}(U_{\cdot},\Omega)$，令 $F_{V_{n}}\in K^{+}_{\mathrm{tot}+}(V_{\cdot n},\Omega)$ 是它到 $V_{\cdot n}$ 的拉回。用各项为 $\phi_{2\cdot}F_{V_{n}}$ 的 Čech 双复形的全复形，定义 $W_{n}$ 上的 $\Omega$-复形 $f_{\cdot}F_{n}$。

% [Tutor Note: 原书 p.260 上一行作 F^{V_n}，此处作 F_{V_n}（下标），指同一对象；底本如此，译文依原样。]

这些 $\Omega$-复形构成 $K^{+}_{\mathrm{tot}+}(W_{\cdot},\Omega)$ 的对象 $f_{\cdot}F$。函子 $f_{\cdot}:K^{+}_{\mathrm{tot}+}(U_{\cdot},\Omega)\to K^{+}_{\mathrm{tot}}(W_{\cdot},\Omega)$ 保持 $\mathcal{D}$-拟同构，因此给出函子 $f_{*}:D(\mathcal{Y})^{+}\to D(\mathcal{Z})^{+}$。验证 $f_{*}$ 的构造不依赖 $U$、$W$、$V$ 的辅助选取，且与 $f$ 的复合相容，留给读者。

% [Tutor Note: 原书 p.260 第二个范畴作 K^+_{tot}(W_·,Ω)，漏了下标 +，底本如此，译文依原样。]

光滑叠的光滑态射 $f:\mathcal{Y}\to\mathcal{Z}$ 给出 t-正合函子 $f^{\dagger}=f^{\dagger}_{\Omega}:D(\mathcal{Z})\to D(\mathcal{Y})$。具体地说，选取 $\mathcal{Y}$ 的容许超覆叠 $U_{\cdot}$、$\mathcal{Z}$ 的容许超覆叠 $W_{\cdot}$，以及与 $f$ 相容的态射 $f_{\cdot}:U_{\cdot}\to W_{\cdot}$。函子 $f^{\cdot}_{\Omega}:K_{\mathrm{tot}\pm}(W_{\cdot},\Omega)\to K_{\mathrm{tot}\pm}(U_{\cdot},\Omega)$ 保持 $\mathcal{D}$-拟同构，故给出导出范畴之间的函子 $f^{\cdot}_{\Omega}$。我们把这个定义的检验留给读者：\SourcePage{261}它不依赖辅助选取，我们的拉回函子与 $f$ 的复合相容，且当 $f$ 拟紧时函子 $f^{\cdot}_{\Omega}$ 是 $f_{*}$ 的左伴随。

\bdp{7.5.7}

% [Tutor Note: 原书 p.261 此编号段后无正文（仅印出编号），底本如此。]

\bdp{7.5.8} \emph{注.} (i) 也可以尝试用适当的非拟凝聚 $\Omega$-复形来定义 $D(\mathcal{Y})$，方式类似于 [LMB93]6.3 中 $\mathcal{O}$-模导出范畴的定义。这样的定义很可能给出同一个范畴 $D^{+}(\mathcal{Y})$。

(ii) 导出范畴的「局部」构造在 $\mathcal{O}$-模的设定中也很方便。例如，它有助于把代数叠的余切复形定义为导出范畴中的真正对象（而不只是像 [LMB93]9.2 那样取其截断的投射极限），也有助于处理 Grothendieck-Serre 对偶。

(iii) 把 $\mathcal{D}$-模换成反常层，就得到局部有限型代数叠上可构造层导出范畴的一个方便定义。

\bdsec{7.6. 等变设定}

\bdp{7.6.1} 我们来解释（有限维）Hecke 模式的 7.1.1 (a)、(b) 两部分。设 $G$ 是代数群，$K\subset G$ 是代数子群。为简单起见，假定 $K$ 仿射；此时下面的叠满足 7.3.1 的条件 (310)。置$^{*)}$\ $\mathcal{H}^{c}:=C(K\backslash G/K,\Omega)$，$\mathcal{H}:=D(K\backslash G/K)$。\footnote{$^{*)}$此处上标「c」表示我们处理的是真正的 DG 复形范畴，而非导出范畴。}分别称这两个范畴为\emph{预 Hecke} 范畴与 \emph{Hecke} 范畴。它们带有如下定义的典范幺半结构。

考虑叠的态射
\begin{equation}\tag{318}
(K\backslash G/K)\times(K\backslash G/K)\xleftarrow{p}K\backslash G\times_{K}G/K\xrightarrow{\bar m}K\backslash G/K.
\end{equation}

\SourcePage{262}

这里 $G\times_{K}G$ 是 $G\times G$ 模去 $K$-作用 $k(g_{1},g_{2})=(g_{1}k^{-1},kg_{2})$ 的商，$p$ 是显然的投影，$\bar m$ 是乘法映射。对 $F_{1},F_{2}\in\mathcal{H}^{c}$，置 $F_{1}\overset{c}{\circledast}F_{2}:=\bar m_{\cdot}p^{\cdot}_{\Omega}(F_{1}\boxtimes F_{2})$。\emph{卷积张量积} $\overset{c}{\circledast}$ 满足显然的结合约束，故 $\mathcal{H}^{c}$ 上有幺半结构。把卷积张量积 $\circledast:\mathcal{H}\times\mathcal{H}\longrightarrow\mathcal{H}$ 定义为 $\overset{c}{\circledast}$ 的右导出函子。有 $F_{1}\circledast F_{2}=\bar m_{*\cdot}p^{\cdot}_{\Omega}(F_{1}\boxtimes F_{2})$；若 $\Omega$-复形 $F_{1}$、$F_{2}$ 是松的（见 7.3.7），则 $F_{1}\circledast F_{2}=F_{1}\overset{c}{\circledast}F_{2}$。于是 $\circledast$ 的结合约束可由 $\overset{c}{\circledast}$ 的结合约束推出，故 $\mathcal{H}$ 是幺半三角范畴。$\mathcal{H}^{c}$ 与 $\mathcal{H}$ 有单位对象 $E$：有 $E_{G}=i_{K\cdot}\Omega_{K}$（这里 $i_{K}:K\hookrightarrow G$ 是嵌入）。

设 $Y$ 是带 $G$-作用的光滑簇。考虑叠 $\mathcal{B}:=K\backslash Y$。$\mathcal{D}(\mathcal{B})$ 上的 Hecke 作用来自图表
\begin{equation}\tag{319}
(K\backslash G/K)\times\mathcal{B}\xleftarrow{p_{Y}}K\backslash(G\times_{K}Y)\xrightarrow{\bar m_{Y}}\mathcal{B}.
\end{equation}

即：对 $F\in\mathcal{H}^{c}$、$T\in C(\mathcal{B},\Omega)$，置 $F\overset{c}{\circledast}T:=\bar m_{Y\cdot}p^{\cdot}_{Y\Omega}(F\boxtimes T)$。与上面一样，$\overset{c}{\circledast}$ 满足显然的结合约束，故 $C(\mathcal{B},\Omega)$ 是幺半 $\mathcal{H}^{c}$-模。把 $\circledast:\mathcal{H}\times D(\mathcal{B})\longrightarrow D(\mathcal{B})$ 定义为 $\overset{c}{\circledast}$ 的右导出函子。有 $F\circledast T=\bar m_{Y*\cdot}p^{\cdot}_{Y\Omega}(F\boxtimes T)$；若 $F$、$T$ 松（见 7.3.7），则 $F\circledast T=F\overset{c}{\circledast}T$。因此 $D(\mathcal{B})$ 是 $\mathcal{H}$-模。

\bdp{7.6.2} \emph{注.} (i) 上面的定义中之所以能考虑无界导出范畴，是因为投影 $\bar m$、$\bar m_{Y}$ 可表（见 7.3.11(ii)）。

(ii) 若 $f:Z\to Y$ 是带 $G$-作用的光滑簇的态射，则 $f_{*}:D(K\backslash Z)\to D(K\backslash Y)$ 是 $\mathcal{H}$-模的态射。

\bdp{7.6.3} 设 $Y$ 是带仿射代数群 $K$ 作用的光滑簇。考虑叠 $\mathcal{B}:=K\backslash Y$。在 7.6 的其余部分，我们将用 $Y$ 上适当的等变复形来描述 $D(\mathcal{B})$。我们还将引入某个导出范畴 $D(K^{\backslash}_{\backslash}Y)$，它介于 $D(K\backslash Y)$ 与 $D(Y)$ 之间，在 7.7 中会用到。

\SourcePage{263}

置 $K_{\Omega}=(K,\Omega_{K})$，$K^{\cdot}_{\Omega}=(K,\Omega^{\cdot}_{K})$（于是 $K^{\cdot}_{\Omega}$ 就是略去 de Rham 微分后的 $K_{\Omega}$）。它们分别是 DG 环化空间范畴与分次环化空间范畴中的群对象。记 $\mathfrak{k}$、$\mathfrak{k}_{\Omega}$、$\mathfrak{k}^{\cdot}_{\Omega}$ 分别为 $K$、$K_{\Omega}$、$K^{\cdot}_{\Omega}$ 的李代数。作为普通复形，$\mathfrak{k}_{\Omega}$ 等于 $\mathrm{id}_{\mathfrak{k}}$ 的锥，故 $\mathfrak{k}^{0}_{\Omega}=\mathfrak{k}=\mathfrak{k}^{-1}_{\Omega}$。由于 $K$ 是 $K_{\Omega}$ 与 $K^{\cdot}_{\Omega}$ 的子群，我们得到相应的 Harish-Chandra 对 $(\mathfrak{k}_{\Omega},K)$、$(\mathfrak{k}^{\cdot}_{\Omega},K)$。注意，$K_{\Omega}$-模与 DG $(\mathfrak{k}_{\Omega},K)$-模相同，$K^{\cdot}_{\Omega}$-模与分次 $(\mathfrak{k}^{\cdot}_{\Omega},K)$-模相同。

% [Tutor Note: 原书 p.263「the category of DG ringed spaced」多一字母 d（应为 spaces），底本如此。]

$Y$ 上的 $K$-作用给出 $K_{\Omega}$ 在 $Y_{\Omega}=(Y,\Omega)$ 上的作用，从而给出 $K^{\cdot}_{\Omega}$ 在 $Y^{\cdot}_{\Omega}=(Y,\Omega^{\cdot})$ 上的作用。对分次 $\Omega^{\cdot}_{Y}$-模 $F^{\cdot}_{Y}$，$F^{\cdot}_{Y}$ 上的 $K^{\cdot}_{\Omega}$-作用等同于 $(\mathfrak{k}^{\cdot}_{\Omega},K)$-作用。具体地说，它是 $F^{\cdot}_{Y}$ 上的 $K$-作用，连同 $K$-等变态射 $\mathfrak{k}\otimes F^{\cdot}_{Y}\to F^{\cdot-1}_{Y}$，$\xi\otimes f\mapsto i_{\xi}(f)$（假定 $K$ 以伴随方式作用在 $\mathfrak{k}$ 上），使得对任意 $\xi\in\mathfrak{k}$ 与 $\nu\in\Omega^{1}_{Y}$ 有 $i_{\xi}(\nu f)=\langle\xi,\nu\rangle f+\nu i_{\xi}(f)$，$i^{2}_{\xi}=0$。

\bdp{7.6.4} 设 $F_{Y}$ 是 $Y$ 上的 $\Omega$-复形。$F_{Y}$ 上的 \emph{$K$-作用} 是分次 $\mathcal{O}_{Y}$-模 $F^{\cdot}_{Y}$ 上的 $K$-作用，使得对任意 $k\in K$，平移 $k^{*}F^{\cdot}_{Y}\xrightarrow{\sim}F^{\cdot}_{Y}$ 是 $\Omega$-复形的态射（即它与微分交换）。$F_{Y}$ 上的 \emph{$K_{\Omega}$-作用} 是把 $F_{Y}$ 看作 $Y_{\Omega}$ 上的 DG 模时 $K_{\Omega}$ 的作用。换言之，它是分次 $\Omega^{\cdot}_{Y}$-模 $F^{\cdot}_{Y}$ 上的 $K^{\cdot}_{\Omega}$-作用，使得 $K$ 按 $\Omega$-复形作用于 $F_{Y}$，而 $\mathfrak{k}_{\Omega}$ 按 DG 李代数作用于 $F_{Y}$。后一条件意味着对任意 $\xi\in\mathfrak{k}$ 有 $di_{\xi}+i_{\xi}d=\mathrm{Lie}_{\xi}$（这里 $\mathrm{Lie}$ 是由 $K$-作用得到的 $\mathfrak{k}$ 在 $F^{\cdot}_{Y}$ 上的作用）。带 $K$-作用的 $\Omega$-复形称为\emph{弱 $K$-等变} $\Omega$-复形，带 $K_{\Omega}$-作用的称为 \emph{$K_{\Omega}$-等变} $\Omega$-复形。

显然，对叠 $\mathcal{B}:=K\backslash Y$ 上的任意 $\Omega$-复形 $F$，$\Omega$-复形 $F_{Y}$ 自动带有 $K_{\Omega}$-作用。

\bdp{7.6.5} \emph{引理.} 函子 $C(K^{\backslash}_{\backslash}Y,\Omega)\longrightarrow$（$Y$ 上的 $K_{\Omega}$-等变 $\Omega$-复形）是 DG 范畴的等价。\hfill$\Box$

\SourcePage{264}

\bdp{7.6.6} \emph{注.} 假定处于 7.6.1 的情形。设 $m:K\backslash G\times G/K\to K\backslash G/K$ 是乘法映射。置 $F_{1}\widetilde{\circledast}F_{2}=m_{\cdot}(F_{1K\backslash G}\boxtimes F_{2G/K})$；这是 $K\backslash G/K$ 上的 $\Omega$-复形。沿投影 $G\times G\to G\times_{K}G$ 的纤维的 $K$-作用给出 $F_{1}\widetilde{\circledast}F_{2}$ 上的 $K_{\Omega}$-作用（取 $K\backslash G/K$ 上的平凡 $K$-作用）。其不变量与 $F_{1}\overset{c}{\circledast}F_{2}$ 一致。类似地，考虑映射 $m_{Y}:(K\backslash G)\times Y\to\mathcal{B}$；置 $F\widetilde{\circledast}T:=m_{Y\cdot}(F_{K\backslash G}\boxtimes T)$。$(K\backslash G)\times Y$ 上显然的 $K$-作用给出该 $\Omega$-复形上的 $K_{\Omega}$-作用，其不变量与 $F\overset{c}{\circledast}T$ 一致。

% [Tutor Note: 原书 p.264 末句「coinside」拼写如此（应为 coincide），底本如此。]

\bdp{7.6.7} 把 $Y$ 上弱 $K$-等变 $\Omega$-复形的范畴记为 $C(K^{\backslash}_{\backslash}Y,\Omega)$，相应的同伦范畴与 $\mathcal{D}$-导出范畴记为 $K(K^{\backslash}_{\backslash}Y,\Omega)$、$D(K^{\backslash}_{\backslash}Y,\Omega)$（弱等变 $\Omega$-复形的态射若作为普通 $\Omega$-复形是 $\mathcal{D}$-拟同构，则称为 $\mathcal{D}$-拟同构）。

\bdp{7.6.8} \emph{注.} (i) 遗忘函子 $C(\mathcal{B},\Omega)\to C(K^{\backslash}_{\backslash}Y,\Omega)$ 有左右伴随函子 $c^{l},c^{r}:C(K^{\backslash}_{\backslash}Y,\Omega)\to C(\mathcal{B},\Omega)$，$c^{l}(F_{Y})=U(\mathfrak{k}_{\Omega})\otimes_{U(\mathfrak{k})}F_{Y}$，$c^{r}(F_{Y})=\mathrm{Hom}_{U(\mathfrak{k})}(U(\mathfrak{k}_{\Omega}),F_{Y})$。这些函子保持拟同构，故给出导出范畴之间的伴随函子。

(ii) 遗忘函子 $C(K^{\backslash}_{\backslash}Y,\Omega)\to C(Y,\Omega)$ 有右伴随函子 $\mathrm{Ind}:C(Y,\Omega)\to C(K^{\backslash}_{\backslash}Y,\Omega)$，$\mathrm{Ind}(T^{\cdot}_{Y})=p_{*}m^{*}(T^{\cdot}_{Y})$，其中 $m,p:K\times Y\rightrightarrows Y$ 是作用映射与投影映射。这些函子保持拟同构，故给出导出范畴之间的伴随函子。复合 $c^{r}\mathrm{Ind}$ 是投影 $Y\to\mathcal{B}$ 的前推函子。

(iii) 注 7.6.6 (ii) 对弱等变 $\Omega$-复形仍然有效。

% [Tutor Note: 原书 p.264 引「Remark 7.6.6 (ii)」，但 7.6.6 并无 (ii)，底本如此，译文依原样。]

(iv) 设 $f:Z\to Y$ 是带 $K$-作用的光滑簇的态射。7.3.6 中正像函子的构造无需改动即可搬到弱等变设定，故有函子 $f_{*}=Rf_{\cdot}:D(K^{\backslash}_{\backslash}Z,\Omega)\to D(K^{\backslash}_{\backslash}Y,\Omega)$。函子 $f_{*}$ 与\SourcePage{265}(i)、(ii) 中所述的函子交换。对 7.2.8、7.3.6 中的拉回函子 $f^{\cdot}_{\Omega}$ 也一样。

(v) 下面是 7.6.1 的弱等变版本。设 7.6.1 中的 $Y$ 还带有仿射代数群 $G'$ 的作用，且与 $G$-作用交换（我们把它写成右作用）。考虑 $Y$ 上带有交换的 $K_{\Omega}$-作用与 $G$-作用的 $\Omega$-复形范畴 $C(K^{\backslash}_{\backslash}Y_{/}^{/}G',\Omega)=C(\mathcal{B}_{/}^{/}G',\Omega)$。则相应的导出范畴 $D(\mathcal{B}_{/}^{/}G',\Omega)$ 是 $\mathcal{H}$-模。$\mathcal{H}$-作用的定义方式与 7.6.1 相同。注 7.6.6 仍然有效。

% [Tutor Note: 原书 p.265 此处作「commuting K_Ω- and G-actions」，按上下文应为 G′-作用；底本如此，译文依原样作 G-作用。]

\bdp{7.6.9} 下面描述上述等变范畴的 $\mathcal{D}$-模对应物（细节见 [BL]）。对 $Y$ 上的 $\mathcal{D}$-模 $M$，$M$ 上的\emph{弱 $K$-作用} 是把 $M$ 看作 $\mathcal{O}_{Y}$-模时的 $K$-作用，使得对任意 $k\in K$，平移 $k^{*}M\xrightarrow{\sim}M$ 是 $\mathcal{D}$-模的态射。带弱 $K$-作用的 $\mathcal{D}$-模称为\emph{弱 $K$-等变} $\mathcal{D}$-模；其范畴记为 $\mathcal{M}(K^{\backslash}_{\backslash}Y)$（按惯例用 $\mathcal{M}^{\ell}$ 或 $\mathcal{M}^{r}$ 指明左、右 $\mathcal{D}$-模）。记号 $C(K^{\backslash}_{\backslash}Y,\mathcal{D})$、$K(K^{\backslash}_{\backslash}Y,\mathcal{D})$、$D(K^{\backslash}_{\backslash}Y,\mathcal{D})=D(K^{\backslash}_{\backslash}Y)$ 的含义是清楚的（参看 7.2）。

7.2.2 中的函子 $\mathcal{D}$ 与 $\Omega$ 把弱等变复形送到弱等变复形，于是有伴随 DG 函子
\begin{equation}\tag{320}
\mathcal{D}:C(K^{\backslash}_{\backslash}Y,\Omega)\to C(K^{\backslash}_{\backslash}Y,\mathcal{D}),\qquad
\Omega:C(K^{\backslash}_{\backslash}Y,\mathcal{D})\to C(K^{\backslash}_{\backslash}Y,\Omega),
\end{equation}
以及三角范畴之间互逆的等价
\begin{equation}\tag{321}
D(K^{\backslash}_{\backslash}Y,\mathcal{D})\leftrightarrows D(K^{\backslash}_{\backslash}Y,\Omega).
\end{equation}

按惯例，把这些等同起来的范畴记为 $D(K^{\backslash}_{\backslash}Y)$。

\bdp{7.6.10} \emph{注.} 对弱 $K$-等变 $\mathcal{D}$-模 $M$，$Y$ 上的 $\mathfrak{k}$-作用以两种方式提升到 $\mathcal{O}$-模 $M$：或者是由 $M$ 上的 $K$-作用定义的无穷小作用，或者经由 $Y$ 上的 $\mathfrak{k}$-作用 $\sigma:\mathfrak{k}\to\Theta_{Y}$ 与 $M$ 上的 $\mathcal{D}$-模结构。分别记这两个作用为 $\xi,m\mapsto\mathrm{Lie}_{\xi}m$ 与 $\sigma_{\xi}m$。置 $\xi^{\natural}m:=\mathrm{Lie}_{\xi}m-\sigma_{\xi}m$。则 $\xi^{\natural}\in\mathrm{End}_{\mathcal{D}}M$，并且

% 尾部接 p.266（原书 p.265 末句未完）。
%% end units/ch07_u03.tex
%% begin units/ch07_u04.tex
% ch07_u04 —— 原书 p.266--p.275（7.6.11 -- 7.8.6 前半）
% 本单元自 p.266 首句起（承 p.265 末句的谓词），至 p.275 末行止
% （该句跨页未完，尾部接 p.276）。

\SourcePage{266}

$\natural:\mathfrak{k}\to\operatorname{End}_{\mathcal{D}}M$ 是 $M$ 上的
$\mathfrak{k}$-作用。注意 $\natural$ 平凡，当且仅当 $M$ 是 $K$-等变
$\mathcal{D}$-模，即 $M\in\mathcal{M}(\mathcal{B})$。

\bdp{7.6.11} $Y$ 上的 $K$-等变 $\mathcal{D}$-复形，指由弱 $K$-等变
$\mathcal{D}$-模构成的复形 $N$，并带态射
$\mathfrak{k}\otimes N^{\cdot}\to N^{\cdot-1}$，$\xi\otimes n\mapsto i_\xi n$，
使得对任意 $\xi\in\mathfrak{k}$ 有 $i_\xi^2=0$，
$di_\xi+i_\xi d=\xi^{\natural}$。按记号滥用，把这种复形的 DG 范畴记作
$C(\mathcal{B},\mathcal{D})$。注意任意 $K$-等变 $\mathcal{D}$-模按显然方式
就是 $K$-等变 $\mathcal{D}$-复形；任意 $K$-等变 $\mathcal{D}$-复形的上同调层
都是 $K$-等变 $\mathcal{D}$-模。于是有上同调函子
$H:C(\mathcal{B},\mathcal{D})\to\mathcal{M}(\mathcal{B})$。把
$C(\mathcal{B},\mathcal{D})$ 的同伦范畴对 $H$-拟同构作局部化，得到三角范畴
$D(\mathcal{B},\mathcal{D})$。容易看出它是 $t$-范畴，其核心为
$\mathcal{M}(\mathcal{B})$。
% [Tutor Note: 原书 p.266 此处印作「such that for any xi in k our has」，
% 其中「our has」应为「one has」之误；底本如此，译文按「有」处理。]

对任意 $F\in C(\mathcal{B},\Omega)$，$\mathcal{D}$-复形 $\mathcal{D}F$
带上算子 $i_\xi^{\mathcal{D}F}=i_\xi^F\otimes\operatorname{id}_{\mathcal{D}Y}$，
是 $K$-等变的。对任意 $N\in C(\mathcal{B},\mathcal{D})$，$\Omega$-复形
$\Omega N$ 带算子 $i_\xi^{\Omega N}$；后者在
$N^i\otimes\Lambda^{-j}\Theta_Y$ 上按
$n\otimes\tau\mapsto i_\xi n\otimes\tau+(-1)^in\otimes\sigma(\xi)\wedge\tau$
作用，故 $\Omega N$ 是 $K_\Omega$-等变的 $\Omega$-复形。于是有伴随函子
$\mathcal{D}$、$\Omega$：
\begin{equation}\tag{322}
C(\mathcal{B},\Omega)\rightleftarrows C(\mathcal{B},\mathcal{D})
\end{equation}
以及三角范畴之间互逆的等价
\begin{equation}\tag{323}
D(\mathcal{B},\Omega)\rightleftarrows D(\mathcal{B},\mathcal{D}).
\end{equation}
后一个等价把 $D(\mathcal{B},\mathcal{D})$ 上上述 $t$-结构与 7.3.2 中定义的
$D(\mathcal{B},\Omega)$ 上的 $t$-结构等同起来。当所论的叠是光滑簇对群作用的商时，
这给出 7.3.4 的另一个证明。

\bdsec{7.7. Harish-Chandra 模及其导出范畴}

\SourcePage{267}

\bdp{7.7.1} 设 $G$ 是仿射代数群，$K\subset G$ 是代数子群，于是有
Harish-Chandra 对 $(\mathfrak{g},K)$。考虑范畴
$\mathcal{M}(K\backslash G/^{/}G)=\mathcal{M}((K\backslash G)/^{/}G)$，
其对象是 $G$ 上的 $\mathcal{D}$-模，带彼此交换的 $K$-作用与弱 $G$-作用
（其中 $K$ 与 $G$ 分别经左、右平移作用在 $G$ 上）。对
$M\in\mathcal{M}(K\backslash G/^{/}G)$，置
$\gamma(M)=\gamma^r(M):=\Gamma(G,M_G)^G$；这里把 $M_G$ 看作 $G$ 上的右
$\mathcal{D}$-模。这是 $(\mathfrak{g},K)$-模：$\mathfrak{g}$ 经右 $G$-平移
不变的向量场作用在 $\gamma(M)$ 上（按 $M$ 上的 $\mathcal{D}$-模结构），
$K$ 经左 $K$-平移作用。

\bdp{7.7.2} \emph{引理.} 函子
$\gamma:\mathcal{M}(K\backslash G/^{/}G)\to\mathcal{M}(\mathfrak{g},K)$
是范畴等价。

\emph{证明.} 留给读者（或见 [Kas]）。\hfill$\Box$

\bdp{7.7.3} \emph{注记.} (i) 置 $\gamma^l(M):=\Gamma(G,M^l_G)^G$，其中
$M^l_G$ 是 $M$ 的左 $\mathcal{D}$-模实现。由与上面相同的理由，这是
$(\mathfrak{g},K)$-模；有显然的等同
$\gamma^l(M)=\gamma^r(M)\otimes\det\mathfrak{g}$。

(ii) 有向量空间的典范同构 $\gamma^l(M)\cong M^l_{G,1}
=M^l_{K\backslash G,1}$，它把 $G$-不变截面送到该截面在 $1\in G$ 处的值。
$M^l_{K\backslash G,1}$ 上的 $(\mathfrak{g},K)$-模结构可描述如下。$K$-作用
来自右 $K$-平移在 $K\backslash G$ 上的（弱）作用（注意 $K$ 是
$1\in K\backslash G$ 的稳定化子），$\mathfrak{g}$-作用来自与弱 $G$-作用
相对应的 $\mathfrak{g}$ 的 $\natural$-作用（见 7.6.10）。

(iii) 设 $P$ 是 $K$-模，$\mathcal{P}$ 是 $K\backslash G$ 上相应的 $G$-等变
向量丛，其纤维为 $\mathcal{P}_1=P$。有
$\mathcal{D}\mathcal{P}=\mathcal{P}\otimes\mathcal{D}_{K\backslash G}
\in\mathcal{M}((K\backslash G)/^{/}G)$，且
$\gamma(\mathcal{D}\mathcal{P})=U(\mathfrak{g})\otimes_{U(\mathfrak{k})}
(P\otimes\det\mathfrak{k}^*)$。

\bdp{7.7.4} 上述引理给出 $(\mathfrak{g},K)$-模的导出范畴
$D(\mathfrak{g},K)$ 上典范的 $\mathcal{H}$-作用，这正是 7.1.1(c) 所许诺的。
事实上，由 7.6.8(v)（以及 7.6.9）知 $D(K\backslash G/^{/}G)$ 是
$\mathcal{H}$-模；而 7.7.2 把 $D(\mathfrak{g},K)$ 与这个范畴等同起来。
下面在 7.8.2 中给出该作用的另一种描述。它与当前定义的等价性在 7.8.9、
7.8.10(i) 中确立。

\SourcePage{268}

本节其余部分（7.7.5--7.7.11）是 Harish-Chandra 情形下
$\mathcal{D}$-$\Omega$ 等价的一个插叙；作为附带收获，7.7.12 给出
Bernstein--Lunts 定理 [BL]1.3 的一个简单证明。读者可以跳过这部分，
直接看 7.8。

\bdp{7.7.5} 下面是 7.7.2 对 $\Omega$-复形的版本。

设 $\Omega_{\mathfrak{g}}$ 是 $\mathfrak{g}$ 的 Chevalley 上链 DG-代数，故
$\Omega_{\mathfrak{g}}^{\cdot}=\Lambda^{\cdot}\mathfrak{g}^*$。它带
$K_\Omega$ 的典范「伴随」作用（记号见 7.6.3）。具体说，$K$ 以余伴随方式
作用在 $\Omega_{\mathfrak{g}}^{\cdot}$ 上；$\xi\in\mathfrak{k}
=\mathfrak{k}_\Omega^{-1}$ 作为 $\Omega_{\mathfrak{g}}^{\cdot}$ 的导子
$i_\xi$ 作用，把 $\nu\in\mathfrak{g}^*=\Omega_{\mathfrak{g}}^1$ 送到
$\langle\nu,\xi\rangle$。

$\Omega_{(\mathfrak{g},K)}$-\emph{复形}指一个 DG
$(\Omega_{\mathfrak{g}},K_\Omega)$-模；也就是说，它是一个复形，带
$\Omega_{\mathfrak{g}}$-作用与 $K_\Omega$-作用，二者关于 $K_\Omega$ 在
$\Omega_{\mathfrak{g}}$ 上的作用相容。对 $\Omega_{(\mathfrak{g},K)}$-复形
$T$，把 $\nu\in\mathfrak{g}^*=\Omega_{\mathfrak{g}}^1$、
$\xi\in\mathfrak{k}=\mathfrak{k}_\Omega^{-1}$ 在 $T$ 上的作用分别记作
$a_\nu$、$i_\xi$。把 $\Omega_{(\mathfrak{g},K)}$-复形的 DG 范畴记作
$C\Omega_{(\mathfrak{g},K)}$，其同伦范畴记作
$K\Omega_{(\mathfrak{g},K)}$。

对 $F\in C(K\backslash G/^{/}G,\Omega)$，置
$\gamma(F):=\Gamma(G,F_G)^G$。这是 $\Omega_{(\mathfrak{g},K)}$-复形。
事实上，$\Omega_{\mathfrak{g}}^{\cdot}$ 经与「$G$ 上对右 $G$-平移不变的
微分形式 DG 代数」的通常等同作用在它上面；而 $K_\Omega$ 作用在
$\gamma(F)$ 上，因为它作用在 $F_G$ 上（见 7.6.4、7.6.5）。

\bdp{7.7.6} \emph{引理.} 函子
$\gamma:C(K\backslash G/^{/}G,\Omega)\to C\Omega_{(\mathfrak{g},K)}$
是 DG 范畴的等价。

\emph{证明.} 留给读者。\hfill$\Box$

\bdp{7.7.7} 我们把 $(\mathfrak{g},K)$-复形与
$\Omega_{(\mathfrak{g},K)}$-复形等同于 $K\backslash G$ 上的弱 $G$-等变复形。
下面写出 Harish-Chandra 情形下的标准函子 $\mathcal{D}$ 与 $\Omega$。
方便起见，引入 DG Harish-Chandra 对
$(\mathfrak{k}_\Omega\times\mathfrak{g},K)$（结构嵌入
$\operatorname{Lie}K\hookrightarrow\mathfrak{k}_\Omega\times\mathfrak{g}$
是对角映射）。

设 $DR_{\mathfrak{g}}$ 是 $\mathfrak{g}$ 的、系数在 $U\mathfrak{g}$ 中的
Chevalley 上链复形（把 $U\mathfrak{g}$ 看作左 $U\mathfrak{g}$-模），故
$DR_{\mathfrak{g}}^i=\Lambda^i\mathfrak{g}^*\otimes U\mathfrak{g}$。
于是 $DR_{\mathfrak{g}}$ 是 $\Omega_{\mathfrak{g}}$-复形，也是
$(\mathfrak{k}_\Omega\times\mathfrak{g},K)$-复形；这两个作用相容。

\SourcePage{269}

（这里 $(\mathfrak{k}_\Omega\times\mathfrak{g},K)$ 经投影
$(\mathfrak{k}_\Omega\times\mathfrak{g},K)\to(\mathfrak{k}_\Omega,K)$
作用在 $\Omega_{\mathfrak{g}}$ 上，见 7.7.5）。具体说，对
$\nu\in\Omega_{\mathfrak{g}}$，
$\varepsilon=(\varepsilon_l,\varepsilon_r)\in\mathfrak{k}\times\mathfrak{g}
=\mathfrak{k}_\Omega^0\times\mathfrak{g}$，
$\xi\in\mathfrak{k}=\mathfrak{k}_\Omega^{-1}$，$k\in K$ 与
$a=\alpha\otimes v\in DR_{\mathfrak{g}}$，有
\[
\nu a=\nu\alpha\otimes v,\qquad
\varepsilon a=\operatorname{Ad}_{\varepsilon_l}(\alpha)\otimes v
+\alpha\otimes(\varepsilon_l v-v\varepsilon_r),
\]
\[
\xi a=i_\xi(\alpha)\otimes v,\qquad
ka=\operatorname{Ad}_k(\alpha)\otimes\operatorname{Ad}_k(v).
\]

对 $(\mathfrak{g},K)$-模的复形（简称 $(\mathfrak{g},K)$-复形）$V$，置
$\Omega V:=\operatorname{Hom}_{\mathfrak{g}}(DR_{\mathfrak{g}},V)$；
按显然方式这是 $\Omega_{(\mathfrak{g},K)}$-复形。对
$\Omega_{(\mathfrak{g},K)}$-复形 $T$，置
$\mathcal{D}T=\mathcal{D}_{(\mathfrak{g},K)}T:=
T\otimes_{\Omega_{\mathfrak{g}},\mathfrak{k}_\Omega}DR_{\mathfrak{g}}
=(T\otimes_{\Omega_{\mathfrak{g}}}DR_{\mathfrak{g}})_{\mathfrak{k}_\Omega}$；
这是 $(\mathfrak{g},K)$-复形。于是有伴随 DG 函子
\begin{equation}\tag{324}
\mathcal{D}=\mathcal{D}_{(\mathfrak{g},K)}:
C\Omega_{(\mathfrak{g},K)}\longrightarrow C(\mathfrak{g},K),\qquad
\Omega:C(\mathfrak{g},K)\longrightarrow C\Omega_{(\mathfrak{g},K)}.
\end{equation}

\emph{注.} 对上述 $T$，令
$\overline{T}^{\cdot}\subset T^{\cdot}$ 为所有算子 $i_\xi$
（$\xi\in\mathfrak{k}$）的核。这是 $T^{\cdot}$ 的 $K$-子模与
$\Lambda^{\cdot}(\mathfrak{g}/\mathfrak{k})^*$-子模（这里
$\Lambda^{\cdot}(\mathfrak{g}/\mathfrak{k})^*\subset
\Lambda^{\cdot}\mathfrak{g}^*=\Omega_{\mathfrak{g}}^{\cdot}$），
且显然的态射
\begin{equation}\tag{325}
\Omega_{\mathfrak{g}}^{\cdot}
\otimes_{\Lambda^{\cdot}(\mathfrak{g}/\mathfrak{k})^*}
\overline{T}^{\cdot}\longrightarrow T^{\cdot},\qquad
\overline{T}^{\cdot}\otimes_{U\mathfrak{k}}U\mathfrak{g}
\longrightarrow\mathcal{D}T^{\cdot}
\end{equation}
是同构。

\bdp{7.7.8} 回到几何情形。有显然的等同
$\Gamma(G,DR_G)^G=DR_{\mathfrak{g}}$（记号见 7.2.2；$G$ 经右平移作用在
自己上）。对 $M\in C((K\backslash G)/^{/}G,\mathcal{D})$，有
$\Omega_{(\mathfrak{g},K)}$-复形的典范同构
$\gamma(\Omega M)\cong\Omega(\gamma M)$，它定义为复合
\[
\Gamma(G,\mathit{Hom}_{\mathcal{D}_G}(DR_G,M_G))^G
=\operatorname{Hom}_{\mathcal{D}_G}(DR_G,M_G)^G
=\operatorname{Hom}_{U\mathfrak{g}}(DR_{\mathfrak{g}},\gamma M).
\]
对 $F\in C(K\backslash G/^{/}G,\Omega)$，有类似的典范同构
$\gamma\mathcal{D}F\cong\mathcal{D}\gamma F$，其定义留给读者。

\bdp{7.7.9} 对 $\Omega_{(\mathfrak{g},K)}$-复形 $T$，置
$H_{\mathfrak{g}}^{\cdot}T=H^{\cdot}\mathcal{D}T
\in\mathcal{M}(\mathfrak{g},K)$。于是
$H_{\mathfrak{g}}^{\cdot}:K\Omega_{(\mathfrak{g},K)}\to
\mathcal{M}(\mathfrak{g},K)$ 是上同调函子。把
$K\Omega_{(\mathfrak{g},K)}$ 中诱导 $H_{\mathfrak{g}}^{\cdot}$ 之间同构的
态射称为 $\mathfrak{g}$-\emph{拟同构}。$\mathfrak{g}$-拟同构构成局部化族；
把 $K\Omega_{(\mathfrak{g},K)}$ 的相应局部化记作
$D\Omega_{(\mathfrak{g},K)}$。函子 $\mathcal{D}$、$\Omega$ 给出互逆的
导出范畴等价

\SourcePage{270}

\begin{equation}\tag{326}
D\Omega_{(\mathfrak{g},K)}\rightleftarrows D(\mathfrak{g},K),
\end{equation}
其中 $D(\mathfrak{g},K):=D\mathcal{M}(\mathfrak{g},K)$。等价 $\gamma$
给出导出范畴的等价
\begin{equation}\tag{327}
D(K\backslash G/^{/}G,\Omega)\cong D\Omega_{(\mathfrak{g},K)},\qquad
D((K\backslash G)/^{/}G,\mathcal{D})\cong D(\mathfrak{g},K).
\end{equation}

\bdp{7.7.10} \emph{注记.} (i) 任意 $\mathfrak{g}$-拟同构都是拟同构；
反之未必。

(ii) 任意 $\Omega_{(\mathfrak{g},K)}$-复形 $T$ 都可看作
$\Omega_{\mathfrak{g}}=\Omega_{(\mathfrak{g},1)}$-复形（忘掉
$K_\Omega$-作用），于是有相应的 $\mathfrak{g}$-模复形
$\mathcal{D}_{\mathfrak{g}}T:=T\otimes_{\Omega_{\mathfrak{g}}}
DR_{\mathfrak{g}}$。显然的投影 $\mathcal{D}_{\mathfrak{g}}T\to
\mathcal{D}_{(\mathfrak{g},K)}T$ 是拟同构。这意味着：
$\Omega_{(\mathfrak{g},K)}$-复形的态射是 $\mathfrak{g}$-拟同构，
当且仅当它是 $\Omega_{\mathfrak{g}}$-复形的 $\mathfrak{g}$-拟同构。

\bdp{7.7.11} 7.7.7、7.7.9 的格式有如下版本。回忆 $DR_{\mathfrak{g}}$ 是
$(\mathfrak{k}_\Omega\times\mathfrak{g},K)$-复形。于是上述
$\mathcal{D}_{\mathfrak{g}}T$ 是
$(\mathfrak{k}_\Omega\times\mathfrak{g},K)$-复形；而对
$(\mathfrak{k}_\Omega\times\mathfrak{g},K)$-复形 $V$，复形
$\Omega V:=\operatorname{Hom}_{\mathfrak{g}}(DR_{\mathfrak{g}},V)$ 是
$\Omega_{(\mathfrak{g},K)}$-复形。函子
\begin{equation}\tag{328}
\mathcal{D}_{\mathfrak{g}}:C\Omega_{(\mathfrak{g},K)}\longrightarrow
C(\mathfrak{k}_\Omega\times\mathfrak{g},K),\qquad
\Omega:C(\mathfrak{k}_\Omega\times\mathfrak{g},K)\longrightarrow
C\Omega_{(\mathfrak{g},K)}
\end{equation}
是伴随的，同伦范畴之间的相应函子亦然。过渡到导出范畴，它们成为
（用 7.7.10(ii)）互逆的等价
\begin{equation}\tag{329}
D\Omega_{(\mathfrak{g},K)}\rightleftarrows
D(\mathfrak{k}_\Omega\times\mathfrak{g},K).
\end{equation}
投影 $(\mathfrak{k}_\Omega\times\mathfrak{g},K)\to(\mathfrak{g},K)$ 给出
全忠实嵌入 $C(\mathfrak{g},K)\to
C(\mathfrak{k}_\Omega\times\mathfrak{g},K)$，从而有正合函子
\begin{equation}\tag{330}
D(\mathfrak{g},K)\longrightarrow
D(\mathfrak{k}_\Omega\times\mathfrak{g},K).
\end{equation}
下述定理属于 Bernstein 与 Lunts [BL]1.3\footnote{[BL] 的作者只考虑有界导出范畴。}：

\SourcePage{271}

\bdp{7.7.12} \emph{定理.} 函子（330）是范畴等价。

\emph{证明.} (328) 中的函子 $\Omega$ 限制到 $C(\mathfrak{g},K)$ 上，
与 (324) 中的 $\Omega$ 一致。于是 7.7.12 由（326）与（329）推出。逆函子
$D(\mathfrak{k}_\Omega\times\mathfrak{g},K)\to D(\mathfrak{g},K)$ 把 $V$
送到 $\mathcal{D}_{(\mathfrak{g},K)}\Omega V$。\hfill$\Box$

\bdsec{7.8. Hecke 作用与局部化函子}

\bdp{7.8.1} 我们将描述 Harish-Chandra 模的导出范畴上典范的 Hecke 作用。
我们考虑扭情形，即 $\mathfrak{g}$ 的中心扩张的表示。以下是特征标的清单。

设 $G'$ 是 $G$ 被 $\mathbb{G}_m$ 的中心扩张，并带分裂 $K\to G'$。于是
$K$ 的原像 $K'\subset G'$ 与 $K\times\mathbb{G}_m$ 等同。置
$\mathfrak{g}':=\operatorname{Lie}G'$，
$\mathfrak{k}':=\operatorname{Lie}K'
=\mathfrak{k}\times\mathbb{C}$。有 Harish-Chandra 对
$(\mathfrak{g}',K')$，以及配套的 DG 对
$(\mathfrak{k}_\Omega\times\mathfrak{g}',K')$（这里结构嵌入
$\mathfrak{k}'\hookrightarrow
\mathfrak{k}_\Omega\times\mathfrak{g}'$ 的第一个分量是投影
$\mathfrak{k}'\to\mathfrak{k}$）。

设 $\mathcal{M}(\mathfrak{g},K)'$ 是由 $(\mathfrak{g}',K')$-模构成的
范畴，其中 $\mathbb{G}_m\subset K'$ 按标准特征标作用；其对象称为
$(\mathfrak{g},K)'$-模，或简称 Harish-Chandra 模。这是 Abel 范畴。
类似地，设 $C(\mathfrak{k}_\Omega\times\mathfrak{g},K)'$ 是由那些
$(\mathfrak{k}_\Omega\times\mathfrak{g}',K')$-复形构成的范畴，其中
$\mathbb{G}_m$ 按标准特征标作用；其对象称为
$(\mathfrak{k}_\Omega\times\mathfrak{g},K)'$-复形，或简称
Harish-Chandra 复形。这是 DG 范畴，带显然的上同调函子，取值于
$\mathcal{M}(\mathfrak{g},K)'$。把相应的导出范畴记作
$D(\mathfrak{g},K)'$；这是以 $\mathcal{M}(\mathfrak{g},K)'$ 为核心的
$t$-范畴。

\emph{注.} 由 Bernstein--Lunts 定理的扭版本，
$D(\mathfrak{g},K)'$ 等价于 $\mathcal{M}(\mathfrak{g},K)'$ 的导出范畴
\footnote{扭 Bernstein--Lunts 由非扭版本（见 7.7.12）施于 Harish-Chandra
对 $(\mathfrak{g}',K')$ 得到。}。下文不用这一事实，因为 Hecke 作用是
自然用 $(\mathfrak{k}_\Omega\times\mathfrak{g},K)'$-复形定义的。

\SourcePage{272}

\bdp{7.8.2} 现在定义 $D(\mathfrak{g},K)'$ 上典范的 $\mathcal{H}$-作用。
先定义预 Hecke 幺半 DG 范畴 $\mathcal{H}^c:=C(K\backslash G/K,\Omega)$
在 $C(\mathfrak{k}_\Omega\times\mathfrak{g},K)'$ 上的作用；Hecke 作用
在过渡到导出范畴之后才出现。

记 $\mathcal{L}_G$ 为 $G$ 上对应于 $\mathbb{G}_m$-torsor $G'\to G$ 的
线丛。$G$ 在自己上的左、右平移作用典范地提升为 $\mathcal{L}_G$ 上的
$G'$-作用。于是 $\mathcal{L}_G$ 的截面等同于 $G'$ 上的函数 $\phi$，
使得对 $c\in\mathbb{G}_m$、$g'\in G'$ 有
$\phi(cg')=c^{-1}\phi(g')$。因此 $\mathbb{G}_m\subset G'$ 经右平移在
$\mathcal{L}_G$ 的截面上作用，作用的特征标与标准特征标互逆。

取 Harish-Chandra 复形
$V\in C(\mathfrak{k}_\Omega\times\mathfrak{g},K)'$。置
$\mathcal{V}_G:=\mathcal{L}_G\otimes V$。则 $\mathcal{V}_G$ 是 $G$ 上
左 $\mathcal{D}$-模的复形。事实上，$\mathfrak{g}'$ 在 $\mathcal{L}_G$
上的无穷小右平移作用与 $V$ 上的 $\mathfrak{g}'$-作用张量起来，给出
$\mathcal{V}_G$ 上的 $\mathfrak{g}$-作用。$\mathcal{V}_G$ 上的左
$\mathcal{D}$-模结构使得左不变向量场经上述 $\mathfrak{g}$-作用作用在
$\mathcal{V}_G$ 上。$\mathcal{D}$-复形 $\mathcal{V}_G$ 关于左 $G'$-平移
弱等变：这些平移按 $\mathcal{L}_G$ 上的相应作用与 $V$ 上的平凡作用之
张量作用。因此由 7.6.10，它带典范的 $\mathfrak{g}'$-作用 $\natural$。

\emph{注.} 对 $\theta\in\mathfrak{g}'$，考虑函数
$\theta^{\natural}:G\to\mathfrak{g}'$，
$\theta^{\natural}(g):=\operatorname{Ad}_g(\theta)$。则对 $v\in V$、
$l\in\mathcal{L}_G$ 有
$\theta^{\natural}(l\otimes v)=l\otimes\theta^{\natural}(v)$。

取 $F\in\mathcal{H}^c$。则 $F_G\otimes\mathcal{V}_G$ 是 $G$ 上的
$\Omega$-复形（见 7.2.3(ii)）。它关于右 $K$-平移是 $K_\Omega$-等变的。
具体说，$K$ 按 $F$、$\mathcal{L}_G$ 上的相应作用与 $V$ 上的结构作用之
张量作用；算子 $i_\xi$ 按 $F$ 上右平移作用的相应算子与 $V$ 的结构算子
之和作用。把 $G/K$ 上相应的 $\Omega$-复形记作
$(F\otimes\mathcal{V})_{G/K}$。$F_G\otimes\mathcal{V}_G$ 上来自
$\mathcal{V}_G$ 上作用 $\natural$ 的 $\mathfrak{g}'$-作用与这个
$K_\Omega$-作用交换，故它定义了 $(F\otimes\mathcal{V})_{G/K}$ 上的
$\mathfrak{g}'$-作用。我们也把它记作 $\natural$。

\emph{注.} 若 $V$ 是 $(\mathfrak{g},K)'$-模的复形，则 $\mathcal{V}_G$ 是
左 $\mathcal{D}_G$-模的复形，关于右 $K$-平移强等变。

\SourcePage{273}

设 $\mathcal{V}_{G/K}$ 是 $G/K$ 上相应的左 $\mathcal{D}$-模复形。有
$(F\otimes\mathcal{V})_{G/K}=F_{G/K}\otimes\mathcal{V}_{G/K}$。

置 $F\overset{\sim}{\circledast}V:=
\Gamma(G,F_G\otimes\mathcal{V}_G)$，以及
\begin{equation}\tag{331}
F\overset{c}{\circledast}V=\Gamma(G/K,(F\otimes\mathcal{V})_{G/K})
=(F\overset{\sim}{\circledast}V)^{K_\Omega}.
\end{equation}

这是 $(\mathfrak{k}_\Omega\times\mathfrak{g},K)'$-复形。事实上，
$\mathfrak{g}'$ 按 $\natural$ 作用，$K$ 按 $F$ 与 $\mathcal{V}$ 的左平移
作用之张量作用，而算子 $i_\xi$ 是 $F$ 的相应算子。Harish-Chandra
相容性留给读者验证。

现在 $\overset{c}{\circledast}$ 定义了
$C(\mathfrak{k}_\Omega\times\mathfrak{g},K)'$ 上的
$\mathcal{H}^c$-模结构。事实上，结合约束
$(F_1\overset{c}{\circledast}F_2)\overset{c}{\circledast}V=
F_1\overset{c}{\circledast}(F_2\overset{c}{\circledast}V)$ 由显然的等同
\[
\Gamma(G,(F_1\overset{c}{\circledast}F_2)^{\cdot}\otimes\mathcal{L}_G)
=[\Gamma(G,F_1^{\cdot}\otimes\mathcal{L}_G)
\otimes\Gamma(G,F_2^{\cdot}\otimes\mathcal{L}_G)]^{K_\Omega}
\]
推出，其中 $K_\Omega$ 按右、左平移作用之张量作用（见 7.6.5）。我们把
Hecke 作用 $\circledast:\mathcal{H}\times D(\mathfrak{g},K)'\to
D(\mathfrak{g},K)'$ 定义为 $\overset{c}{\circledast}$ 的右导出函子。
若 $F$ 松，则 $F\circledast V=F\overset{c}{\circledast}V$，故
$\circledast$ 的结合约束由 $\overset{c}{\circledast}$ 的结合约束推出。

\emph{注.} 由上一则注，对
$M\in\mathcal{M}(K\backslash G/K)\subset\mathcal{H}$、
$V\in\mathcal{M}(\mathfrak{g},K)'$ 有
\begin{equation}\tag{332}
H^{\cdot}M\circledast V
=H_{DR}^{\cdot}(G/K,M\otimes\mathcal{V}_{G/K}).
\end{equation}

\bdp{7.8.3} \emph{注.} 设我们的扭是平凡的，故
$G'=G\times\mathbb{G}_m$。有显然的等价
$\mathcal{M}(\mathfrak{g},K)'=\mathcal{M}(\mathfrak{g},K)$ 与
$D(\mathfrak{g},K)=D(\mathfrak{g},K)'$（见 7.7.11）。于是我们定义了
$D(\mathfrak{g},K)$ 上的 Hecke 作用。7.8.9 将表明，这个作用确实与
7.7.4 中的作用一致。

回到一般情形。设 $U'$ 是 $\mathfrak{g}$ 的扭包络代数；用 $\mathfrak{z}$
记其由 $\operatorname{Ad}G$-不变元素构成的子代数。交换代数 $\mathfrak{z}$
按显然方式作用在任意 Harish-Chandra 复形上，故
$C(\mathfrak{k}_\Omega\times\mathfrak{g},K)'$ 从而 $D(\mathfrak{g},K)$
是 $\mathfrak{z}$-范畴。
% [Tutor Note: 原书 p.273 末句此处印作「hence D(g,K)」（g 为 fraktur 体），
% 无撇号；按上下文应为带撇的 D(g,K)。底本如此，保留原样。]

\SourcePage{274}

\bdp{7.8.4} \emph{引理.} $C(\mathfrak{k}_\Omega\times\mathfrak{g},K)'$、
$D(\mathfrak{g},K)'$ 上的 Hecke 作用是 $\mathfrak{z}$-线性的。

\emph{证明.} 用 7.8.2 中第一则注。\hfill$\Box$

\bdp{7.8.5} \emph{例.}（5 中将用到）令 $Vac':=U'/U'\mathfrak{k}$ 为扭
真空模。我们来显式计算 $F\circledast Vac'$。记号沿用 7.8.2。于是按 7.8.2
中第二则注，有 $G/K$ 上的左 $\mathcal{D}$-模 $\mathcal{V}_{G/K}$，它关于
左 $G$-平移弱等变，且 $\mathcal{V}_G=\mathcal{L}_G\otimes Vac'$。嵌入
$\mathbb{C}\subset Vac'$ 给出嵌入
$\mathcal{L}_{G/K}\subset\mathcal{V}_{G/K}$。容易看出，相应的左
$\mathcal{D}_{G/K}$-模态射 $\mathcal{D}_{G/K}
\otimes_{\mathcal{O}_{G/K}}\mathcal{L}_{G/K}\to\mathcal{V}_{G/K}$
是弱 $G$-等变 $\mathcal{D}$-模的同构。

\emph{注.} $\mathcal{D}_{G/K}\otimes\mathcal{L}_{G/K}$ 上与 $\natural$
对应的 $\mathfrak{g}'$-作用由下列公式给出：$\alpha'(\psi\otimes l)
=\psi\otimes\alpha'(l)-\psi\cdot\alpha\otimes l$，其中
$\alpha'\in\mathfrak{g}'$，$\alpha$ 是 $G/K$ 上相应的左平移向量场，
$\alpha'(l)$ 是 $l\in\mathcal{L}_{G/K}$ 的无穷小左平移。

于是对 $F\in\mathcal{H}^c$，有
$(F\otimes\mathcal{V})_{G/K}=F_{G/K}\otimes\mathcal{D}_{G/K}
\otimes_{\mathcal{O}_{G/K}}\mathcal{L}_{G/K}
=\mathcal{D}(F_{G/K})\otimes_{\mathcal{O}_{G/K}}\mathcal{L}_{G/K}$。
因此
\begin{equation}\tag{333}
F\overset{c}{\circledast}Vac'
=\Gamma(G/K,\mathcal{D}(F_{G/K})\otimes\mathcal{L}_{G/K}).
\end{equation}

这里 $\Gamma(G/K,\mathcal{D}(F_{G/K})\otimes\mathcal{L}_{G/K})$ 上的
$(\mathfrak{k}_\Omega\times\mathfrak{g},K)'$-作用定义如下。$\mathfrak{g}'$-
作用来自上面注中所述 $\mathcal{D}(F_{G/K})\otimes\mathcal{L}_{G/K}$ 上的
$\mathfrak{g}'$-作用；$K$-作用是左平移作用；算子 $i_\xi$ 来自
$F_{G/K}$ 上相应的算子。

过渡到导出函子（相当于在上述公式中考虑松的 $F$），得到
\begin{equation}\tag{334}
F\circledast Vac'=R\Gamma(G/K,
\mathcal{D}(F_{G/K})\otimes\mathcal{L}_{G/K}).
\end{equation}
特别地，对 $M\in\mathcal{M}(K\backslash G/K)$ 有
\begin{equation}\tag{335}
M\circledast Vac'=R\Gamma(G/K,
M_{G/K}\otimes\mathcal{L}_{G/K}).
\end{equation}

\SourcePage{275}

这里右端的 $\mathfrak{g}'$-作用来自
$M_{G/K}\otimes\mathcal{L}_{G/K}$ 上的 $\mathfrak{g}'$-作用，由公式
$\alpha'(m\otimes l)=m\otimes\alpha'(l)-m\alpha\otimes l$ 给出。

\bdp{7.8.6} 我们来解释 7.1.1 中「Hecke 模式」的 (d) 部分。先定义局部化
函子 $\Delta$。记号沿用 7.8.1。设 $Y$ 是 $G$ 作用其上的光滑簇，
$\mathcal{L}=\mathcal{L}_Y$ 是 $Y$ 上的线丛。设 $\mathcal{L}$ 带
$G'$-作用，该作用提升 $Y$ 上的 $G$-作用，并使
$\mathbb{G}_m\subset G'$ 按与标准特征标相反的特征标作用。线丛
$\omega_Y\otimes\mathcal{L}$ 带类似的作用。

我们如下定义一个 DG 函子
\begin{equation}\tag{336}
\Delta_\Omega=\Delta_{\Omega\mathcal{L}}:
C(\mathfrak{k}_\Omega\times\mathfrak{g},K)'\longrightarrow
C(K\backslash Y,\Omega)
\end{equation}
注意 $(\mathfrak{g}',K')$ 从而
$(\mathfrak{k}_\Omega\times\mathfrak{g}',K')$ 作用在
$\omega_Y\otimes\mathcal{L}$ 上（因为 $G'$ 如此）。对 Harish-Chandra
复形 $V$，考虑 $\mathcal{O}$-模的复形
$\omega_Y\otimes\mathcal{L}\otimes V$。
$\omega_Y\otimes\mathcal{L}$ 与 $V$ 上
$(\mathfrak{k}_\Omega\times\mathfrak{g}',K')$-作用的张量积给出
$\omega_Y\otimes\mathcal{L}\otimes V$ 上的
$(\mathfrak{k}_\Omega\times\mathfrak{g},K)$-作用。置
\[
\Delta_\Omega(V):=\operatorname{Hom}_{\mathfrak{g}}
(DR_{\mathfrak{g}},\omega_Y\otimes\mathcal{L}\otimes V)[-\dim K]
\]
（记号见 7.7.7）。换言之，$\Delta_\Omega(V)$ 是 $\mathfrak{g}$ 的、系数在
$\omega_Y\otimes\mathcal{L}\otimes V$ 中的平移后的 Chevalley 链复形。
这是 $Y$ 上的 $\Omega$-复形。由于 $DR_{\mathfrak{g}}$ 与
$\omega_Y\otimes\mathcal{L}\otimes V$ 都是
$(\mathfrak{k}_\Omega\times\mathfrak{g},K)$-复形，故
$\Delta_\Omega(V)$ 是 $K_\Omega$-等变的，即
$\Delta_\Omega(V)\in C(K\backslash Y,\Omega)$。

注意 $\Delta_\Omega(V)$ 带典范的递增有限滤过，其相邻商等于
$\Lambda^i\mathfrak{g}\otimes\omega_Y\otimes\mathcal{L}\otimes V
[i-\dim K]$。因此 $\Delta_\Omega$ 把拟同构送到 $\mathcal{D}$-拟同构。
于是它给出三角函子
\begin{equation}\tag{337}
L\Delta=L\Delta_{\mathcal{L}}:
D(\mathfrak{g},K)'\longrightarrow D(K\backslash Y)
\end{equation}
上述注还表明，$L\Delta$ 是右 $t$-正合函子。核心之间相应的右正合函子是
$\Delta_{\mathcal{L}}:\mathcal{M}(\mathfrak{g},K)'\to$
% 尾部接 p.276（原书此句跨页，p.275 末行止于箭头）。
%% end units/ch07_u04.tex
%% begin units/ch07_u05.tex
% ch07_u05 —— 原书 p.276–p.285（7.8.7 -- 7.10.1 开头）
% 本单元自 p.275 末句「$\Delta_{\mathcal{L}}:\mathcal{M}(\mathfrak{g},K)'\to$」的断处续起，
% 至 7.10.1 的「Set」止；尾部接 p.286。

\SourcePage{276}

$\mathcal{M}^{\ell}(K\backslash Y)$ 把 $(\mathfrak{g},K)'$-模 $V$ 送到 $K$-等变左
$\mathcal{D}_{Y}$-模 $(\mathcal{D}_{Y}\otimes\mathcal{L})\otimes_{U(\mathfrak{g}')}V$。
更一般地，
\[
H^{i}_{\mathcal{D}}L\Delta_{\mathcal{L}}(V)=H_{-i}(\mathfrak{g},\mathcal{D}_{Y}\otimes\mathcal{L}\otimes V).
\]

\bdp{7.8.7} \emph{注.} (i) 上述构造只用到 $(\mathfrak{g}',K')$ 在 $(Y,\mathcal{L})$ 上的作用
（不需要整个 $G'$-作用）。
% [Tutor Note: 原书 p.276 此处 "construction" 误拼作「construcion」，底本如此，译文无法保留拼写，照常翻译。]

(ii) 可以证明，$L\Delta_{\mathcal{L}}$ 是 $\Delta_{\mathcal{L}}$ 的左导出函子
（见 7.8.1 的注）。

(iii) 设 $(\mathfrak{g}',K')$ 是 $(\mathfrak{g},K)$ 的平凡扩张，于是
$(\mathfrak{g},K)'$-模与 $(\mathfrak{g},K)$-模相同，而 $\mathcal{L}$ 是
$\mathcal{O}_{Y}$ 配上 $(\mathfrak{g}',K')$ 的显然作用。则
\[
\Delta_{\mathcal{L}}(V)=\mathcal{D}_{Y}\otimes_{U(\mathfrak{g})}V,
\]
即 $\Delta_{\mathcal{L}}$ 与 1.2.4 的函子 $\Delta$ 一致。

\bdp{7.8.8} \emph{命题.} 函子 $L\Delta_{\mathcal{L}}:D(\mathfrak{g},K)'\to D(K\backslash Y)$
是 $\mathcal{H}$-模的态射。

\emph{证明.} 只需证明函子
$\Delta_{\Omega\mathcal{L}}:C(\mathfrak{k}_{\Omega}\times\mathfrak{g},K)'\to C(K\backslash Y,\Omega)$
是 $\mathcal{H}^{c}$-模的态射。

取 7.8.2 中的 $F$、$V$。我们要定义与结合约束相容的 $\Omega$-复形的典范等同
\[
\alpha:\Delta_{\Omega}(F\overset{c}{\circledast}V)\xrightarrow{\ \sim\ }F\overset{c}{\circledast}\Delta_{\Omega}(V).
\]
我们将建立一个与 $K_{\Omega}$-作用相容的典范同构
\[
\tilde{\alpha}:\Delta_{\Omega}(F\widetilde{\circledast}V)\xrightarrow{\ \sim\ }F\widetilde{\circledast}\Delta_{\Omega}(V)
\]
（记号见 7.6.6、7.8.2）。取 $K_{\Omega}$-不变量即得 $\alpha$。

设 $m,p:G\times Y\to Y$ 为作用映射与投影映射，$i:G\times Y\to G\times Y$
为对称态射 $i(g,x)=(g,gx)$；有 $pi=m$。$\mathcal{L}$ 上的 $G'$-作用给出线丛的
$i$-同构 $\tilde{i}:\mathcal{O}_{G}\boxtimes\mathcal{L}_{Y}\xrightarrow{\sim}\mathcal{L}_{G}\boxtimes\mathcal{L}_{Y}$。

下面，对 $\mathfrak{g}$-复形 $P$，记 $C(P)$ 为系数在 $P$ 中、平移 $\dim K$
的李代数链 Chevalley 复形。于是 $C(P)^{\cdot}=C^{\cdot}\otimes P^{\cdot}$，
其中 $C^{a}:=\Lambda^{\dim K-a}\mathfrak{g}$。考虑 $\Omega$-复形
\[
F_{G}\boxtimes\Delta_{\Omega}(V)=F_{G}\boxtimes C(\omega_{Y}\otimes\mathcal{L}_{Y}\otimes V)
\]
与
\[
C((F_{G}\otimes V_{G})\boxtimes(\omega_{Y}\otimes\mathcal{L}_{Y}))
=C((F_{G}\otimes(\mathcal{L}_{G}\otimes V))\boxtimes(\omega_{Y}\otimes\mathcal{L}_{Y}));
\]
这里 $(F_{G}\otimes V_{G})\boxtimes(\omega_{Y}\otimes\mathcal{L}_{Y})$ 上的
$\mathfrak{g}$-作用，是 $\mathfrak{g}'$-作用 $\natural$ 与
$\omega_{Y}\otimes\mathcal{L}_{Y}$ 上的标准 $\mathfrak{g}'$-作用之张量积（见 7.8.2）。

\SourcePage{277}

存在 $\Omega$-复形的典范 $i$-同构
\[
\tilde{\alpha}':F_{G}\boxtimes\Delta_{\Omega}(V)
\xrightarrow{\ \sim\ }C((F_{G}\otimes V_{G})\boxtimes(\omega_{Y}\otimes\mathcal{L}_{Y})),
\]
其定义如下。对 $f\in F_{G}$、$\lambda\in C^{\cdot}$、
$l\in\omega_{Y}\otimes\mathcal{L}_{Y}$、$v\in V$，有
\[
\tilde{\alpha}'(f\otimes\lambda\otimes l\otimes v)
=a(\lambda)\otimes f\otimes\tilde{i}(l)\otimes v;
\]
这里 $a(\lambda)\in\mathcal{O}_{G\times Y}\otimes C^{\cdot}$ 是函数
$a(\lambda)(g,y)=\operatorname{Ad}_{g}(\lambda)$。
$\alpha$ 与微分交换，这一点留给读者验证（用 7.8.2 的注）。

现在有显然的等同
$m\cdot(F_{G}\boxtimes\Delta_{\Omega}(V))=F\widetilde{\circledast}\Delta_{\Omega}(V)$
与
$p\cdot C((F_{G}\otimes V_{G})\boxtimes(\omega_{Y}\otimes\mathcal{L}_{Y}))
=\Delta_{\Omega}(F\widetilde{\circledast}V)$。
于是 $\tilde{\alpha}'$ 给出所要的典范同构 $\tilde{\alpha}$。
$\tilde{\alpha}$ 与 $K_{\Omega}$-作用及结合约束的相容性，留给读者验证。\hfill$\square$

\bdp{7.8.9} 考虑 $Y=G$、$G$ 以左平移作用的情形，并取
$\mathcal{L}=\mathcal{L}_{Y}$ 为 $\mathcal{L}_{G}$ 的对偶线丛（见 7.8.2），
配上左平移给出的显然 $G'$-作用。右 $G'$-平移作用于我们的数据。
因此 $\Omega$-复形 $\Delta_{\Omega}(V)$ 关于 $G'$ 的右平移作用是弱 $G'$-等变的。

设 $C(K\backslash G/^{/}G,\Omega)'\subset C(K\backslash G/^{/}G',\Omega)$
为下述弱 $G'$-等变 $\Omega$-复形 $T$ 构成的子范畴：$\mathbb{G}_{m}\subset G'$
按标准特征标作用于 $T$。记 $D(K\backslash G/^{/}G)'$ 为相应的
$\mathcal{D}$-导出范畴。复形 $\Delta_{\Omega}(V)$ 落在该子范畴中，故有三角函子
$L\Delta:D(\mathfrak{g},K)'\to D(K\backslash G/^{/}G)'$。这些范畴都是
$\mathcal{H}$-模（后者见 7.6.8(v)、7.6.9）。由 7.8.8，$L\Delta$ 是
$\mathcal{H}$-模的态射。7.7.6 与 7.7.11 的一个变体表明，$L\Delta$ 是
t-范畴的等价。
% [Tutor Note: 原书 p.277 此处作「This categories are $\mathcal{H}$-Modules」，主谓数不一致；底本如此，照译为「这些范畴都是」。]

\bdp{7.8.10} \emph{注.} i) 若 $G'$ 是 $G$ 的平凡扩张，则
$D(\mathfrak{g},K)'=D(\mathfrak{g},K)$，且 $L\Delta$ 与 7.7.2 中函子
$\gamma^{-1}$ 所定义的等价一致。这表明 7.7.4 与 7.8.3 中的 Hecke 作用确实一致。

(ii) 设扩张是任意的。则拉回函子
$r:D(K\backslash G/K)\to D(K'\backslash G'/K')$ 是幺半范畴的态射，

\SourcePage{278}

而全忠实嵌入 $D(\mathfrak{g},K)'\hookrightarrow D(\mathfrak{g}',K')$ 是
$\mathcal{H}$-模的 $r$-态射。于是在 $(\mathfrak{g}',K')$ 的情形，扭的图景本质上
等价于非扭的图景。不过在应用中，方便的做法是把扭（又称水平，又称中心荷）
单独保留。
% [Tutor Note: 原书 p.278 此处作「is $r$-Morphism of Hecke Modules」，英文语法不规范（应为 a morphism）；底本如此，照译为「是 $\mathcal{H}$-模的 $r$-态射」。]

\bdp{7.8.11} 下面解释 7.1.1 中「Hecke 模式」(d) 的 $\Gamma$ 部分。
本文主体不需要这一题材，读者可以跳过本节余下部分。我们处理扭的版本，
故处于 7.8.6 的情形。对 $T\in C(K\backslash Y,\Omega)$，$Y$ 上的
$\mathcal{D}$-复形 $\mathcal{D}T_{Y}$ 是 $K$-等变的（见 7.6.11）。
把 $\mathcal{D}T_{Y}$ 看作带 $(\mathfrak{k}_{\Omega}\times\mathfrak{g},K)$-作用的
$\mathcal{O}$-复形。置
\[
\Gamma_{\mathcal{L}}(T):=\Gamma(Y,\mathcal{D}T_{Y}\otimes(\omega_{Y}\otimes\mathcal{L}_{Y})^{*}).
\]
这是 Harish-Chandra 复形（回忆 $(\mathfrak{g}',K)$ 作用于
$\omega_{Y}\otimes\mathcal{L}_{Y}$），于是有 DG 函子
$\Gamma_{\mathcal{L}}:C(K\backslash Y,\Omega)\to C(\mathfrak{k}_{\Omega}\times\mathfrak{g},K)'$。
设
\[
R\Gamma_{\mathcal{L}}:D(K\backslash Y)\to D(\mathfrak{g},K)'
\]
为其右导出函子。若 $T$ 是松的，则
$\Gamma_{\mathcal{L}}(T)=R\Gamma_{\mathcal{L}}(T)$，故 $R\Gamma_{\mathcal{L}}$
的定义是正当的。

注意 $R\Gamma_{\mathcal{L}}$ 是左 t-正合函子；记
$\Gamma_{\mathcal{L}}:\mathcal{M}(K\backslash Y)\to\mathcal{M}(\mathfrak{g},K)'$
为相应的左正合函子。有
$\Gamma_{\mathcal{L}}(M)=\Gamma(Y,M\otimes(\omega_{Y}\otimes\mathcal{L}_{Y})^{*})$。
若处于 7.8.7(iii) 的情形，则在右 $\mathcal{D}$-模与左 $\mathcal{D}$-模的标准
等同之下，该函子与 1.2.4 的函子 $\Gamma$ 一致。

\bdp{7.8.12} \emph{引理.} 函子 $R\Gamma_{\mathcal{L}}$ 是 $\mathcal{H}$-模的态射。

\emph{证明.} 只需证明 $\Gamma_{\mathcal{L}}$ 是 $\mathcal{H}^{c}$-模的态射，
即对 $F\in\mathcal{H}^{c}$ 与上述 $T$，定义与结合约束相容的典范同构
$\beta:\Gamma_{\mathcal{L}}(F\overset{c}{\circledast}T)\xrightarrow{\sim}
F\overset{c}{\circledast}\Gamma_{\mathcal{L}}(T)$。我们先写出与
$K_{\Omega}$-作用相容的典范同构
$\tilde{\beta}:\Gamma_{\mathcal{L}}(F\widetilde{\circledast}T)
\xrightarrow{\sim}F\widetilde{\circledast}\Gamma_{\mathcal{L}}(T)$；
取 $K_{\Omega}$-不变量即得 $\beta$。

$\mathcal{L}$ 上的 $G'$-作用给出同构
$m_{Y}^{*}((\omega_{Y}\otimes\mathcal{L}_{Y})^{*})
=\mathcal{L}_{G}\otimes(\omega_{Y}\otimes\mathcal{L}_{Y})^{*}$，
而 $\mathcal{D}_{Y}$ 上的 $G$-作用（作为左 $\mathcal{O}_{Y}$-模）给出
% [Tutor Note: 原书 p.278 末行「(as on a left $\mathcal{O}_Y$-module」的右括号缺失；底本如此，译文补足了右括号。]

\SourcePage{279}

同构 $m_{Y}^{*}(\mathcal{D}_{Y})=\mathcal{O}_{G}\boxtimes\mathcal{D}_{Y}$。
这些同构把 $\Gamma_{\mathcal{L}}(F\widetilde{\circledast}T)^{\cdot}$ 与
\[
\Gamma(G\times Y,(F^{\cdot}\otimes\mathcal{L}_{G})\boxtimes
(\mathcal{D}T_{Y}^{\cdot}\otimes(\omega_{Y}\otimes\mathcal{L}_{Y})^{*}))
\]
等同起来。该向量空间等于
$\Gamma(G,F^{\cdot}\otimes\mathcal{L}_{G})\otimes
\Gamma(Y,\mathcal{D}T_{Y}^{\cdot}\otimes(\omega_{Y}\otimes\mathcal{L}_{Y})^{*})$，
即 $(F\widetilde{\circledast}\Gamma_{\mathcal{L}}(T))^{\cdot}$。
我们的 $\tilde{\beta}$ 就是这些等同的复合。$\tilde{\beta}$ 是 Harish-Chandra
复形的同构且与 $K_{\Omega}$-作用相容，这一点留给读者验证。\hfill$\square$

\bdsec{7.9. 额外的对称性与参数}

\bdp{7.9.1} 在本文主体（即 5.4）中，我们用到 7.1.1 的 Hecke 模式的一个
等变版本。确切地说，给定一个额外的 Harish-Chandra 对 $(\mathfrak{l},P)$，
它作用于 $(G,K)$，而我们要找 7.1.1(a)--(d) 的 $(\mathfrak{l},P)$-等变版本。
下面简要说明这一设定；全部细节见本节余下部分。Hecke 范畴 $\mathcal{H}$ 是
$K\backslash G/K$ 上弱 $(\mathfrak{l},P)$-等变 $\mathcal{D}$-模范畴的导出版本。
这是幺半三角范畴（即 7.1.1(a) 在当前设定下的类似物）。$\mathcal{H}$ 作用于
$(\mathfrak{l}\bowtie\mathfrak{g},P\bowtie K)$-模的适当导出范畴 $D_{HC}$；
这是类似于 7.1.1(c) 的 Harish-Chandra 对应物。几何对应物如下。设 $X$ 是带
$(\mathfrak{l},P)$-结构 $X^{\wedge}$ 的「参数」空间（见 2.6.4）。考虑以
$X^{\wedge}$ 为参数、带 $G$-作用的光滑簇族 $Y^{\wedge}$。设 $X^{\wedge}$ 上的
$(\mathfrak{l},P)$-作用提升到 $Y^{\wedge}$，且与 $G$-作用相容。则 $\mathcal{H}$
作用于 $X$-叠 $\mathcal{B}=(P\bowtie K)\backslash Y^{\wedge}$ 的
$\mathcal{D}$-模导出范畴 $D(\mathcal{B})$（这是 7.1.1(b) 的版本）。
有适当的局部化函子 $L\Delta:D_{HC}\to D(\mathcal{B})$，它与 Hecke 作用交换
（这就是 7.1.1(d)）。对带 $(\mathfrak{l},P)$-作用的代数 $A$，上述构造有
$A$-线性版本：考察带 $A$-作用的 Harish-Chandra 模与带 $A_{X}$-作用的
$\mathcal{D}$-模（$A_{X}$ 的定义见 2.6.6）。相应的三角范畴记作
$\mathcal{H}_{A}$、$D_{HCA}$ 与 $D(\mathcal{B},A_{X})$。

\SourcePage{280}

这些构造本质上是前几节构造的直截了当的修改；我们写下它们，是为在 5.4 中
直接引用。

\emph{注.} 等变 Hecke 模式不能归约为把 $G$ 换成对应于 Harish-Chandra 对
$(\mathfrak{l}\bowtie\mathfrak{g},P\bowtie G)$ 的群归纳概形后的普通模式。
事实上，我们的 $\mathcal{H}$ 比相应的「普通」Hecke 范畴大得多：后者由
$K\backslash G/K$ 上强 $P$-等变 $\mathcal{D}$-模构成。特别地，$\mathcal{H}$
以 $(\mathfrak{l},P)$-模的张量范畴为张量子范畴。上述纤维化结构 $Y/X$ 正是
以使整个 $\mathcal{H}$ 作用于 $D(\mathcal{B})$ 所需。
% [Tutor Note: 原书 p.280 此处作「much larger then」，应为「than」；底本如此，照译为「大得多」。]

\bdp{7.9.2} 于是考虑作用于 $(G,K)$ 的 Harish-Chandra 对 $(\mathfrak{l},P)$。
这里 $P$ 可以是任意仿射群概形（不必有限型），但我们假设
$\operatorname{Lie}P$ 在 $\mathfrak{l}$ 中余维有限。考虑 DG 范畴
$\mathcal{H}^{c}$，其对象是 $K\backslash G/K$ 上的 $\Omega$-复形 $F$，
带 $(\mathfrak{l},P)$-作用，且该作用提升 $G/K$ 上的 $(\mathfrak{l},P)$-作用。
这样的 $F$ 等同于 $G$ 上
$(\mathfrak{l},P)\bowtie(K_{\Omega}\times K_{\Omega})$-等变的 $\Omega$-复形。
称 $\mathcal{H}^{c}$ 为 $(\mathfrak{l},P)$-等变预 Hecke 范畴。
$\mathcal{H}^{c}$ 的同伦范畴中那些作为普通 $\Omega$-复形的
$\mathcal{D}$-拟同构的态射构成局部化族。$(\mathfrak{l},P)$-等变 Hecke 范畴
$\mathcal{H}$ 就是相应的局部化。于是 $\mathcal{H}$ 是 t-范畴，其核心等于
$G/K$ 上带弱 $(\mathfrak{l}\bowtie\mathfrak{k},P\bowtie K)$-作用的
$\mathcal{D}$-模范畴（这里 $K$ 以左平移作用于 $G/K$），且 $K$ 的作用实际上是
强的。

现在 $\mathcal{H}^{c}$ 是 DG 幺半范畴，$\mathcal{H}$ 是幺半三角范畴。
事实上，7.6.1 的全部定义在当前情形都适用。

\emph{注.} 取 Harish-Chandra 模 $V\in\mathcal{M}(\mathfrak{l},P)$。
把它对应到 $G/K$ 的指定点处的摩天层，看作集中在 $0$ 次、带平凡
$K_{\Omega}$-作用的 $\Omega$-复形。这是 $\mathcal{H}^{c}$ 的对象。函子
$\mathcal{M}(\mathfrak{l},P)\to\mathcal{H}^{c},\mathcal{H}$ 是全忠实幺半函子。
注意 $\mathcal{M}(\mathfrak{l},P)$ 以典范方式属于（预）Hecke 幺半范畴的中心，
即对任意上述 $V$、

\SourcePage{281}

$F\in\mathcal{H}$，有与 $F$ 的张量积及 $V$ 的张量积都相容的典范同构
$V\circledast{}F\xrightarrow{\sim}F\circledast{}V$。事实上，两个对象都等于
$V\otimes F$。

\bdp{7.9.3} 为定义 $\mathcal{D}$-模上的 Hecke 作用，需要先作些准备。

设 $X$ 是光滑簇，$Y$ 是 $\mathcal{D}_{X}$-概形。$Y$ 上的
$\mathcal{D}_{X}\Omega_{/X}$-复形是带 $\mathcal{D}_{X}$-结构的
DG $\Omega_{Y/X}$-模（$\mathcal{D}_{X}$-结构 $:=$ 沿 $Y/X$ 上结构联络的叶的
平坦联络）。确切地说，$Y$ 上的 $\mathcal{D}_{X}$-结构在
$\Omega_{Y/X}(\mathcal{D}_{X}):=\mathcal{D}_{X}\otimes_{\mathcal{O}_{X}}\Omega_{Y/X}$
上定义了结合 DG 代数结构。于是 $Y$ 上的 $\mathcal{D}_{X}\Omega_{/X}$-复形是
左 DG $\Omega_{Y/X}(\mathcal{D}_{X})$-模，且作为 $\mathcal{O}_{Y}$-模是拟凝聚的。

$Y$ 上 $\mathcal{D}_{X}\Omega_{/X}$-复形的 DG 范畴
$C(Y,\mathcal{D}_{X}\Omega_{/X})$ 是张量范畴（张量积取在 $\Omega_{Y/X}$ 上）。
拉回函子
\[
C(\mathcal{M}^{\ell}(X))\longrightarrow C(Y,\mathcal{D}_{X}\Omega_{/X}),
\qquad
M\mapsto\Omega_{Y/X}\otimes_{\mathcal{O}_{X}}\Omega_{Y/X},
\]
是张量函子。
% [Tutor Note: 原书 p.281 此处的拉回函子印作 $M\mapsto\Omega_{Y/X}\otimes_{\mathcal{O}_X}\Omega_{Y/X}$，第二个因子应为 $M$；底本如此，保留原样。]
特别地 $C(Y,\mathcal{D}_{X}\Omega_{/X})$ 是 $\mathcal{M}^{\ell}(X)$-模
（有 $M\circledast F=M\otimes_{\mathcal{O}_{X}}F$）。

注意，对 $Y$ 上的 $\mathcal{D}_{X}\Omega_{/X}$-复形 $F$，有绝对 $\Omega$-复形
$\Omega_{X}F$，它定义为沿 $X$ 的 de Rham 复形，系数在 $F$ 中
\footnote{如同 7.2，函子 $\Omega_{X}$ 有左伴随函子 $\mathcal{D}_{X}$。}。
于是若 $Y$ 是光滑簇，就有 $\mathcal{D}_{X}\Omega_{/X}$-复形的
$\mathcal{D}$-拟同构概念。$C(Y,\mathcal{D}_{X}\Omega_{/X})$ 的同伦范畴相应的
局部化记作 $D(Y,\mathcal{D}_{X}\Omega_{/X})$。函子
$\Omega_{X}:D(Y,\mathcal{D}_{X}\Omega_{/X})\longrightarrow D(Y,\Omega)$
是三角范畴的等价。

\bdp{7.9.4} 现在设 $X$ 是带 $(\mathfrak{l},P)$-结构 $X^{\wedge}$ 的光滑簇
（见 2.6.4）。设 $Y^{\wedge}$ 是带 $(\mathfrak{l},P)\bowtie G$-作用的概形，
且有与作用相容的光滑态射 $p^{\wedge}:Y^{\wedge}\to X^{\wedge}$（于是 $G$ 沿
纤维作用，$p^{\wedge}$ 与 $(\mathfrak{l},P)$ 的作用交换）。置
$Y:=P\backslash Y^{\wedge}$。这是带光滑投影 $p:Y\to X$ 的光滑簇。
$Y^{\wedge}$ 上的 $(\mathfrak{l},P)$-作用在 $Y$ 上定义了 $\mathcal{D}_{X}$-概形
结构。$Y^{\wedge}$ 上的 $G$-作用在 $Y$ 上给出水平的 $G_{X}$-作用（群
$\mathcal{D}_{X}$-概形 $G_{X}$ 的定义见 2.6.6）。
% [Tutor Note: 原书 p.281 此处 "horizontal" 误拼作「horisontal」（p.283 同）；底本如此，照译为「水平的」。]

\SourcePage{282}

考虑 $X$ 上纤维化的叠 $\mathcal{B}:=K_{X}\backslash Y=(P\bowtie K)\backslash Y^{\wedge}$，
于是有相应的左 $\mathcal{D}$-模范畴 $\mathcal{M}^{\ell}(\mathcal{B})$ 与
$\mathcal{B}$ 上 $\Omega$-复形的 t-范畴 $D(\mathcal{B})$。该 t-范畴还有另一种
用 $\mathcal{D}_{X}\Omega_{/X}$-复形的实现，下面描述。

考虑 DG 群 $\mathcal{D}_{X}$-概形
$G_{\Omega X}:=(G_{X},\Omega_{G_{X}/X})$、$K_{\Omega X}$。
$Y$ 上 $\mathcal{D}_{X}\Omega_{/X}$-复形的 $K_{\Omega X}$-作用按 7.6.4 定义。
现在有 $Y$ 上 $K_{\Omega X}$-等变 $\mathcal{D}_{X}\Omega_{/X}$-复形的 DG 范畴
$C(K_{X}\backslash Y,\mathcal{D}_{X}\Omega_{/X})$。把它的同伦范畴按
$\mathcal{D}$-拟同构局部化，得到三角范畴
$D(K_{X}\backslash Y,\mathcal{D}_{X}\Omega_{/X})$。de Rham 函子 $\Omega_{X}$
把它与 $D(\mathcal{B})$ 等同。

现在可以定义 $D(\mathcal{B})$ 上的 Hecke 作用。先构造 $\mathcal{H}^{c}$ 在
$C(K_{X}\backslash Y,\mathcal{D}_{X}\Omega_{/X})$ 上的作用
$\overset{c}{\circledast}$。事实上，对 $F\in\mathcal{H}^{c}$，有 $G_{X}$ 上的
$\mathcal{D}_{X}\Omega_{/X}$-复形 $F_{X}$，它关于左、右平移是
$K_{\Omega X}$-等变的。于是对
$T\in C(K_{X}\backslash Y,\mathcal{D}_{X}\Omega_{/X})$，在 $\mathcal{D}_{X}$-概形
$G_{X}\times Y$（$G_{X}$ 与 $Y$ 在 $X$ 上的纤维积）上有
$\mathcal{D}_{X}\Omega_{/X}$-复形 $F\boxtimes T$。它关于 $G_{X}\times Y$ 上所有
$K_{X}$-作用是 $K_{\Omega X}$-等变的。故 $F\boxtimes T$ 下降到
$G_{X}\times_{K_{X}}Y$。定义
$F\overset{c}{\circledast}T\in C(K_{X}\backslash Y,\mathcal{D}_{X}\Omega_{/X})$
为上述复形沿作用映射 $G_{X}\times_{K_{X}}Y\to Y$ 的前推。Hecke 作用
$\circledast:H\times D(\mathcal{B})\longrightarrow D(\mathcal{B})$ 是
$\overset{c}{\circledast}$ 的右导出函子；按惯例，可以用松的
$\mathcal{D}_{X}\Omega_{/X}$-复形来计算它。

\emph{注.} 对 $W\in\mathcal{M}(\mathfrak{l},P)\subset\mathcal{H}^{c}$ 与上述
$T$，有 $W\overset{c}{\circledast}T=W\circledast{}T=W_{X}\otimes T$
（$\mathcal{D}_{X}$-模 $W_{X}$ 的定义见 2.6.6）。

\bdp{7.9.5} 下面定义 Harish-Chandra 范畴。设 $G'$ 如 7.8.1，并设给定
$G$ 上 $(\mathfrak{l},P)$-作用到 $G'$ 上的提升，它保持 $K\subset G'$ 且固定
$\mathbb{G}_{m}\subset G'$。于是有 Harish-Chandra 对
$(\mathfrak{l},P)\bowtie(\mathfrak{g}',K')$。设 $C_{HC}$ 为
$(\mathfrak{l},P)\bowtie(\mathfrak{k}_{\Omega}\times\mathfrak{g},K)'$-复形的
范畴，即带相容 $(\mathfrak{l},P)$-作用的
$(\mathfrak{k}_{\Omega}\times\mathfrak{g},K)'$-复形（记号见 7.8.1）。设
$D_{HC}$ 为相应的导出范畴。这是 t-范畴，核心为
$\mathcal{M}_{HC}=\mathcal{M}(\mathfrak{l}\bowtie\mathfrak{g},P\bowtie K)'$。
下面

\SourcePage{283}

我们把 $C_{HC}$、$D_{HC}$ 的对象简称为 Harish-Chandra 复形，把
$\mathcal{M}_{HC}$ 的对象简称为 Harish-Chandra 模。

预 Hecke 范畴 $\mathcal{H}^{c}$ 作用于 $C_{HC}$。事实上，7.8.2 的构造在
当前情形完全说得通（$(\mathfrak{l},P)$ 通过结构搬运作用于
$F\overset{c}{\circledast}V$）。$D_{HC}$ 上的 $\mathcal{H}$-作用
$\circledast$ 是 $\overset{c}{\circledast}$ 的右导出函子。7.8.4--7.8.5 的
结果原样移植到当前设定。

\emph{注.} 对 $W\in\mathcal{M}(\mathfrak{l},P)\subset\mathcal{H}^{c}$ 与
Harish-Chandra 复形 $V$，有 Harish-Chandra 复形的典范同构
$W\overset{c}{\circledast}V=W\circledast{}V=W\otimes V$。

\bdp{7.9.6} 下面转向局部化函子。7.8.6 的构造移植到我们的设定如下。
从 7.9.4 的 $Y^{\wedge}$ 出发。设它带线丛 $\mathcal{L}_{Y^{\wedge}}$，且
$Y^{\wedge}$ 上的 $(\mathfrak{l},P)\bowtie G$-作用提升为
$\mathcal{L}_{Y^{\wedge}}$ 上的 $(\mathfrak{l},P)\bowtie G'$-作用，使
$\mathbb{G}_{m}\subset G'$ 按与标准特征标相反的特征标作用。设
$\mathcal{L}_{Y}$ 是 $\mathcal{L}_{Y^{\wedge}}$ 沿 $P$-作用下降到 $Y$ 的线丛。
该线丛带典范 $\mathcal{D}_{X}$-结构，它来自 $\mathcal{L}_{Y^{\wedge}}$ 上的
$\mathfrak{l}$-作用。它还带水平的 $G'_{X}$-作用。

有 DG 函子
\begin{equation}\tag{338}
\Delta_{\Omega}=\Delta_{\Omega\mathcal{L}}:C_{HC}\longrightarrow
C(K_{X}\backslash Y,\mathcal{D}_{X}\Omega_{/X}),
\end{equation}
$\Delta_{\Omega}(V)=\operatorname{Hom}_{\mathfrak{g}_{X}}
(DR_{\mathfrak{g}_{X}},\omega_{Y/X}\times\mathcal{L}_{Y}\times V)[-\dim K]$
（对照（336））。
% [Tutor Note: 原书 p.283 此处印作 $\omega_{Y/X}\times\mathcal{L}_Y\times V$（用 $\times$），而 7.8.6 的（336）式用 $\otimes$；底本如此，保留原样。]
如 7.8.6，该函子把拟同构送到 $\mathcal{D}$-拟同构，故给出三角函子
\begin{equation}\tag{339}
L\Delta=L\Delta_{\mathcal{L}}:D_{HC}\longrightarrow D(\mathcal{B}),
\end{equation}
它是右 t-正合的。两个核心之间的相应右正合函子
$\Delta_{\mathcal{L}}:\mathcal{M}_{HC}\longrightarrow\mathcal{M}^{\ell}(\mathcal{B})$
把 $V$ 送到 $K_{X}$-等变左 $\mathcal{D}_{Y}$-模
\[
(\mathcal{D}_{Y/X}\otimes\mathcal{L}_{Y})\otimes_{U(\mathfrak{g}'_{X})}V_{X}.
\]

函子 $\Delta_{\Omega}$、$L\Delta$ 与 Hecke 作用交换。事实上，7.8.8 的证明
逐字照搬到我们的设定。特别地，对任意 $W\in\mathcal{M}(\mathfrak{l},P)$、
$V\in D_{HC}$，有 $L\Delta(W\otimes V)=W_{X}\otimes L\Delta(V)$。

\SourcePage{284}

\bdp{7.9.7} \emph{$A$-线性版本.} 设此外还给定带 $(\mathfrak{l},P)$-作用的
交换代数 $A$。把它接到上述模式上如下。

(i) 记 $\mathcal{H}^{c}_{A}$ 为满足如下条件的对象 $F\in\mathcal{H}^{c}$ 的
DG 范畴：$F$ 带 $A$-作用，$A$ 与 $(\mathfrak{l},P)$ 的作用相容，且 $F$ 是
$A$-平坦的。设 $\mathcal{H}_{A}$ 为相应的 $\mathcal{D}$-导出范畴。卷积积按
7.9.2 定义（张量积取在 $A$ 上），故 $\mathcal{H}^{c}_{A}$ 与
$\mathcal{H}_{A}$ 是幺半范畴。设 $\mathcal{M}(\mathfrak{l},P)^{fl}_{A}$ 为
带 $(\mathfrak{l},P)$-作用的平坦 $A$-模的张量范畴。如 7.9.2 的注，有典范
全忠实幺半函子
$\mathcal{M}(\mathfrak{l},P)^{fl}_{A}\longrightarrow\mathcal{H}^{c}_{A},\mathcal{H}_{A}$，
它们把 $\mathcal{M}(\mathfrak{l},P)^{fl}_{A}$ 送到 Hecke 范畴的中心。

(ii) 设处于 7.9.4 的情形。考虑 $\mathcal{B}$ 上带 $A_{X}$-作用的左
$\mathcal{D}$-模范畴 $\mathcal{M}^{\ell}(\mathcal{B},A_{X})$（$\mathcal{D}_{X}$-代数
$A_{X}$ 的定义见 2.6.6）。设 $C(\mathcal{B},A_{X}\otimes\Omega)$ 为
$\mathcal{B}$ 上带 $A_{X}$-作用的 $\Omega$-复形的 DG 范畴，
$D(\mathcal{B},A_{X})$ 为相应同伦范畴关于 $\mathcal{D}$-拟同构的局部化。
这是 t-范畴，核心为 $\mathcal{M}^{\ell}(\mathcal{B},A_{X})$。如 7.9.4，也可以
用 $\mathcal{D}_{X}\Omega_{/X}$-复形来定义该 t-范畴。确切地说，设
$C(K_{X}\backslash Y,A_{X}\mathcal{D}_{X}\Omega_{/X})$ 为
$C(K_{X}\backslash Y,\mathcal{D}_{X}\Omega_{/X})$ 中带 $A_{X}$-作用（与
$K_{\Omega X}$-作用交换）的对象构成的 DG 范畴。按 $\mathcal{D}$-拟同构
局部化，得到三角范畴
$D(K_{X}\backslash Y,A_{X}\mathcal{D}_{X}\Omega_{/X})$。de Rham 函子
$\Omega_{X}$ 把它与 $D(\mathcal{B},A_{X})$ 等同。

$A$-线性设定下的 Hecke 作用完全按 7.9.4 定义。7.9.4 的注仍然成立
（张量积取在 $A_{X}$ 上）。

(iii) 设处于 7.9.5 的情形。定义 $C_{HCA}$ 为带相容 $A$-作用的
Harish-Chandra 复形范畴（于是 $A$ 与
$(\mathfrak{k}_{\Omega}\times\mathfrak{g},K)'$ 的作用交换）。设 $D_{HCA}$ 为
相应的导出范畴。这是 t-范畴，核心 $\mathcal{M}_{HCA}$ 等于带相容 $A$-作用的
$(\mathfrak{l}\bowtie\mathfrak{g},P\bowtie K)'$-模范畴。所有

\SourcePage{285}

关于 Hecke 作用的构造与结果都原样成立。在 7.9.5 的注中取
$W\in\mathcal{M}(\mathfrak{l},P)^{fl}_{A}$；张量积 $W\otimes V$ 取在 $A$ 上。
局部化函子的 $A$-线性设定无需改动。

\emph{注.} 有显然的函子（用 $A$ 作张量），它们把上述普通范畴送到附上 $A$
的范畴。这些函子与我们考虑的所有结构相容。忘掉 $A$-作用的函子
$D(\mathcal{B},A_{X})\to D(\mathcal{B})$、$D_{HCA}\to D_{HC}$ 是
$\mathcal{H}$-模的态射。它们与局部化函子交换。

\bdp{7.9.8} \emph{变体.} 设除 $A$ 之外还给定与 $(\mathfrak{l},P)$-作用相容的
交换代数态射 $e:\mathfrak{z}\to A$。这里
$\mathfrak{z}:=U(\mathfrak{g})^{\prime\operatorname{Ad}G}$（故若 $G$ 连通，
则 $\mathfrak{z}$ 是 $U(\mathfrak{g})'$ 的中心）。于是 $\mathfrak{z}$ 以两种
方式作用于 $\mathcal{M}_{HCA}$ 或 $C_{HCA}$ 的任一对象。记
$\mathcal{M}^{e}_{HCA}$、$C^{e}_{HCA}$ 为 $\mathfrak{z}$ 的两种作用一致的
对象构成的范畴；设 $D^{e}_{HC}$ 为相应的导出范畴。$\mathcal{H}^{c}_{A}$ 在
$C_{HCA}$ 上的作用是 $\mathfrak{z}$-线性的（见 7.8.4），故它保持
$C^{e}_{HCA}$。于是有 $\mathcal{H}_{A}$ 在 $D^{e}_{HCA}$ 上的作用。显然函子
$D^{e}_{HCA}\to D_{HCA}$ 是 $\mathcal{H}_{A}$-模的态射。

\emph{注.} 若 $e$ 是满射，则 $\mathcal{M}^{e}_{HCA}$ 是 $\mathcal{M}_{HC}$ 中
由被 $\operatorname{Ker}e$ 零化的 Harish-Chandra 模构成的满子范畴。
$C^{e}_{HCA}$ 亦然。

\bdsec{7.10. $\mathcal{D}$-晶体}

下文概述 $\mathcal{D}$-模理论的晶体方法。与通常的形式体系不同，它不区分
光滑与不光滑的概形。

本节中，「概形」指「局部有限型的 $\mathbb{C}$-概形」。代数空间与叠亦然。
形式概形或代数空间假设为局部归纳有限型的
\footnote{即：任一闭子概形都是有限型的。}。

\bdp{7.10.1} 设 $f:Y\to X$ 是概形的拟有限态射。则 Grothendieck 的函子
$Rf^{!}:D^{b}(X,\mathcal{O})\to D^{b}(Y,\mathcal{O})$ 左 t-正合。置

% 尾部接 p.286：$f^{!}:=H^{0}Rf^{!}:\mathcal{M}(X,\mathcal{O})\to\mathcal{M}(Y,\mathcal{O})$，……
%% end units/ch07_u05.tex
%% begin units/ch07_u06.tex
% ch07_u06 —— 原书 p.286--p.295（7.10.1 末段至 7.10.20 证明前半）
% 承接 p.285 末尾 7.10.1 的断句「Set」，本单元自 $f^{!}:=H^{0}Rf^{!}$ 续译；
% 尾部停在 p.295 末行，7.10.20 的证明接 p.296。

\SourcePage{286}

$f^{!}:=H^{0}Rf^{!}:\mathcal{M}(X,\mathcal{O})\to\mathcal{M}(Y,\mathcal{O})$；这是左正合函子。因此，诸范畴 $\mathcal{M}(X,\mathcal{O})$ 连同函子 $f^{!}$，构成概形与拟有限态射的范畴上的纤维范畴。

下面给出 $f^{!}$ 的显式描述。由 Zariski 主定理，任意拟有限态射都是有限态射与开浸入的复合；我们在这两种情形下分别描述 $f^{!}$。若 $f$ 是开浸入（或更一般地，$f$ 平展），则 $f^{!}=f^{*}$。若 $f$ 有限，则 $f^{!}$ 是函子 $f_{*}:\mathcal{M}(Y,\mathcal{O})\to\mathcal{M}(X,\mathcal{O})$ 的右伴随。确切地说，$f_{*}\mathcal{O}_{Y}$ 是有限 $\mathcal{O}_{X}$-代数；函子 $f_{*}$ 把 $\mathcal{M}(Y,\mathcal{O})$ 与这样的范畴等同：其对象是作为 $\mathcal{O}_{X}$-模为拟凝聚的 $f_{*}\mathcal{O}_{Y}$-模。现在设 $M$ 是 $X$ 上的 $\mathcal{O}$-模，则相应的 $f_{*}\mathcal{O}_{Y}$-模 $f_{*}f^{!}M$ 是 $\mathcal{H}om_{\mathcal{O}_{X}}(f_{*}\mathcal{O}_{Y},M)$。特别地，若 $f$ 是闭浸入，则 $f^{!}M\subset M$ 是（概形意义下）支撑于 $Y$ 上的截面构成的子模。

上述图景按如下方式推广到形式概形（或代数空间）的设定。对形式概形 $\hat X$，记 $\mathcal{M}(\hat X,\mathcal{O})$ 为离散拟凝聚 $\mathcal{O}_{\hat X}$-模的范畴\footnote{该范畴是阿贝尔范畴。更一般的设定见 7.11.4。}。例如，设 $\hat X$ 是概形 $V$ 沿其闭子概形 $X$ 的形式完备化，则 $\mathcal{M}(\hat X,\mathcal{O})$ 与 $V$ 上集合论意义下支撑于 $X$ 的 $\mathcal{O}$-模范畴相同。若 $\hat X$ 仿射，则对任意 $M\in\mathcal{M}(\hat X,\mathcal{O})$ 有 $M=\bigcup M_{X'}$；这里 $X'$ 取遍 $\hat X$ 的闭子概形构成的（有向）集合，而 $M_{X'}\in\mathcal{M}(X',\mathcal{O})$ 是概形意义下支撑于 $X'$ 的截面构成子模。拉回函子 $f^{!}$ 以唯一的方式\footnote{我们假设它们与 $f$ 的复合相容。}推广到形式代数空间的拟有限态射的设定。事实上，设 $\hat f:\hat Y\to\hat X$ 是这样的态射。为定义 $\hat f^{!}:\mathcal{M}(\hat X,\mathcal{O})\to\mathcal{M}(\hat Y,\mathcal{O})$，可设 $\hat X$、$\hat Y$ 仿射。此时 $\hat f^{!}M=\bigcup\hat f^{!}|_{Y'}M_{X'}$，其中 $Y'$ 是 $\hat Y$ 的闭子概形，且 $\hat f(Y')\subset X'$。若 $\hat f$ 是 ind-有限\footnote{$:=Y_{\mathrm{red}}\to X_{\mathrm{red}}$ 是有限态射。}，则 $\hat f^{!}$ 的显式描述留给读者。

\SourcePage{287}

\bdp{7.10.2} 对概形或代数空间 $X$，记 $X_{cr}$ 为如下图表构成的范畴：$X\overset{j}{\longleftarrow}S\overset{i}{\hookrightarrow}\hat S$，其中 $j$ 是拟有限态射，$i$ 是仿射概形的闭浸入，且相应的理想 $\mathcal{I}\subset\mathcal{O}_{\hat S}$ 幂零。我们通常把 $X_{cr}$ 的这个对象写作 $(S,\hat S)$，或简写作 $\hat S$。$X_{cr}$ 中的态射 $(S,\hat S)\to(S',\hat S')$ 是概形态射 $\phi:\hat S\to\hat S'$，使得 $\phi(S)\subset S'$，且 $\phi|_{S}:S\to S'$ 是 $X$-概形的态射。

注意，对如上任意 $\phi$，态射 $\phi:\hat S\to\hat S'$ 都拟有限。因此，诸范畴 $\mathcal{M}(\hat S,\mathcal{O})$ 连同拉回函子 $\phi^{!}$，构成 $X_{cr}$ 上的纤维范畴 $\mathcal{M}^{!}(X_{cr},\mathcal{O})$。

有时考虑更大的范畴 $X_{\widehat{cr}}$ 较为方便：它由与上述类似的图表构成，但允许 $\hat S$ 是形式概形（于是 $\mathcal{I}$ 是投射幂零理想，即 $\hat S_{\mathrm{red}}=S_{\mathrm{red}}$）。与上面一样，有 $X_{\widehat{cr}}$ 上的纤维范畴 $\mathcal{M}^{!}(X_{\widehat{cr}},\mathcal{O})$。

\bdp{7.10.3} \emph{定义.} $X$ 上的 $\mathcal{D}$-晶体是 $\mathcal{M}^{!}(X_{cr},\mathcal{O})$ 的笛卡儿截面。$X$ 上的 $\mathcal{D}$-晶体构成 $\mathbb{C}$-范畴 $\mathcal{M}_{\mathcal{D}}(X)$。

确切地说，$\mathcal{D}$-晶体 $M$ 是这样一条规则：对任意 $(S,\hat S)\in X_{cr}$，指定 $\hat S$ 上的 $\mathcal{O}$-模 $M_{\hat S}=M_{(S,\hat S)}$；对态射 $\phi:(S,\hat S)\to(S',\hat S')$，指定等同 $\alpha_{\phi}:M_{\hat S}\xrightarrow{\sim}\phi^{!}M_{\hat S'}$，且这些等同与 $\phi$ 的复合相容。

特别地，若 $\phi$ 是由理想 $\mathcal{I}\subset\mathcal{O}_{\hat S'}$ 定义的闭浸入，则 $M_{\hat S}$ 是 $M_{\hat S'}$ 中被 $\mathcal{I}$ 零化的截面构成子模。

在定义 7.10.3 中可以把 $X_{cr}$ 换成 $X_{\widehat{cr}}$：得到的 $\mathcal{D}$-晶体范畴相同。事实上，对 $(S,\hat S)\in X_{\widehat{cr}}$ 有 $M_{\hat S}=\bigcup M_{(S,\hat S')}$，其中 $\hat S'$ 取遍所有满足 $S\subset\hat S'\subset\hat S$ 的子概形。

\bdp{7.10.4} \emph{变体.} 设 $X_{cr}^{(i)},\ldots,X_{cr}^{(iv)}$ 是 $X_{cr}$ 的完全子范畴，其对象 $(S,\hat S)$ 分别满足下列条件之一（在 (ii)--(iv) 中设 $X$ 是概形）：
\begin{enumerate}[label=(\roman*),leftmargin=2.2em]
\item $S\to X$ 平展。
\item $S\to X$ 是开浸入。
\item （设 $X$ 仿射）$S\xrightarrow{\sim}X$。
\SourcePage{288}
\item $S\to X$ 是局部闭浸入。
\end{enumerate}

记 $\mathcal{M}^{(i)}_{\mathcal{D}}(X),\ldots,\mathcal{M}^{(iv)}_{\mathcal{D}}(X)$ 为 $\mathcal{M}^{!}(X_{cr},\mathcal{O})$ 在相应子范畴 $X_{cr}^{(a)}$ 上的笛卡儿截面构成的范畴。有显然的限制函子 $\mathcal{M}_{\mathcal{D}}(X)\to\mathcal{M}^{(a)}_{\mathcal{D}}(X)$。这些函子是范畴等价，其验证留给读者\footnote{只需注意：若把 $\mathcal{M}_{\mathcal{D}}(X)$ 换成 $\mathcal{M}^{(a)}_{\mathcal{D}}(X)$，7.10.6、7.10.7、7.10.8 连同其证明逐字仍然成立。}。

\emph{注.} 范畴 $X_{cr}^{(ii)}$ 就是 $X$ 的 Grothendieck 晶体 site（的底范畴），故按 Grothendieck 的术语，$\mathcal{D}$-晶体就是纤维范畴 $\mathcal{M}^{!}(X_{cr}^{(ii)},\mathcal{O})$ 的晶体。我们把 $X_{cr}$ 当作基本设定，因为它直接推广到归纳概形的设定（见 7.11.6）。

\bdp{7.10.5} 设 $f:Y\to X$ 是拟有限态射。它给出忠实函子 $Y_{cr}\to X_{cr}$，把 $Y\overset{j}{\longleftarrow}S\overset{i}{\hookrightarrow}\hat S$ 送到 $Y\overset{fj}{\longleftarrow}S\overset{i}{\hookrightarrow}\hat S$。%
% [Tutor Note: 原书 p.288 此处第二个图表仍自 $Y$ 起头，按定义（进而给出 $X_{cr}$ 的对象）应自 $X$ 起头；底本如此，保留原样。]
于是得到相应的「限制」函子 $f^{!}:\mathcal{M}_{\mathcal{D}}(X)\to\mathcal{M}_{\mathcal{D}}(Y)$。它与 $f$ 的复合相容。

特别地，当 $U$ 在 $X$ 上平展时，诸范畴 $\mathcal{M}_{\mathcal{D}}(U)$ 构成小平展 site $X_{\text{ét}}$ 上的纤维范畴，记作 $\mathcal{M}_{\mathcal{D}}(X_{\text{ét}})$。

\bdp{7.10.6} \emph{引理.} $\mathcal{D}$-晶体是平展拓扑下的局部对象，即 $\mathcal{M}_{\mathcal{D}}(X_{\text{ét}})$ 是范畴的层。\hfill$\Box$

\bdp{7.10.7} 下面给出 $\mathcal{D}$-晶体的一个方便的「具体」描述。

设有闭浸入 $X\hookrightarrow V$，其中 $V$ 是形式光滑\footnote{见 7.11.1。}的形式代数空间，且 $X_{\mathrm{red}}=V_{\mathrm{red}}$\footnote{即 $X$ 在 $\mathcal{O}_{V}$ 中的理想投射幂零。}。若 $X$ 仿射，这样的对象总存在：可把 $X$ 嵌入光滑概形 $W$，再取 $V$ 为 $W$ 沿 $X$ 的形式完备化。

对 $n\geq1$，记 $V^{<n>}$ 为 $V^{n}$ 沿对角线 $V\subset V^{n}$（等价地，沿 $X\subset V^{n}$）的形式完备化。投影 $p_{1},p_{2}$：\SourcePage{289}$V^{<2>}\to V$，$p_{12},p_{23},p_{13}:V^{<3>}\to V^{<2>}$ 都是 ind-有限，于是有函子 $p^{!}_{i}:\mathcal{M}(V,\mathcal{O})\to\mathcal{M}(V^{<2>},\mathcal{O})$、$p^{!}_{ij}:\mathcal{M}(V^{<2>},\mathcal{O})\to\mathcal{M}(V^{<3>},\mathcal{O})$。由于 $V$ 形式光滑，这些函子正合。

记 $\mathcal{M}_{\mathcal{D}V}(X)$ 为满足下述条件的对 $(M_{V},\tau)$ 构成的范畴：$M_{V}\in\mathcal{M}(V,\mathcal{O})$，而 $\tau:p^{!}_{1}M_{V}\xrightarrow{\sim}p^{!}_{2}M_{V}$ 是同构，且
\begin{equation}\tag{340}
p^{!}_{23}(\tau)p^{!}_{12}(\tau)=p^{!}_{13}(\tau).
\end{equation}

\bdp{7.10.8} \emph{命题.} 范畴 $\mathcal{M}_{\mathcal{D}}(X)$ 与 $\mathcal{M}_{\mathcal{D}V}(X)$ 典范等价。

\emph{证明.} 所论对象是局部的，故可设 $X$ 仿射。对 $M\in\mathcal{M}_{\mathcal{D}}(X)$，有 $M_{V}=M_{(X,V)}\in\mathcal{M}(V,\mathcal{O})$。由于 $p^{!}_{i}M_{V}=M_{V^{<2>}}$，所得的 $\tau$ 显然满足 (340)。反之，设给定 $(M_{V},\tau)\in\mathcal{M}_{\mathcal{D}V}(X)$；我们来定义相应的 $\mathcal{D}$-晶体 $M$。对 $(S,\hat S)\in X_{cr}$，选取延拓结构态射 $j:S\to X$ 的 $j':\hat S\to V$；这样的 $j'$ 存在，因为 $V$ 形式光滑。考虑 $\mathcal{O}_{\hat S}$-模 $j'^{!}M_{V}$。若 $j'':\hat S\to V$ 是 $j$ 的另一延拓，则有典范同构 $\nu_{j'j''}:j'^{!}M_{V}\xrightarrow{\sim}j''^{!}M_{V}$。确切地说，$(j',j'')$ 把 $\hat S$ 映入 $V^{<2>}$，故 $j'^{!}M_{V}=(j',j'')^{!}p^{!}_{1}M_{V}$；再用 $j''^{!}M_{V}$ 的类似描述，令 $\nu_{j'j''}:=(j',j'')^{!}(\tau)$。由 (340)，这些等同可传递，故 $j'^{!}M_{V}$ 不依赖于 $j'$ 的选取。这就是 $M_{(S,\hat S)}$。$M$ 的结构同构 $\alpha_{\phi}$ 的定义是显然的。\hfill$\Box$

\bdp{7.10.9} \emph{推论.} (i) 对任意 $X$，范畴 $\mathcal{M}_{\mathcal{D}}(X)$ 是阿贝尔范畴。

(ii) 对 $\hat S\in X_{cr}$，函子 $\mathcal{M}_{\mathcal{D}}(X)\to\mathcal{M}(\hat S,\mathcal{O})$，$M\mapsto M_{\hat S}$ 左正合。

(iii) 对拟有限态射 $j:Y\to X$，函子 $j^{!}:\mathcal{M}_{\mathcal{D}}(X)\to\mathcal{M}_{\mathcal{D}}(Y)$ 左正合；若 $j$ 平展，则 $j^{!}$ 正合。

\emph{证明.} 命题 (i) 在 $X$ 仿射时成立。事实上，按 7.10.7 选取 $X\hookrightarrow V$。范畴 $\mathcal{M}_{\mathcal{D}V}(X)$ 是阿贝尔范畴，因为函子 $p^{!}_{i}$、$p^{!}_{ij}$ 正合，故由 7.10.8 得证。

\SourcePage{290}

若 $j:U\to X$ 是仿射概形的平展态射，则函子 $j^{!}:\mathcal{M}_{\mathcal{D}}(X)\to\mathcal{M}_{\mathcal{D}}(U)$ 正合。事实上，设 $U\hookrightarrow V_{U}$ 是 $X\hookrightarrow V$ 的 $U$-局部化（故 $V_{U}$ 在 $V$ 上平展）；则 $j^{!}$ 与平展局部化函子 $\mathcal{M}_{\mathcal{D}V}(X)\to\mathcal{M}_{DV_{U}}(U)$ 一致，而后者显然正合。

于是 (i) 由 7.10.6 推出。其余留给读者。\hfill$\Box$

\bdp{7.10.10} \emph{引理.} 对平展态射 $p:U\to X$，函子 $p^{!}$ 有右伴随函子 $p_{*}:\mathcal{M}_{\mathcal{D}}(U)\to\mathcal{M}_{\mathcal{D}}(X)$。若 $p$ 是开浸入，则 $p^{!}p_{*}$ 是恒等函子。

\emph{证明.} 下面给出 $p_{*}$ 的显式构造。对 $(S,\hat S)\in X_{cr}$，令 $S_{U}:=S\times_{X}U$；设 $\hat p_{S}:\hat S_{U}\to\hat S$ 是平展态射，它到 $S\hookrightarrow\hat S$ 的拉回是投影 $S_{U}\to S$。于是 $(S_{U},\hat S_{U})\in U_{cr}$，并有函子 $X_{cr}\to U_{cr}$，$(S,\hat S)\mapsto(S_{U},\hat S_{U})$。

现在对 $N\in\mathcal{M}_{\mathcal{D}}(U)$ 令 $(p_{*}N)_{\hat S}:=(\hat p_{S})_{\cdot}N_{\hat S_{U}}$。等同 $\alpha_{\phi}$ 来自基变换同构 $\phi^{!}\hat p_{S'}{}_{\cdot}=\hat p_{S}{}_{\cdot}\phi^{!}_{U}$。\hfill$\Box$

现在设 $i:Y\hookrightarrow X$ 是闭浸入，$j:U:=X\setminus Y\hookrightarrow X$ 是补开浸入。记 $\mathcal{M}_{\mathcal{D}}(X)_{Y}$ 为 $\mathcal{M}_{\mathcal{D}}(X)$ 的完全子范畴，由满足 $j^{!}M=0$ 的 $\mathcal{D}$-晶体 $M$ 构成。

\bdp{7.10.11} \emph{引理.} (i) 函子 $i^{!}$ 有左伴随函子 $i_{*}:\mathcal{M}_{\mathcal{D}}(Y)\to\mathcal{M}_{\mathcal{D}}(X)$。

(ii) $i_{*}$ 把 $\mathcal{M}_{\mathcal{D}}(Y)$ 映入 $\mathcal{M}_{\mathcal{D}}(X)_{Y}$，且
\[
i_{*}:\mathcal{M}_{\mathcal{D}}(Y)\to\mathcal{M}_{\mathcal{D}}(X)_{Y},\qquad i^{!}:\mathcal{M}_{\mathcal{D}}(X)_{Y}\to\mathcal{M}_{\mathcal{D}}(Y)
\]
是互逆的范畴等价。

(iii) 设 $p:Z\to X$ 是拟有限态射；令 $Y_{Z}:=Y\times_{Z}X$，于是有 $i_{Z}:Y_{Z}\hookrightarrow Z$ 与 $p_{Y}:Y_{Z}\to Y$。则有函子的典范等同 $p^{!}i_{*}=i_{Y*}p^{!}_{Y}:\mathcal{M}_{\mathcal{D}}(Y)\to\mathcal{M}_{\mathcal{D}}(Z)$。

\emph{证明.} 下面给出 $i_{*}$ 的显式构造。取 $Y$ 上的 $\mathcal{D}$-晶体 $N$。对 $(S,\hat S)\in X_{cr}$ 令 $S_{Y}:=S\times_{X}Y$，于是 $S_{Y}$ 是 $S$ 的闭子概形，\SourcePage{291}从而是 $\hat S$ 的闭子概形。投影 $S_{Y}\to Y$ 拟有限，故 $N$ 给出 $S_{Y}$ 上的 $\mathcal{D}$-晶体。把 $(i_{*}N)_{(S,\hat S)}$ 定义为 $\hat S$ 上相应的 $\mathcal{O}$-模（见 7.10.3）。$i_{*}N$ 的结构同构 $\alpha_{\phi}$ 按显然的方式来自 $N$ 的相应同构。

$i_{*}$ 的伴随性质以及性质 (ii)、(iii) 都是显然的。\hfill$\Box$

\bdp{7.10.12} \emph{命题.} 若 $X$ 光滑，则 $\mathcal{M}_{\mathcal{D}}(X)$ 典范等价于 $X$ 上 $\mathcal{D}$-模的范畴 $\mathcal{M}(X)$。

\emph{证明.} 对 $V=X$ 使用 $\mathcal{M}_{\mathcal{D}}(X)$ 的描述 7.10.7。于是 $\mathcal{D}$-晶体 $M$ 等同于对 $(M_{X},\tau)$，其中 $M_{X}\in\mathcal{M}(X,\mathcal{O})$，而 $\tau:p^{!}_{1}M_{X}\xrightarrow{\sim}p^{!}_{2}M_{X}$ 是 $X^{<2>}$ 上 $\mathcal{O}$-模的同构，且满足 (340)。我们来说明，这样的 $\tau$ 与 $M_{X}$ 上的右 $\mathcal{D}$-模结构是同一回事。

把 $\mathcal{D}_{X}$ 看作 $\mathcal{M}(X^{<2>},\mathcal{O})$ 的对象（经由 $\mathcal{O}_{X}$-双模结构）。有典范同构 $\mathcal{D}_{X}\xrightarrow{\sim}p^{!}_{1}\mathcal{O}_{X}$，它把 $\partial\in\mathcal{D}_{X}$ 与截面 $(f\otimes g\mapsto f\partial(g))\in\mathcal{H}om_{\mathcal{O}_{X}}(\mathcal{O}_{X^{<2>}},\mathcal{O}_{X})=p^{!}_{1}\mathcal{O}_{X}$ 对应。因此 $M_{X}\otimes_{\mathcal{O}_{X}}\mathcal{D}_{X}\xrightarrow{\sim}M_{X}\otimes_{\mathcal{O}_{X}}p^{!}_{1}\mathcal{O}_{X}\xrightarrow{\sim}p^{!}_{1}M_{X}$。于是由伴随性，
\begin{equation}\tag{341}
\operatorname{Hom}(p^{!}_{1}M_{X},p^{!}_{2}M_{X})=\operatorname{Hom}(p_{2\cdot}p^{!}_{1}M_{X},M_{X})=\operatorname{Hom}(M_{X}\otimes\mathcal{D}_{X},M_{X}).
\end{equation}
这里把 $M_{X}\otimes\mathcal{D}_{X}$ 经由 $\mathcal{D}_{X}$ 上的右 $\mathcal{O}$-模结构看作 $\mathcal{O}_{X}$-模。于是 $\tau:p^{!}_{1}M_{X}\to p^{!}_{2}M_{X}$ 等同于态射 $M_{X}\otimes\mathcal{D}_{X}\to M_{X}$。可以验证，对 $\tau$ 的条件恰好意味着这个箭头是 $\mathcal{D}_{X}$ 在 $M_{X}$ 上的右幺作用。关于评注与细节，见下一个注记。\hfill$\Box$

\bdp{7.10.13} \emph{注.} 我们在更一般的设定下讨论 7.10.12 的若干要点。由于 $\mathcal{O}_{X^{<2>}}$ 是 $\mathcal{O}_{X}\otimes_{\mathbb{C}}\mathcal{O}_{X}$ 的完备化，可以把 $\mathcal{M}(X^{<2>},\mathcal{O})$ 的对象看作 $X$ 上某些 $\mathcal{O}_{X}$-双模层，称为 Diff-双模\footnote{[BB93] 中用的是「微分双模」这一术语；细节请参看那里。}。若 $A$、$B$ 是 Diff-双模，则 $A\otimes_{\mathcal{O}_{X}}B$ 也是（故 $\mathcal{M}(X^{<2>},\mathcal{O})$ 是幺半范畴）。注意 $A\otimes_{\mathcal{O}_{X}}B$ 实际上（以显然的方式）是\SourcePage{292}$\mathcal{M}(X^{<3>},\mathcal{O})$ 的对象。

由伴随性，对任意 $C\in\mathcal{M}(X^{<2>},\mathcal{O})$，Diff-双模的态射 $A\otimes_{\mathcal{O}_{X}}B\to C$ 等同于 $\mathcal{M}(X^{<3>},\mathcal{O})$ 中的态射 $A\otimes_{\mathcal{O}_{X}}B\to p^{!}_{13}C$。于是对 Diff-代数\footnote{即 Diff-双模的幺半范畴中的代数。} $A$，其乘法等同于 $\mathcal{M}(X^{<3>},\mathcal{O})$ 中的态射 $m:A\otimes_{\mathcal{O}_{X}}A\to p^{!}_{13}A$（结合律如何用这些术语写出，留给读者）。类似地，对（右）$A$-模 $M_{X}$，可以把 $A$-作用写成 $\mathcal{M}(X^{<2>},\mathcal{O})$ 中的态射 $a:M_{X}\otimes_{\mathcal{O}_{X}}A\to p^{!}_{2}M_{X}$；作用（结合）性质恰好是说，由 $m$ 与 $a$ 得到的 $\mathcal{M}(X^{<3>},\mathcal{O})$ 中两个态射 $M_{X}\otimes_{\mathcal{O}_{X}}A\otimes_{\mathcal{O}_{X}}A\to p^{!}_{3}M$ 相同。现在设 $A=\mathcal{D}_{X}$，或更一般地，$A$ 是 tdo。则 $m:A\otimes_{\mathcal{O}_{X}}A\to p^{!}_{13}A$ 是同构\footnote{这个性质或许刻画了 tdo。}。若 $M_{X}$ 是（可能非幺的）$A$-模，则 $a:M_{X}\otimes_{\mathcal{O}_{X}}A\to p^{!}_{2}M_{X}$ 是同构，当且仅当该模是幺模。

\bdp{7.10.14} 在光滑设定下，把 7.10.5、7.10.10、7.10.11(i) 中的函子 $f^{!}$、$p_{*}$、$i_{*}$ 分别与标准的 $\mathcal{D}$-模函子等同起来，这件事留给读者。

结合 7.10.12 与 7.10.11(ii) 可见：若 $X$ 是任意代数空间，则 $X$ 上的 $\mathcal{D}$-晶体等同于 [Sa91] 意义下 $X$ 上的 $\mathcal{D}$-模\footnote{Saito 更愿意处理解析设定，但他的定义有明显的代数版本（而上述定义有明显的解析版本）。}。

\bdp{7.10.15} 本节其余部分是 tdo 与扭 $\mathcal{D}$-模的晶体设定的概要。先讨论晶体 $\mathcal{O}^{*}$-gerbe。在光滑概形的情形，这样的 gerbe 等同于「tdo 模去线丛的扭」这一概念的平展局部化版本。然后对晶体 $\mathcal{O}^{*}$-gerbe $\mathcal{C}$，定义相应的扭 $\mathcal{D}$-晶体的阿贝尔范畴 $\mathcal{M}_{\mathcal{C}}(X)$。

\bdp{7.10.16} 与前面一样，$X$ 是任意代数空间。范畴 $X_{cr}$ 带有 site 的结构（平展晶体拓扑）：覆盖是态射族\SourcePage{293}$\{(S_{i},\hat S_{i})\to(S,\hat S)\}$，使得 $\{\hat S_{i}\to\hat S\}$ 是 $\hat S$ 的平展覆盖。其上带有环层 $\mathcal{O}_{cr}$，其中 $\mathcal{O}_{cr}(S,\hat S)=\mathcal{O}(\hat S)$。于是有相应的可逆元层 $\mathcal{O}^{*}_{cr}$。

\bdp{7.10.17} \emph{定义.} $X$ 上的晶体 $\mathcal{O}^{*}$-gerbe 是 $X_{cr}$ 上的 $\mathcal{O}^{*}_{cr}$-gerbe\footnote{即 [De-Mi] 术语中 $X_{cr}$ 上带 $\mathcal{O}^{*}_{cr}$ 的 gerbe。}。

确切地说，意思如下。考虑 $X_{cr}$ 上的 Picard 群胚层 $\mathcal{P}ic_{cr}$，其中 $\mathcal{P}ic_{cr}(S,\hat S):=\mathcal{P}ic(\hat S)$（即 $\hat S$ 上线丛的 Picard 群胚）。于是 $X$ 上的晶体 $\mathcal{O}^{*}$-gerbe 是 $X_{cr}$ 上的 $\mathcal{P}ic_{cr}$-torsor $\mathcal{C}$（即 $\mathcal{C}$ 是 $X_{cr}$ 上的纤维范畴，带有 $\mathcal{P}ic_{cr}$ 的作用，使得每个纤维 $\mathcal{C}(\hat S)=\mathcal{C}(S,\hat S)$ 是 $\mathcal{P}ic(\hat S)$-torsor），且在 $X_{cr}$ 上局部地 $\mathcal{C}(S,\hat S)$ 非空。

晶体 $\mathcal{O}^{*}$-gerbe 构成 Picard 2-群胚 $\mathcal{G}_{cr}(X)$。gerbe 的等价类群是 $H^{2}(X_{cr},\mathcal{O}^{*}_{cr})$。对一对 gerbe $\mathcal{C}$、$\mathcal{C}'$，态射 $\phi:\mathcal{C}\to\mathcal{C}'$ 构成 $\mathcal{P}ic(X_{cr})$-torsor。这里 $\mathcal{P}ic(X_{cr})$ 是 $X_{cr}$ 上 $\mathcal{O}^{*}_{cr}$-torsor 的 Picard 群胚\footnote{若 $X$ 光滑，则这样的 torsor 等同于 $X$ 上带平坦联络的线丛。}。

\bdp{7.10.18} \emph{注.} (i) 设 $X_{\text{ét}\,cr}$ 是 $X$ 的小平展晶体 site（作为范畴它等于 7.10.4 中的 $X^{(i)}_{cr}$，拓扑由 $X_{cr}$ 诱导）。$X$ 上的晶体 $\mathcal{O}^{*}$-gerbe 经限制给出 $X_{\text{ét}\,cr}$ 上的 $\mathcal{O}^{*}_{cr}$-gerbe。我们得到 gerbe 的 Picard 2-群胚的等价，其验证留给读者\footnote{我们把 $X_{cr}$ 当作基本设定，因为它直接推广到归纳概形的情形，见 7.11.6。}。

(ii) 我们的 $\mathcal{G}_{cr}(X)$ 是与复形 $\tau_{\leq2}R\Gamma(X_{cr},\mathcal{O}^{*}_{cr})=\tau_{\leq2}R\Gamma(X_{\text{ét}\,cr},\mathcal{O}^{*}_{cr})$ 相联系的 Picard 2-群胚。为计算 $R\Gamma$，考察由 $(\mathcal{O}_{cr}/\mathcal{I}_{cr})(S,\hat S)=\mathcal{O}(S)$ 定义的典范理想 $\mathcal{I}_{cr}\subset\mathcal{O}_{cr}$。有环化拓扑的典范态射 $i:X_{\text{ét}}\to X_{\text{ét}\,cr}$，$i^{-1}(S,\hat S)=S$，而 $\mathcal{I}_{cr}$ 适合短正合列 $0\to\mathcal{I}_{cr}\to\mathcal{O}_{cr}\to i_{\cdot}\mathcal{O}_{X}\to 0$。过渡\SourcePage{294}到可逆元层，得到短正合列
\begin{equation}\tag{342}
0\longrightarrow\mathcal{I}_{cr}\xrightarrow{\exp}\mathcal{O}^{*}_{cr}\longrightarrow i_{\cdot}\mathcal{O}^{*}_{X}\longrightarrow 0,
\end{equation}
其中 $\exp$ 是指数映射（由于每个 $\mathcal{I}_{cr}(S,\hat S)$ 都是幂零理想，我们的 $\exp$ 定义得当）。由于 $R\Gamma(X_{\text{ét}\,cr},i_{\cdot}\mathcal{O}^{*}_{X})=R\Gamma(X_{\text{ét}},\mathcal{O}^{*})$。可以用 (343) 来计算 $R\Gamma(X_{cr},\mathcal{O}^{*}_{cr})$。%
% [Tutor Note: 原书 p.294 此处句末印作句号而非逗号，且所引编号印作 (343)——讨论的其实是本页的 (342)，(343) 在 p.295；底本如此，保留原样。]
例如，由于 $H^{0}(X_{cr},\mathcal{I}_{cr})=0$，群 $H^{0}(X_{cr},\mathcal{O}^{*}_{cr})$ 就是 $X$ 上局部常值可逆函数构成的群 $\mathcal{O}^{*}(X)_{con}$。

(iii) 设 $X$ 光滑。令 $\Omega^{\geq1}_{X}:=(0\to\Omega^{1}_{X}\to\Omega^{2}_{X}\ldots)$。按 Grothendieck，有 $R\Gamma(X_{cr},\mathcal{O}_{cr})=R\Gamma(X,\Omega_{X})$ 与 $R\Gamma(X_{cr},\mathcal{I}_{cr})=R\Gamma(X_{cr},\mathit{Cône}(\mathcal{O}_{cr}\to i_{\cdot}\mathcal{O}_{X})[-1])=R\Gamma(X,\Omega^{\geq1}_{X})$。于是 (342) 给出长上同调序列
\[
\begin{gathered}
0\longrightarrow\mathcal{O}^{*}(X)_{con}\longrightarrow\mathcal{O}^{*}(X)\xrightarrow{\ dlog\ }\Omega^{1cl}(X)\longrightarrow H^{1}(X_{cr},\mathcal{O}^{*}_{cr})\longrightarrow\\
\longrightarrow\mathcal{P}ic(X)\xrightarrow{\ c_{1}\ }H^{2}(X,\Omega^{\geq1}_{X})\longrightarrow H^{2}(X_{cr},\mathcal{O}^{*}_{cr})\longrightarrow Br(X)\longrightarrow0.
\end{gathered}
\]
这里 $H^{1}(X_{cr},\mathcal{O}^{*}_{cr})$ 是 $X$ 上带平坦联络的线丛的同构类群。右端的 0 成立，因为 $H^{2}(X_{\text{ét}},\mathcal{O}^{*})=Br(X)$ 是扭群，而 $H^{3}(X_{cr},\mathcal{I}_{cr})$ 是 $\mathbb{C}$-向量空间。

(iv) 若 $X$ 是概形，则可考虑更弱的拓扑 $X_{Zar\,cr}$（作为范畴它等于 7.10.4 中的 $X^{(ii)}_{cr}$）。得到 $X_{Zar\,cr}$ 上 $\mathcal{O}^{*}_{cr}$-gerbe 的相应 Picard 2-群胚 $\mathcal{G}_{Zar\,cr}(X)$。由平展下降，拉回函子 $\mathcal{G}_{Zar\,cr}(X)\to\mathcal{G}_{cr}(X)$ 是 Picard 2-群胚的全忠实态射；即 $\mathcal{G}_{Zar\,cr}(X)$ 是 Zariski 局部平凡晶体 $\mathcal{O}^{*}$-gerbe 的 2-群胚。容易看出\footnote{参看 7.10.22。}，$\mathcal{C}\in\mathcal{G}_{cr}(X)$ 属于 $\mathcal{G}_{Zar\,cr}(X)$，当且仅当 $X_{\text{ét}}$ 上的 $\mathcal{O}^{*}$-gerbe $i_{\cdot}\mathcal{C}$ 是 Zariski 局部平凡的。例如，若 $X$ 光滑，则 $H^{2}(X_{Zar},\mathcal{O}^{*})=0$，故 $\mathcal{G}_{cr}(X)/\mathcal{G}_{Zar\,cr}(X)=Br(X)$。

\SourcePage{295}

\bdp{7.10.19} 下面给出（适当刚性化的）晶体 $\mathcal{O}^{*}$-gerbe 的一个方便的「具体」描述。

设有如 7.10.7 中的 $X\hookrightarrow V$。对 $\mathcal{C}\in\mathcal{G}_{cr}(X)$ 与 $X$ 在 $V$ 中的无穷小邻域 $X'$，有 $\mathcal{P}ic(X')$-torsor $\mathcal{C}(X')$。令 $\mathcal{C}(V):=\varprojlim\mathcal{C}(X')$（$:=$ $\mathcal{C}$ 在诸 $X'$ 构成的有向集上的笛卡儿截面群胚）；这是 $\mathcal{P}ic(V)$-torsor。

考虑这样的对 $(\mathcal{C},\mathcal{E}_{V})$：$\mathcal{C}\in\mathcal{G}_{cr}(X)$，$\mathcal{E}_{V}\in\mathcal{C}(V)$。这样的对象构成 Picard 群胚 $\mathcal{G}^{V}_{cr}(X)$。确切地说，态射 $(\mathcal{C},\mathcal{E}_{V})\to(\mathcal{C}',\mathcal{E}'_{V})$ 是一对 $(F,\nu)$，其中 $F:\mathcal{C}\to\mathcal{C}'$ 是态射，$\nu:F(\mathcal{E}_{V})\xrightarrow{\sim}\mathcal{E}'_{V}$\footnote{注意这样的对没有对称性，故 $\mathcal{G}^{V}_{cr}(X)$ 是通常的群胚（而 $\mathcal{G}_{cr}(X)$ 是 2-群胚）。}。我们来描述 $\mathcal{G}^{V}_{cr}(X)$。

沿用 7.10.7 的记号。设 $\mathcal{R}$ 是 $V^{<2>}$ 上的线丛，$\beta:p^{*}_{12}\mathcal{R}\otimes p^{*}_{23}\mathcal{R}\xrightarrow{\sim}p^{*}_{13}\mathcal{R}$ 是 $V^{<3>}$ 上线丛的同构，使得下述 $V^{<4>}$ 上线丛同构的图表交换（结合条件）：
\begin{equation}\tag{343}
\begin{tikzcd}
\mathcal{R}_{12}\otimes\mathcal{R}_{23}\otimes\mathcal{R}_{34} \arrow[r] \arrow[d] & \mathcal{R}_{13}\otimes\mathcal{R}_{34} \arrow[d]\\
\mathcal{R}_{12}\otimes\mathcal{R}_{24} \arrow[r] & \mathcal{R}_{14}
\end{tikzcd}
\end{equation}
这里 $\mathcal{R}_{ij}$ 是 $\mathcal{R}$ 经投影 $p_{ij}:V^{<4>}\to V^{<2>}$ 的拉回，而箭头来自 $\beta$。

这样的对 $(\mathcal{R},\beta)$ 构成 Picard 群胚 $G(V)$（关于张量积）。

\bdp{7.10.20} \emph{命题.} Picard 群胚 $\mathcal{G}^{V}_{cr}(X)$ 与 $G(V)$ 典范等价。

\emph{证明.} 对 $(\mathcal{C},\mathcal{E}_{V})\in\mathcal{G}^{V}_{cr}(X)$ 令 $\mathcal{R}:=\mathcal{H}om(p^{*}_{1}\mathcal{E}_{V},p^{*}_{2}\mathcal{E}_{V})\in\mathcal{P}ic(V)$，并把 $\beta$ 定义为复合同构；显然 $(\mathcal{R},\beta)\in G(V)$。于是得到 Picard 群胚的态射 $\mathcal{G}^{V}_{cr}(X)\to G(V)$。

% 尾部接 p.296：7.10.20 的证明在 p.296 上续（「The inverse Morphism assigns to $(R,\beta)$…」），本单元至此停笔。
%% end units/ch07_u06.tex
%% begin units/ch07_u07.tex
% ch07_u07 —— 原书 p.296--p.305（7.10.20 证明续 -- 7.11.6）
% 本单元自 7.10.20 证明的续文起（承上单元 p.295 末句），至 7.11.6 末句止；
% 尾部接 p.306。

\SourcePage{296}

逆向态射把 $(\mathcal{R},\beta)$ 送到由平凡 gerbe 借助
$(\mathcal{R},\beta)$ 粘成的对 $(\mathcal{C},\mathcal{E}_V)$。具体地说，
$(\mathcal{C},\mathcal{E}_V)$ 的定义如下。由于 $V$ 形式光滑，结构态射
$j:S\to X$ 可延拓为 $j':\hat S\to V$。此时 $\mathcal{C}(\hat S)$ 是一个
$\mathcal{P}ic(\hat S)$-torsor，并带有如下附加结构：

\begin{enumerate}[label=(\roman*),leftmargin=2.2em]
\item 对任一如上之 $j'$，给定 $\mathcal{C}(\hat S)$ 的一个对象，记作
$j'^*\mathcal{E}_V$。
\item 若 $j'':\hat S\to V$ 是 $j$ 的另一延拓，则我们有一个线丛之间的等同
$\theta_{j''j'}:\operatorname{Hom}(j'^*\mathcal{E}_V,j''^*\mathcal{E}_V)
\xrightarrow{\ \sim\ }(j'',j')^*\mathcal{R}$。
\end{enumerate}

我们要求 (ii) 把 $\operatorname{Hom}$ 的复合与 $\beta$ 所定义的同构等同起来。
容易看出，这样的 $\mathcal{C}(\hat S)$ 存在且唯一（相差唯一的等价）。诸纤维
$\mathcal{C}(\hat S)$ 以显然的方式粘成一个晶体 $\mathcal{O}^*$-gerbe。按构造有
$\mathcal{E}_V\in\mathcal{C}(V)$。\hfill$\square$

\bdp{7.10.21} \emph{注.} 设 $\mathcal{E}'_V$ 是 $\mathcal{C}(V)$ 的另一对象，
而 $(\mathcal{R}',\beta')\in G(V)$ 是对应于 $(\mathcal{C},\mathcal{E}'_V)$ 的对。
置 $\mathcal{L}:=\operatorname{Hom}(\mathcal{E}_V,\mathcal{E}'_V)
\in\mathcal{P}ic(V)$。则
\[
\mathcal{R}'=\operatorname{Ad}_{\mathcal{L}}\mathcal{R}
:=(p_2^*\mathcal{L})\otimes\mathcal{R}\otimes(p_1^*\mathcal{L})^{\otimes-1},
\qquad
\beta'=\operatorname{Ad}_{\mathcal{L}}\beta .
\]
现在设 $\mathcal{C}$ 是 $X$ 上任一晶体 $\mathcal{O}^*$-gerbe，并设我们有如上的
$X\hookrightarrow V$。要用 7.10.20 描述 $\mathcal{C}$，必须保证
$\mathcal{C}(V)$ 非空。

% [Tutor Note: 原书 p.296 此处正文只写「$\mathcal{C}(V)$ is non-empty
% if$^{*)}\ldots$」，脚注 \*) 仅作「and, certainly, only if」；正文的 if 与脚注
% 合读才是「当且仅当」。译文按原书拆分照录，不合并。]
\bdp{7.10.22} \emph{引理.} 设 $X$ 仿射，且 $V$ 是可数多个子概形之并。则
$\mathcal{C}(V)$ 非空，只要 $\mathcal{C}(X,X)$ 非空%
\footnote{并且当然，仅当。}。

\emph{证明.} 设 $X'\subset V$ 是 $X$ 的一个无穷小邻域。则任一
$\mathcal{E}_X\in\mathcal{C}(X,X)$ 都可延拓为
$\mathcal{E}_{X'}\in\mathcal{C}(X,X')$，且所有这样的延拓彼此同构。现在有 $X$
的无穷小邻域序列 $X\subset X^{(1)}\subset X^{(2)}\cdots$，使得
$V=\varinjlim X^{(n)}$。归纳地定义序列
$\mathcal{E}_X(n)\in\mathcal{C}(X,X^{(n)})$ 以及等同
$\mathcal{E}_X(n+1)|_{X^{(n)}}=\mathcal{E}_X(n)$。这就是
$\mathcal{E}_V\in\mathcal{C}(V)$。\hfill$\square$

\SourcePage{297}

\bdp{7.10.23} \emph{注.} (i) 考虑 $X_{\text{ét}}$ 上的 $\mathcal{O}^*$-gerbe
$i\mathcal{C}$（故 $i\mathcal{C}(U)=\mathcal{C}(U,U)$）。则
$\mathcal{C}(X,X)\neq\varnothing$ 当且仅当 $i\mathcal{C}$ 是平凡 gerbe，即它在
$H^2(X_{\text{ét}},\mathcal{O}^*)=Br(X)$ 中的类消没%
\footnote{这个类是 $\mathcal{C}$ 的类在映射
$H^2(X_{cr},\mathcal{O}^*_{cr})\to
H^2(X_{cr},i\cdot\mathcal{O}^*)=Br(X)$ 下的像。}。

(ii) 对任意代数空间 $X$ 与 $\mathcal{C}\in\mathcal{G}_{cr}(X)$，可在
$X_{\text{ét}}$ 局部用 7.10.20 描述 $\mathcal{C}$。具体地说，存在 $X$ 的平展
覆盖 $U_i$，使 $U_i$ 仿射且 $\mathcal{C}(U_i,U_i)\neq\varnothing$。把 $U_i$
嵌入光滑概形，并取相应的形式完备化 $V_i$。于是由 7.10.22，可用 7.10.20、
7.10.21 描述 $\mathcal{C}_{U_i}$。

\bdp{7.10.24} \emph{定义.} 设 $\mathcal{C}\in\mathcal{G}_{cr}(X)$。$X$ 上的
$\mathcal{C}$-扭 $\mathcal{D}$-晶体是一个 Cartesian 函子
$M:\mathcal{C}\to\mathcal{M}^{!}(X_{cr},\mathcal{O})$，使得对任一
$\varepsilon\in\mathcal{C}(\hat S)$ 与 $f\in\mathcal{O}^*(\hat S)$，有
$M(f\varepsilon)=f\cdot\operatorname{id}_{M(\varepsilon)}$。

$\mathcal{C}$-扭 $\mathcal{D}$-晶体构成一个 $\mathbb{C}$-范畴
$\mathcal{M}_{\mathcal{C}}(X)$。它以函子性的方式依赖于 $\mathcal{C}$
（对态射 $\mathcal{C}\to\mathcal{C}'$，相应地有范畴等价
$\mathcal{M}_{\mathcal{C}}(X)\xrightarrow{\ \sim\ }
\mathcal{M}_{\mathcal{C}'}(X)$，等等）。

诸范畴 $\mathcal{M}_{\mathcal{C}}(U)=\mathcal{M}_{\mathcal{C}_U}(U)$，
$U\in X_{\text{ét}}$，以显然的方式构成 $X_{\text{ét}}$ 上的范畴层
$\mathcal{M}_{\mathcal{C}}(X_{\text{ét}})$。

设 $\mathcal{C}^{triv}$ 是平凡化 gerbe，故
$\mathcal{C}^{triv}(\hat S)=\mathcal{P}ic(\hat S)$。以
$\mathcal{C}^{triv}$ 扭的 $\mathcal{D}$-晶体与通常的 $\mathcal{D}$-晶体相同。
具体地说，把 $M\in\mathcal{M}_{\mathcal{C}^{triv}}(X)$ 与 $\mathcal{D}$-晶体
$M_{\hat S}:=M(\mathcal{O}_{\hat S})$ 等同起来。

\emph{注.} 在上述定义中，可把 $X_{cr}$ 换成
$X_{\text{ét}\,cr}$。若 $X$ 是概形且
$\mathcal{C}\in\mathcal{G}_{Zar\,cr}(X)$，则可把 $X_{cr}$ 换成 $X_{Zar\,cr}$。
得到的是同一个范畴 $\mathcal{M}_{\mathcal{C}}(X)$。

\bdp{7.10.25} 这里是 7.10.7、7.10.8 的扭版本。设我们处于情形 7.10.19，即
有 $(\mathcal{C},\mathcal{E}_V)\in\mathcal{G}^V_{cr}(X)$ 以及相应的
$(\mathcal{R},\beta)\in G(V)$（见 7.10.20）。范畴
$\mathcal{M}_{\mathcal{C}}(X)$ 可描述如下。设
$\mathcal{M}_{\mathcal{R}}(X)=\mathcal{M}_{\mathcal{R}\beta}(X)$ 是二元对
$(M_V,\tau)$ 的范畴，其中 $M_V\in\mathcal{M}(V,\mathcal{O})$，而
$\tau:(p_1^{!}M_V)\otimes\mathcal{R}\xrightarrow{\ \sim\ }p_2^{!}M_V$
是 $\mathcal{M}(V^{<2>},\mathcal{O})$\SourcePage{298}中的同构，且满足%
\footnote{我们用 $\beta$ 来等同等式左右两端所在的模。}
\begin{equation}\tag{344}
p_{23}^{!}(\tau)p_{12}^{!}(\tau)=p_{13}^{!}(\tau).
\end{equation}

\bdp{7.10.26} \emph{引理.} 范畴 $\mathcal{M}_{\mathcal{C}}(X)$ 与
$\mathcal{M}_{\mathcal{R}}(X)$ 典范等价。

\emph{证明.} 对 $M\in\mathcal{M}_{\mathcal{C}}(X)$，置
$M_V=M(\mathcal{E}_V):=\bigcup M(\mathcal{E}_{(X,X')})$，并把 $\tau$ 定义为诸同构
\[
(p_1^{!}M_V)\otimes\mathcal{R}=M(p_1^*\mathcal{E}_V)\otimes\mathcal{R}
=M((p_1^*\mathcal{E}_V)\otimes\mathcal{R})
=M(p_2^*\mathcal{E}_V)=p_2^{!}M_V
\]
的复合。其余部分是 7.10.8 的证明的直接改写。\hfill$\square$

\bdp{7.10.27} \emph{引理.} 对任意 $X$ 与 $\mathcal{C}\in\mathcal{G}_{cr}(X)$，
范畴 $\mathcal{M}_{\mathcal{C}}(X)$ 是阿贝尔范畴。

\emph{证明.} 对 7.10.9 的证明作显然的修改。用到 7.10.23(ii)、7.10.22、
7.10.26。\hfill$\square$

\bdp{7.10.28} 从现在起，我们设 $X$ 是光滑代数空间。我们要把上述图景与 tdo
以及扭 $\mathcal{D}$-模的通常框架作比较。首先，把晶体
$\mathcal{O}^*$-gerbe 与 tdo 联系起来\footnote{关于 tdo 的基本事实，例如参看
[BB93] 2.1。}。

考察 $V=X$ 时的 7.10.19。考虑由对 $(\mathcal{C},\mathcal{E}_X)$ 构成的
Picard 群胚 $\mathcal{G}^{\sim}_{cr}(X):=\mathcal{G}^{V}_{cr}(X)$，其中
$\mathcal{C}$ 是 $X$ 上的晶体 $\mathcal{O}^*$-gerbe，而
$\mathcal{E}_X\in\mathcal{C}(X)$。

$\mathcal{G}^{\sim}_{cr}(X)$ 有一个方便的解释。考虑 $X$ 上的
$\mathcal{I}_{cr}$-gerbe（即 $X_{cr}$ 上的 $\mathcal{I}_{cr}$-gerbe）。由于
$H^0(X_{cr},\mathcal{I}_{cr})=0$，这些 gerbe 构成（平移的）Picard 群胚
$\mathcal{G}\mathcal{I}_{cr}(X)$。指数态射
$\mathcal{I}_{cr}\hookrightarrow\mathcal{O}^*_{cr}$ 给出函子
$\exp:\mathcal{G}\mathcal{I}_{cr}(X)\to\mathcal{G}_{cr}(X)$。由于
$\mathcal{I}_{(X,X)}=0$，对任一 $\mathcal{I}_{cr}$-gerbe $\mathcal{B}$，群胚
$\mathcal{B}_X$ 平凡，故群胚 $(\exp\mathcal{B})_X$ 有特异对象
$\mathcal{E}_{\mathcal{B}X}$（在典范同构意义下确定）。这样我们就定义了 Picard
群胚的态射
\begin{equation}\tag{345}
\exp:\mathcal{G}\mathcal{I}_{cr}(X)\longrightarrow \mathcal{G}^{\sim}_{cr}(X),
\end{equation}

\SourcePage{299}

$\mathcal{B}\mapsto(\exp\mathcal{B},\mathcal{E}_{\mathcal{B}X})$。这是 Picard
群胚的等价（这由 (342) 推出）。

\emph{例.} (342) 的「边界映射」给出 Picard 群胚的态射
$c:\mathcal{P}ic(X)\to\mathcal{G}\mathcal{I}_{cr}(X)$（即晶体 Chern 类）。
用 (345) 的术语，它把 $\mathcal{L}\in\mathcal{P}ic(X)$ 送到对
$(\mathcal{C}^{triv},\mathcal{L})$。

\bdp{7.10.29} \emph{命题.} $\mathcal{G}^{\sim}_{cr}(X)$ 与 $X$ 上 tdo 构成的
Picard 群胚 $\mathcal{T}\mathcal{D}\mathcal{O}(X)$ 典范等价。

\emph{证明.} 按 $V=X$ 时的 7.10.20，把我们的 $\mathcal{G}^{\sim}_{cr}(X)$ 与
$G(X)$ 等同起来。现设 $(\mathcal{R},\beta)\in G(X)$，相应的 tdo
$\mathcal{D}_{\mathcal{R}}=\mathcal{D}_{(\mathcal{R},\beta)}$ 如下定义。我们沿用
7.10.13 的记号。把 $\mathcal{D}_X$ 看作 Diff-双模（即
$\mathcal{M}(X^{<2>},\mathcal{O})$ 的一个对象）。置
\[
\mathcal{D}_{\mathcal{R}}:=
\mathcal{D}_X\otimes_{\mathcal{O}_{X^{<2>}}}\mathcal{R}.
\]
乘法态射
$m_{\mathcal{R}}:\mathcal{D}_{\mathcal{R}}\otimes_{\mathcal{O}_X}
\mathcal{D}_{\mathcal{R}}\to p_{13}^{!}\mathcal{D}_{\mathcal{R}}$ 是
$\mathcal{D}_X$ 的相应态射与 $\beta$ 的张量积。容易验证
$\mathcal{D}_{\mathcal{R}}$ 是 tdo，且 $G(X)\to\mathcal{T}\mathcal{D}\mathcal{O}(X)$、
$(\mathcal{R},\beta)\mapsto\mathcal{D}_{\mathcal{R}}$ 是 Picard 群胚的态射。

逆向态射把 $X$ 上的 tdo $A$ 送到对象 $(\mathcal{R},\beta)$，其中
$\mathcal{R}:=\operatorname{Hom}_{\mathcal{O}_{X^{<2>}}}(\mathcal{D}_X,A)$，而
$\beta$ 由比较 $\mathcal{D}_X$ 与 $A$ 的乘法态射 $m$ 来定义。细节留给读者。
\hfill$\square$

\bdp{7.10.30} \emph{注.} 上述等价还有另一种（等价的）写法。由 (345)，
$\mathcal{G}^{\sim}_{cr}(X)$ 等价于 $\mathcal{G}\mathcal{I}_{crys}(X)$，即等价于
与复形 $\tau_{\le1}(R\Gamma(X_{crys},\mathcal{I}_{Xcrys})[1])$ 相伴的 Picard
群胚。按 [BB93] 2.1.6、2.1.4，$\mathcal{T}\mathcal{D}\mathcal{O}(X)$ 是与复形
$\tau_{\le1}(R\Gamma(X,\Omega_X^{\ge1})[1])$ 相伴的 Picard 群胚。现在上述两个复形
典范拟同构（见 7.10.18(iii)）。

\bdp{7.10.31} 这里是 7.10.12 的扭版本。对
$(\mathcal{C},\mathcal{E}_X)\in\mathcal{G}^{\sim}_{cr}(X)$，考虑相应的
$(\mathcal{R},\beta)\in G(X)$ 与 tdo $\mathcal{D}_{\mathcal{R}}$。取
$M\in\mathcal{M}_{\mathcal{C}}(X)$。按 7.10.26，可把 $M$ 看作
$\mathcal{M}_{\mathcal{R}}(X)$ 中的对 $(M_X,\tau)$。\SourcePage{300}由于%
\footnote{参看 7.10.12 与 7.10.29 的证明。}
$p_1^{!}M_X=M_X\otimes_{\mathcal{O}_X}\mathcal{D}_X$ 且
$\mathcal{D}_{\mathcal{R}}=
\mathcal{D}_X\otimes_{\mathcal{O}_{X^{<2>}}}\mathcal{R}$，我们可以把 $\tau$
改写成 $\mathcal{M}(X^{<2>},\mathcal{O})$ 中的同构
\begin{equation}\tag{346}
M_X\otimes_{\mathcal{O}_X}\mathcal{D}_{\mathcal{R}}
\xrightarrow{\ \sim\ }p_2^{!}M_X .
\end{equation}
由伴随性，可把 (346) 看作 $\mathcal{O}_X$-模的态射
\begin{equation}\tag{347}
M_X\otimes\mathcal{D}_{\mathcal{R}}\to M_X .
\end{equation}
记 $\mathcal{M}^{r}(X,\mathcal{D}_{\mathcal{R}})$ 为 $X$ 上右
$\mathcal{D}_{\mathcal{R}}$-模的范畴。

\bdp{7.10.32} \emph{命题.} 态射 (347) 是 $\mathcal{D}_{\mathcal{R}}$ 在
$M_X$ 上的右幺作用。函子
$\mathcal{M}_{\mathcal{C}}(X)\to\mathcal{M}^{r}(X,\mathcal{D}_{\mathcal{R}})$、
$M\mapsto M_X$ 是范畴等价。

\emph{证明.} 留给读者（见 7.10.12、7.10.13）。\hfill$\square$

\bdsec{7.11. 归纳概形上的 $\mathcal{D}$-模}

本节讨论形式光滑归纳概形上的 $\mathcal{D}$-模理论。注意 $\mathcal{D}$-晶体
图景（见 7.10）在归纳概形的框架下立刻有意义；而通常的进路（微分算子等）
要费些篇幅才写得下来。

\bdp{7.11.1} 严格意义下的归纳概形 $X$ 是一个「空间」（即交换
$\mathbb{C}$-代数范畴上的集合值函子
$A\mapsto X(A)=X(\operatorname{Spec}A)$），它可以表为
$\varinjlim X_\alpha$，其中 $\{X_\alpha\}$ 是拟紧概形的有向族，使得所有映射
$i_{\alpha\beta}:X_\alpha\to X_\beta$（$\alpha\le\beta$）都是闭嵌入。若 $X$
能如上表出，且指标 $\alpha$ 的集合可数，则称 $X$ 是
$\aleph_0$-归纳概形\footnote{并非所有自然的归纳概形例子都是
$\aleph_0$-归纳概形；例如对每个无穷维向量空间 $V$，函子
$A\mapsto\operatorname{End}_A(V\otimes A)$ 是归纳概形，但不是
$\aleph_0$-归纳概形。}。若 $P$\SourcePage{301}是概形的一个在取闭子概形下
保持的性质，则称 $X$ 满足归纳-$P$ 性质，如果每个 $X_\alpha$ 都满足 $P$。

置 $X_{red}:=\varinjlim X_{\alpha red}$；若 $X_{red}=X$，则称归纳概形 $X$
是既约的。

形式概形是这样的归纳概形 $X$：$X_{red}$ 是概形（关于形式概形的这一定义与
EGA 中定义的关系，参看 7.12.17 中的讨论）。$\aleph_0$-形式概形是
$\aleph_0$-归纳概形的形式概形。归纳概形 $Z$ 沿闭子概形 $Y\subset Z$ 的
完备化，是满足 $Y'_{red}=Y_{red}$ 的闭子概形 $Y'\subset Z$ 的正向极限。
对形式概形，我们不写「归纳仿射」而写作「仿射」。形式概形 $X$ 仿射，当且仅当
$X_{red}$ 仿射。

仿照 Grothendieck（[Gr64]、[Gr67]），若对每个 $A$ 与每个幂零理想
$I\subset A$，映射 $X(A)\to X(A/I)$ 都是满射，则称 $X$ 形式光滑。容易看出，
对归纳有限型的归纳概形，形式光滑性是局部性质（参看 [Gr67] 中命题 17.1.6
的证明）\footnote{我们不知道对非归纳有限型的归纳概形这是否成立。对概形，
答案是肯定的。这可由 [Gr68a] 中的注 9.5.8 与下述出人意料的结果（[RG]，
p.~82，3.1.4）推出：投射模这一性质对 Zariski 拓扑是局部的，甚至对 fpqc
拓扑也是局部的（无需任何有限性假设！）。}。若对任意如上之 $A$、$I$，从
$X(A)$ 到 $Y(A)$ 与 $X(A/I)$ 在 $Y(A/I)$ 上的纤维积的映射是满射，则称
态射 $X\to Y$ 形式光滑。

设 $X$ 是归纳概形。若对任意满足 $Y\subset Z$ 的闭子概形 $Z\subset X$，$Y$
在 $\mathcal{O}_Z$ 中的理想都有限生成，则称闭拟紧子概形 $Y\subset X$ 是
合理的。若 $X$ 是其合理子概形之并，即 $X$ 可表为 $\varinjlim X_\alpha$
且所有 $X_\alpha$ 都合理，则称 $X$ 是合理的。合理归纳概形的闭子空间是合理
归纳概形；两个合理归纳概形之积合理。

\SourcePage{302}

\emph{注.} 把上述定义中的「概形」一词换成「代数空间」，就得到归纳代数空间
的概念。全部讨论都自动适用于归纳代数空间的情形。

\bdp{7.11.2} \emph{例.} (i) 归纳仿射的归纳概形 $X$ 等同于一个投射代数，即
交换代数范畴中的一个投射对象 $R$，它可表为 $\varprojlim R_\alpha$，使得映射
$R_\beta\to R_\alpha$（$\beta\ge\alpha$）都是满射。我们写
$X=\operatorname{Spf}R:=\varinjlim\operatorname{Spec}R_\alpha$。拓扑由开理想
$I_\alpha\subset R$ 定义的完备拓扑交换代数 $R$ 可看作投射代数（置
$R_\alpha:=R/I_\alpha$）。并非所有投射代数都是这种类型：若指标 $\alpha$ 的
集合不可数，则从 $R_\alpha$ 的集合论投射极限到 $R_{\alpha_0}$ 的映射未必是
满射\footnote{即使映射 $R_\beta\to R_\alpha$（$\beta\ge\alpha$）都是满射
（正如我们所设）。}。当然，归纳仿射的 $\aleph_0$-归纳概形，等同于拓扑由
$R$ 的可数或有限个开理想系定义的完备拓扑代数。

(ii) 设 $V$ 是 Tate 向量空间（见 4.2.13）。则 $V$（更确切地说，函子
$A\mapsto V\hat\otimes A$）是合理的归纳仿射归纳概形。事实上，$V$ 中每个
c-格都是仿射概形。有 $V=\operatorname{Spf}R$，其中
$R=\varprojlim\operatorname{Sym}(U_\alpha^*)$，$U_\alpha$ 跑遍 $V$ 中
c-格的集合。

若 $X$ 是合理归纳概形，则对 $x\in X(\mathbb{C})$，$X$ 在 $x$ 处的切空间
$\Theta_x$ 是 Tate 向量空间：$\Theta_x$ 的拓扑由 $X$ 的含 $x$ 的合理子概形
在 $x$ 处的切空间定义。于是若 $H$ 是合理群归纳概形，则其李代数
$\operatorname{Lie}H$ 是 Tate 向量空间范畴中的李代数。

(iii) 对如上之 $V$，记 $Gr(V)$ 为 $V$ 中 c-格的「空间」。更确切地说，
$Gr(V)$ 是把交换代数 $A$ 送到 $V\hat\otimes A$ 中 c-格（按 4.2.14 的意义）
的集合的函子。显然 $Gr(V)$ 是归纳固有、形式光滑的归纳概形（事实上，它是
所有 c-格对 $U_1\subset U_2\subset V$ 对应的 $U_2/U_1$ 的 Grassmann 概形
之并）。

\SourcePage{303}

(iv) 设 $K$ 是局部域，$O\subset K$ 是相应的局部环（故
$K\simeq\mathbb{C}((t))$，$O\simeq\mathbb{C}[[t]]$）。对任意「空间」$Y$，
我们有「空间」$Y(O)\subset Y(K)$，其定义为 $Y(O)(A):=
Y(A\hat\otimes O)$、$Y(K)(A):=Y(A\hat\otimes K)$（这里
$A\hat\otimes O=A[[t]]$，$A\hat\otimes K=A((t))$）。设 $Y$ 是仿射概形。则
$Y(O)$ 也是仿射概形，而 $Y(K)$ 是归纳仿射的 $\aleph_0$-归纳概形。若 $Y$
有限型，则 $Y(K)$ 合理。若 $Y$ 光滑，则 $Y(O)$ 与 $Y(K)$ 都形式光滑。

设 $G$ 是仿射代数群，$\mathfrak{g}$ 是其李代数。考虑群归纳概形 $G(K)$。
有 $\operatorname{Lie}(G(K))=\mathfrak{g}(K)=\mathfrak{g}\otimes K$，
$\operatorname{Lie}(G(O))=\mathfrak{g}(O)=\mathfrak{g}\otimes O$。

(v) 设 $G$ 是合理群归纳概形，且 $G_{red}$ 是仿射群概形。函子
$G\mapsto(\operatorname{Lie}G,G_{red})$ 是如上诸 $G$ 的范畴与 Harish-Chandra
对的范畴之间的等价。对归纳概形 $X$，$G$ 在 $X$ 上的作用，等同于
$(\operatorname{Lie}G,G_{red})$ 在 $X$ 上的作用。类似地，$G$-模等同于
$(\operatorname{Lie}G,G_{red})$-模，等等。

\bdp{7.11.3} 在归纳概形的框架下，有两种不同的方式定义 $\mathcal{O}$-模；
相应的对象分别称为 $\mathcal{O}^{p}$-模与 $\mathcal{O}^{!}$-模。我们从较直接
（但较不重要）的 $\mathcal{O}^{p}$-模概念开始\footnote{这里「$p$」代表
「投射极限」。}，它对任意「空间」$X$ 都有意义（见 7.11.1）。

$X$ 上的 $\mathcal{O}^{p}$-模 $P$ 是这样一条规则：它把交换代数 $A$ 与点
$\phi\in X(A)$ 送到 $A$-模 $P_\phi$，并把代数的态射 $f:A\to B$ 送到
$B$-模的等同 $f_P:B\otimes_f P_\phi\xrightarrow{\ \sim\ }P_{f\phi}$。这些等同
须与 $f$ 的复合相容。若 $X=\varinjlim X_\alpha$ 是归纳概形，则这样的 $P$
等同于：$X_\alpha$ 上（拟凝聚）$\mathcal{O}$-模 $P_{X_\alpha}$ 的族，连同对
$\alpha\le\beta$ 的等同 $i^*_{\alpha\beta}P_{X_\beta}=P_{X_\alpha}$，且满足
显然的传递性。若每个 $P_\phi$（或每个 $P_{X_\alpha}$）都平坦，则称 $P$ 平坦。
$X$ 上可逆 $\mathcal{O}^{p}$-模（又称线丛）的定义类似。

\SourcePage{304}

记 $X$ 上 $\mathcal{O}^{p}$-模的范畴为 $\mathcal{M}^{p}(X,\mathcal{O})$。这是
张量 $\mathbb{C}$-范畴。$\mathcal{M}^{p}(X,\mathcal{O})$ 中的单位对象是函数的
「层」$\mathcal{O}_X$。注意 $\mathcal{M}^{p}(X,\mathcal{O})$ 未必是阿贝尔
范畴。平坦 $\mathcal{O}^{p}$-模的范畴
$\mathcal{M}^{pfl}(X,\mathcal{O})$ 是（Quillen 意义下的）正合范畴。

对任意 $P,P'\in\mathcal{M}^{p}(X,\mathcal{O})$，向量空间
$\operatorname{Hom}(P,P')$ 带有显然的拓扑；态射的复合是连续的。特别地，
$\Gamma(X,P):=\operatorname{Hom}(\mathcal{O}_X,P)$ 是拓扑向量空间，且是拓扑环
$\Gamma(X,\mathcal{O}_X)$ 上的模。

\emph{注.} (i) 若把 $\mathcal{O}$-模换成任一在仿射概形范畴上纤维化的范畴，
上述定义仍有意义。例如可考虑左 $\mathcal{D}$-模（又称带可积联络的
$\mathcal{O}$-模）；归纳概形上相应的对象称为（左）$\mathcal{D}^{p}$-模。

(ii) 若 $X$ 是归纳仿射的 $\aleph_0$-归纳概形，
$X=\operatorname{Spf}R=\varinjlim\operatorname{Spec}R/I_\alpha$
（见 7.11.2(i)），则 $X$ 上的 $\mathcal{O}^{p}$-模等同于完备且分离的拓扑
$R$-模 $P$，使得 $I_\alpha P\subset P$ 的闭包构成拓扑的基。

\bdp{7.11.4} 下面转向 $\mathcal{O}^{!}$-模。这里必须设 $X$ 是合理归纳概形。
$X$ 上的 $\mathcal{O}^{!}$-模 $M$ 是这样一条规则：它把合理子概形 $Y\subset X$
送到拟凝聚 $\mathcal{O}_Y$-模 $M_{(Y)}$，并对 $Y\subset Y'$ 给出态射
$M_{(Y)}\to M_{(Y')}$。这些态射把 $M_{(Y)}$ 与
$i^{!}_{YY'}M_{(Y')}:=\operatorname{Hom}_{\mathcal{O}_{Y'}}
(\mathcal{O}_Y,M_{(Y')})$ 等同起来，且满足显然的传递性条件%
\footnote{我们需要考虑合理子概形，才能保证 $i^{!}$ 保持拟凝聚性。}。若写
$X=\varinjlim X_\alpha$ 且诸 $X_\alpha$ 合理，则只需考虑诸 $X_\alpha$，而不必
考虑所有合理子概形。$X$ 上的 $\mathcal{O}^{!}$-模构成阿贝尔范畴
$\mathcal{M}(X,\mathcal{O})$。注意对任意闭子概形 $Y\subset X$，范畴
$\mathcal{M}(Y,\mathcal{O})$ 是 $\mathcal{M}(X,\mathcal{O})$ 的全子范畴，且对
子商封闭；又对任一 $\mathcal{O}^{!}$-模 $M$，有
$M=\varinjlim M_{(X_\alpha)}$。

范畴 $\mathcal{M}(X,\mathcal{O})$ 是张量范畴
$\mathcal{M}^{p}(X,\mathcal{O})$ 上的模。具体地说，对
$M\in\mathcal{M}(X,\mathcal{O})$、$P\in\mathcal{M}^{p}(X,\mathcal{O})$，其张量
积 $M\otimes P\in$
\SourcePage{305}$\mathcal{M}(X,\mathcal{O})$ 是
$\varinjlim_\alpha M_{(X_\alpha)}
\otimes_{\mathcal{O}_{X_\alpha}}P_{X_\alpha}$。函子
$\otimes:\mathcal{M}(X,\mathcal{O})\times
\mathcal{M}^{pfl}(X,\mathcal{O})\to\mathcal{M}(X,\mathcal{O})$ 双正合。

对 $\mathcal{O}^{!}$-模 $M$，我们把其整体截面的空间定义为
$\Gamma(X,M):=\varinjlim_\alpha\Gamma(X_\alpha,M_{(X_\alpha)})$。函子
$\Gamma(X,\cdot)$ 左正合。

\emph{注.} (i) 诸范畴 $\mathcal{M}(Y,\mathcal{O})$ 连同函子
$i^{!}_{YY'}$ 一起，构成 $X$ 的合理子概形范畴（有序集）上的纤维化范畴；
而 $\mathcal{M}(X,\mathcal{O})$ 是其 Cartesian 截面的范畴。

(ii) 若 $X=\operatorname{Spf}R$ 且投射代数 $R$ 是拓扑代数（见 7.11.2），则
$X$ 上的 $\mathcal{O}^{!}$-模等同于离散 $R$-模（这里「离散」是指 $R$-作用
关于 $M$ 上的离散拓扑连续）。

(iii) 若 $P$ 平坦，则
$(M\otimes P)_{(X_\alpha)}=M_{(X_\alpha)}\otimes P_{X_\alpha}$。

\bdp{7.11.5} 设群归纳概形（或任意群「空间」）$K$ 作用在 $X$ 上。则对每个
交换代数 $A$，群 $K(A)$ 作用在 $\operatorname{Spec}A\times X$ 上。对
$M\in\mathcal{M}(X,\mathcal{O})$，$K$ 在 $M$ 上的作用由
$\mathcal{O}_{\operatorname{Spec}A}\boxtimes M\in
\mathcal{M}(\operatorname{Spec}A\times X,\mathcal{O})$ 上的 $K(A)$-作用给出，
且要求对任意态射 $A\to A'$ 相应的作用相容。记 $X$ 上 $K$-等变
$\mathcal{O}^{!}$-模的范畴为
$\mathcal{M}(K\backslash X,\mathcal{O})$。$K$-等变 $\mathcal{O}^{p}$-模的
定义留给读者。

\bdp{7.11.6} 7.10 的所有基本定义与结果（拓扑 $X_{cr}$、$\mathcal{D}$-晶体、
晶体 $\mathcal{O}^*$-torsor、扭 $\mathcal{D}$-晶体、基本函子性这些概念的定义）对任意
归纳有限型的归纳概形 $X$ 都显然有意义。于是从 $\mathcal{D}$-晶体的观点看，
$\mathcal{D}$-模理论自动推广到归纳概形的框架。

本节剩余部分要讨论的是 $\mathcal{D}$-模的通常进路（微分算子环等），它在我们
的归纳概形形式光滑时适用。比较 $\mathcal{D}$-晶体与 $\mathcal{D}$-模两种
框架的结果 7.10.12、7.10.29、7.10.32，对形式光滑的归纳概形逐字成立。

% 尾部接 p.306（本单元末句完结于 7.11.6 的「remain literally true for
% formally smooth ind-schemes」。）
%% end units/ch07_u07.tex
%% begin units/ch07_u08.tex
% ch07_u08 —— 原书 p.306–p.315（7.11.7–7.11.15(ii)）

\SourcePage{306}

以下不再提 D-晶体。本书正文中我们使用通常的 D-模（我们遇到的归纳概形是仿射 Grassmann 概形，它们形式光滑）。不过要注意，正是 D-晶体的进路让下述事实成为显然（当 $Y$ 是 Schubert 胞腔时我们会用到它）：设 $Y$ 是有限型概形，$i:Y\hookrightarrow X$ 是 $Y$ 到上述形式光滑 $X$ 的闭嵌入。则 $X$ 上支撑在 $Y$ 上（集合论意义下）的 D-模范畴只依赖于 $Y$，而不依赖于 $i$ 与 $X$。事实上，该范畴与 $X$ 上 D-晶体的范畴典范等同。

\bdp{7.11.7} 我们来解释归纳概形情形下什么是微分算子。设 $X$ 是 ind-finite type 的归纳概形。对 $X$ 上 $\mathcal{O}^!$-模 $M$，置
\begin{equation}\tag{348}
\operatorname{Der}(\mathcal{O}_X,M):=\varinjlim\operatorname{Der}(\mathcal{O}_Y,M_{(Y)})=\varinjlim\operatorname{Hom}(\Omega_Y,M_{(Y)}).
\end{equation}
这里 $Y$ 是 $X$ 的闭子概形。把 $\operatorname{Der}(\mathcal{O}_X,M)$ 视为 $X$ 上 $\mathcal{O}^!$-模。类似地，置
\begin{equation}\tag{349}
\mathcal{D}(M)=\operatorname{Diff}(\mathcal{O}_X,M):=\varinjlim\operatorname{Diff}(\mathcal{O}_Y,M_{(Y)}).
\end{equation}
把微分算子层 $\operatorname{Diff}(\mathcal{O}_Y,M_{(Y)})$ 视为 [BB93] 意义下的「微分 $\mathcal{O}_Y$-双模」，即 $Y\times Y$ 上支撑在对角线上的 $\mathcal{O}$-模（集合论意义下）。
% [Tutor Note: 原书 p.306 此处引号两侧都印成右引号（」differential O_Y-bimodule」），底本如此，译文改用中文引号。]

于是 $\mathcal{D}(M)$ 是 $X\times X$ 上支撑在对角线上的 $\mathcal{O}^!$-模（集合论意义下）。按两个 $\mathcal{O}_X$-模结构中的任一个，都可把它视为 $X$ 上 $\mathcal{O}^!$-模。注意 $\mathcal{D}(M)$ 带典范递增滤链 $\mathcal{D}_\cdot(M)$，其中 $\mathcal{D}_i(M)$ 是支撑在对角线第 $i$ 个无穷小邻域上的截面子模；等价地，$\mathcal{D}_i(M)=\varinjlim\operatorname{Diff}_i(\mathcal{O}_Y,M_{(Y)})$ 是阶 $\le i$ 的微分算子子模。有 $\mathcal{D}_0(M)=M$、$\bigcup\mathcal{D}_i(M)=\mathcal{D}(M)$，且 $\operatorname{gr}_\cdot\mathcal{D}(M)$ 上两个 $\mathcal{O}^!$-模结构重合。有显然的嵌入 $\operatorname{Der}(\mathcal{O}_X,M)\subset\mathcal{D}_1(M)$。

\SourcePage{307}

现设 $X$ 形式光滑。下一个命题中，我们按左 $\mathcal{O}$-模结构把 $\mathcal{D}(M)$ 视为 $X$ 上 $\mathcal{O}^!$-模。

\bdp{7.11.8} \emph{命题.} (i) 函子 $\operatorname{Der}(\mathcal{O}_X,\cdot)$、$\mathcal{D}$、$\mathcal{D}_i$ 正合，且与正向极限交换。故存在平坦 $\mathcal{O}^p$-模 $\Theta_X$、$\mathcal{D}_X$，以及 $\mathcal{D}_X$ 的一列平坦子模 $\mathcal{D}_{iX}$ 构成的滤链，使
\[
\operatorname{Der}(\mathcal{O}_X,M)=M\otimes\Theta_X,\qquad\mathcal{D}(M)=M\otimes\mathcal{D}_X,\qquad\mathcal{D}_i(M)=M\otimes\mathcal{D}_{iX}.
\]
(ii) 有典范等同 $\operatorname{gr}_\cdot\mathcal{D}_X=\operatorname{Sym}^\cdot\Theta_X$。

\emph{注.} 7.12.12 中将证明 $\Theta_X$、$\mathcal{D}_X$、$\mathcal{D}_{iX}$ 这些 $\mathcal{O}^p$-模是 Raynaud--Gruson 意义下的 Mittag-Leffler 模（见 7.12.1、7.12.2、7.12.9）。若 $X$ 是 $\aleph_0$-归纳概形，则这些 $\mathcal{O}^p$-模限制到 $X$ 的子概形上局部自由（更精确的陈述见 7.12.13）。

\emph{证明.} (i) 这些函子显然左正合，且与正向极限交换。$\operatorname{Der}(\mathcal{O}_X,\cdot)$ 的右正合性由 $X$ 的形式光滑性推出（把导子 $\mathcal{O}_X\to M$ 标准地解读为态射 $\operatorname{Spec}(\operatorname{Sym}^\cdot M/\operatorname{Sym}^{\ge 2}M)\to X$）。于是得到 $\Theta_X\in M^{pfl}(X,\mathcal{O})$。

(ii) 我们定义典范同构\footnote{在一般情形（基域特征可以非零）下，须把 $\operatorname{Sym}^\cdot$ 换成 $\Gamma^\cdot$：对任意平坦 $A$-模 $P$，把 $\Gamma^n(P)$ 定义为 $P^{\otimes n}$ 中的 $S_n$-不变量。注意（因 $P$ 是投射模的归纳极限）$\Gamma^n(P)$ 平坦，且对任意 $A$-模 $M$ 有 $(M\otimes P^{\otimes n})^{S_n}=M\otimes\Gamma^n(P)$。}
\begin{equation}\tag{350}
\sigma_\cdot:\operatorname{gr}_\cdot\mathcal{D}(M)\xrightarrow{\sim}M\otimes\operatorname{Sym}^\cdot\Theta_X.
\end{equation}
这显然推出命题。

注意对任意 $n\ge 0$，显然的态射 $M\otimes\Theta_X^{\otimes n}\to\varinjlim\operatorname{Hom}(\Omega_Y^{\otimes n},M_{(Y)})$ 是同构（用到 $\Omega_Y$ 凝聚这一事实）。因此 (350) 等价于等同
\begin{equation}\tag{351}
\sigma_n:\operatorname{gr}_n\mathcal{D}(M)\xrightarrow{\sim}\varinjlim\operatorname{Hom}(\operatorname{Sym}^n\Omega_Y,M_{(Y)}).
\end{equation}

\SourcePage{308}

我们的 $\sigma_n$ 是下列映射的归纳极限：
\[
\sigma_{nY}:\operatorname{gr}_n\operatorname{Diff}(\mathcal{O}_Y,M_{(Y)})\to\operatorname{Hom}(\operatorname{Sym}^n\Omega_Y,M_{(Y)}),
\]
它们定义如下。有 $\operatorname{Diff}_n(\mathcal{O}_Y,M_{(Y)})=\operatorname{Hom}_{\mathcal{O}_Y}(\mathcal{O}_{Y\times Y}/\mathcal{I}^{n+1},M_{(Y)})$，其中 $\mathcal{I}\subset\mathcal{O}_{Y\times Y}$ 是对角线理想（源按某个投影映射视为 $\mathcal{O}_Y$-模）。又 $\mathcal{I}/\mathcal{I}^2=\Omega_Y$，故 $\mathcal{I}^n/\mathcal{I}^{n+1}$ 是 $\operatorname{Sym}^n\Omega_Y$ 的商，而 $\sigma_{nY}$ 来自映射 $\operatorname{Sym}^n\Omega_Y\to\mathcal{I}^n/\mathcal{I}^{n+1}\subset\mathcal{O}_{Y\times Y}/\mathcal{I}^{n+1}$。

余下要证 $\sigma_n$ 是同构；可设 $n\ge 1$。显然 $\sigma_{nY}$ 单射，故 $\sigma_n$ 单射。为看出 $\sigma_n$ 满，考虑概形 $Z:=\operatorname{Spec}(\operatorname{Sym}^\cdot\Omega_Y/\operatorname{Sym}^{\ge n+1}\Omega_Y)$。其子概形 $\operatorname{Spec}(\operatorname{Sym}^\cdot\Omega_Y/\operatorname{Sym}^{\ge 2}\Omega_Y)=\operatorname{Spec}(\mathcal{O}_{Y\times Y}/\mathcal{I}^2)\subset Y\times Y\subset Y\times X$ 的嵌入，借 $X$ 的形式光滑性延拓为 $Y$ 上的态射 $i:Z\to Y\times X$。容易看出 $i$ 是闭嵌入。存在闭子概形 $Y'\subset X$，使 $Y\subset Y'$ 且 $Z\subset Y\times Y'$。于是 $Z$ 是 $Y'\times Y'$ 中对角线第 $n$ 个无穷小邻域的子概形。因此得到嵌入 $\operatorname{Hom}(\operatorname{Sym}^n\Omega_Y,M_{(Y)})\subset\operatorname{Hom}_{\mathcal{O}_Y}(\mathcal{O}_Z,M_{(Y)})\subset\operatorname{Diff}_n(\mathcal{O}_{Y'},M_{(Y')})$。把它们与 $\sigma_{nY'}$ 复合，得到嵌入 $\operatorname{Hom}(\operatorname{Sym}^n\Omega_Y,M_{(Y)})\subset\operatorname{Hom}(\operatorname{Sym}^n\Omega_{Y'},M_{(Y')})$。这意味着 $\sigma_n$ 满。\hfill$\Box$

\bdp{7.11.9} 为解释归纳概形上的 D-模是什么，用微分双模的语言比较方便。

设 $X$ 是任意合理归纳概形。$X$ 上的 Diff-双模 $D$（参看 [BB93]）是这样的规则：对任意合理子概形 $Y\subset X$，指派 $Y\times X$ 上一个 $\mathcal{O}^!$-模 $D_Y$，它在 $Y\times X$ 中支撑在（集合论意义下的）对角线 $Y$ 上；对 $Y\subset Y'$ 有等同 $D_{Y'}\otimes\mathcal{O}_Y\xrightarrow{\sim}D_Y$，这些等同按显然意义传递。

Diff-双模的范畴 $M^{di}(X,\mathcal{O})$ 是幺半 $\mathbb{C}$-范畴。即对 $D,D'\in M^{di}(X,\mathcal{O})$，其张量积 $D\otimes D'$ 由
\[
(D\otimes D')_Y:=\varinjlim(D_Y)_{(Y\times Y')}\otimes_{\mathcal{O}_{Y'}}D'_{Y'}
\]
定义。我们的 $\mathcal{O}_X$ 是 $M^{di}(X,\mathcal{O})$ 的单位对象\SourcePage{309}（见下面的注 (i)）。范畴 $M(X,\mathcal{O})$ 是右 $M^{di}(X,\mathcal{O})$-模：对 $\mathcal{O}^!$-模 $M$ 有 $M\otimes D=\varinjlim M_{(Y)}\otimes D_Y$，其中把 $M_{(Y)}\otimes D_Y$ 按 $D_Y$ 上的右 $\mathcal{O}^!$-模结构视为 $X$ 上 $\mathcal{O}^!$-模。
% [Tutor Note: 原书 p.309 作 「a right M^{di}(X,O)-Module」，Module 无故大写；底本排版异常，译文按普通名词处理。]

\emph{注.} (i) $X$ 上 $\mathcal{O}^p$-模等同于支撑在（概形论意义下的）对角线上的微分 $\mathcal{O}_X$-双模。故有幺半范畴的全忠实嵌入 $M^p(X,\mathcal{O})\subset M^{di}(X,\mathcal{O})$。它与 7.11.4、7.11.9 中在 $M(X,\mathcal{O})$ 上的作用相容。
% [Tutor Note: 原书 p.309 此处作 「the Actions on M(X,O)」，Actions 无故大写；底本排版异常，译文按普通名词处理。]

(ii) 遗忘\footnote{忘掉右 $\mathcal{O}$-模结构。}函子 $M^{di}(X,\mathcal{O})\to M^p(X,\mathcal{O})$ 忠实，故可把 Diff-双模看作带额外结构的 $\mathcal{O}^p$-模。称 Diff-双模平坦，指它作为 $\mathcal{O}^p$-模平坦。平坦 Diff-双模的范畴是正合范畴（参看 7.11.3）。

$X$ 上的 Diff-代数指幺半范畴 $M^{di}(X,\mathcal{O})$ 中的幺结合代数 $D$。$X$ 上的 $D^!$-模指 $M(X,\mathcal{O})$ 中的（必为右）$D$-模 $M$。常把这样的 $M$ 简称为 D-模。$D$-模范畴记作 $M(X,D)$，它是阿贝尔范畴。

\emph{注.} (i) 遗忘函子 $M(X,D)\to M(X,\mathcal{O})$ 有左伴随函子，即 $M\mapsto M\otimes D$。

(ii) 范畴 $M^p(X,\mathcal{O})$ 以显然方式成为左 $M^{di}(X,\mathcal{O})$-模。故可考虑 $D^p$-模，即在 $M^p(X,\mathcal{O})$ 中的左 $D$-模。

对 $D\in M^{di}(X,\mathcal{O})$，置 $\Gamma(X,D):=\varprojlim\Gamma(Y\times X,D_Y)$；这是拓扑向量空间。有显然的连续映射 $\Gamma(X,D)\otimes\Gamma(X,D')\to\Gamma(X,D\otimes D')$。对 $M\in M(X,\mathcal{O})$ 有类似的映射 $\Gamma(X,M)\otimes\Gamma(X,D)\to\Gamma(X,M\otimes D)$。因此对 Diff-代数 $D$，$\Gamma(X,D)$ 是拓扑环；且对任意 $D$-模 $M$，向量空间 $\Gamma(X,M)$ 是离散 $\Gamma(X,D)$-模。

设群归纳概形（或任意群「空间」）$K$ 作用在 $X$ 上。$K$ 在 Diff-代数 $D$ 上的弱\footnote{关于强作用见 [BB93]。}作用定义如下。\SourcePage{310}对任意交换代数 $A$，群 $K(A)$ 作用在 $\operatorname{Spec}A\times X$ 上。于是 $K$ 在 $D$ 上的弱作用是下述规则：对每个 $A$，把该作用提升为 $\operatorname{Spec}A\times X$ 上 Diff-代数 $\mathcal{O}_{\operatorname{Spec}A}\boxtimes D$ 上的作用。对任意态射 $A\to A'$，相应作用须按显然方式相容。若 $M$ 是 $D$-模，则 $K$ 在 $M$ 上的弱作用指 $K$ 在 $M$ 上作为 $\mathcal{O}^!$-模的作用（见 7.11.4），使 $D$-作用态射 $M\otimes D\to M$ 与 $K$-作用相容。弱 $K$-等变 $D$-模的范畴记作 $M(K\backslash X,D)$。

\bdp{7.11.10} 对合理 $\aleph_0$-归纳概形 $X$ 上的微分双模与代数，还有更具体的「层论」看法。\footnote{$\aleph_0$ 假设使我们能用拓扑代数代替投射代数；见 7.11.2(i)。}我们分两步解释。
% [Tutor Note: 原书 p.310 的引号两侧都印成右引号（」sheaf-theoretic」），底本如此，译文改用中文引号。]

(i) 设 $X_{red}$ 是概形，即 $X$ 是形式概形\footnote{形式光滑仿射 $\aleph_0$-形式概形（ind-finite type）的描述见 7.12.22、7.12.23。}。于是 $X$ 的底层拓扑空间有定义，$\mathcal{O}_X$ 是拓扑代数的层。任意 Diff-双模 $D$ 给出一个拓扑 $\mathcal{O}_X$-双模的层 $\varprojlim D_{X_\alpha}$，仍滥用记号记作 $D$。它满足下列性质：
\begin{itemize}[label=--]
\item $D$ 上拓扑的基由 $\mathcal{I}\cdot D$ 的闭包构成，其中 $\mathcal{I}\subset\mathcal{O}_X$ 是开理想；该拓扑完备且分离。
\item 商 $D/\mathcal{I}\cdot D$ 是 $X\times X$ 上支撑在（集合论意义下的）对角线处的 $\mathcal{O}^!$-模。
\end{itemize}
显然 $M^{di}(X,\mathcal{O})$ 等价于这类拓扑 $\mathcal{O}_X$-双模层的范畴。注意 $D\otimes D'=D\hat{\otimes}_{\mathcal{O}_X}D'$。因此 $X$ 上的 Diff-代数等同于 $X$ 上拓扑代数的层 $D$，连同代数层的连续态射 $\epsilon:\mathcal{O}_X\to D$，且要求 $D$ 上的 $\mathcal{O}_X$-双模结构满足上述条件。

\SourcePage{311}

$X$ 上的 D-模等同于离散右 D-模的层，它作为 $\mathcal{O}_X$-模拟凝聚（即它是 $X$ 上 $\mathcal{O}^!$-模）。

(ii) 设 $X$ 是任意合理 $\aleph_0$-归纳概形。对合理子概形 $Y\subset X$，记 $Y^\wedge$ 为 $X$ 沿 $Y$ 的完备化。这是 (i) 中的形式概形。对 $X$ 上的 Diff-双模 $D$，令 $D_{Y^\wedge}$ 为 $D$ 到 $Y^\wedge$ 的（$\mathcal{O}^p$-模）拉回。它是 $Y^\wedge$ 上的 Diff-双模，故如 (i) 可看作 $\mathcal{O}_{Y^\wedge}$-双模的层。若 $Y'\subset X$ 是含 $Y$ 的另一个合理子概形，则有 $\mathcal{O}_{Y'^\wedge}$-双模层的连续态射 $D_{Y'^\wedge}\to D_{Y^\wedge}$；它把 $D_{Y^\wedge}$ 与 $D_{Y'^\wedge}$ 关于由 $\mathcal{I}\cdot D_{Y'^\wedge}$ 的闭包生成的拓扑的完备化等同，其中 $\mathcal{I}\subset\mathcal{O}_{Y'^\wedge}$ 是使 $\operatorname{Spec}(\mathcal{O}/\mathcal{I})_{red}=Y_{red}$ 的开理想。这些态射满足显然的传递性。显然 $X$ 上的 Diff-双模等同于这类数据。

因此 $X$ 上的 Diff-代数 $D$ 可看作如下数据：
\begin{itemize}[label=--]
\item 一簇拓扑代数的层 $D_{Y^\wedge}$，同时对任意合理子概形 $Y\subset X$ 定义满足 (i) 中条件的态射 $\epsilon_{Y^\wedge}:\mathcal{O}_{Y^\wedge}\to D_{Y^\wedge}$；
\item 对 $Y\subset Y'$ 有连续态射 $r_{YY'}:D_{Y'^\wedge}\to D_{Y^\wedge}$，它如上把 $D_{Y^\wedge}$ 与 $D_{Y'^\wedge}$ 的完备化等同。要求相容性 $r_{YY'}\epsilon_{Y'^\wedge}=\epsilon_{Y^\wedge}$、$r_{YY''}=r_{YY'}r_{Y'Y''}$。
\end{itemize}
用这种语言描述 D-模，留给读者。

\emph{注.} 对 Diff-代数 $D$，拓扑代数 $\Gamma(X,D)$ 是拓扑代数 $\Gamma(Y,D_{Y^\wedge})$ 的射影极限。

\bdp{7.11.11} \emph{关键例子.} 设 $X$ 是 ind-finite type 的形式光滑归纳概形。考虑 7.11.8(i) 中定义的 $\mathcal{O}^p$-模 $\mathcal{D}_X$。于是对子概形 $Y\subset X$，$\mathcal{O}_Y$-模 $(\mathcal{D}_X)_Y$ 是 $\mathcal{D}(\mathcal{O}_Y):=\varinjlim\operatorname{Diff}(\mathcal{O}_{Y'},\mathcal{O}_Y)$，带其左 $\mathcal{O}_Y$-模结构。我们的 $\mathcal{D}_X$ 带显然的 Diff-双模结构。微分算子的复合使 $\mathcal{D}_X$ 成为 $X$ 上的 Diff-代数。由 7.11.8，$\mathcal{D}_X$ 带典范环滤链 $\mathcal{D}_{iX}$，使 $\operatorname{gr}_\cdot\mathcal{D}_X=\operatorname{Sym}^\cdot\Theta_X$。\SourcePage{312}拓扑代数 $\Gamma(X,\mathcal{D}_X)$ 称为 $X$ 上整体微分算子的环。

$\mathcal{D}_X$-模范畴记作 $M(X,\mathcal{D})$，或简记 $M(X)$。若群「空间」$K$ 作用在 $X$ 上，则 $\mathcal{D}_X$ 带典范弱 $K$-作用（由结构的平移定义）。于是有弱 $K$-等变 D-模的范畴 $\mathcal{M}(K\backslash X,\mathcal{D}_X)=\mathcal{M}(K\backslash X)$。

\emph{扭版本.} 本书正文中还需要扭微分算子环（alias tdo）、这类环的族及其上的模。相应定义是有限维情形下通常定义的直接修改（例如参看 [BB93]）。下面具体描述我们需要的 tdo 特例。

设 $X$ 如上，$\mathcal{L}$ 是 $X$ 上的线丛（见 7.11.3）。

a. 作用在 $\mathcal{L}$ 上的微分算子的 Diff-代数 $\mathcal{D}_{\mathcal{L}}$ 的定义与 $\mathcal{D}_X$ 完全相同，只须在 (349) 中把 $\mathcal{D}(M)$ 换成
\[
\mathcal{D}_{\mathcal{L}}(M)=\operatorname{Diff}(\mathcal{L},M\otimes\mathcal{L}):=\varinjlim\operatorname{Diff}(\mathcal{L}_Y,M_{(Y)}\otimes\mathcal{L}_Y);
\]
命题 7.11.8（及其证明）无须任何改动依然成立。等价地，$\mathcal{D}_{\mathcal{L}}=\mathcal{L}\otimes\mathcal{D}_X\otimes\mathcal{L}^{\otimes-1}$。

b. 如下定义 $X$ 上的 Diff-代数 $\mathcal{D}_{\mathcal{L}^h}$。设 $\pi:X^\sim\to X$ 是对应于 $\mathcal{L}$ 的 $\mathbb{G}_m$-torsor（故 $X^\sim$ 是 $\mathcal{L}$ 去掉零截面）。考虑 $X$ 上的 Diff-代数 $\mathcal{D}^\sim:=\pi_*\mathcal{D}_{X^\sim}$（故对子概形 $Y\subset X$ 有 $(\mathcal{D}^\sim)_Y:=\pi_*((\mathcal{D}_{X^\sim})_{\pi^{-1}Y})$）。$\mathcal{D}_{X^\sim}$ 上的弱 $\mathbb{G}_m$-作用给出 $\mathcal{D}^\sim$ 上的弱 $\mathbb{G}_m$-作用（相对 $X$ 上的平凡 $\mathbb{G}_m$-作用）。我们的 $\mathcal{D}_{\mathcal{L}^h}$ 是 $\mathcal{D}^\sim$ 中 $\mathbb{G}_m$-不变量的子代数。

记 $h$ 为 $\mathcal{D}_{\mathcal{L}^h}$ 中对应于 $-t\frac{d}{dt}\in\operatorname{Lie}\mathbb{G}_m$ 的作用的整体截面。则 $\mathcal{D}_{\mathcal{L}^h}$ 是 $h$ 在 $\mathcal{D}^\sim$ 中的中心化子。注意对任意子概形 $Y\subset X$，$\mathcal{L}_{Y^\wedge}$ 的平凡化（它在 $Y$ 上局部存在）给出等同
\[
\mathcal{D}_{\mathcal{L}^hY^\wedge}\xrightarrow{\sim}\mathcal{D}_{Y^\wedge}\otimes_{\mathbb{C}}\mathbb{C}[h].
\]

\emph{注.} (i) 考虑 $\mathcal{O}^p$-模 $\pi_*(\mathcal{O}_{X^\sim})=\oplus\mathcal{L}^{\otimes n}$。它带 $\mathcal{D}_{\mathcal{L}^h}$ 的作用，且该作用保持分次。$\mathcal{D}_{\mathcal{L}^h}$ 在 $\mathcal{L}^{\otimes n}$ 上的作用把 $\mathcal{D}_{\mathcal{L}^h}/(h-n)\mathcal{D}_{\mathcal{L}^h}$ 与 $\mathcal{D}_{\mathcal{L}^{\otimes n}}$ 等同。

\SourcePage{313}

(ii) 设 $M^\sim$ 是 $X^\sim$ 上弱 $\mathbb{G}_m$-等变的 D-模。置 $M:=(\pi_*M^\sim)^{\mathbb{G}_m}$；它是 $\mathcal{D}_{\mathcal{L}^h}$-模。函子 $\mathcal{M}(\mathbb{G}_m\backslash X^\sim)\to\mathcal{M}(X,\mathcal{D}_{\mathcal{L}^h})$，$M^\sim\mapsto M$，是范畴等价。

\bdp{7.11.12} 我们来解释 D-Ω 复形在归纳概形情形下的相互作用。先定义 Ω-复形。这里设 $X$ 是任意合理归纳概形。

对任意合理子概形 $Y\subset Y'$，有交换 DG 代数的满态射 $\Omega_{Y'}\to\Omega_Y$。$X$ 上的 $\Omega^!$-复形 $F$（或简称 Ω-复形）是这样的规则：对合理子概形 $Y\subset X$ 指派 DG $\Omega_Y$-模 $F_{[Y]}$，并对 $Y\subset Y'$ 给出 $\Omega_{Y'}$-模的态射 $F_{[Y]}\to F_{[Y']}$，这些态射把 $F_{[Y]}$ 与 $i^!_{\Omega YY'}F_{[Y']}:=\operatorname{Hom}_{\Omega_{Y'}}(\Omega_Y,F_{[Y']})$ 等同，且满足显然的传递条件。假定 $F^i_{[Y]}$ 作为 $\mathcal{O}_Y$-模拟凝聚。如 7.11.4，只须考虑 $X_\alpha$，不必考虑所有合理 $Y$。如 7.2.1 的注，这样的 $F$ 等同于一个 $\mathcal{O}^!$-模复形，其微分是阶 $\le 1$ 的微分算子。记 $C(X,\Omega)$ 为 $\Omega^!$-复形的 DG 范畴。

若 $f:Y\to X$ 是归纳概形的可表示拟紧态射（故 $Y=\varinjlim Y_\alpha$，其中 $Y_\alpha:=f^{-1}(X_\alpha)$），则有拉回函子
\[
f^\cdot_\Omega:C(X,\Omega)\to C(Y,\Omega),\qquad f^\cdot_\Omega(F):=\varinjlim\Omega_{Y_\alpha}\otimes_{f^{-1}\Omega_{X_\alpha}}F_\alpha.
\]
若 $f$ 满且形式光滑，则 $f^\cdot_\Omega$ 满足下降性质。

设群「空间」$K$ 作用在 $X$ 上。$K$ 在 $X$ 上 Ω-复形 $F$ 上的作用定义如下：对任意 $g\in K(A)$，把 $g$ 在 $\operatorname{Spec}A\times X$ 上的作用提升为 $\mathcal{O}_{\operatorname{Spec}A}\otimes F\in C(\operatorname{Spec}A\times X,\Omega)$ 上的作用；且须满足显然的相容性。相应范畴记作 $C(K\backslash X,\Omega)$。

\emph{注.} (i) 设 $K$ 是群归纳概形，于是有李代数 $\operatorname{Lie}K$。7.6.4 中用算子 $i_\xi$ 给出的 $K_\Omega$-作用的定义，立即适用于当前情形。$K_\Omega$-等变 Ω-复形的范畴记作 $C(K\backslash X,\Omega)$。

\SourcePage{314}

(ii) 若我们的 $K$ 是仿射群概形，则 $K_\Omega$-等变 Ω-复形等同于这样的 Ω-复形 $F$：它带 $K\times X$ 上 Ω-复形的同构 $m^\cdot_\Omega F=p^\cdot_X F$，且满足通常的条件（见 7.6.5）。

\bdp{7.11.13} 设 $X$ 是 ind-finite type 的形式光滑归纳概形。记 $C(X,\mathcal{D})$ 为 $X$ 上 D-模复形（简称 D-复形）的 DG 范畴。有 DG 函子
\begin{equation}\tag{352}
\mathcal{D}:C(X,\Omega)\to C(X,\mathcal{D}),
\end{equation}
它把 Ω-复形 $F$ 送到 D-复形 $\mathcal{D}F$，其分量为 $(\mathcal{D}F)^n:=\mathcal{D}(F^n)=F^n\otimes\mathcal{D}_X$（见 7.11.8），微分由公式 $d(a):=d_F\circ a$ 定义（这里 $a\in\mathcal{D}(F^n)=\operatorname{Diff}(\mathcal{O}_X,F^n)$）。该函子有右伴随函子
\begin{equation}\tag{353}
\Omega:C(X,\mathcal{D})\to C(X,\Omega),
\end{equation}
它可以具体描述如下。对子概形 $Y\subset X$，有 D-复形 $DR_Y:=\mathcal{D}(\Omega_Y)$。它也是左 DG $\Omega_Y$-模。于是对 D-复形 $M$ 有
\[
\Omega M=\varinjlim\operatorname{Hom}(DR_Y,M)=\bigcup\operatorname{Hom}(DR_Y,M).
\]
引理 7.2.4 连同其证明仍然成立。如 7.2.5，有上同调函子 $H^\cdot_D:C(X,\Omega)\to\mathcal{M}(X)$，$H^\cdot_D(F)=H^\cdot(\mathcal{D}F)$，以及相应的 D-拟同构概念。$\mathcal{D}$、$\Omega$ 的伴随态射是拟同构与 D-拟同构\footnote{D-模的 de Rham 复形不从下界有界，这一事实并不妨碍上述图景。}。

\bdp{7.11.14} 称一个 $\mathcal{O}^!$-复形或 $\mathcal{O}^!$-模具有\emph{拟紧支撑}，若它在某个闭子概形的补集上为零。同一定义适用于 D-复形与 Ω-复形。我们用下标「c」标记相应的范畴。函子 $\mathcal{D}$ 与 $\Omega$ 保持相应的满 DG 子范畴 $C_c(X,\Omega)\subset C(X,\Omega)$、$C_c(X,\mathcal{D})\subset C(X,\mathcal{D})$。

\SourcePage{315}

为确保我们的导出范畴正是所要的（即它们具有好的函子性质），我们额外假设\emph{对角线态射 $X\to X\times X$ 是仿射的}（参看 7.3.1）。例如，只要假设 $X$ 分离就够。

记 $D(X,\mathcal{O})$ 为 $C_c(X,\mathcal{O})$ 的同伦范畴关于拟同构的局部化；这是以 $M_c(X,\mathcal{O})$ 为核心的 t-范畴。类似地定义 $D(X,\mathcal{D})$（假设 $X$ 是 ind-finite type 的形式光滑概形）；这是以 $M_c(X)$ 为核心的 t-范畴。设 $D(X,\Omega)$ 是 $C_c(X,\Omega)$ 的同伦范畴关于 D-拟同构的局部化。函子 $\mathcal{D}$ 与 $\Omega$ 给出 $D(X,\mathcal{D})$ 与 $D(X,\Omega)$ 的典范等同。故照例把这些等同后的范畴简记为 $D(X)$\footnote{为得到以 $\mathcal{M}(X)$ 为核心的 t-范畴，可以考虑那些为具有拟紧支撑的子复形之并的复形；但要保证该范畴有好的函子性质，须假设 $X$ 满足某些额外条件（例如存在形式光滑满态射 $Y\to X$，其中 $Y$ 归纳仿射）。由所有复形构成的范畴具有不好的同调与函子性质。注意通常的补救办法——只考虑从下界有界的 Ω-复形——在这里不起作用（D-模的 de Rham 复形不满足该条件）。}。

称具有拟紧支撑的 $\mathcal{O}^!$-模 $F$ 是\emph{松弛}的，若对任意使得 $F$ 支撑在 $Y^\wedge$ 上的闭子概形 $Y\subset X$，以及 $Y^\wedge$ 上任意平坦 $\mathcal{O}^p$-模 $P$，都有 $H^a(X,P\otimes F)=0$（$a>0$）。称 $\mathcal{O}^!$-、D- 或 Ω-复形 $F$ 松弛，若每个 $\mathcal{O}^!$-模 $F^i$ 都松弛。有下述与 7.3.8 平行的引理：

\bdp{7.11.15} \emph{引理.} i) 对任意 $F'\in C_c(X,\Omega)$，存在 D-拟同构 $F'\to F$，使 $F$ 松弛，且 $F$ 与 $F'$ 的支撑相同。
% [Tutor Note: 原书 p.315 此项标号印作 i)，下一项印作 (ii)，底本如此，照录。]

(ii) 若 $f:X\to X'$ 是归纳概形的形式光滑仿射态射，则函子
\[
f^\cdot_\Omega:C_c(X',\Omega)\to C_c(X,\Omega),\qquad f^\cdot:C_c(X,\Omega)\to C_c(X',\Omega)
\]
把松弛复形送到松弛复形。

% 尾部接 p.316（承接 7.11.15 的 (iii)）
%% end units/ch07_u08.tex
%% begin units/ch07_u09.tex
% ch07_u09 —— 原书 p.316--p.325（7.11.15(iii) -- 7.12.17 中段）
% 承接 ch07_u08：本单元自 7.11.15 的 (iii) 续起，不写 \bdchap。
% 尾部接 p.326（7.12.17 中（356）式之后的句子未完）。

\SourcePage{316}

(iii) 若 $F_1,F_2$ 分别是 $X_1,X_2$ 上的松弛复形，则 $F_1\boxtimes F_2$ 是 $X_1\times X_2$ 上的松弛 Ω-复形。

\emph{证明.} 按显然的方式修改 7.3.8 的证明。\hfill$\square$

由此可见，导出范畴 $D(X)$ 可以用松弛复形来定义。

\bdp{7.11.16} 归纳概形的任意态射 $f:X\to Y$ 给出推前函子
\[
f_{\cdot}:C(X,\Omega)\to C(Y,\Omega),
\]
它保持子范畴 $C_c$。$f_{\cdot}$ 把具拟紧支撑的松弛复形之间的 D-拟同构仍送到 D-拟同构，这一点留给读者核验（参看 7.3.9、7.3.11(ii)）。于是右导出函子 $Rf_{\cdot}=f_{*}:D(X)\to D(Y)$ 良定义：若 $F$ 是具拟紧支撑的松弛复形，则 $f_{*}F=f_{\cdot}F$。由于 $f_{\cdot}$ 把松弛复形送到松弛复形，$f_{*}$ 与 $f$ 的复合相容。

对 $M\in D(X,\mathcal{D})$，记 $M_{\mathcal{O}}\in D(X,\mathcal{O})$ 为把同一个 $M$ 看作 $\mathcal{O}^!$-模复形。有典范的积分态射
\[
i_f:Rf_{\cdot}(M_{\mathcal{O}})\to(f_{*}M)_{\mathcal{O}},
\]
它位于 $D(Y,\mathcal{O})$ 中，定义同 7.2.11；该态射与 $f$ 的复合相容。

\bdp{7.11.17} 按 7.6.1 定义 Hecke 幺半范畴 $\mathcal{H}$。出发点是 ind-affine 群归纳概形 $G$ 及其仿射群子概形 $K\subset G$。假设 $G/K$（关于 fpqc 拓扑的层商）是 ind-finite type 的归纳概形；它自动形式光滑，且对角态射仿射。显然 $G$ 是合理归纳概形，$K$ 是它的合理子概形。考虑 $G$ 上具拟紧支撑的 $(K\times K)_{\Omega}$-等变 $\Omega^!$-复形构成的 DG 范畴 $\mathcal{H}^c$（见 7.11.12 注 (i)）。由下降，这样的复形等同于 $G/K$ 上或 $K\backslash G$ 上的 $K_{\Omega}$-等变容许 $\Omega^!$-复形。相应的 D-拟同构概念彼此等价。我们的 $\mathcal{H}$ 就是相应的 D-导出范畴。
% [Tutor Note: 原书 p.316 此处印作 "is a ind-scheme of ind-finite type"，冠词用 a 而不用 an；底本如此，中文无对应形式。]

\SourcePage{317}

7.6.1 的构造在我们的情形下完全适用。故 $\mathcal{H}^c$ 是 DG 幺半范畴，$\mathcal{H}$ 是三角幺半范畴。

\bdp{7.11.18} 设概形 $Y$ 带 $G$-作用，且存在开拟紧子概形的递增族 $U_0\subset U_1\subset\ldots$，使 $Y=\bigcup U_i$，并使得对某个合理群子概形 $K_i\subset G$，$K_i$ 在 $U_i$ 上的作用自由，而 $K_i\backslash U_i$ 是光滑概形（特别地有限型）。于是叠 $\mathcal{B}=K\backslash Y$ 光滑（它被诸概形 $(K_i\cap K)\backslash U_i$ 覆盖），其对角态射仿射，故可用 7.3.12 的 $D(\mathcal{B})$ 定义。

为在 $D(\mathcal{B})$ 上定义 $\mathcal{H}$-作用，照 7.6.1 的办法做，并作如下修改（这些修改源于 $Y$ 与 $G$ 可能非拟紧）。可设上述 $U_i$ 都 $K$-不变；置 $\mathcal{B}_i=K\backslash U_i\subset\mathcal{B}$。取松弛 Ω-复形 $F=\bigcup F_n\in C_a(K\backslash G/K,\Omega)$（故 $F_n$ 的支撑 $S_n$ 拟紧）与 $T\in C(\mathcal{B}_{\cdot},\Omega)$。取 $j(n,i)$ 为关于 $n$、$i$ 都递增的序列，使 $S_n^{-1}\cdot U_i\subset U_{j(n,i)}$。考虑 $\mathcal{B}_i$ 上的 Ω-复形
\[
(F_n\circledast T)_i:=\bar m_{U_i}\cdot p^{\cdot}_{U_i\Omega}(F_n\boxtimes T_{j(n,i)})|_{\mathcal{B}_i}
\]
与
\[
(F_n\circledast T)'_i:=\bar m_{U_i}\cdot p^{\cdot}_{U_i\Omega}(F_n\boxtimes T_{j(n+1,i)})|_{\mathcal{B}_i}.
\]
有显然的态射 $(F_n\circledast T)'_i\to(F_{n+1}\circledast T)_i$ 与 $(F_n\circledast T)'_i\to(F_n\circledast T)_i$，后者是拟同构。置
\[
(F\circledast T)_i:=\operatorname{Cone}\big(\oplus(F_n\circledast T)_i\to\oplus(F_n\circledast T)_i\big),
\]
其中箭头是上述态射的（逐分量）差。
% [Tutor Note: 原书 p.317 此处 Cone 的两个直和项都印成 $(F_n\circledast T)_i$；按上下文第二项应为 $(F_n\circledast T)'_i$，底本如此，保留原样。]

这些 $(F\circledast T)_i$ 按显然的方式构成 $C(\mathcal{B},\Omega)$ 的对象 $F\circledast T$。留给读者核验：作为 $D(\mathcal{B})$ 的对象，$F\circledast T$ 不依赖辅助数据（$U_i$ 与 $j(n,i)$）的选取；且 $\circledast$ 是 $D(\mathcal{B})$ 上的 $\mathcal{H}$-作用。
% [Tutor Note: 原书 p.317 本节两处把「作用」印作「Action」（H-Action），无故大写；底本如此，译文照常作「作用」。]

\bdsec{7.12. 归纳概形与 Mittag-Leffler 模}
% [Tutor Note: 原书 p.317 的 7.12 节标题与正文同行（run-in 粗体）；底本如此，译文将标题单独排出。]

Raynaud 与 Gruson [RG] 引入了 Mittag-Leffler 模这一非凡的概念。本节说明：平坦 Mittag-Leffler 模这一概念，在某种意义上正是 ind-finite type 形式光滑归纳概形概念的线性化版本（见 7.12.12、7.12.14、7.12.15）。
\SourcePage{318}%
利用可数生成的平坦 Mittag-Leffler 模必投射这一事实，我们描述 ind-finite type 的形式光滑仿射 $\aleph_0$-形式概形（见 7.12.22、7.12.23）。

读者可以跳过本节：它的结果在本书其余部分用不到（我们收录它们，只为澄清形式光滑归纳概形这一概念）。

7.11 中我们假定「归纳概形」意即「$\mathbb{C}$ 上的归纳概形」（这其实无关紧要）。本节我们宁可放弃这个假设。

\bdp{7.12.1} 设 $A$ 是环\footnote{我们假设 $A$ 交换，但在 7.12.1--7.12.8 中这并不必要（只需按显然的方式在「模」字前加上「左」或「右」）。}。记 $\mathcal{C}$ 为有限表现 $A$-模的范畴。据 [RG] p.69，称 $A$-模 $M$ 为 \emph{Mittag-Leffler 模}，若任意态射 $f:F\to M$（$F\in\mathcal{C}$）都可分解为 $F\xrightarrow{u}G\to M$（$G\in\mathcal{C}$），使得对 $f$ 的任意分解 $F\xrightarrow{u'}G'\to M$（$G'\in\mathcal{C}$），存在态射 $\varphi:G'\to G$ 使 $u=\varphi u'$。

\bdp{7.12.2} 设 $M=\varinjlim M_i$，$i\in I$，其中 $I$ 是有向有序集，$M_i\in\mathcal{C}$。据 loc.cit.，$M$ 是 Mittag-Leffler 模当且仅当：对每个 $i\in I$ 存在 $j\ge i$，使得对每个 $k\ge i$，态射 $u_{ij}:M_i\to M_j$ 都可分解为 $\varphi_{ijk}u_{ik}$，其中 $\varphi_{ijk}:F_k\to F_j$。
% [Tutor Note: 原书 p.318 此处印作 $\varphi_{ijk}:F_k\to F_j$，而上文用的是 $M_i$ 记号；底本如此，保留原样。]

$I$ 是滤过范畴时也一样；若 $I$ 是「$\mathcal{C}$ 的对象到 $M$ 的全部态射」构成的范畴，而 $F_i\in\mathcal{C}$ 是态射 $i$ 的源，则上述陈述就是同义反复。

\bdp{7.12.3} 上述归纳系统 $(M_i)$ 的这一性质，在把 $\mathcal{C}$ 换成任意范畴 $\mathcal{C}'$ 后仍有意义。若 $\mathcal{C}'$ 是集合范畴的对偶，即有集合的射影系统 $(E_i,u_{ij}:E_j\to E_i)$，则由 EGA $0_{\mathrm{III}}$ 13.1.2 得到下述 Mittag-Leffler 条件：对每个 $i\in I$ 存在 $j\ge i$，使得 $u_{ij}(E_j)=u_{ik}(E_k)$ 对所有 $k\ge j$ 成立。该条件成立当且仅当射影系统 $(E_i,u_{ij})$ 等价于射影系统 $(\tilde E_{\alpha},\tilde u_{\alpha\beta})$，其中映射 $\tilde u_{\alpha\beta}$ 都是%
\SourcePage{319}%
满的。证「仅当」只需对充分大的 $j$ 置 $\tilde E_i:=u_{ij}(E_j)$。

\bdp{7.12.4} 设 $M=\varinjlim M_i$，$M_i\in\mathcal{C}$。据 [RG]，$M$ 是 Mittag-Leffler 模当且仅当：对任意从 $\mathcal{C}$ 到集合范畴的反变函子 $\Phi$，射影系统 $(\Phi(M_i))$ 都满足 Mittag-Leffler 条件（证「若」时考虑函子 $\Phi(N)=\operatorname{Hom}(N,\prod_i M_i)$ 或 $\tilde\Phi(N)=\coprod_i\operatorname{Hom}(N,M_i)$）。设 $M$ 平坦，置 $M^{*}_i=\operatorname{Hom}(M_i,A)$。据 [RG]，$M$ 是 Mittag-Leffler 模当且仅当射影系统 $(M^{*}_i)$ 满足 Mittag-Leffler 条件。$M_i$ 投射时这显然。一般情形由 Lazard 引理得出（存在由有限生成投射模构成、与 $(M_i)$ 等价的归纳系统）。

\bdp{7.12.5} 考虑从 $A$-模范畴到阿贝尔群范畴的两类函子：

1) 对 $A$-模 $M$ 有函子
\begin{equation}\tag{354}
L\mapsto L\otimes_A M;
\end{equation}

2) 对 $A$-模的射影系统 $N_i$（$i$ 属于某个有向有序集）有函子
\begin{equation}\tag{355}
L\mapsto\varinjlim_i\operatorname{Hom}(N_i,L).
\end{equation}
% [Tutor Note: 原书 p.319 的（354）（355）与 p.325 的（356）编号都排在左缘；底本如此，译文按通行的右缘排法排印。]

\bdp{7.12.6} \emph{命题.}
\begin{enumerate}[label=(\roman*),leftmargin=2.2em]
\item 函子 (354) 同构于形如 (355) 的函子，当且仅当 $M$ 平坦；
\item 函子 (354) 同构于射影系统 $(N_i)$ 对应的函子 (355)（转移映射 $N_j\to N_i$，$i\le j$，都满），当且仅当 $M$ 是平坦 Mittag-Leffler 模；
\item 对应满转移映射射影系统的函子 (355) 同构于形如 (354) 的函子，%
\SourcePage{320}%
当且仅当 (355) 正合且各 $N_i$ 有限生成。
\end{enumerate}

\emph{证明.} 若 (354) 与 (355) 同构，则 (354) 左正合，故 $M$ 平坦。若 $M$ 平坦，则由 Lazard 引理 $M=\varinjlim P_i$，其中 $P_i$ 投射且有限生成；于是对应 $N_i=P^{*}_i$ 的 (355) 同构于 (354)。(i) 得证。由 (i) 推出 (ii)：对上述 $P_i$，射影系统 $(P^{*}_i)$ 等价于转移映射满的 $(N_i)$，当且仅当 $(P^{*}_i)$ 满足 Mittag-Leffler 条件（见 7.12.3）。证 (iii)：形如 (354) 的函子正是右正合且与无限直和交换（从而与归纳极限交换）的加性函子。形如 (355) 的函子右正合当且仅当它正合。若各 $N_i$ 有限生成，则 (355) 与无限直和交换。若转移映射满且 (355) 与归纳极限交换，则各 $N_i$ 有限生成。\hfill$\square$

\bdp{7.12.7} 由 7.12.6，平坦 Mittag-Leffler 模「等同于」有限生成模的射影系统 $(N_i)$ 的等价类（转移映射 $N_j\to N_i$（$i\le j$）都满，且函子 (355) 正合）。更确切地说，$M=\varinjlim_i\operatorname{Hom}(N_i,A)$（此时 (354) 与 (355) 同构）。

\bdp{7.12.8} \emph{定理}(Raynaud--Gruson).（D. Lazard 呢？据 [RG] p.73，此想法可溯至 D. Lazard 学位论文第一章的定理 3.1 与 3.2，Bull. Soc. Math. France 第 97 卷 (1969), 81--128；另见他在同刊第 95 卷 (1967), 95--108 的工作。）下列条件等价：
\begin{enumerate}[label=(\roman*),leftmargin=2.2em]
\item $M$ 是平坦 Mittag-Leffler 模；
\SourcePage{321}%
\item $M$ 的每个有限或可数子集都含于某个可数生成的投射子模 $P\subset M$，且 $M/P$ 平坦；
\item $M$ 的每个有限子集都含于某个投射子模 $P\subset M$，且 $M/P$ 平坦。
\end{enumerate}

特别地，投射模都是 Mittag-Leffler 模，而可数生成\footnote{可数生成的假设不可少；见 7.12.24。}的平坦 Mittag-Leffler 模都是投射模。

(iii)$\Rightarrow$(i) 容易。（只须证：若 $F,F'$ 有限表现，$\varphi:F\to F'$，$\psi:F'\to M$，且 $\psi\varphi(F)\subset P$，则存在 $\tilde\psi:F'\to M$ 使 $\tilde\psi(F')\subset P$ 且 $\tilde\psi\varphi=\psi\varphi$；这要用到 $\operatorname{Hom}(L,M)\to\operatorname{Hom}(L,M/P)$ 对有限表现 $L$ 满，特别取 $L=\operatorname{Coker}\varphi$。）

(i)$\Rightarrow$(ii) 见 [RG] p.73--74。关键论证如下。设序列 $P_1\to P_2\to\ldots$，各 $P_i$ 有限生成投射，且射影系统 $(P^{*}_i)$ 满足 Mittag-Leffler 性质。要证 $P:=\varinjlim P_i$ 投射，只须证：对任意正合列 $0\to N'\to N\to N''\to 0$，映射 $\operatorname{Hom}(P,N)\to\operatorname{Hom}(P,N'')$ 满。对每个 $i$，序列 $0\to P^{*}_i\otimes N'\to P^{*}_i\otimes N\to P^{*}_i\otimes N''\to 0$ 正合，问题归结为这些序列的射影极限正合。据 EGA $0_{\mathrm{III}}$ 13.2.2，这由射影系统 $(P^{*}_i\otimes N')$ 的 Mittag-Leffler 性质推出。

\emph{注.} 若指标集 $i$ 不可数，我们就不能\footnote{EGA $0_{\mathrm{III}}$ 13.2.2 的论证基于下述事实：由可数集 $I$ 参数化的非空集射影系统 $(Y_i)_{i\in I}$，若满足 Mittag-Leffler 条件，则其射影极限非空。不可数情形下这不对。例如取不可数集 $S$，对每个有限 $F\subset A$ 记 $Y_F$ 为单射 $F\to\mathbb{N}$ 的集合；映射 $Y_{F'}\to Y_F$（$F'\supset F$）都满，但 $\varprojlim_F Y_F=\emptyset$。}用 EGA $0_{\mathrm{III}}$ 13.2.2。
% [Tutor Note: 原书 p.321 脚注前文设的是不可数集 $S$，此处印作「有限 $F\subset A$」，$A$ 应为 $S$；底本如此，保留原样。]

\SourcePage{322}

$P$ 投射性的另一证明如下（其实是同一个证明的另一版本）。记 $f_i:P_i\to P_{i+1}$。Mittag-Leffler 性质说，把序列 $\{P_i\}$ 换成它的子序列后，存在 $g_i:P_{i+1}\to P_i$ 使 $g_{i+1}f_{i+1}f_i=f_i$。置 $\mathcal{P}:=\oplus_i P_i$。记 $f:\mathcal{P}\to\mathcal{P}$、$g:\mathcal{P}\to\mathcal{P}$ 为 $f_i$、$g_i$ 诱导的算子。则 $gf^2=f$。有正合列
\[
0\to\mathcal{P}\xrightarrow{1-f}\mathcal{P}\to P\to 0.
\]
因 $\mathcal{P}$ 投射，只须证该列分裂，即有 $h:\mathcal{P}\to\mathcal{P}$ 使 $h(1-f)=1$。事实上置 $h=1-(1-g)^{-1}gf$，并用等式 $gf^2=f$。\footnote{D. Arinkin 注意到，先验地显然：若 $f,g$ 是一个（非交换）环 $R$ 的元素，$gf^2=f$，且 $1-g$ 有左逆，则 $1-f$ 有左逆。事实上，记 $\mathbf{1}$ 为 $1$ 在 $R/R(1-f)$ 中的像，则 $f\mathbf{1}=\mathbf{1}$、$gf^2\mathbf{1}=g\mathbf{1}$，故 $g\mathbf{1}=\mathbf{1}$，从而 $\mathbf{1}=0$。}

\bdp{7.12.9} \emph{命题.} 设 $B$ 是 $A$-代数。若 $M$ 是 Mittag-Leffler $A$-模，则 $B\otimes_A M$ 是 Mittag-Leffler $B$-模。若 $B$ 在 $A$ 上忠实平坦，则反之亦然。

这在 [RG] 中有证明。证明容易：把 $M$ 表为有限表现模的归纳极限，再用 7.12.2。

于是，概形上 Mittag-Leffler $\mathcal{O}$-模的概念、归纳概形上 Mittag-Leffler $\mathcal{O}^p$-模的概念，都清楚了。

\bdp{7.12.10} \emph{命题.} 诺特概形 $S$ 上可数型的平坦 Mittag-Leffler $\mathcal{O}$-模 $\mathcal{F}$ 局部自由。若 $S$ 仿射且连通，而 $\mathcal{F}$ 无限型，则 $\mathcal{F}$ 自由。

这是 7.12.8 与下述结果的直接推论。

\bdp{7.12.11} \emph{定理.} 若 $R$ 诺特且 $\operatorname{Spec}R$ 连通，则每个不是有限生成的投射 $R$-模都自由。

本定理由 Bass 证明（见 [Ba63] 的推论 4.5）。

\SourcePage{323}

\bdp{7.12.12} \emph{命题.} 设 $X$ 是域上 ind-finite type 的形式光滑归纳概形。则 $\mathcal{O}^p$-模 $\Theta_X$、$\mathcal{D}_X$、$\mathcal{D}_{iX}$（见 7.11.8）都是平坦 Mittag-Leffler 模。

\emph{证明.} 我们先证 $\mathcal{D}_X$ 到闭子概形 $Y\subset X$ 的限制是平坦 Mittag-Leffler $\mathcal{O}_Y$-模（对 $\Theta_X$ 与 $\mathcal{D}_{iX}$ 的论证相同）。可设 $Y$ 仿射（否则用 $X\setminus F$ 代替 $X$，其中 $F\subset Y$ 适当闭）。据 7.12.6，只须证：
\begin{enumerate}[label=(\roman*),leftmargin=2.2em]
\item $\mathcal{O}_Y$-模范畴上定义的函子 $L\mapsto L\otimes\mathcal{D}_X$ 正合；
\item 它具有 (355) 的形式，其中 $\mathcal{O}_Y$-模 $N_i$ 凝聚。
\end{enumerate}
按定义，$L\otimes\mathcal{D}_X$ 就是 (349) 所定义的层 $\mathcal{D}(L)$，故 (ii) 显然。(i) 我们已在 7.11.8 证明。\hfill$\square$

\bdp{7.12.13} \emph{命题.} 设 $X$ 是域上 ind-finite type 的形式光滑 $\aleph_0$-归纳概形，$Y\subset X$ 是局部闭子概形。则 $\Theta_X$ 到 $Y$ 的限制局部自由。若 $Y$ 仿射且连通，而 $\Theta_X$ 到 $Y$ 的限制无限型，则它自由。

这由 7.12.12 与 7.12.10 推出。

\bdp{7.12.14} \emph{命题.} 设 $A$ 是环，$M$ 是 $A$-模。由 $F_M(R)=M\otimes R$ 定义「$A$-空间」$F_M$（即从 $A$-代数范畴到集合范畴的函子）。则 $F_M$ 是归纳概形当且仅当 $M$ 是平坦 Mittag-Leffler 模；此时 $F_M$ 在 $A$ 上形式光滑，且在 $A$ 上 ind-finite type。

\emph{证明.} 若 $M$ 是平坦 Mittag-Leffler 模，则由 7.12.6(ii) $F_M$ 是归纳概形，由 7.12.6(iii) 它在 $A$ 上 ind-finite type。形式光滑由定义得出。反之，设 $F_M$ 是归纳概形。把 $F_M$ 表为 $\varinjlim\mathcal{S}_i$，其中 $\mathcal{S}_i$ 是 $F_M$ 的闭子概形，都含零截面 $0\in F_M(A)$。记 $N_i$ 为 $\mathcal{S}_i$ 的余切层%
\SourcePage{324}%
到 $0:\operatorname{Spec}A\hookrightarrow\mathcal{S}_i$ 的限制。则函子 (355) 同构于 (354)，故由 7.12.6(ii) $M$ 是平坦 Mittag-Leffler 模。\hfill$\square$
% [Tutor Note: 原书 p.323 此处 "closed subshemes" 系 "closed subschemes" 之误；底本如此，译文照常作「闭子概形」。]

\emph{注.} 若 $M$ 是任意平坦 $A$-模，则它是有限生成投射 $A$-模 $M_i$ 的有向族的归纳极限，故 $F_M=\varinjlim F_{M_i}$ 是广义意义下的归纳概形（态射 $F_{M_i}\to F_{M_j}$ 不必是闭嵌入）。容易看出：若 $F_M$ 是广义意义下的归纳概形，则 $M$ 平坦。

\bdp{7.12.15} \emph{命题.} 设 $(N_i)_{i\in I}$ 是由有向集 $I$ 参数化的有限生成 $A$-模射影系统，其所有转移映射 $N_j\to N_i$（$j\ge i$）都满。置 $\mathbb{A}(N_i):=\operatorname{Spec}\operatorname{Sym}(N_i)$，$\mathcal{S}:=\varinjlim_i\mathbb{A}(N_i)$。归纳概形 $\mathcal{S}$ 在 $A$ 上形式光滑，当且仅当 $\mathcal{S}$ 同构于 7.12.14 中对应某个平坦 Mittag-Leffler 模 $M$ 的归纳概形 $F_M$。

\emph{证明.} $\mathcal{S}$ 形式光滑，当且仅当函子 (355) 正合（把形式光滑的定义用于形如 $A\oplus J$ 的 $A$-代数，其中 $A\cdot J\subset J$、$J^2=0$）。再用 7.12.6(iii)。\hfill$\square$

\bdp{7.12.16} \emph{命题.} 设 $M$ 是平坦 Mittag-Leffler 模，$F_M$ 是 7.12.14 中的归纳概形。下列条件等价：
\begin{enumerate}[label=(\roman*),leftmargin=2.2em]
\item 对应 $F_M$ 的投射代数（见 7.11.2(i)）是拓扑代数；
\item $M$ 是 [RG] 意义下的严格 Mittag-Leffler 模。
\end{enumerate}

据 [RG] p.74，称 $M$ 为\emph{严格 Mittag-Leffler 模}，若对每个 $f:F\to M$（$F\in\mathcal{C}$）存在 $u:F\to G$（$G\in\mathcal{C}$），使 $f=gu$ 且 $u=hf$，其中 $g:G\to M$、$h:M\to G$（回忆 $\mathcal{C}$ 是有限表现模的范畴）。若 $M=\varinjlim M_i$，$M_i\in\mathcal{C}$，而 $u_{ij}:M_i\to M_j$、$u_i:M_i\to M$ 是典范映射，则 $M$ 严格 Mittag-Leffler 当且仅当：对每个 $i$ 存在 $j\ge i$ 使 $u_{ij}=\varphi_{ij}u_j$，其中 $\varphi_{ij}:M\to M_j$。显然，投射模都是严格 Mittag-Leffler 模，%
\SourcePage{325}%
而严格 Mittag-Leffler 模都是 Mittag-Leffler 模。反之一般不真（见 7.12.24）。
% [Tutor Note: 原书 p.324 此处 "strictly" 误拼作「stritly」；底本如此，译文无法保留拼写，照常译作「严格」。]

\emph{证明.} 把 $M$ 表为 $\varinjlim P_i$，其中 $P_i$ 有限生成投射。置 $N_i:=\operatorname{Im}(P^{*}_j\to P^{*}_i)$，$j$ 充分大。考虑下述条件：
\begin{enumerate}[label=(\alph*),leftmargin=2.2em]
\item 映射 $\varphi_i:\varprojlim_r\operatorname{Sym}(N_r)\to\operatorname{Sym}(N_i)$ 满；
\item $\operatorname{Im}\varphi_i\supset N_i$ 对每个 $i$ 成立；
\item 映射 $\varprojlim_r N_r\to N_i$ 对每个 $i$ 满；
\item 对每个 $i$ 存在 $j\ge i$，使 $\operatorname{Hom}(M,A)$ 与 $\operatorname{Hom}(P_j,A)$ 在 $\operatorname{Hom}(P_i,A)$ 中的像相等。
\end{enumerate}
显然 (i)$\Leftrightarrow$(a)$\Leftrightarrow$(b)$\Leftrightarrow$(c)$\Leftrightarrow$(d)。对 $i\le j$ 考虑映射 $u_{ij}:P_i\to P_j$ 与 $u_i:P_i\to M$。为证 (d)$\Leftrightarrow$(ii)，只须证：$\operatorname{Hom}(M,A)$ 与 $\operatorname{Hom}(P_j,A)$ 在 $\operatorname{Hom}(P_i,A)$ 中的像相等，当且仅当 $u_{ij}=\varphi u_j$，其中 $\varphi:M\to P_j$。证「仅当」：$\operatorname{Hom}(M,P_j)$ 与 $\operatorname{Hom}(P_j,P_j)$ 在 $\operatorname{Hom}(P_i,P_j)$ 中的像相等，故 $\mathrm{id}\in\operatorname{Hom}(P_j,P_j)$ 在 $\operatorname{Hom}(P_i,P_j)$ 中的像，是某个 $\varphi\in\operatorname{Hom}(M,P_j)$ 的像。\hfill$\square$

\bdp{7.12.17} 转入形式光滑仿射 $\aleph_0$-归纳概形的结构之前，先讨论 7.11.1 的形式概形定义与 Grothendieck 的定义（见 EGA I）之间的关系。两者即使在仿射情形也不等价。我们意义下的形式仿射概形，是可表为 $\varinjlim\operatorname{Spec}R_\alpha$ 的归纳概形 $X$，其中 $(R_\alpha)$ 是环的射影系统，映射 $u_{\alpha\beta}:R_\beta\to R_\alpha$（$\beta\ge\alpha$）都满，且 $\operatorname{Ker}u_{\alpha\beta}$ 的元素幂零。Grothendieck 则要求可把 $X$ 表为 $\varinjlim\operatorname{Spec}R_\alpha$，使得映射
\begin{equation}\tag{356}
\varprojlim_\beta R_\beta\to R_\alpha
\end{equation}
% 尾部接 p.326（此句在 p.325 末断开，续文为「是满的」一段）。
%% end units/ch07_u09.tex
%% begin units/ch07_u10.tex
% ch07_u10 —— 原书 p.326–p.335（7.12.18 命题 -- 7.12.24 例 (d)）
% 底本：A. Beilinson, V. Drinfeld, Quantization of Hitchin's Integrable System
% and Hecke Eigensheaves，第 7 章。
% 承接：上单元（p.316–p.325）末句断在 "...are surjective"，本单元正文自「是满射」续起。
% 尾部接 p.336。

\SourcePage{326}

是满射\footnote{这比 $u_{\alpha\beta}$ 的满射性更强；例如，若 $M$ 是平坦 Mittag-Leffler $A$-模而非严格 Mittag-Leffler 模，则 7.12.6 的论证表明，$F_M$ 沿零截面的完备化不能表示为 $\varinjlim\operatorname{Spec}R_{\alpha}$，使映射 (356) 为满射。}，且理想 $\operatorname{Ker}u_{\alpha\beta}$ 幂零。
我们意义下的合理 $\aleph_0$-形式概形，就是 EGA I 意义下的形式概形。Grothendieck 意义下的拟紧形式概形，若具有由「defining ideals (English?)」（「Idéaux de définition」；见 EGA I 10.5.1）构成的基本系统，则是我们意义下的形式概形；特别地，EGA I 意义下的 noetherian 形式概形都如此。
% [Tutor Note: 原书 p.326 正文即印作「defining ideals (English?)」，括号中的 (English?) 是作者对英文术语的夹注，并非排版残留；底本如此，保留原样。]

由于我们主要关心域上 ind-finite type 的仿射 $\aleph_0$-形式概形，我们的定义与 EGA I 的定义之间的差别无关紧要。

\bdp{7.12.18} \emph{命题.} 设 $X$ 是 $A$ 上 ind-finite type 的形式光滑 $\aleph_0$-归纳概形，$S\subset X$ 是闭子概形，且 $S\to\operatorname{Spec}A$ 是同构。设 $X_{red}=S_{red}$（特别地，$X$ 是形式概形）。以 $M$ 记相对切层 $\Theta_{X/A}$ 限制到 $S$ 后的整体截面所成的 $A$-模。则 $M$ 是可数生成的投射模，且 $(X,S)$ 同构于归纳概形 $F_M$（见 7.12.14）沿零截面的完备化 $\widehat F_M$。

\emph{注.} 在 $A$ 上 ind-finite type 的形式光滑归纳概形 $X$ 上，$\mathcal{O}^{p}$-模 $\Theta_{X/A}$ 的定义与 $A=\mathbb{C}$ 的情形完全相同（见 7.11.8、7.11.7）。

\emph{证明.} 与 7.12.12 中一样，可证 $M$ 是平坦 Mittag-Leffler 模。$\aleph_0$ 假设蕴涵 $M$ 可数生成。由 7.12.8，$M$ 投射。

把 $X$ 表示为 $\varinjlim X_n$（$n\in\mathbb{N}$），其中 $X_n$ 是 $X$ 的含 $S$ 的闭子概形，且 $X_n\subset X_{n+1}$。设 $X^{(1)}$ 是 $S$ 在 $X$ 中的第一次无穷小邻域，即 $X^{(1)}$ 是 $S$ 在 $X_n$ 中的第一次无穷小邻域之并，
\SourcePage{327}%
$n\in\mathbb{N}$。显然 $X^{(1)}=F^{(1)}_{M}:=0\in F_M$ 的第一次无穷小邻域。

嵌入 $X^{(1)}\to F_M$ 可延拓为态射 $\varphi:X\to\widehat F_M$（为构造 $\varphi$，先构造 $\varphi_n:X_n\to\widehat F_M$，使 $\varphi_n|_{X_{n-1}}=\varphi_{n-1}$，且 $\varphi_n$ 在 $X_n\cap X^{(1)}$ 上的限制是典范嵌入 $X_n\cap X^{(1)}\hookrightarrow F^{(1)}_{M}$；这之所以可行，是因为 $\widehat F_M$ 在 $A$ 上形式光滑）。完全类似地，把嵌入 $F^{(1)}_{M}=X^{(1)}\hookrightarrow X$ 延拓为态射 $\psi:\widehat F_M\to X$。因 $\varphi$ 与 $\psi$ 诱导 $F^{(1)}_{M}$ 与 $X^{(1)}$ 之间的同构，故 $\varphi$ 与 $\psi$ 都是 ind-闭嵌入，且 $\varphi\psi$ 是同构。于是 $\varphi$ 与 $\psi$ 是同构。

\bdp{7.12.19} \emph{例.} 我们将构造一对 $(X,S)$，它满足 7.12.18 中除 $\aleph_0$ 假设以外的全部条件，且 $(X,S)$ 不 $A$-同构于形如 $\widehat F_M$ 的形式概形。

设给定平坦 Mittag-Leffler 模的一个非平凡扩张
\begin{equation}\tag{357}
0\to N'\to N\to L\to 0.
\end{equation}
这样的扩张对「大多数」环 $A$ 都存在；见 7.12.24(b, a$''$, d)。把 (357) 与 $A[t]$ 张量，得到扩张 $0\to N'[t]\to N[t]\to L[t]\to 0$。把此扩张乘以 $t$，得到 $0\to N'[t]\to Q\to L[t]\to 0$。归纳概形 $F_Q$ 在 $A[t]$ 上形式光滑，从而在 $A$ 上形式光滑。设 $S\subset F_Q$ 是零截面复合 $\operatorname{Spec}A\to\operatorname{Spec}A[t]\to F_Q$ 的像。以 $X$ 记 $F_Q$ 沿 $S$ 的完备化。

在证明 $(X,S)$ 具有所需性质之前，我们先更具体地描述 $X$。对 $A$-代数 $R$，$F_Q$ 的一个 $R$-点由两部分构成：一个 $A$-态射 $A[t]\to R$，以及 $Q\otimes_{A[t]}R$ 的一个元素。换言之，$F_Q$ 的 $R$-点由一个三元组 $(n,l,t)$ 给出，其中 $n\in N\otimes_A R$，$l\in L\otimes_A R$，$t\in R$，且
\begin{equation}\tag{358}
\pi(n)=tl,
\end{equation}
这里 $\pi$ 是投影 $N\otimes_A R\to L\otimes_A R$。

\SourcePage{328}

于是 $F_Q$ 是 $F_N\times F_L\times\mathbb{A}^1$ 的闭归纳子概形，由方程 (358) 定义。因此 $X\subset\widehat F_N\times\widehat F_L\times\widehat{\mathbb{A}}^1$ 由同一方程 (358) 定义（这里 $\widehat{\mathbb{A}}^1$ 是 $\mathbb{A}^1$ 在 $0\in\mathbb{A}^1$ 处的完备化）。

现设 $(X,S)$ 与 $\widehat F_M$ 是 $A$-同构的。则 $M$ 就是 $\Theta_{X/A}$ 限制到 $S$ 后的整体截面所成的模。把 (358) 线性化，得到
\begin{equation}\tag{359}
M=N'\oplus L\oplus A\subset N\oplus L\oplus A.
\end{equation}
复合
\begin{equation}\tag{360}
\widehat F_M\xrightarrow{\sim}X\hookrightarrow\widehat F_N\times\widehat F_L\times\mathbb{A}^1
\end{equation}
由「Taylor 级数」$\sum_{n=1}^{\infty}\varphi_n$ 给出，其中 $\varphi_n$ 是次数为 $n$ 的齐次多项式映射 $M\to N\oplus L\oplus A$；显然 $\varphi_1$ 是嵌入 (359)。置 $f=\operatorname{pr}_N\circ\varphi_2$，其中 $\operatorname{pr}_N$ 是投影 $N\oplus L\oplus A\to N$。因 $M=N'\oplus L\oplus A$，二次映射 $M\to N$ 构成的模以双线性映射 $L\times A\to N$（即 $\operatorname{Hom}(L,N)$）构成的模为直和项。$f$ 在 $\operatorname{Hom}(L,N)$ 中的像给出 (357) 的一个分裂（用下述事实：态射 (360) 经过由方程 (358) 定义的归纳子概形 $X\subset\widehat F_N\times\widehat F_L\times\mathbb{A}^1$ 分解）。于是得到矛盾。

\bdp{7.12.20} \emph{命题.} 设 $X$ 是环 $A$ 上形式光滑的归纳概形。设下列两个假设之一成立：(i) $X$ 是 ind-affine 的；(ii) $A$ noetherian，且 $X$ 在 $A$ 上 ind-finite type。则 $X$ 是一族在 $A$ 上形式光滑的 ind-闭 $\aleph_0$-归纳概形的有向并。

\emph{证明.} 只需证明：对闭子概形的任意递增序列 $Y_n\subset X$，存在在 $A$ 上形式光滑的 ind-闭 $\aleph_0$-归纳概形 $Y\subset X$，使 $Y\supset Y_n$ 对所有 $n$ 成立。

设 $X$ 是 ind-affine 的。则每个 $Y_n$ 都是仿射的。把 $Y_n$ 表示为 $A$ 上形式光滑概形 $V_n$ 的闭子概形（例如，把 $Y_n$ 的坐标环表示为 $A$ 上多项式代数的商）。

\SourcePage{329}

设 $Y'_n\subset V_n$ 是 $Y_n$ 在 $V_n$ 中的第一次无穷小邻域。因 $X$ 形式光滑，态射 $Y_n\hookrightarrow X$ 可延拓为态射 $Y'_n\to Z_n\subset X$，其中 $Z_n\subset X$ 是某个闭子概形。置 $Y^{(2)}_n:=Z_1\cup\ldots\cup Z_n$。现把上述构造用于 $(Y^{(2)}_n)$，得到新序列 $(Y^{(3)}_n)$，如此类推。所有 $Y^{(k)}_n$ 之并在 $A$ 上形式光滑。

若 $X$ ind-quasicompact 但非 ind-affine，则上述构造的显然修改给出含所有 $Y_n$ 的 ind-闭 $\aleph_0$-归纳概形 $Y\subset X$，使得对 $A$ 上任意仿射概形 $S$ 与由理想 $\mathcal{I}\subset\mathcal{O}_S$（$\mathcal{I}^2=0$）定义的闭子概形 $S_0\subset S$，每个 $A$-态射 $S_0\to Y$ 都局部地延拓为态射 $S\to Y$。若假设 (ii) 成立，则这蕴涵整体延拓的存在。

\bdp{7.12.21} 我们将描述域 $C$ 上 ind-finite type 的形式光滑仿射 $\aleph_0$-形式概形（据 7.12.20，一般情形在某种意义上可归到 $\aleph_0$ 情形）。首先有如下例子。

(0) 置 $R_{mn}:=C[x_1,\ldots,x_m][[x_{m+n},\ldots,x_{m+n}]]$。则 $\operatorname{Spf}R_{mn}$ 是 $C$ 上形式光滑的仿射 $\aleph_0$-形式概形。
% [Tutor Note: 原书 p.329 此处印作 $C[x_1,\ldots,x_m][[x_{m+n},\ldots,x_{m+n}]]$，上下标重复「m+n」；据上下文及原书 p.330 的同一记号，应为 $[[x_{m+1},\ldots,x_{m+n}]]$。底本如此，保留原样。]

(i) 设 $I\subset R_{mn}$ 是理想，$A:=R_{mn}/I$。以 $\mathcal{I}$ 记 $\operatorname{Spf}R_{mn}$ 上对应于 $I$ 的理想层。当然，$\operatorname{Spf}A$ 是 $C$ 上 ind-finite type 的仿射 $\aleph_0$-形式概形。它形式光滑，当且仅当对每个 $u\in\operatorname{Spf}A$，$\mathcal{I}$ 在 $u$ 处的茎由某些 $f_1,\ldots,f_r\in I$ 生成，且 Jacobi 矩阵 $\bigl(\frac{\partial f_i}{\partial x_j}(u)\bigr)$ 的秩为 $r$。

(ii) 设 $A$ 如 (i)，且 $\operatorname{Spf}A$ 形式光滑。则 $\operatorname{Spf}A[[y_1,y_2,\ldots]]$ 是 $C$ 上 ind-finite type 的形式光滑仿射 $\aleph_0$-形式概形。

在 7.12.22 与 7.12.23 中我们将证明：域上每个连通的、ind-finite type 的形式光滑仿射 $\aleph_0$-形式概形，都同构于例 (i) 或 (ii) 中的形式概形。

\SourcePage{330}

\bdp{7.12.22} \emph{命题.} 设 $X$ 是域 $C$ 上 ind-finite type 的形式光滑的仿射形式概形，且 $\Theta_X$ 凝聚（即 $\Theta_X$ 到 $X$ 的每个闭子概形的限制都有限生成）。则 $X$ 同构于例 7.12.21(i) 中的形式概形。

\emph{证明.} 把 $X$ 表示为 $\varinjlim\operatorname{Spec}A_i$，使当 $i\le j$ 时态射 $A_j\to A_i$ 是核幂零的满射。各代数 $A_i$ 都是有限型的。可设指标集 $i$ 有最小元 $0$。置 $I_i:=\operatorname{Ker}(A_i\to A_0)$。

\emph{引理.} 对每个 $k\in\mathbb{N}$，存在 $i_1$，使态射 $A_i/I_i^k\to A_{i_1}/I_{i_1}^k$ 对一切 $i\ge i_1$ 都是双射。

假定引理，对足够大的 $i$ 置 $A_{(k)}:=A_i/I_i^k$，$I_{(k)}:=\operatorname{Ker}(A_{(k)}\to A_0)$。显然 $A_{(1)}=A_0$，$A_{(k)}=A_{(k+1)}/I^k_{(k+1)}$，$I_{(k)}=I_{(k+1)}/I^k_{(k+1)}$。有 $X=\operatorname{Spf}A$，其中 $A:=\varprojlim A_{(k)}$。取 $A_{(1)}=A_0$ 的生成元 $\bar x_1,\ldots,\bar x_m$，以及 $A_0$-模 $I_{(2)}$ 的生成元 $\bar x_{m+1},\ldots,\bar x_{m+n}$。把 $\bar x_1,\ldots,\bar x_{m+n}$ 提升为 $\tilde x_1,\ldots,\tilde x_{m+n}\in A$。置
\[
R_{mn}:=C[x_1,\ldots,x_m][[x_{m+1},\ldots,x_{m+n}]].
\]
存在唯一的连续同态 $f:R_{mn}\to A$，使 $x_i\mapsto\tilde x_i$。显然 $f$ 是满射。此外，$f$ 诱导满射 $\mathfrak{a}^k\twoheadrightarrow\operatorname{Ker}(A\to A_{(k)})$，其中 $\mathfrak{a}\subset R_{mn}$ 是由 $x_{m+1},\ldots,x_{m+n}$ 生成的理想。故 $f$ 是开映射。因此 $f$ 诱导 $A$ 与 $R_{mn}$ 的某个商之间的拓扑同构。命题得证。

余下证明引理。存在 $i_0$，使对每个 $i\ge i_0$，态射 $\operatorname{Spec}A_{i_0}\to\operatorname{Spec}A_i$ 诱导切空间之间的同构（事实上，因 $\Theta_X$ 到 $\operatorname{Spec}A_0$ 的限制有限生成，对应于 $A_0$-模 $N_i:=\Omega_i\otimes_{A_i}A_0$ 的函子 (355) 同构于函子 $L\mapsto\operatorname{Hom}(Q,L)$，其中 $Q$ 是某个 $A_0$-模；故存在 $i_0$，使 $N_i=N_{i_0}$ 对 $i\ge i_0$ 成立）。可设 $i_0=0$。置 $Y_i:=\operatorname{Spec}A_i/I_i^k$（特别地，$Y_0=\operatorname{Spec}A_0$）。态射 $Y_0\to Y_i$ 诱导切空间之间的同构。

\SourcePage{331}

把 $A_0$ 表示为 $C[x_1,\ldots,x_n]/J$，并置 $\tilde Y_0:=\operatorname{Spec}C[x_1,\ldots,x_n]/J^k$。因 $X$ 形式光滑，态射 $Y_0\hookrightarrow X$ 可延拓为态射 $\tilde Y_0\to X$。其像含于某个 $Y_{i_1}$ 中。我们来证明：对 $i\ge i_1$，嵌入 $\nu:Y_{i_1}\hookrightarrow Y_i$ 是同构。我们有态射 $f:\tilde Y_0\to Y_{i_1}$。另一方面，态射 $Y_0\hookrightarrow\tilde Y_0$ 延拓为 $g:Y_i\to\tilde Y_0$。复合 $\nu fg:Y_i\to Y_i$ 在 $Y_0$ 上诱导恒等。故 $\nu fg$ 有限，且诱导切空间之间的同构。因此 $\nu fg$ 是闭嵌入。因 $Y_i$ noetherian，闭嵌入 $Y_i\to Y_i$ 是同构。故 $\nu fg$ 是同构，从而 $\nu$ 是同构。

\bdp{7.12.23} \emph{命题.} 设 $X$ 是域 $C$ 上连通的、ind-finite type 的形式光滑仿射 $\aleph_0$-形式概形，且 $\Theta_X$ 不凝聚（即 $\Theta_X$ 到 $X_{red}$ 的限制是无限型的）。则 $X$ 同构于例 7.12.21(ii) 中的形式概形。

\emph{证明.} 我们将构造形式光滑态射
\[
X\to\operatorname{Spf}C[[y_1,y_2,\ldots]],
\]
它在 $0\in\operatorname{Spf}C[[y_1,y_2,\ldots]]$ 上的纤维是 7.12.21(i) 中的形式概形。

把 $X$ 表示为 $\varinjlim\operatorname{Spec}A_n$（$n\in\mathbb{N}$），使对每个 $n$，态射 $A_{n+1}\to A_n$ 是核幂零的满射。各代数 $A_n$ 都是有限型的。由 7.12.13，$\Theta_X$ 到 $\operatorname{Spec}A_n$ 的限制自由；其秩可数。这意味着：对每个 $n$，投射系统 $(\Omega_{A_i}\otimes_{A_i}A_n)$（$i\ge n$）等价于投射系统
\[
\ldots\to A_n^3\to A_n^2\to A_n
\]
（这里映射 $A_n^{k+1}\to A_n^k$ 是到前 $k$ 个坐标的投影）。故把序列 $(A_n)$ 换成其子序列后，得到图解
\[
\ldots\twoheadrightarrow\Omega_{A_3}\twoheadrightarrow F_2\twoheadrightarrow\Omega_{A_2}\twoheadrightarrow F_1\twoheadrightarrow\Omega_{A_1}
\]
\SourcePage{332}
其中 $F_n$ 是有限生成的自由 $A_n$-模，且 $A_n$-模 $G_n:=\operatorname{Ker}(F_{n+1}\otimes_{A_{n+1}}A_n\to F_n)$ 也自由。对每个 $n$ 取 $G_n$ 的一组基 $e_{n1},\ldots,e_{nk_n}$。把 $e_{ni}$ 提升为 $\tilde e_{ni}\in\operatorname{Ker}(\Omega_{A_{n+2}}\otimes_{A_{n+2}}A_n\to F_n)\subset\operatorname{Ker}(\Omega_{A_{n+2}}\otimes_{A_{n+2}}A_n\to\Omega_{A_n})$，并把 $\tilde e_{ni}$ 表示为 $df_{ni}$，其中 $f_{ni}\in\operatorname{Ker}(A_{n+2}\to A_2)$。
% [Tutor Note: 原书 p.332 此处印作 $f_{ni}\in\operatorname{Ker}(A_{n+2}\to A_2)$，按上下文应为 $A_n$（$\tilde e_{ni}$ 是 $\Omega_{A_{n+2}}\otimes_{A_{n+2}}A_n$ 中的元素，$df_{ni}$ 的 $f_{ni}$ 应在 $A_{n+2}$ 中而像在 $A_n$ 中）。底本如此，保留原样。]
最后把 $f_{ni}$ 提升为 $\tilde f_{ni}\in A:=\varprojlim_m A_m$，并把 $f_{ni}$（$n\in\mathbb{N}$，$i\le k_n$）排成一个序列 $\varphi_1,\varphi_2,\ldots$。此序列收敛到 $0$，故有连续同态 $C[[y_1,y_2,\ldots]]\to A$，使 $y_i\mapsto\varphi_i$。它诱导态射
\begin{equation}\tag{361}
f:X\to Y:=\operatorname{Spf}C[[y_1,y_2,\ldots]].
\end{equation}
由构造可见，微分
\begin{equation}\tag{362}
df:\Theta_X\to f^*\Theta_Y
\end{equation}
是满射，且其核凝聚（事实上，(362) 到 $\operatorname{Spec}A_1\subset X$ 的限制显然具有这些性质，故到 $\operatorname{Spec}A_n$（$n\in\mathbb{N}$）的限制也具有这些性质）。

\emph{引理.} 形式光滑、ind-finite type 的归纳概形之间的态射 $f:X\to Y$ 形式光滑，当且仅当其微分 (362) 是满射。此时 $\Theta_{X/Y}$ 是 (362) 的核。

假定引理，可见 (361) 形式光滑，且 $\Theta_{X/Y}$ 凝聚。故 (361) 在 $0\in Y$ 上的纤维 $X_0$ 满足命题 7.12.22 的条件。因此 $X_0$ 同构于例 7.12.21(i) 中的形式概形。我们来证明 $X$ 同构于 $\tilde X:=X_0\times Y$。因 $X$ 在 $Y$ 上形式光滑，嵌入 $X_0\hookrightarrow X$ 延拓为 $Y$-态射 $\alpha:\tilde X\to X$。因 $\tilde X$ 在 $Y$ 上形式光滑，嵌入 $X_0\hookrightarrow\tilde X$ 延拓为 $Y$-态射 $\beta:X\to\tilde X$。$\alpha$ 与 $\beta$ 都是 ind-闭嵌入（若有限型概形的态射 $\nu:Y\to Z$ 诱导同构 $Y_{red}\to Z_{red}$，且 $\nu$ 的每个几何纤维既约，则 $\nu$ 是闭嵌入）。$Y$-态射 $\beta\alpha:X_0\times Y\to X_0\times Y$
\SourcePage{333}%
在 $0\in Y$ 上诱导恒等，故 $\beta\alpha$ 是同构。因此 $\alpha$ 与 $\beta$ 是同构，命题得证。

引理的证明是标准的。关于 $\Theta_{X/Y}$ 的陈述由定义推出。为证第一个陈述，取仿射概形 $S$ 及理想 $\mathcal{I}\subset\mathcal{O}_S$（$\mathcal{I}^2=0$），以 $S_0\subset S$ 记对应于 $\mathcal{I}$ 的子概形。对态射 $\psi:S_0\to X$，以 $E_X(S,\mathcal{I},\psi)$（相应地 $E_Y(S,\mathcal{I},\psi)$）记 $\psi$（相应地 $f\psi$）到态射 $S\to X$（相应地 $S\to Y$）的延拓之集。$f$ 的形式光滑性指：对上述所有 $S,\mathcal{I},\psi$，$f_*:E_X(S,\mathcal{I},\psi)\to E_Y(S,\mathcal{I},\psi)$ 是满射。因 $X$ 与 $Y$ 形式光滑，$E_X(S,\mathcal{I},\psi)$ 与 $E_Y(S,\mathcal{I},\psi)$ 非空。据 [Gr67] 的 16.5.14，它们分别是 $V_X(S,\mathcal{I},\psi)$ 与 $V_Y(S,\mathcal{I},\psi)$ 上的 torsor（即非空的仿射空间），其中 $V_X(S,\mathcal{I},\psi):=\operatorname{Hom}(\psi^*\Omega_X,\mathcal{I})=\Gamma(S_0,\psi^*\Theta_X\otimes\mathcal{I})$，$V_Y(S,\mathcal{I},\psi)=\Gamma(S_0,\psi^*f^*\Theta_Y\otimes\mathcal{I})$。映射 $f_*$ 是仿射的，而相应的线性映射 $\Gamma(S_0,\psi^*\Theta_X\otimes\mathcal{I})\to\Gamma(S_0,\psi^*f^*\Theta_Y\otimes\mathcal{I})$ 由 (362) 诱导。故引理的第一个陈述显然成立。

\bdp{7.12.24} \emph{Mittag-Leffler 模的例子.}

(a) 据 [RG] p.77 的 2.4.1，对每个 noetherian $A$ 与投射 $A$-模 $P$，$A$-模 $P^*:=\operatorname{Hom}_A(P,A)$ 是严格 Mittag-Leffler 且平坦的。为证 $P^*$ 严格 Mittag-Leffler，可如下论证：对任意 $f:F\to P^*$（$F$ 有限型），$f^*:P\to F^*$ 的像由某些 $l_1,\ldots,l_n\in F^*$ 生成；这些 $l_i$ 给出 $u:F\to A^n$，使 $f=gu$ 且 $u=hf$，其中 $g:A^n\to P^*$，$h:P^*\to A^n$。

特别地，若 $A$ noetherian，则对每个集合 $I$，$A$-模 $A^I$ 是严格 Mittag-Leffler 且平坦的。

(a$'$) 众所周知，若 $A$ 是 Dedekind 环而非域，则对无限集 $I$，$A^I$ 不投射。事实上，可设 $I$ 可数。固定非零素理想 $\mathfrak{p}\subset A$，考虑由满足 $a_i\to 0$（在 $\mathfrak{p}$-进拓扑下）的元素 $a=(a_i)\in A^I$ 构成的子模 $M$。若 $A^I$ 投射，则局部化 $M_{\mathfrak{p}}$ 自由。因 $M/\mathfrak{p}M$ 的维数可数，$M_{\mathfrak{p}}$ 的秩应可数。但 $M$ 含一个同构于 $A^I$ 的子模，故 $(A^I)_{\mathfrak{p}}$ 的秩应可数。然而 $(A^I)_{\mathfrak{p}}/\mathfrak{p}\cdot(A^I)_{\mathfrak{p}}=(A/\mathfrak{p})^I$ 的维数不可数，矛盾。

\SourcePage{334}

(a$''$) 设 $A$ 在 $\mathbb{Z}$ 上或在域上有限生成\footnote{我们不知道是否只需假定 $A$ noetherian。}。若 $A$ 不是 Artin 环，且 $I$ 无限，则 $A^I$ 不投射：用 (a$'$) 以及 $A$ 上有限 Dedekind 环 $B$ 的存在性。

(b) 若 $L$ 是非投射的平坦 Mittag-Leffler 模，则存在非分裂正合列 $0\to N'\to N\to L\to 0$，其中 $N$ 与 $N'$ 是平坦 Mittag-Leffler 模。事实上，若 $N$ 是投射模，$N\to L$ 是满射，则它不分裂，且 $\operatorname{Ker}(N\to L)$ 是 Mittag-Leffler 的（[RG] p.71 的 2.1.6）。

(c) [RG] 中指出：若
\begin{equation}\tag{363}
0\to A\xrightarrow{f}M'\to M\to 0
\end{equation}
是 $A$-模的非分裂正合列，且 $M$ 平坦、Mittag-Leffler，则 $M'$ 是 Mittag-Leffler 的，但不是严格 Mittag-Leffler 的。事实上，若 $M'$ 严格 Mittag-Leffler，则存在有限表现模 $G$ 与态射 $u:A\to G$，使 $f=gu$ 且 $u=hf$，其中 $g:G\to M'$，$h:M'\to G$。因 $M$ 是有限生成投射模的直极限，可设 $\operatorname{Im}g\subset\operatorname{Im}f$。于是 $gh$ 给出 (363) 的分裂，即得矛盾。

这里还有另一个论证。$F_{M'}$ 在 $0\in F_M$ 上的纤维是 $F_{M'}$ 的闭子概形，典范同构于 $\operatorname{Spec}A\times\mathbb{A}^1$；若 (363) 非分裂，则投影 $\operatorname{Spec}A\times\mathbb{A}^1\to\mathbb{A}^1$ 不能延拓为函数 $F_{M'}\to\mathbb{A}^1$，故由 7.12.16，$M'$ 不是严格 Mittag-Leffler 的。

\SourcePage{335}

(d) 设 $A$ 是 Dedekind 环，既非域也非完备局部环。则据 [RG] p.76，存在非分裂正合列 (363)，其中 $M$ 是平坦的严格 Mittag-Leffler $A$-模。下面是它的构造。以 $K$ 记 $A$ 的分式域。固定非零素理想 $\mathfrak{p}\subset A$，考虑完备化 $\widehat A_{\mathfrak{p}}$、$\widehat K_{\mathfrak{p}}$；则 $\widehat A_{\mathfrak{p}}\ne A$，$\widehat K_{\mathfrak{p}}\ne K$。以 $M$ 记满足 $a_n\in\mathfrak{p}^{-n}$ 且在 $\widehat K_{\mathfrak{p}}$ 中收敛的序列 $(a_n)$ 构成的模；有态射 $\lim:M\to\widehat K_{\mathfrak{p}}$。注意 $M$ 是严格 Mittag-Leffler 模\footnote{$M$ 是 Mittag-Leffler 模这一事实是显然的：$A$ 是 Dedekind 环，$M$ 平坦，且对每个有限维子空间 $V\subset M\otimes K$，模 $V\cap M$ 有限生成}。事实上，据上面的 (a)，$\prod_{n=1}^{\infty}\mathfrak{p}^{-n}$ 是严格 Mittag-Leffler 的，而 $\bigl(\prod_{n=1}^{\infty}\mathfrak{p}^{-n}\bigr)/M$ 平坦，故 $M$ 严格 Mittag-Leffler。我们断言 $\operatorname{Ext}(M,A)\ne 0$，即态射 $\varphi:\operatorname{Hom}(M,K)\to\operatorname{Hom}(M,K/A)$ 不满。更确切地说，设 $l:M\to K/A$ 是 $\lim:M\to\widehat K_{\mathfrak{p}}$ 与态射 $\widehat K_{\mathfrak{p}}\to\widehat K_{\mathfrak{p}}/\widehat A_{\mathfrak{p}}\hookrightarrow K/A$ 的复合。我们将证明 $l\notin\operatorname{Im}\varphi$。

设 $l$ 来自 $\tilde l:M\to K$。$\tilde l$ 在 $\mathfrak{p}^{-n}\subset M$ 上的限制定义 $c_n\in\operatorname{Hom}(\mathfrak{p}^{-n},A)=\mathfrak{p}^n$。则 $\tilde l=\tilde l'$，其中 $\tilde l':M\to K_{\mathfrak{p}}$ 把 $(a_n)\in M$ 送到
% [Tutor Note: 原书 p.335 此处印作 $\tilde l':M\to K_{\mathfrak{p}}$（$K$ 上无 hat），而前文 $\lim:M\to\widehat K_{\mathfrak{p}}$；按上下文应为 $\widehat K_{\mathfrak{p}}$。底本如此，保留原样。]
\begin{equation}\tag{364}
\sum_{n=1}^{\infty}c_na_n+\lim_{n\to\infty}a_n.
\end{equation}
事实上，$\tilde l'-\tilde l$ 是态射 $M/M_0\to\widehat A_{\mathfrak{p}}$，其中 $M_0$ 是使 $a_n=0$ 对足够大的 $n$ 成立的 $(a_n)\in M$ 之集；另一方面，$\operatorname{Hom}(M/M_0,\widehat A_{\mathfrak{p}})=0$，因为 $M/M_0$ 是 $\mathfrak{p}$-可除的（即 $\mathfrak{p}M+M_0=M$）。因 $\tilde l'=\tilde l$，表达式 (364) 对每个序列 $(a_n)\in M$ 都属于 $K\subset\widehat K_{\mathfrak{p}}$。这是不可能的（分别考虑非零 $c_n$ 的个数有限与无限两种情形）。

% 尾部接 p.336
%% end units/ch07_u10.tex
%% begin units/ch07_u11.tex
% ch07_u11 —— 原书 p.336–p.345（第 7 章 7.13 BRST 基础）
% 底本：A. Beilinson, V. Drinfeld, Quantization of Hitchin's Integrable System
% and Hecke Eigensheaves

\SourcePage{336}

\emph{注.} 在 (d) 中我们不得不排除 $A$ 为完备局部环的情形。真正的原因由下述结果说明：
\begin{enumerate}[label=\arabic*),leftmargin=2.6em]
\item 据 [J]，若 $A$ 是完备局部 Noether 环，$M$ 是平坦 $A$-模，$N$ 是有限生成 $A$-模，则 $\operatorname{Ext}(M,N)=0$；
\item 据 [RG]（p.76，2.3.3 的注 4），若 $A$ 是 Artin 环的投射极限（这是 [RG] 中「线性紧」一词的含义吗？），则每个（平坦？）Mittag-Leffler $A$-模都是严格 Mittag-Leffler 的。（[RG] 中没有平坦性假设，但去掉这一假设后他们的论证还正确吗？例如，[RG] 中的 $F_i$ 为何是线性紧的？）
\end{enumerate}

\bdsec{7.13. BRST 基础}

BRST 构造是哈密顿约化的精致版本，它在无穷维情形尤为重要。本文正文两次调用 BRST：先是用它定义 Feigin--Frenkel 同构，随后用它构造 Hecke 性质证明中所用的局部化函子 $L\Delta$。本节简要叙述一般的 BRST 构造；函子 $L\Delta$ 将在下一节研究。

BRST 的常用数学参考文献是 [F84]、[FGZ86]、[KS] 与 [Ak]。我们试图写一份不依赖冗余结构的叙述（如 $\mathbb{Z}$-分次、正规序等）。
% [Tutor Note: 原书 p.336 此处印作 "redundand structures"，"redundand" 系 "redundant" 之误，底本如此，保留原样。]

我们先从有限维情形开始。在关于 Tate 中心扩张的一段插叙之后，再解释无穷维版本。

\bdp{7.13.1} 设 $F$ 是有限维向量空间。以 $\mathrm{Cl}^{\cdot}=\mathrm{Cl}^{\cdot}_F$ 记 $F\oplus F^{*}$ 的 Clifford 代数，其分次使 $F$ 的度为 $-1$、$F^{*}$ 的度为 $1$。我们把 $\mathrm{Cl}^{\cdot}$ 视为分次向量空间的张量范畴中的代数\footnote{带「超」交换约束。}。置 $\mathrm{Cl}^{\cdot}_i:=\Lambda^{\le i}F\cdot\Lambda F^{*}\subset\mathrm{Cl}^{\cdot}$。则 $\mathrm{Cl}^{\cdot}_0=\Lambda^{\cdot}F^{*}\subset\mathrm{Cl}^{\cdot}_1\subset\ldots$ 是 $\mathrm{Cl}^{\cdot}$ 上的环滤链。\emph{经典 Clifford 代数} $\mathcal{C}l^{\cdot}=\mathcal{C}l^{\cdot}_F:=\operatorname{gr}\mathrm{Cl}^{\cdot}$ 是交换的（作为分次代数），故它是\SourcePage{337}%
通常意义下的泊松代数。

置 $\mathcal{C}l^{\cdot}_i:=\operatorname{gr}_i\mathrm{Cl}^{\cdot}$。交换分次代数 $\mathcal{C}l^{\cdot}$ 由 $F=\mathcal{C}l^{-1}_1$ 与 $F^{*}=\mathcal{C}l^{0}_1$ 自由生成。泊松括号 $\{\,,\,\}$ 在 $F$ 与 $F^{*}$ 上为零，而对 $f\in F$、$f^{*}\in F^{*}$ 有 $\{f,f^{*}\}=f^{*}(f)$。

子空间 $\mathcal{C}l^{0}_1$ 是 $\mathcal{C}l^{\cdot}$ 的李子代数；它使 $F$ 与 $F^{*}$ 正规化，相应的伴随作用把它与 $\operatorname{End}_F$ 及 $\operatorname{End}_{F^{*}}$ 等同。令 $E^{Lie}=\operatorname{End}^{Lie}_F$ 表示把 $\operatorname{End}_F$ 视为李代数。则 $E^{\flat}=\operatorname{End}^{\flat}_F:=\mathrm{Cl}^{0}_1$ 是 $E^{Lie}$ 被 $\mathbb{C}$ 的中心扩张。

\emph{注.} (i) $\mathrm{Cl}$ 在 $\Lambda F^{*}\xrightarrow{\sim}\mathrm{Cl}/\mathrm{Cl}\cdot F$ 上的作用把它与「奇」向量空间 $F^{odd}$ 上的微分算子代数等同。$\mathrm{Cl}$ 上的滤链就是通常按微分算子阶数给出的滤链，故 $\mathcal{C}l^{\cdot}$ 是 $F^{odd}$ 的余切丛上函数的泊松代数。

(ii)（只在有限维情形有效）扩张 $\operatorname{End}^{\flat}_F$ 分裂（分裂方式不唯一）。事实上，有分裂 $s',s'':E^{Lie}\to E^{\flat}$，它们分别把 $E^{Lie}$ 与 $F^{*}\cdot F$ 及 $F\cdot F^{*}$ 等同。其余任一分裂都等于 $s_{\lambda}=\lambda s'+(1-\lambda)s''$，其中 $\lambda\in\mathbb{C}$。例如 $s_{1/2}$ 是「酉」分裂，它也可如下定义。注意 $\mathrm{Cl}$ 带有典范反自同构（作为\emph{分次}代数），它在 $F$ 与 $F^{*}$ 上是恒等。它保持 $\mathrm{Cl}^{0}_1$，而「酉」分裂就是 $-1$ 特征子空间。

\bdp{7.13.2} 这里是 BRST 构造的「经典」版本。设 $\mathfrak{n}$ 是有限维李代数，$\mathcal{R}$ 是泊松代数，$l^{c}:\mathfrak{n}\to\mathcal{R}$ 是李代数态射\footnote{「c」表示「classical」（经典）。}。置 $\mathcal{C}l^{\cdot}:=\mathcal{C}l^{\cdot}_{\mathfrak{n}}$。$\mathfrak{n}$ 的伴随作用给出李代数态射 $a^{c}:\mathfrak{n}\to\mathcal{C}l^{0}_1$。置 $\mathcal{A}^{\cdot}:=\mathcal{C}l^{\cdot}\otimes\mathcal{R}$；这是泊松分次代数。它还带有附加分次 $\mathcal{A}^{\cdot}_{(i)}:=\mathcal{C}l^{\cdot}_i\otimes\mathcal{R}$，它与乘法相容（但与泊松括号不相容）。我们有李代数态射 $\mathcal{L}ie:\mathfrak{n}\to\mathcal{A}^{0}$，$n\mapsto\mathcal{L}ie_n:=a^{c}(n)\otimes 1+1\otimes l^{c}(n)$。下文对 $n\in\mathfrak{n}$，以 $i^{c}_n$ 记 $\mathcal{C}l^{-1}_1\subset\mathcal{A}^{-1}_{(1)}$ 中相应的元素。有 $\{\mathcal{L}ie_{n_1},i^{c}_{n_2}\}=i^{c}_{[n_1,n_2]}$。

\SourcePage{338}

下列关键引理及其「量子」版本 7.13.7 本质上属于 Akman [Ak]。

\bdp{7.13.3} \emph{引理.} 存在唯一元素 $Q^{c}=Q^{c}_{\mathcal{A}}\in\mathcal{A}^{1}$，使得对任意 $n\in\mathfrak{n}$ 有 $\{Q^{c},i^{c}_n\}=\mathcal{L}ie_n$。事实上 $Q^{c}\in\mathcal{A}^{1}_{(\le1)}$。有 $\{Q^{c},Q^{c}\}=0$。

\emph{证明.} 把 $\mathcal{A}$ 视为 $\Lambda\mathfrak{n}$-模，其中 $n\in\mathfrak{n}=\Lambda^{1}\mathfrak{n}$ 作用为 $\operatorname{Ad}_{i^{c}_n}=\{i^{c}_n,\cdot\}$。被所有 $\operatorname{Ad}_{i^{c}_n}$ 零化的元素构成的子空间（即 $\mathfrak{n}\subset\mathcal{A}^{-1}_{(1)}$ 的中心化子）等于 $\Lambda\mathfrak{n}\otimes\mathcal{R}$。它是 $\mathcal{A}^{\le0}$ 的子空间，故 $Q^{c}$ 的唯一性显然。我们的 $\Lambda\mathfrak{n}$-模是自由的，故 $Q^{c}$ 的存在性由下述事实推出：映射 $n_1,n_2\mapsto\{\mathcal{L}ie_{n_1},i^{c}_{n_2}\}$ 是斜对称的。因为 $\mathcal{L}ie_n\in\mathcal{A}^{0}_{(\le1)}$，所以 $Q^{c}\in\mathcal{A}^{1}_{(\le1)}$。最后，因 $\{Q^{c},Q^{c}\}\in\mathcal{A}^{2}$，要验证它为零，只需证明对任意 $n,n'\in\mathfrak{n}$ 有 $\operatorname{Ad}_{i^{c}_n}\operatorname{Ad}_{i^{c}_{n'}}(\{Q^{c},Q^{c}\})=0$。事实上，
\[
\operatorname{Ad}_{i^{c}_n}\operatorname{Ad}_{i^{c}_{n'}}(\{Q^{c},Q^{c}\})
=2\operatorname{Ad}_{i^{c}_n}(\{\mathcal{L}ie_{n'},Q^{c}\})
=2\{i^{c}_{[n,n']},Q^{c}\}+2\{\mathcal{L}ie_{n'},\mathcal{L}ie_n\}=0.
\]
\hfill$\Box$

\emph{注.} 以 $\mathfrak{n}_{\heartsuit}$ 记下述李分次代数：其非零分量为 $\mathfrak{n}^{-1}_{\heartsuit}=\mathfrak{n}$、$\mathfrak{n}^{0}_{\heartsuit}=\mathfrak{n}$、$\mathfrak{n}^{1}_{\heartsuit}=\mathbb{C}=\mathbb{C}\cdot Q$；$\mathfrak{n}^{0}_{\heartsuit}$ 上的李括号与 $\mathfrak{n}$ 的相同；$\mathfrak{n}^{0}_{\heartsuit}$ 在 $\mathfrak{n}^{-1}_{\heartsuit}$ 上的伴随作用是 $\mathfrak{n}$ 的伴随作用；算子 $\operatorname{Ad}_{Q}:\mathfrak{n}^{-1}_{\heartsuit}\to\mathfrak{n}^{0}_{\heartsuit}$ 是 $\operatorname{id}_{\mathfrak{n}}$。于是带微分 $\operatorname{Ad}_{Q}$ 的 $\mathfrak{n}_{\heartsuit}$ 是李 DG 代数\footnote{注意 $\mathfrak{n}_{\heartsuit}/\mathfrak{n}^{1}_{\heartsuit}$ 就是 7.6.3 中的李 DG 代数 $\mathfrak{n}_{\Omega}$。}。此时 7.13.3 说明：存在李分次代数的典范态射 $\mathcal{L}ie:\mathfrak{n}_{\heartsuit}\to\mathcal{A}^{\cdot}$，其分量分别为 $n\mapsto i^{c}_n$、$n\mapsto\mathcal{L}ie_n$、$Q\mapsto Q^{c}$。

\bdp{7.13.4} 置 $d:=\operatorname{Ad}_{Q^{c}}=\{Q^{c},\cdot\}$。这是 $\mathcal{A}^{\cdot}$ 的度 $1$ 的导子，且平方为零。于是 $\mathcal{A}$ 是泊松 DG 代数；它称为 $\mathcal{R}$ 的\emph{BRST 约化}。态射 $\mathcal{L}ie:\mathfrak{n}_{\heartsuit}\to\mathcal{A}$ 是李 DG 代数态射。

若 $H^{i}\mathcal{A}=0$（$i\ne0$），则称 BRST 约化是\emph{正则的}。

易见 $Q^{c}=Q_1+Q_0$，其中 $Q_1\in\mathcal{A}^{1}_{(1)}=\mathfrak{n}\otimes\Lambda^{2}\mathfrak{n}^{*}\otimes\mathcal{R}$ 与 $Q_0\in\mathcal{A}^{1}_{(0)}=\mathfrak{n}^{*}\otimes\mathcal{R}$ 分别是 $\frac{1}{2}a^{c}\in\operatorname{Hom}(\mathfrak{n},\mathcal{C}l^{0}_1)=\mathfrak{n}^{*}\otimes\mathcal{C}l^{0}_1\subset\mathcal{A}^{1}\otimes\mathcal{A}^{0}$ 在乘法映射下的像，而 $l\in\operatorname{Hom}(\mathfrak{n},\mathcal{R})=\mathcal{A}^{1}_{(1)}$。
% [Tutor Note: 原书 p.338 此处印作 "$l\in\operatorname{Hom}(\mathfrak{n},\mathcal{R})=\mathcal{A}^{1}_{(1)}$"；因 $\operatorname{Hom}(\mathfrak{n},\mathcal{R})=\mathfrak{n}^{*}\otimes\mathcal{R}=\mathcal{A}^{1}_{(0)}$，末项 $(1)$ 疑为 $(0)$ 之误，底本如此，保留原样。]
把微分按双分次分解，\SourcePage{339}%
可见 $\mathcal{A}$ 是以 $d':\mathcal{A}^{i}_{(j)}\to\mathcal{A}^{i+1}_{(j)}$、$d'':\mathcal{A}^{i}_{(j)}\to\mathcal{A}^{i+1}_{(j-1)}$ 为双微分的双复形的全复形。

BRST 微分保持滤链 $\mathcal{A}_{(\le i)}$。特别地 $\mathcal{A}_{(0)}=C(\mathfrak{n},\mathcal{R})$ 是 $\mathcal{A}$ 的 DG 子代数，于是有分次代数的典范态射
\begin{equation}\tag{365}
H^{\cdot}(\mathfrak{n},\mathcal{R})\to H^{\cdot}\mathcal{A}.
\end{equation}

注意 $(\mathcal{A}^{\cdot}_{(-)},d'')$ 是关于 $l^{c}:\mathfrak{n}\to\mathcal{R}$ 的 Koszul 复形 $P:=\Lambda^{-\cdot}\mathfrak{n}\otimes\mathcal{R}$。于是 $\mathcal{A}$ 是 $\mathfrak{n}$ 的以 $P$ 为系数的李代数上链的 Chevalley 复形 $C^{\cdot}(\mathfrak{n},P)$。显然的投影 $P\to\mathcal{R}/\mathcal{R}l^{c}(\mathfrak{n})$ 给出 DG 代数同构 $\mathcal{A}/\mathcal{I}\xrightarrow{\sim}C(\mathfrak{n},\mathcal{R}/\mathcal{R}l^{c}(\mathfrak{n}))$，其中 $\mathcal{I}\subset\mathcal{A}$ 是由元素 $i^{c}_n$（$n\in\mathfrak{n}$）生成的 DG 理想。过渡到上同调，得到分次代数的典范态射
\begin{equation}\tag{366}
H^{\cdot}\mathcal{A}\to H^{\cdot}(\mathfrak{n},\mathcal{R}/\mathcal{R}l^{c}(\mathfrak{n})).
\end{equation}

若 $H_i(P)=0$（$i\ne0$），则称 $l^{c}$ 是\emph{正则的}。

\bdp{7.13.5} \emph{引理.} 若 $l^{c}$ 正则，则 (366) 是同构。

\emph{证明.} 正则性意味着投影 $P\to\mathcal{R}/\mathcal{R}l^{c}(\mathfrak{n})$ 是拟同构。故 $\mathcal{A}\to C^{\cdot}(\mathfrak{n},\mathcal{R}/\mathcal{R}l^{c}(\mathfrak{n}))$ 也是拟同构。\hfill$\Box$

于是对负的 $i$ 有 $H^{i}\mathcal{A}=0$，且 $H^{0}\mathcal{A}\xrightarrow{\sim}[\mathcal{R}/\mathcal{R}l^{c}(\mathfrak{n})]^{\mathfrak{n}}$，这正是 $\mathcal{R}$ 关于哈密顿作用 $l^{c}$ 的通常哈密顿约化。

\bdp{7.13.6} 现在转向 BRST 的「量子」版本。设 $\mathfrak{n}$ 是有限维李代数。置 $\mathrm{Cl}^{\cdot}:=\mathrm{Cl}^{\cdot}_{\mathfrak{n}}$。以 $\mathfrak{n}^{\flat}$ 记 $\mathfrak{n}$ 被 $\mathbb{C}$ 的中心扩张，它定义为 $\operatorname{End}^{\flat}_{\mathfrak{n}}$ 沿伴随作用态射 $\mathfrak{n}\to\operatorname{End}_{\mathfrak{n}}$ 的拉回（记号见 7.13.1 末）。换言之，$\mathfrak{n}^{\flat}$ 是 $\mathfrak{n}$ 被 $\mathbb{C}$ 的中心扩张，配有李代数映射 $a:\mathfrak{n}^{\flat}\to\mathrm{Cl}^{0}$，使 $a(1_{\mathfrak{n}^{\flat}})=1$\footnote{这里 $1_{\mathfrak{n}^{\flat}}$ 是 $\mathbb{C}\subset\mathfrak{n}^{\flat}$ 的生成元。}，\SourcePage{340}%
且由 $\mathfrak{n}\oplus\mathfrak{n}^{*}$ 上的伴随作用在 $\mathrm{Cl}$ 上诱导的 $\mathfrak{n}$-作用与 $a$ 的伴随作用相同。

设 $R$ 是结合代数，$l:\mathfrak{n}^{\flat}\to R$ 是李代数态射，且 $l(1_{\mathfrak{n}^{\flat}})=-1$。置 $A^{\cdot}:=\mathrm{Cl}^{\cdot}\otimes R$；这是结合分次代数。我们有李代数态射 $\operatorname{Lie}:=a+l:\mathfrak{n}\to A^{0}$，$n\mapsto\operatorname{Lie}_n:=a(n^{\flat})+l(n^{\flat})$，其中 $n^{\flat}$ 是 $n$ 到 $\mathfrak{n}^{\flat}$ 的任一提升。下文对 $n\in\mathfrak{n}$，以 $i_n$ 记 $\mathrm{Cl}^{-1}_1\subset A^{-1}$ 中相应的元素。有 $[\operatorname{Lie}_{n_1},i_{n_2}]=i_{[n_1,n_2]}$。

\bdp{7.13.7} \emph{引理.} 存在唯一元素 $Q=Q_A\in A^{1}$，使得对任意 $n\in\mathfrak{n}$ 有 $[Q,i_n]=\operatorname{Lie}_n$。事实上 $Q\in\mathrm{Cl}^{1}_1\otimes R$。有 $Q^{2}=0$。

\emph{证明.} 与「经典」版本 7.13.3 的证明相同。\hfill$\Box$

置 $d:=\operatorname{Ad}_Q$\footnote{当然，我们取「超」意义下的 $\operatorname{Ad}$，故对 $v\in A^{odd}$ 有 $dv=Qv+vQ$。}；这是 $A$ 的度 $1$ 的导子，且平方为零。于是 $A$ 是结合 DG 代数，称为 $R$ 的\emph{BRST 约化}。如同 7.13.3 后的注与 7.13.4 那样，我们有李 DG 代数的典范态射 $\operatorname{Lie}:\mathfrak{n}_{\heartsuit}\to A$，其分量为 $n\mapsto i_n$、$n\mapsto\operatorname{Lie}_n$、$Q\mapsto Q_A$。

若 $H^{i}A=0$（$i\ne0$），则称 BRST 约化是\emph{正则的}。

以 $C(\mathfrak{n},R)$ 记 $\mathfrak{n}$ 的以 $R$ 为系数的李代数上链的 Chevalley DG 代数（关于作用 $\operatorname{Ad}_l$）。作为分次代数它等于 $\Lambda^{\cdot}\mathfrak{n}^{*}\otimes R$，故它是 $A^{\cdot}$ 的子代数。

\bdp{7.13.8} \emph{引理.} 嵌入 $C(\mathfrak{n},R)\subset A$ 与微分相容。

\emph{证明.} 只需说明：在 $R,\mathfrak{n}^{*}\subset A$ 上，我们的微分分别等于 $\mathfrak{n}$-作用映射 $R\to\mathfrak{n}^{*}\otimes R$ 的对偶与括号映射 $\mathfrak{n}^{*}\to\Lambda^{2}\mathfrak{n}^{*}$ 的对偶。如同 $Q$ 的唯一性证明，只需验证对任意 $n,n_1,n_2\in\mathfrak{n}$、$n^{*}\in\mathfrak{n}^{*}$、$r\in R$ 有 $[i_n,[Q,r]]=[l(n),r]$ 与 $[i_{n_1},[i_{n_2},[Q,n^{*}]]]=n^{*}([n_1,n_2])$；这是直接的计算。\hfill$\Box$
% [Tutor Note: 原书 p.340 此处印作 "It suuffices to show"，"suuffices" 系 "suffices" 之误，底本如此，保留原样。]

\emph{注.} 我们看到 $d$ 保持环滤链 $\mathrm{Cl}^{\cdot}\otimes R$。在 $\mathrm{Cl}_i\otimes R/\mathrm{Cl}_{i-1}\otimes R=\Lambda^{\cdot+i}\mathfrak{n}^{*}\otimes\Lambda^{\cdot}\mathfrak{n}\otimes R=C^{\cdot+i}(\mathfrak{n},\Lambda^{\cdot}\mathfrak{n}\otimes R)$ 上，它与 Chevalley 微分相同。

\SourcePage{341}

DG 代数嵌入 $C(\mathfrak{n},R)\subset A$ 给出分次代数态射
\begin{equation}\tag{367}
H^{\cdot}(\mathfrak{n},R)\to H^{\cdot}A.
\end{equation}
特别地，因 $R$ 的中心 $\mathfrak{z}$ 落在 $R^{\mathfrak{n}}$ 中，得到态射
\begin{equation}\tag{368}
\mathfrak{z}\to H^{0}A.
\end{equation}

\bdp{7.13.9} \emph{注.}（只在有限维情形有效）设 $I$ 是由元素 $i_n$（$n\in\mathfrak{n}$）生成的 $A$ 的左 DG 理想。商复形 $A/I$ 可如下计算。设 $\mathfrak{n}\hookrightarrow\mathfrak{n}^{\flat}$ 是由 7.13.1 的注 (ii) 中的分裂 $s'$ 定义的分裂。则 $I$ 作为普通理想由元素 $i_n$ 与 $l(n)$（$n\in\mathfrak{n}$）生成。把投影 $A\to A/I$ 限制到 $C(\mathfrak{n},R)$，得到复形的同构 $A/I\xrightarrow{\sim}C(\mathfrak{n},R/Rl(\mathfrak{n}))$，它给出态射
\begin{equation}\tag{369}
H^{\cdot}A\to H^{\cdot}(\mathfrak{n},R/Rl(\mathfrak{n})).
\end{equation}

\bdp{7.13.10} \emph{注.} 设 $C^{\cdot}$ 是不可约分次 $\mathrm{Cl}^{\cdot}$-模（这样的 $C^{\cdot}$ 在同构与分次平移的意义下唯一）。若 $M=(M^{\cdot},d_M)$ 是 $R$-复形（即 $R$-模的复形），则 $M\otimes C^{\cdot}:=(M^{\cdot}\otimes C^{\cdot},d)$，其中 $d:=d_M\otimes\operatorname{id}_{C^{\cdot}}+Q\cdot$，是 $A$-复形（即 DG $A$-模）。函子 $\cdot\otimes C^{\cdot}$：($R$-复形) $\to$ ($A$-复形) 是范畴等价。

\bdp{7.13.11} 我们来比较「量子」与「经典」两种情形。设处于 7.13.6 的情形。设 $R_0\subset R_1\subset\ldots$ 是 $R$ 上的递增环滤链，使 $\cup R_i=R$ 且 $\mathcal{R}:=\operatorname{gr}R$ 交换。则 $\mathcal{R}$ 按通常方式是泊松代数。给 $A$ 配上滤链 $A_{\cdot}$，它等于滤链 $\mathrm{Cl}^{\cdot}$ 与 $R_{\cdot}$ 的张量积。则 $\mathcal{A}:=\operatorname{gr}A$ 作为泊松分次代数等于 $\mathcal{C}l^{\cdot}\otimes\mathcal{R}$。置 $\mathcal{A}_i:=\operatorname{gr}_iA$。

设 $l(\mathfrak{n}^{\flat})\subset R_1$；以 $l^{c}$ 记相应的态射 $\mathfrak{n}\to\mathcal{R}_1$。则 $(\mathcal{R},l^{c})$ 是定义 7.13.2 中「经典」BRST 构造的数据。由 7.13.3 有相应的「经典」BRST 元素 $Q^{c}$。易见 $Q\in A_1$，且 $Q^{c}$ 等于 $Q$ 在 $\mathcal{A}_1$ 中的像。

\SourcePage{342}

因此滤链 $A_{\cdot}$ 关于微分稳定，且 $\operatorname{gr}A$ 作为泊松 DG 代数与相应的「经典」$\mathcal{A}$ 相同。于是有收敛到 $H^{\cdot}A$ 的谱序列，其首项为
\[
E^{p,q}_1=H^{p+q}\mathcal{A}_{-p}.
\]

\bdp{7.13.12} \emph{引理.} (i) 设 $l^{c}$ 正则。则对 $i<0$ 有 $H^{i}\mathcal{A}=0$，且 $\operatorname{gr}H^{0}\mathcal{A}\subset[\mathcal{R}/\mathcal{R}l^{c}(\mathfrak{n})]^{\mathfrak{n}}$。

(ii) 若此外还有 $H^{i}(\mathfrak{n},\mathcal{R}/\mathcal{R}l^{c}(\mathfrak{n}))=0$（$i>0$），则对 $i\ne0$ 有 $H^{i}\mathcal{A}=0$，且 $\operatorname{gr}H^{0}\mathcal{A}\xrightarrow{\sim}[\mathcal{R}/\mathcal{R}l^{c}(\mathfrak{n})]^{\mathfrak{n}}$。

\emph{证明.} 看谱序列与 7.13.5。\hfill$\Box$

\bdp{7.13.13} 在下述情形可以显式算出代数 $H^{0}A$。设处于 7.13.11 的情形，且 $l:\mathfrak{n}^{\flat}\to R_1$ 是单射。以 $\mathfrak{b}'$ 记 $l(\mathfrak{n}^{\flat})$ 在 $R_1$ 中的正规化子。于是 $\mathfrak{b}'$ 是含 $\mathfrak{n}^{\flat}$ 的李代数，且有延拓 $l$ 的李代数嵌入 $l^{\mathfrak{b}}:\mathfrak{b}'\to R_1$。置 $\mathfrak{b}:=\mathfrak{b}'/\mathbb{C}$，于是 $\mathfrak{b}'$ 是 $\mathfrak{b}$ 被 $\mathbb{C}$ 的中心扩张。$\mathfrak{b}$ 的伴随作用给出李代数态射 $\mathfrak{b}\to\operatorname{End}_{\mathfrak{n}}$；以 $\mathfrak{b}^{\flat}$ 记中心扩张 $\operatorname{End}^{\flat}$ 的拉回（见 7.13.1）。则 $\mathfrak{n}^{\flat}$ 是 $\mathfrak{b}^{\flat}$ 的李子代数，且有延拓 $a$ 的李代数态射 $a^{\mathfrak{b}}:\mathfrak{b}^{\flat}\to\mathrm{Cl}^{0}_1$。

设 $\mathfrak{b}^{\natural}$ 是扩张 $\mathfrak{b}'$ 与 $\mathfrak{b}^{\flat}$ 的 Baer 和。按构造有典范分裂 $s:\mathfrak{n}\to\mathfrak{b}^{\natural}$。它关于 $\mathfrak{b}$ 的伴随作用不变，故 $s(\mathfrak{n})$ 是 $\mathfrak{b}^{\natural}$ 中的理想。置 $\mathfrak{h}^{\natural}:=\mathfrak{b}^{\natural}/s(\mathfrak{n})$；这是 $\mathfrak{h}:=\mathfrak{b}/\mathfrak{n}$ 被 $\mathbb{C}$ 的中心扩张。

置 $\operatorname{Lie}^{\mathfrak{b}}:=a^{\mathfrak{b}}\otimes1+1\otimes l^{\mathfrak{b}}:\mathfrak{b}^{\natural}\to A^{0}_1$。这是李代数态射，它在 $\mathbb{C}\subset\mathfrak{b}^{\natural}$ 上等于 $\operatorname{id}_{\mathbb{C}}$。其像与 $Q$ 交换（因为我们的构造都是自然的），即它属于 $\operatorname{Ker}d$。有 $\operatorname{Lie}^{\mathfrak{b}}\circ s=\operatorname{Lie}=d\circ i:\mathfrak{n}\to A^{0}$，故 $\operatorname{Lie}^{\mathfrak{b}}$ 给出典范态射 $\operatorname{Lie}^{\mathfrak{h}}:\mathfrak{h}^{\natural}\to H^{0}A$。以 $U^{\natural}\mathfrak{h}$ 记 $\mathfrak{h}$ 的对应于 $\mathfrak{h}^{\natural}$ 的扭包络代数。我们的 $\operatorname{Lie}^{\mathfrak{h}}$ 给出结合代数的典范态射\SourcePage{343}%
\begin{equation}\tag{370}
h:U^{\natural}\mathfrak{h}\to H^{0}A.
\end{equation}

此态射有显然的「经典」版本 $h^{c}:\operatorname{Sym}\mathfrak{h}\to H^{0}\mathcal{A}$。它与投影 $H^{0}A\to[\mathcal{R}/l^{c}(\mathfrak{n})\mathcal{R}]^{\mathfrak{n}}$ 的复合（见 (366)）是显然的态射 $\operatorname{Sym}\mathfrak{h}\to[\mathcal{R}/l^{c}(\mathfrak{n})\mathcal{R}]^{\mathfrak{n}}$，其在 $\mathfrak{h}$ 上的限制是 $l^{\mathfrak{b}}$ 与投影 $R_1\to R_1/R_0$ 的复合。

\bdp{7.13.14} \emph{引理.} 设 $l^{c}$ 正则，且态射 $\operatorname{Sym}\mathfrak{h}\to[\mathcal{R}/l^{c}(\mathfrak{n})\mathcal{R}]^{\mathfrak{n}}$ 是同构。则 (370) 是同构。

\emph{证明.} 用 7.13.12(i)。\hfill$\Box$

\bdp{7.13.15} \emph{例.}（参看 [Ko78]）
\begin{enumerate}[label=(\roman*),leftmargin=2.2em]
\item 我们用 7.13.13 的记号。设 $\mathfrak{g}$ 是（有限维）半单李代数，$\mathfrak{b}\subset\mathfrak{g}$ 是 Borel 子代数，$\mathfrak{n}:=[\mathfrak{b},\mathfrak{b}]$。置 $R:=U\mathfrak{g}$，并令 $R_{\cdot}$ 是 $R$ 上的标准滤链，故 $\mathcal{R}=\operatorname{Sym}\mathfrak{g}$。因 $\mathfrak{n}$ 的伴随作用是幂零的，扩张 $\mathfrak{n}^{\flat}$ 典范地平凡化。令 $l:\mathfrak{n}\to\mathfrak{g}\subset R$ 是显然的嵌入。则 $\mathfrak{b}'$ 等于 $\mathfrak{b}\oplus\mathbb{C}$，故此扩张成为平凡的。我们借 7.13.1 的注 (ii) 中的分裂 $s'$ 平凡化扩张 $\mathfrak{b}^{\flat}$。因此我们分裂了扩张 $\mathfrak{b}^{\natural}$，从而 $U^{\natural}\mathfrak{h}=\operatorname{Sym}\mathfrak{h}$。

7.13.14 的条件成立。事实上，$l^{c}$ 显然是正则的，且显然的嵌入 $i^{c}:\operatorname{Sym}\mathfrak{h}\hookrightarrow[\operatorname{Sym}(\mathfrak{g}/\mathfrak{n})]^{\mathfrak{n}}$ 是同构，因为 $\mathfrak{n}$ 沿投影 $(\mathfrak{g}/\mathfrak{n})^{*}\to\mathfrak{h}^{*}$ 的一般纤维单可迁地作用。因此 $h:\operatorname{Sym}\mathfrak{h}\xrightarrow{\sim}H^{0}A$。

下面说明典范态射 (368) $\mathfrak{z}\to H^{0}A=\operatorname{Sym}\mathfrak{h}$ 就是通常的 Harish-Chandra 态射。显然的嵌入 $i:\operatorname{Sym}\mathfrak{h}\xrightarrow{\sim}[R/Rl(\mathfrak{n})]^{\mathfrak{n}}$ 是同构；按定义，Harish-Chandra 态射是嵌入 $\mathfrak{z}\hookrightarrow R^{\mathfrak{n}}$ 与此同构的逆的复合。考虑 (369) 中的映射 $p:H^{0}A\to[R/Rl(\mathfrak{n})]^{\mathfrak{n}}$。由 $p$ 的定义可得 $ph=i$，这蕴含我们的断言。

\SourcePage{344}

\item 现在设 $\psi:\mathfrak{n}\to\mathbb{C}$ 是 $\mathfrak{n}$ 的非退化特征标（我们用 7.13.15(i) 的记号）。置 $R_t:=R[t]$，$l_t:=l+t\psi:\mathfrak{n}\to R_t$。
% [Tutor Note: 原书 p.344 此处作「7.13.15 (i)」，而本条自身即 7.13.15(ii)，所指记号实在 (i) 中给出，编号疑有误（或当为 7.13.13），底本如此，保留原样。]
\end{enumerate}

\bdp{7.13.16} 现在过渡到无穷维情形。我们需要固定一些 Clifford 代数的记号。设 $F$ 是 Tate 向量空间，于是有归纳概形 $Gr(F)$（见 7.11.2(iii)）。归纳概形 $Gr(F)\times Gr(F)$ 带有「相对行列式」的典范线丛 $\lambda$。这是分次线丛，带有典范同构
\begin{equation}\tag{371}
\lambda_{(P,P'')}=\lambda_{(P,P')}\otimes\lambda_{(P',P'')}
\end{equation}
以及对 $P'\subset P$ 的等同 $\lambda_{(P,P')}=\det(P/P')$，它们满足显然的相容性；这里我们假定 $\det(P/P')$ 位于度 $-\dim(P/P')$。

考虑带标准对称形式的 Tate 向量空间 $F\oplus F^{*}$ 与 Clifford 代数 $\mathrm{Cl}=\mathrm{Cl}_F:=\mathrm{Cl}(F\oplus F^{*})$。设 $C$ 是不可约离散 $\mathrm{Cl}$-模\footnote{这里「离散」意指 $C$ 中任一元素的零化子是 $F\oplus F^{*}$ 的开子空间。}。因 $C$ 在张量一个一维向量空间的意义下唯一\footnote{$C$ 同构于费米 Fock 空间 $\varinjlim_U\Lambda^{\cdot}(F/U)\otimes\det(P/U)^{*}$（参看 (182)），其中 $P$ 是 $F$ 中的 c-格，$U$ 取遍 $P$ 的一切 c-子格。}，相应的射影空间 $\mathbb{P}$ 有典范的定义（这是归纳概形）。对任意 c-格 $P\subset F\widehat{\otimes}A$，以 $\lambda^{C}_P$ 记 $C\otimes A$ 中被来自 $P$ 与 $P^{\perp}\subset F^{*}\widehat{\otimes}A$ 的 Clifford 算子零化的元素之集。$A$-子模 $\lambda^{C}_P\subset C\otimes A$ 是一条「线」（即秩 $1$ 的直和项），故 $\lambda^{C}$ 是 $C\otimes\mathcal{O}_{Gr(F)}$ 的线子丛。它定义典范嵌入 $Gr(F)\hookrightarrow\mathbb{P}$。有典范等同
\begin{equation}\tag{372}
\lambda_{(P,P')}=\lambda^{C}_P\otimes(\lambda^{C}_{P'})^{*}
\end{equation}
它与 (371) 相容：若 $P'\subset P$，则同构 $\lambda_{(P,P')}\otimes\lambda^{C}_{P'}\xrightarrow{\sim}\lambda^{C}_P$ 由显然的映射 $\lambda_{(P,P')}=\det(P/P')\to\mathrm{Cl}_F/\mathrm{Cl}_F\cdot P'$ 诱导。

代数 $\mathrm{Cl}$ 带有典范分次，使 $F\subset\mathrm{Cl}^{-1}$、$F^{*}\subset\mathrm{Cl}^{1}$。设 $C^{\cdot}$ 是 $C$ 上与 $\mathrm{Cl}_F$ 的分次相容的分次；它在相差一个平移的意义下\SourcePage{345}%
唯一。于是 $\lambda^{C}$ 是齐次线，且 (372) 是分次线丛的同构。
% [Tutor Note: 原书 p.345 此处印作 "homogenuous"，系 "homogeneous" 之误，底本如此，保留原样。]

\bdp{7.13.17} 以 $\overline{\mathrm{Cl}}^{\cdot}=\overline{\mathrm{Cl}}^{\cdot}_F$ 记 $\mathrm{Cl}^{\cdot}$（作为分次代数）关于由左理想 $\mathrm{Cl}\cdot U$ 生成的拓扑的完备化，其中 $U\subset F\oplus F^{*}$ 是开子空间。于是 $C$ 是离散 $\overline{\mathrm{Cl}}$-模。$\overline{\mathrm{Cl}}$ 的作用给出拓扑分次代数的同构 $\overline{\mathrm{Cl}}^{\cdot}\xrightarrow{\sim}\operatorname{End}^{\cdot}_{\mathbb{C}}C$。

分次代数 $\mathrm{Cl}^{\cdot}$ 有典范滤链 $\mathrm{Cl}^{\cdot}_0=\Lambda^{\cdot}F^{*}\subset\mathrm{Cl}^{\cdot}_1\subset\ldots$（见 7.13.1）。我们把 $\overline{\mathrm{Cl}}^{\cdot}$ 上的滤链 $\overline{\mathrm{Cl}}^{\cdot}_i$ 定义为 $\mathrm{Cl}^{\cdot}_i$ 的闭包。如同 7.13.1，经典 Clifford 代数 $\overline{\mathcal{C}l}^{\cdot}:=\operatorname{gr}\overline{\mathrm{Cl}}^{\cdot}$ 是泊松分次拓扑代数。它带有附加分次 $\overline{\mathcal{C}l}^{\cdot}_i:=\operatorname{gr}_i\overline{\mathrm{Cl}}^{\cdot}$；有
\[
\overline{\mathcal{C}l}^{a}_{i}=\varprojlim_{U,V}\Lambda^{i}(F/U)\otimes\Lambda^{a+i}(F^{*}/V),
\]
其中 $U,V$ 分别是 $F,F^{*}$ 中的 c-格。

以 $E=E_F$ 记 $F$ 的自同态的结合代数。令 $E^{Lie}$ 表示 $E$ 作为李代数。注意 $\overline{\mathcal{C}l}^{0}_1$ 是 $\overline{\mathcal{C}l}^{\cdot}$ 的李子代数，它使 $\overline{\mathcal{C}l}^{-1}_1$ 正规化。$\overline{\mathcal{C}l}^{0}_1$ 在 $\overline{\mathcal{C}l}^{-1}_1=F$ 上的伴随作用把 $\overline{\mathcal{C}l}^{0}_1$ 与 $E^{Lie}$ 等同\footnote{用上面关于 $\overline{\mathcal{C}l}^{0}_1$ 的显式描述。}。置 $E^{\flat}:=\overline{\mathcal{C}l}^{0}_1$；这是 $\mathrm{Cl}^{\cdot}$ 的李子代数，它是 $\overline{\mathcal{C}l}^{0}_1=E^{Lie}$ 被 $\mathbb{C}$ 的中心扩张。

我们看到 $E^{\flat}$ 以与 Clifford 作用相容的方式作用在 $C$ 上；此作用保持 $C$ 上的分次。

接下来几节 7.13.18--7.13.22 给出 $E^{\flat}$ 及其某些子代数的方便描述。读者可以跳过它们，直接看 7.13.23。

\bdp{7.13.18} 这里给出 $E^{Lie}$ 的中心扩张 $E^{\flat}$ 的显式描述，它本质上属于 Tate [T]。

设 $E_+\subset E$ 是由有界算子（即像有界的算子）组成的（双边）理想，$E_-\subset E$ 是由离散算子（即核为开的算子）组成的理想。有 $E_++E_-=E$；置 $E_{tr}:=E_+\cap E_-$。对任意 $A\in E_{tr}$，其迹 $\operatorname{tr}A$ 有定义（若 $U'\subset U\subset F$ 是 c-格，使得
% 尾部接 p.346（本单元在此断处停止）
%% end units/ch07_u11.tex
%% begin units/ch07_u12.tex
% ch07_u12 —— 原书 p.346–p.355（第 7 章 7.13.18 末–7.14.3）
% 底本：A. Beilinson, V. Drinfeld, Quantization of Hitchin's Integrable System
% and Hecke Eigensheaves
% 承接：上单元（p.336–p.345）末句断在 "...are c-lattices such that"，
% 本单元正文自「$A(F)\subset U$、$A(U')=0$，则得到……」续起。

\SourcePage{346}

$A(F)\subset U$、$A(U')=0$，则得到 $A^{\sim}:U/U'\to U/U'$，并令 $\operatorname{tr}A:=\operatorname{tr}A^{\sim}$）。

泛函 $\operatorname{tr}:E_{tr}\to\mathbb{C}$ 对 $E^{Lie}$ 的伴随作用不变；它在 $[E_+,E_-]\subset E_{tr}$ 上取零值。

我们的扩张 $E^{\flat}$ 带典范分裂 $s_+:E_+\to E^{\flat}$、$s_-:E_-\to E^{\flat}$。即：对 $A\in E_+$，其提升 $s_+(A)$ 由下述性质刻画——$s_+(A)$ 零化 $C$ 中一切被来自 $\operatorname{Im}A\subset\mathfrak{g}$ 的 Clifford 算子零化的元素。同样，$s_-(A)$ 是 $A\in E_-$ 的唯一提升，它零化 $C$ 中一切被来自 $(\operatorname{Ker}A)^{\perp}\subset F^{*}$ 的 Clifford 算子零化的元素。截面 $s_{\pm}$ 与 $E$ 的伴随作用交换；且对 $A\in E_{tr}$ 有 $s_-(A)-s_+(A)=\operatorname{tr}A\in\mathbb{C}\subset E^{\flat}$。
% [Tutor Note: 原书 p.346 此句印作 "from Im $A\subset\mathfrak{g}$"；因 $A\in E_+$ 是 $F$ 的有界自同态，此处 $\mathfrak{g}$ 疑为 $F$ 之误，底本如此，保留原样。]

易见具有这些性质的资料 $(E^{\flat},s_{\pm})$ 唯一确定。确实，考虑 $E$-双模的精确序列
\begin{equation}\tag{373}
0\to E_{tr}\xrightarrow{(-,+)}E_+\oplus E_-\xrightarrow{(+,+)}E\to 0.
\end{equation}
现在 $s=(s_+,s_-)$ 把 $E^{\flat}$ 与扩张 (373) 沿 $\operatorname{tr}:E_{tr}\to\mathbb{C}$ 的前推等同。$E^{Lie}$ 在 $E^{\flat}$ 上的伴随作用来自 $E$-双模 $E_+\oplus E_-$ 上的伴随作用。

\emph{注.} (i) 向量空间 $F\otimes F^{*}$ 上有 4 个自然拓扑，其开子空间基分别由 $U\otimes V$、$U\otimes F^{*}$、$F\otimes V$ 与 $U\otimes F^{*}+F\otimes V$ 构成，其中 $U\subset F$、$V\subset F^{*}$ 是开子空间。相应的完备化分别等于 $E_{tr}$、$E_+$、$E_-$ 与 $E$。迹泛函是典范配对 $F\otimes F^{*}\to\mathbb{C}$ 的连续延拓。

(ii) 置 $(E_-/E_{tr})^{\flat}:=E_-/\operatorname{Ker}\operatorname{tr}$；这是 $(E_-/E_{tr})^{Lie}$ 被 $\mathbb{C}$ 的中心扩张。注意 $E_-/E_{tr}\cong E/E_+$，故有投影 $\pi_-:E^{Lie}\to(E_-/E_{tr})^{Lie}$。它典范地提升为扩张的态射 $\pi^{\flat}_-:E^{\flat}\to(E_-/E_{tr})^{\flat}$，其核为 $s_+(E_+)$。换言之，$E^{\flat}$ 是 $(E_-/E_{tr})^{\flat}$ 沿 $\pi_-$ 的拉回。$\pm$ 互换时同理。

(iii) 设 $F^{i}$ 是 $F$ 由闭子空间给出的有限滤链；以 $B\subset E_{F}$ 记保持该滤链的自同态构成的子代数。我们有\SourcePage{347}%
$B^{Lie}$ 的诱导中心扩张 $B^{\flat}$。另一方面，有显然的投影 $\operatorname{gr}^{i}:B\to E_{\operatorname{gr}^{i}F}$；以 $B^{\flat i}$ 记扩张 $E^{\flat}_{\operatorname{gr}^{i}F}$（即 $E^{Lie}_{\operatorname{gr}^{i}F}$ 的中心扩张）沿该投影的拉回。以 $B^{\flat'}$ 记诸扩张 $B^{\flat i}$ 的 Baer 和。则存在典范（且唯一）的扩张同构 $B^{\flat'}\xrightarrow{\sim}B^{\flat}$。

确实，$B^{\flat'}$ 与 (373) 的精确子序列
\[
0\to B\cap E_{tr}\to(B\cap E_+)\oplus(B\cap E_-)\to B\to 0
\]
所定义的扩张一致（注意对 $e\in B\cap E_{tr}$ 有 $\operatorname{tr}(e)=\sum\operatorname{tr}(\operatorname{gr}^{i}e)$）。特别地我们看到，$B^{\flat}$ 在李子代数 $\operatorname{Ker}\operatorname{gr}^{\cdot}$ 上典范分裂。

\bdp{7.13.19} 置 $K=\mathbb{C}((t))$、$O:=\mathbb{C}[[t]]$。设 $F$ 是带通常拓扑的有限维 $K$-向量空间；这是 Tate $\mathbb{C}$-向量空间。设 $i:\mathcal{D}\hookrightarrow E$ 是作用在 $F$ 上的 $K$-微分算子代数，于是得到 $\mathcal{D}^{Lie}$ 的中心扩张 $\mathcal{D}^{\flat}$。下面（沿用 [BS] 2.4）用几何的语言重述 $\mathcal{D}^{\flat}$ 的 Tate 描述。
% [Tutor Note: 原书 p.347 此处印作 "the agebra of K-differential operators"，"agebra" 系 "algebra" 之误，底本如此，保留原样。]

置 $F':=\operatorname{Hom}_{K}(F,K)$、$F^{\circ}:=F'\otimes_{K}\omega_{K}$。显然 $F^{\circ}$ 与 Tate 对偶 $F^{*}$ 一致（用配对 $f^{\circ},f\mapsto\langle f^{\circ},f\rangle:=\operatorname{Res}(f^{\circ},f)$）。我们的 $F$ 是左 $\mathcal{D}$-模，而 $F^{\circ}$ 带唯一的右 $\mathcal{D}$-模结构，使 $\langle,\rangle$ 为 $\mathcal{D}$-不变配对；注意 $\mathcal{D}$ 以微分算子作用在 $F^{\circ}$ 上，这就是通常几何中的「伴随」作用。

设 $K\widehat{\otimes}K$ 是 $K\otimes K$ 关于以 $(t^{n}O)\otimes(t^{n}O)$ 为基的拓扑之完备化，即 $K\widehat{\otimes}K:=\mathbb{C}[[t_{1},t_{2}]][t_{1}^{-1}][t_{2}^{-1}]$。设 $F\widehat{\otimes}F^{\circ}$ 是 $F\otimes F^{\circ}$ 的类似完备化；这是有限维 $K\widehat{\otimes}K$-模。以 $F\widehat{\otimes}F^{\circ}(\infty\Delta)$ 记 $F\widehat{\otimes}F^{\circ}$ 对 $(t_{1}-t_{2})^{-1}$（即对角线方程）的局部化。考虑标准精确序列
\begin{equation}\tag{374}
0\to F\widehat{\otimes}F^{\circ}\to F\widehat{\otimes}F^{\circ}(\infty\Delta)\xrightarrow{r}\mathcal{D}\to 0
\end{equation}
其中投影 $r$ 把「核」$k=k(t_{1},t_{2})dt_{2}\in F\widehat{\otimes}F^{\circ}(\infty\Delta)$ 映到微分算子 $r(k):F\to F$，$f(t)\mapsto\operatorname{Res}_{t_{2}=t}(k(t,t_{2}),f(t_{2}))dt_{2}$。注意 $F\widehat{\otimes}F^{\circ}$ 以自然方式显然是 $\mathcal{D}$-双模。此双作用\SourcePage{348}%
唯一地延拓为 $F\widehat{\otimes}F^{\circ}(\infty\Delta)$ 上与 $K$-双模结构相容的 $\mathcal{D}$-双作用。易见 (374) 是 $\mathcal{D}$-双模的精确序列。设 $\operatorname{tr}:F^{\circ}\widehat{\otimes}F\to\mathbb{C}$ 为态射 $f\otimes f^{\circ}\mapsto\langle f^{\circ},f\rangle$（即它是对角线上限制的留数）。它关于 $\mathcal{D}^{Lie}$ 的伴随作用不变。以 $D^{\flat'}$ 记 (374) 沿 $\operatorname{tr}$ 的前推。$\mathcal{D}$ 在 $F\widehat{\otimes}F^{\circ}(\infty\Delta)$ 上的伴随作用给出 $D^{\flat'}$ 上的 $D^{Lie}$-模结构。对 $l^{\flat}_{1},l^{\flat}_{2}\in D^{\flat'}$，置 $[l^{\flat}_{1},l^{\flat}_{2}]:=l_{1}(l^{\flat}_{2})$，其中 $l_{1}$ 是 $l^{\flat}_{1}$ 在 $D^{Lie}$ 中的像。
% [Tutor Note: 原书 p.348 此处作 "The adjoint action of  on $F\hat\otimes F^{\circ}(\infty\Delta)$ yields a $D^{Lie}$-module structure on $D^{\flat'}$"，"of" 之后脱漏 $\mathcal{D}$；译文按上下文补出，底本如此，保留原样。]

\bdp{7.13.20} \emph{引理.} 括号 $[\,,\,]$ 斜对称，故它使 $D^{\flat'}$ 成为 $D^{Lie}$ 被 $\mathbb{C}$ 的中心扩张。存在中心扩张的唯一同构
\[
D^{\flat'}\xrightarrow{\sim}D^{\flat}.
\]

\emph{证明.} 只需建立 $D^{Lie}$-模扩张的同构 $D^{\flat'}\xrightarrow{\sim}D^{\flat}$。它来自 $\mathcal{D}$-双模精确序列的典范嵌入 $i^{\sim}:(374)\hookrightarrow(373)$，后者定义如下。态射 $\mathcal{D}\hookrightarrow E$ 是我们的标准嵌入 $i$，而 $i^{\sim}:F\widehat{\otimes}F^{\circ}=F\widehat{\otimes}F^{*}\xrightarrow{\sim}E_{tr}$ 是显然的同构（见 7.13.18 的注 (i)）。映射 $i^{\sim}=(i^{\sim}_+,i^{\sim}_-):F\widehat{\otimes}F^{\circ}(\infty\Delta)\to E_+\oplus E_-$ 把「核」$k$ 送到算子 $i^{\sim}_-(k)$ 与 $i^{\sim}_+(k)$：前者等于 $f\mapsto-\operatorname{Res}_{t_{2}=0}(k(t,t_{2}),f(t_{2}))dt_{2}$，后者等于 $f\mapsto(\operatorname{Res}_{t_{2}=t}+\operatorname{Res}_{t_{2}=0})(k(t,t_{2}),f(t_{2}))dt_{2}$。这里 $f\in F$，且 $(k(t,t_{2}),f(t_{2}))dt_{2}\in F((t_{2}))dt_{2}$。算子 $i^{\sim}_{\pm}(k)$ 属于 $E_{\pm}$ 一事留给读者验证\footnote{对 $i^{\sim}_-(k)$ 这是显然的。要验证 $i^{\sim}_+(k)\in E_+$，可用 Parshin 留数公式（[Pa76]，§1，命题 7）作用于 2-形式 $(k(t_{1},t_{2}),g(t_{1})f(t_{2}))dt_{1}\wedge dt_{2}$，其中 $g$ 属于 $F^{*}$ 中充分小的 c-格。}。因 $i^{\sim}$ 保持迹泛函，它给出所要的 $D^{Lie}$-模同构 $D^{\flat'}\xrightarrow{\sim}D^{\flat}$。\hfill$\Box$

\emph{注.} 设 $D_{i}\subset D$ 是次数 $\le i$ 的微分算子构成的空间。扩张 $D^{\flat}_{i}$ 带由嵌入\SourcePage{349}%
$D^{\flat}_{i}\subset\overline{\mathrm{Cl}}_{F}$ 引出的自然拓扑。这是 Tate 拓扑；$D_{i}$ 上的商拓扑与其作为有限维 $K$-向量空间的自然拓扑一致。

\bdp{7.13.21} \emph{例.} 置 $\mathcal{E}:=\operatorname{End}_{K}F=D_{0}\subset D$，于是得到 $\mathcal{E}^{Lie}$ 的中心扩张 $\mathcal{E}^{\flat}$。设 $\mathcal{L}\subset D^{Lie}$ 是 $\mathcal{E}$ 的正规化子；它以伴随的方式作用在 $\mathcal{E}^{\flat}$ 上。我们将把扩张 $\mathcal{E}^{\flat}$ 描述为 $\mathcal{L}$-模\footnote{因 $\mathcal{E}\subset\mathcal{L}$，这里特别描述了 $\mathcal{E}$ 的伴随作用，它正是 $\mathcal{E}^{\flat}$ 上的李括号。}。

易见 $\mathcal{L}$ 等于阶 $\le1$、其象征属于 $\operatorname{Der}_{K}\cdot\operatorname{id}_{F}$ 的微分算子李代数。换言之，$\mathcal{L}$ 由对 $(\tau,\tau^{\sim})$ 组成，其中 $\tau\in\operatorname{Der}K$、$\tau^{\sim}$ 是 $\tau$ 在 $F$ 上的作用；即 $\mathcal{L}$ 是 $(K,F)$ 的无穷小对称李代数。

同上，置 $\mathcal{E}^{\circ}:=\mathcal{E}\otimes_{K}\omega_{K}$。我们用配对 $\langle,\rangle:\mathcal{E}^{\circ}\times\mathcal{E}\to\mathbb{C}$、$\langle a,b\rangle:=\operatorname{Res}\operatorname{tr}_{K}(ab)$ 把 $\mathcal{E}^{\circ}$ 与 Tate 向量空间对偶 $\mathcal{E}^{*}$ 等同。$\mathcal{L}$ 在 $\mathcal{E}^{\circ}$ 上的伴随作用是 $(\tau,\tau^{\sim})(e\otimes\nu)=[\tau^{\sim},e]\otimes\nu+e\otimes\operatorname{Lie}_{\tau}\nu$。设 $\omega_{K}^{\otimes1/2}$ 是 $\operatorname{Spec}K$ 上的半形式层。它带 $\mathcal{L}$-作用（$(\tau,\tau^{\sim})$ 通过 $\operatorname{Lie}_{\tau}$ 作用），故 $\mathcal{L}$ 作用在 $\otimes\omega_{K}^{\otimes1/2}$ 上。考虑 $F\otimes\omega_{K}^{\otimes1/2}$ 上联络的集合 $\operatorname{Conn}(F\otimes\omega_{K}^{\otimes1/2})$\footnote{它不依赖于 $\omega_{K}^{\otimes1/2}$ 的选取。}。因 $\operatorname{End}_{K}F=\operatorname{End}_{K}(F\otimes\omega_{K}^{\otimes1/2})$，我们的 $\operatorname{Conn}(F\otimes\omega_{K}^{\otimes1/2})$ 是 $\mathcal{E}^{\circ}$-torsor（主齐性空间）；$\mathcal{L}$ 以显然的方式作用于其上。
% [Tutor Note: 原书 p.349 此处印作 "so $L$ acts on $\otimes\omega_{K}^{\otimes1/2}$"，$\omega$ 前多印一个 $\otimes$，底本如此，保留原样。]

\bdp{7.13.22} \emph{引理.} 存在唯一的 $\mathcal{L}$-且 $\mathcal{E}^{\circ}$-不变配对
\[
\langle,\rangle:\operatorname{Conn}(F\otimes\omega_{K}^{\otimes1/2})\times\mathcal{E}^{\flat}\to\mathbb{C}
\]
使对任意 $\nabla\in\operatorname{Conn}(F\otimes\omega_{K}^{\otimes1/2})$ 有 $\langle\nabla,1_{\mathcal{E}^{\flat}}\rangle=1$。

\emph{注.} (i) 元素 $\lambda\in\mathcal{E}^{\circ}$ 按公式 $\nabla\mapsto\nabla+\lambda$ 与 $e^{\flat}\mapsto e^{\flat}+\langle\lambda,e\rangle$ 分别作用在 $\operatorname{Conn}(F\otimes\omega_{K}^{\otimes1/2})$ 与 $\mathcal{E}^{\flat}$ 上（这里 $e:=e^{\flat}\bmod\mathbb{C}_{\mathcal{E}^{\circ}}=\mathcal{E}$）。故 $\langle,\rangle$ 的 $\mathcal{E}^{\circ}$-不变性即指 $\langle\nabla+\lambda,e^{\flat}\rangle=\langle\nabla,e^{\flat}\rangle-\langle\lambda,e\rangle$。

(ii) 显然 $\langle,\rangle$ 把 $\mathcal{E}^{\flat}$ 与 $\operatorname{Conn}(F\otimes\omega_{K}^{\otimes1/2})$ 上连续仿射泛函构成的 $\mathcal{L}$-模等同。这就是所承诺的 $\mathcal{E}^{\flat}$ 的描述。

\SourcePage{350}

\emph{证明.} $\langle,\rangle$ 的唯一性由下述事实推出：$\operatorname{Conn}(F\otimes\omega_{K}^{\otimes1/2})$ 无 $\mathcal{L}$-不变元素。

为定义 $\langle\nabla,e^{\flat}\rangle$，我们选取 $F$ 上的联络 $\nabla_{F}$ 与 $\omega_{K}$ 上的联络 $\nabla_{\omega}$，使得 $\nabla=\nabla_{F}+\frac{1}{2}\nabla_{\omega}$。

a. 联络 $\nabla_{F}$ 把 $F\otimes K$ 与 $K\otimes F$ 在对角线形式邻域上的限制等同，即它给出 $K\widehat{\otimes}K$-模的同构 $\epsilon(\nabla_{F}):F\widehat{\otimes}K\xrightarrow{\sim}K\widehat{\otimes}F$。令 $\epsilon(\nabla_{F}):F\widehat{\otimes}F^{\circ}\to K\widehat{\otimes}\omega_{K}$ 为 $\epsilon(\nabla_{F})\otimes\operatorname{id}_{F^{\circ}}$ 与显然态射 $K\widehat{\otimes}(F\otimes F^{\circ})\to K\widehat{\otimes}\omega_{K}$（由配对 $F\otimes F^{\circ}\to\omega_{K}$ 定义）的复合。用对角线方程局部化 $\epsilon(\nabla_{F})$，便得到态射 $F\widehat{\otimes}F^{\circ}(\infty\Delta)\to K\widehat{\otimes}\omega_{K}(\infty\Delta)$。把它作用于 $e^{\flat}$，得到 1-形式 $\epsilon(\nabla_{F},e^{\flat})\in K\widehat{\otimes}\omega_{K}(\Delta)$，它在那些满足 $\operatorname{Res}_{0}\phi(t,t)dt=0$ 的形式 $\phi(t_{1},t_{2})dt_{2}\in K\widehat{\otimes}\omega_{K}$ 构成的子空间模下良定义。注意对 $\lambda\in\mathcal{E}^{\circ}$ 有 $\epsilon(\nabla_{F}+\lambda,e^{\flat})=\epsilon(\nabla_{F},e^{\flat})-\operatorname{tr}_{K}(\lambda\cdot e)$（这里 $\operatorname{tr}_{K}(\lambda\cdot e)\in\omega_{K}=K\widehat{\otimes}\omega_{K}/(t_{1}-t_{2})K\widehat{\otimes}\omega_{K}$）。

b. 设 $\nu\in\omega_{K}\widehat{\otimes}K(\Delta)$ 是在对角线上留数为 1 的形式（即 $\nu$ 模 $\omega_{K}\widehat{\otimes}K$ 等于 $\frac{dt_{1}}{t_{1}-t_{2}}$）。设 $\psi(\nabla_{\omega})$ 是类似的形式，满足 $\psi(\nabla_{\omega})^{\otimes2}=-\nabla^{(1)}_{\omega}\nu$\footnote{这里 $\nabla^{(1)}_{\omega}$ 是沿第一个变量的协变导数。}。注意 $\psi(\nabla_{\omega})$ 在模 $(t_{1}-t_{2})\omega_{K}\widehat{\otimes}K$ 的意义下良定义。对 $l\in\omega_{K}$ 有 $\psi(\nabla_{\omega}+l)=\psi(\nabla_{\omega})-l$（这里把 $l$ 看作 $\omega_{K}\widehat{\otimes}K/(t_{1}-t_{2})\omega_{K}\widehat{\otimes}K$ 的元素）。

c. 考虑 2-形式 $\epsilon(\nabla_{F},e^{\flat})\wedge\nu$。置
\[
\langle\nabla,e^{\flat}\rangle:=\operatorname{Res}_{0}\operatorname{Res}_{\Delta}(\epsilon_{\nabla}(e^{\flat})\wedge\nu)
\]
则 $\langle\nabla,e^{\flat}\rangle$ 良定义（即它不依赖辅助选取），且 $\langle,\rangle$ 是 $\mathcal{E}^{\circ}$-不变的。因所有构造都是自然的，它也 $\mathcal{L}$-不变。\hfill$\Box$
% [Tutor Note: 原书 p.350 此处印作 "Since all the constructions where natural"，"where" 系 "were" 之误，底本如此，保留原样。]

\emph{注.} (i) 设 $e_{\alpha}$ 是 $F$ 的 $F$-基，$e'_{\alpha}$ 是 $F'$ 的对偶基，$\nabla$ 是使 $e'_{\alpha}\cdot(dt)^{-1/2}$ 为水平截面的联络。以 $(e_{\alpha}\cdot e'_{\beta})^{\flat}\in\mathcal{E}^{\flat}$ 记 $e_{\alpha}\otimes e'_{\beta}\frac{dt_{2}}{t_{2}-t_{1}}$ 的像。则 $\langle\nabla,(e_{\alpha}\cdot e_{\beta})^{\flat}\rangle=\delta_{\alpha,\beta}$。
% [Tutor Note: 原书 p.350 末式印作 "$(e_\alpha\cdot e_\beta)^{\flat}$"，第二处 $e_\beta$ 脱漏撇号（上文定义为 $(e_\alpha\cdot e'_\beta)^{\flat}$）；又上句印作 "an F-basis of F"，"F" 疑为 "K" 之误。底本如此，保留原样。]

\SourcePage{351}

(ii) 上述引理是局部 Riemann--Roch 定理的特例；参看 [BS] 的附录。

\bdp{7.13.23} 现在设 $\mathfrak{n}$ 是 Tate 意义下的李代数，即带连续李括号 $[\,,\,]$ 的 Tate 向量空间。下述引理或可使读者更安心。

\emph{引理.} $\mathfrak{n}$ 有由 $\mathfrak{n}$ 的李子代数构成的 $0$ 的邻域基。

\emph{证明.} 取任意 c-格 $P\subset\mathfrak{n}$。我们要找开李子代数 $\mathfrak{k}\subset P$。注意
\begin{equation}\tag{375}
\mathfrak{n}_{P}:=\{\alpha\in\mathfrak{n}:[\alpha,P]\subset P\}
\end{equation}
是开李子代数。置 $\mathfrak{k}:=P\cap\mathfrak{n}_{P}$。\hfill$\Box$

\bdp{7.13.24} 我们对 $F=\mathfrak{n}$ 使用 7.13.17 的记号。于是得到 Clifford 分次拓扑代数 $\overline{\mathrm{Cl}}^{\cdot}=\overline{\mathrm{Cl}}^{\cdot}_{\mathfrak{n}}$、相应的经典 Clifford 分次拓扑代数 $\overline{\mathcal{C}l}^{\cdot}=\operatorname{gr}\overline{\mathrm{Cl}}^{\cdot}$（它是泊松分次拓扑代数）、Tate 向量空间 $\mathfrak{n}$ 的自同态李代数 $E^{Lie}$ 的中心扩张 $E^{\flat}$，以及嵌入 $E^{\flat}\hookrightarrow\overline{\mathrm{Cl}}^{0}$。伴随作用给出态射 $\mathfrak{n}\to E^{Lie}$；以 $\mathfrak{n}^{\flat}$ 记扩张 $E^{\flat}$ 到 $\mathfrak{n}$ 的拉回。于是 $\mathfrak{n}^{\flat}$ 是 $\mathfrak{n}$ 被 $\mathbb{C}$ 的中心扩张。我们把 $\mathfrak{n}^{\flat}$ 赋以最弱的拓扑，使投影 $\mathfrak{n}^{\flat}\to\mathfrak{n}$ 与态射 $\mathfrak{n}^{\flat}\to\overline{\mathrm{Cl}}^{0}$ 都连续。则 $\mathfrak{n}^{\flat}$ 是 Tate 空间，且映射 $\mathfrak{n}^{\flat}/\mathbb{C}\to\mathfrak{n}$ 是同胚\footnote{确实，扩张 $\mathfrak{n}^{\flat}$ 在形如 (375) 的任意子代数上有典范连续分裂（其像由零化 $\lambda_{P}$ 的算子组成）。}。

\bdp{7.13.25} 现在我们准备把 BRST 构造搬到无穷维情形。先从「经典」版本开始。设 $\mathcal{R}$ 是拓扑泊松代数。我们假设 $\mathcal{R}$ 完备、分离，且是拓扑。
% [Tutor Note: 原书 p.351 末句作 "We assume that R is complete and separated and topology."，句子残缺，"and topology" 前疑有脱漏，底本如此，保留原样。]

\SourcePage{352}

\bdp{7.13.26} 以 $\mathcal{M}(\mathfrak{g})^{\flat}$ 记离散 $\mathfrak{g}^{\flat}$-模 $V$ 的范畴，其中 $1\in\mathbb{C}\subset\mathfrak{g}^{\flat}$ 以 $-\operatorname{id}_{V}$ 作用。对这样的 $V$，$\mathfrak{g}^{\flat}$ 在 $C^{\cdot}$ 与 $V$ 上的作用给出 $C^{\cdot}\otimes V$ 上的 $\mathfrak{g}$-模结构。它显然也是 $\mathrm{Cl}_{\mathfrak{g}}$-模，且 $\mathfrak{g}$-作用与 Clifford 作用相容。对 $\alpha\in\mathfrak{g}$，我们把其在 $C^{\cdot}\otimes V$ 上的作用记作 $\operatorname{Lie}_{\alpha}$，把 Clifford 算子 $C^{\cdot}\otimes V\to C^{\cdot-1}\otimes V$ 记作 $i_{\alpha}$。

宜把作用在 $C^{\cdot}\otimes V$ 上的算子改写如下（参看 7.7.5）。设 $\Omega_{\mathfrak{g}}$ 是 $\mathfrak{g}$ 的连续李代数上链构成的 DG 代数。相应的平面分次代数 $\Omega^{\cdot}_{\mathfrak{g}}$ 是 $\mathfrak{g}^{*}$ 的完备外代数。我们把它与完备 Clifford 代数 $\overline{\mathrm{Cl}}_{\mathfrak{g}}$ 中由 $\mathfrak{g}^{*}\subset\mathrm{Cl}_{\mathfrak{g}}$ 生成的闭子代数等同，于是 $\Omega^{\cdot}_{\mathfrak{g}}$ 以 Clifford 算子作用在 $C^{\cdot}\otimes V$ 上。现在设 $\mathfrak{g}_{\Omega}$ 是如下定义的 DG 李代数。其非零分量只有 $\mathfrak{g}^{0}_{\Omega}=\mathfrak{g}^{-1}_{\Omega}=\mathfrak{g}$，微分 $\mathfrak{g}^{-1}_{\Omega}\to\mathfrak{g}^{0}_{\Omega}$ 是 $\operatorname{id}_{\mathfrak{g}}$，$\mathfrak{g}^{0}_{\Omega}$ 上的括号就是 $\mathfrak{g}$ 的括号。回忆 $\mathfrak{g}_{\Omega}$ 作用在 $\Omega_{\mathfrak{g}}$ 上（即 $\mathfrak{g}^{0}_{\Omega}$ 按余伴随方式作用，$\mathfrak{g}^{-1}_{\Omega}$ 按「常数」导子作用）。分次李代数 $\mathfrak{g}^{\cdot}_{\Omega}$ 通过算子 $\operatorname{Lie}_{\alpha}$ 与 $i_{\alpha}$ 作用在 $C^{\cdot}\otimes V$ 上。故 $C^{\cdot}\otimes V$ 是分次 $(\Omega^{\cdot}_{\mathfrak{g}},\mathfrak{g}^{\cdot}_{\Omega})$-模。
% [Tutor Note: 原书 p.352 此处印作 "The corresponding plane graded algebra"，"plane" 疑为 "plain" 之误，底本如此，保留原样。]

\bdp{7.13.27} \emph{命题.} 存在唯一线性映射 $d:C^{\cdot}\otimes V\to C^{\cdot+1}\otimes V$，使对任意 $\alpha\in\mathfrak{g}$ 有 $\operatorname{Lie}_{\alpha}=di_{\alpha}+i_{\alpha}d$。有 $d^{2}=0$，且 $C_{\mathfrak{g}}(V):=(C^{\cdot}\otimes V,d)$ 是 DG $(\Omega_{\mathfrak{g}},\mathfrak{g}_{\Omega})$-模。

\emph{证明.} 唯一性。两个这样的 $d$ 之差是与一切 $i_{\alpha}$ 交换的算子。易见此算子代数与由 $\mathfrak{g}^{-1}_{\Omega}$ 与 $\operatorname{End}V$ 生成的闭子代数一致。因它没有正次数的算子，故毕。

类似论证表明 $(\Omega_{\mathfrak{g}},\mathfrak{g}_{\Omega})$ 的作用与微分相容，且 $d^{2}=0$（先证 $[d,\operatorname{Lie}_{\alpha}]=0$，再证其余性质）。

存在性。我们显式写出 $d$。设 $e_{i}$（$i\in I$）是 $\mathfrak{g}$ 的拓扑基（见 4.2.13），$e^{*}_{i}$ 是 $\mathfrak{g}^{*}$ 的对偶基。对半无限（关于 $\mathfrak{g}$）\SourcePage{353}%
子集 $A\subset I$，以 $\lambda_{A}\subset C^{\cdot}$ 记与由 $e_{a}$（$a\in A$）生成的 c-格相应的齐次线 $\lambda^{C}$（见 7.13.16）。换言之，$\lambda_{A}$ 是被 Clifford 算子 $e_{a}$、$e^{*}_{b}$（$a\in A$，$b\in I\backslash A$）零化的向量构成的子空间。我们的 $C^{\cdot}$ 是诸 $\lambda_{A}$ 的直和。注意对上述 $a,b$ 有 $e^{*}_{a}(\lambda_{A})=\lambda_{A\backslash a}$、$e_{b}(\lambda_{A})=\lambda_{A\cup b}$。
% [Tutor Note: 原书 p.353 此处印作 "the homoguenous line"，"homoguenous" 系 "homogeneous" 之误，底本如此，保留原样。]

置 $V_{A}:=\lambda_{A}\otimes V$；则 $C^{\cdot}\otimes V$ 是诸 $V_{A}$ 的直和。对 $c\in I$ 置 $L_{c}:=\operatorname{Lie}_{e_{c}}$、$i_{c}:=i_{e_{c}}$；对半无限子集 $A,A'$，以 $L^{A,A'}_{c}$、$i^{A,A'}_{c}$ 记这些算子从 $V_{A}$ 到 $V_{A'}$ 的 $A,A'$-分量。

设 $A,B$ 是半无限子集，满足 $|A|-|B|=1$（这里 $|A|-|B|:=|A\backslash(A\cap B)|-|B\backslash(A\cap B)|$）。任取 $a=a_{A,B}\in A\backslash(A\cap B)$（此集合非空）。以 $d^{A,B}$ 记复合 $V_{A}\to V_{B\cup a}\to V_{B}$，其中第一个箭头是 $L^{A,B\cup a}_{a}$，第二个是 Clifford 算子 $e^{*}_{a}$。易见以 $d^{A,B}$ 为分量的算子 $d:C^{\cdot}\otimes V\to C^{\cdot+1}\otimes V$ 定义正确（用到下述事实：对任意 $v\in V$，只有有限多个 $a\in A$ 使 $L_{a}(\lambda_{A}\otimes v)$ 非零）。
% [Tutor Note: 原书 p.353 此处作 "for any $v\in V$ and there is only finitely many $a\in A$ such that..."，"and" 之后疑有脱漏，底本如此，译文略去该 and。]

还需说明我们的 $d$ 满足命题的条件，即对任意 $c\in I$ 有 $[d,i_{c}]=L_{c}$。此事可直接计算验证；要点是 $[L_{a},i_{b}]$ 关于 $a,b$ 的斜对称性。细节留给读者。\hfill$\Box$

\bdp{7.13.28} 若 $V$ 是 $\mathcal{M}(\mathfrak{g})^{\flat}$ 中的复形，则以 $C_{\mathfrak{g}}(V)$ 记双复形 $C(V^{\cdot})$ 的全复形。这是离散 DG $(\Omega_{\mathfrak{g}},\mathfrak{g}_{\Omega})$-模（简称 $(\Omega_{\mathfrak{g}},\mathfrak{g}_{\Omega})$-复形）。函子 $C_{\mathfrak{g}}$ 是 DG 范畴 $C(\mathfrak{g})^{\flat}$（即 $\mathcal{M}(\mathfrak{g})^{\flat}$ 中的复形范畴，我们称之为 $\mathfrak{g}^{\flat}$-复形）与 $(\Omega_{\mathfrak{g}},\mathfrak{g}_{\Omega})$-复形构成的 DG 范畴 $C(\Omega_{\mathfrak{g}},\mathfrak{g}_{\Omega})$ 之间的等价。逆函子把 $F\in C(\Omega_{\mathfrak{g}},\mathfrak{g}_{\Omega})$ 送到复形 $\operatorname{Hom}_{\mathrm{Cl}_{\mathfrak{g}}}(C^{\cdot},F)$。

\bdp{7.13.29} 设 $\mathfrak{k}\subset\mathfrak{g}$ 是开有界李子代数。对 $a\ge0$，以 $C^{\cdot}_{a}\subset C^{\cdot}$ 记被来自 $\mathfrak{k}^{\perp}\subset\mathfrak{g}^{*}$ 的任意 $a+1$ 个 Clifford 算子的乘积零化的元素构成的子空间。则 $0=C^{\cdot}_{-1}\subset C^{\cdot}_{0}\subset C^{\cdot}_{1}\subset\ldots$ 是 $C^{\cdot}=\cup C^{\cdot}_{a}$ 上的递增滤链。$\mathfrak{g}^{*}$ 中任意 Clifford 算子 $\nu$ 保持\SourcePage{354}%
我们的滤链；若 $\nu$ 属于 $\mathfrak{k}^{\perp}$，则它把 $C^{\cdot}_{a}$ 送到 $C^{\cdot+1}_{a-1}$。来自 $\mathfrak{g}$ 的任意 Clifford 算子把 $C^{\cdot}_{a}$ 送到 $C^{\cdot-1}_{a+1}$；若它属于 $\mathfrak{k}$，则保持滤链。于是 $\operatorname{gr}_{*}C^{\cdot}$ 是向量空间 $(\mathfrak{g}/\mathfrak{k})\oplus(\mathfrak{g}/\mathfrak{k})^{*}\oplus\mathfrak{k}\oplus\mathfrak{k}^{*}$（赋以标准「双曲」形式）的 Clifford 代数 $\mathrm{Cl}_{\mathfrak{g},\mathfrak{k}}$ 上的模。

这是不可约 $\mathrm{Cl}_{\mathfrak{g},\mathfrak{k}}$-模；且 $C^{\cdot}_{0}$ 是子代数 $\mathrm{Cl}_{\mathfrak{k}}\subset\mathrm{Cl}_{\mathfrak{g},\mathfrak{k}}$ 上的不可约模。齐次线 $\lambda_{\mathfrak{k}}=\lambda^{(C)}_{\mathfrak{k}}$（见 7.13.16）位于 $C^{\cdot}_{0}$ 中，且 $\operatorname{gr}_{*}C^{\cdot}$ 是由此线生成的、作为子代数 $\Lambda(\mathfrak{g}/\mathfrak{k})\otimes\Lambda\mathfrak{k}^{*}\subset\mathrm{Cl}_{\mathfrak{g},\mathfrak{k}}$ 上的自由模。若 $\lambda_{\mathfrak{k}}\subset C^{0}$（必要时可平移 $\cdot$-滤链来假定这一点），则 $\operatorname{gr}_{a}C^{b}=\Lambda^{a}(\mathfrak{g}/\mathfrak{k})\otimes\Lambda^{b+a}\mathfrak{k}^{*}\otimes\lambda_{\mathfrak{k}}$。
% [Tutor Note: 原书 p.354 此处作 "(we may assume this shifting the · filtration if necessary)"，语句不顺，"this shifting" 疑为 "this by shifting" 之误，又 "homoguenous" 系 "homogeneous" 之误；底本如此，保留原样。]

设 $\mathfrak{k}^{\flat}\subset\mathfrak{g}^{\flat}$ 是 $\mathfrak{k}$ 的原像。这是 $\mathfrak{k}$ 被 $\mathbb{C}$ 的中心扩张，它典范分裂：分裂 $\mathfrak{k}\to\mathfrak{k}^{\flat}$ 的像由那些零化 $\lambda_{\mathfrak{k}}$ 的元素组成（我们考虑 $\mathfrak{k}^{\flat}$ 在 $C^{\cdot}$ 上的李代数作用）。

对 $V\in C(\mathfrak{g})^{\flat}$，子空间 $C^{\cdot}_{a}\otimes V$ 是 $C_{\mathfrak{g}}(V)$ 的子复形；以 $C_{\mathfrak{g}}(V)_{a}$ 记之。我们得到 $C_{\mathfrak{g}}(V)$ 上的滤链，它被来自 $\mathfrak{g}^{*}$ 与 $\mathfrak{k}$ 的 Clifford 算子保持；相继商 $\operatorname{gr}_{a}C_{\mathfrak{g}}(V)$ 是 $(\Omega_{\mathfrak{k}},\mathfrak{k}_{\Omega})$-复形。对 $\mathfrak{k}$-复形 $P$，以 $C_{\mathfrak{k}}(P)$ 记 $\mathfrak{k}$ 的、系数在 $P$ 中的李代数上链构成的 Chevalley 复形；这是 $(\Omega_{\mathfrak{k}},\mathfrak{k}_{\Omega})$-复形。等同
\[
\operatorname{gr}_{a}C_{\mathfrak{g}}(V)^{\cdot}=\Lambda^{\cdot+a}\mathfrak{k}^{*}\otimes(V^{\cdot}\otimes\Lambda^{a}(\mathfrak{g}/\mathfrak{k})\otimes\lambda_{\mathfrak{k}})
\]
是 $(\Omega_{\mathfrak{k}},\mathfrak{k}_{\Omega})$-复形的同构
\begin{equation}\tag{376}
\operatorname{gr}_{a}C_{\mathfrak{g}}(V)\simeq C_{\mathfrak{k}}(V\otimes\Lambda^{a}(\mathfrak{g}/\mathfrak{k})\otimes\lambda_{\mathfrak{k}})[a]
\end{equation}

这里 $\mathfrak{k}$ 以伴随的方式作用在 $\Lambda^{a}(\mathfrak{g}/\mathfrak{k})$ 上。相应的谱序列收敛到 $H^{\cdot}C_{\mathfrak{g}}(V)$；其首项是 $E^{p,q}_{1}=H^{p+q}\operatorname{gr}_{-p}C_{\mathfrak{g}}(V)=H^{q}(\mathfrak{k},\Lambda^{-p}(\mathfrak{g}/\mathfrak{k})\otimes V\otimes\lambda_{\mathfrak{k}})$。

\bdp{7.13.30} \emph{注.} 假设有 $\mathfrak{k}^{\flat}$-子复形 $T\subset V$，使 $V$ 由 $T$ 诱导出，即 $V=U(\mathfrak{g}^{\flat})\otimes_{U(\mathfrak{k}^{\flat})}T$。则嵌入的复合 $C_{\mathfrak{k}}(T\otimes\lambda_{\mathfrak{k}})\subset C_{\mathfrak{g}}(V)_{0}\subset C_{\mathfrak{g}}(V)$ 是拟同构。

\SourcePage{355}

\bdsec{7.14. 无穷维情形下的局部化函子} 现在我们可以解释 7.1.1 中「Hecke 模式」的部分 (c)、(d) 在当前的无穷维情形下是什么意思。

\bdp{7.14.1} 设 $G,K$ 如 7.11.17，$G'$ 是 $G$ 被 $\mathbb{G}_{m}$ 的中心扩张，并带分裂 $K\to G'$（参看 7.8.1）。则 $\mathfrak{g},\mathfrak{g}'$ 是 Tate 意义下的李代数，而 $\mathfrak{k}=\operatorname{Lie}K$ 是 $\mathfrak{g},\mathfrak{g}'$ 的开有界李子代数。7.8.1 中所有范畴在当前情形下都有显然的意义。

按 7.8.2 的方式在范畴 $D(\mathfrak{g},K)'$ 上定义 Hecke 作用。现在线丛 $\mathcal{L}_{G}$ 是 $G$ 上的 $\mathcal{O}^{p}$-模，而 $\mathcal{V}_{G}$ 是左 $\mathcal{D}^{p}$-模构成的复形（见 7.11.3）。7.8.2 的所有构造逐字搬到我们的情形，7.8.4--7.8.5 亦然（在 7.8.4 中，$U'$ 照例应取完备化扭包络代数）。

\bdp{7.14.2} 为定义局部化函子 $L\Delta$，我们需要一些准备。设 $Y$ 是概形，$F$ 是 Tate 向量空间。$Y$ 上的 $\mathrm{Cl}_{F}$-模是 $Y$ 上的 $\mathbb{Z}$-分次 $\mathcal{O}$-模 $C^{\cdot}$，并带分次 Clifford 代数 $\mathrm{Cl}^{\cdot}_{F}$ 的连续作用（见 7.13.16）。对任意 c-格 $P\subset F$，以 $\lambda_{P}(C^{\cdot})$ 记 $C^{\cdot}$ 的如下分次 $\mathcal{O}$-子模：它由被来自 $P\subset F$ 与 $P^{\perp}\subset F^{*}$ 的 Clifford 算子零化的局部截面组成。函子 $\lambda_{P}:\mathcal{C}(Y)\to\{Y$ 上分次 $\mathcal{O}$-模的范畴$\}$ 是范畴等价\footnote{逆函子是与适当的不可约分次 Clifford 模作张量积。}。对两个 c-格 $P_{1},P_{2}$，有典范同构
\begin{equation}\tag{377}
\lambda_{P_{1}}(C^{\cdot})\simeq\lambda_{(P_{1},P_{2})}\otimes\lambda_{P_{2}}(C^{\cdot})
\end{equation}
它满足显然的传递性（见 7.13.16）。对 c-格的 $Y$-族亦然（见 loc. cit.）。

\bdp{7.14.3} 现在假设处于 7.11.18 的情形。则 $Y$ 带典范的 $\mathrm{Cl}_{\mathfrak{g}}$-模 $C^{\cdot}_{Y}$，其定义如下。设 $K\subset G$ 是合理的群子概形，$\mathfrak{k}:=\operatorname{Lie}K$。以 $\omega_{(K\backslash Y)}$ 记典范
% 尾部接 p.356（本单元在此断处停止）
%% end units/ch07_u12.tex
%% begin units/ch07_u13.tex
% ch07_u13 —— 原书 p.356–p.366（第 7 章 7.14.3–7.15.17）
% 底本：A. Beilinson, V. Drinfeld, Quantization of Hitchin's Integrable System
% and Hecke Eigensheaves
% 承接：上单元（p.346–p.355）末句断在「以 $\omega_{(K\backslash Y)}$ 记典范」，
% 本单元正文自「丛 $\omega_{K\backslash Y}=\det\Omega_{K\backslash Y}$ 沿投影…的拉回」续起。

\SourcePage{356}
丛 $\omega_{K\backslash Y}=\det\Omega_{K\backslash Y}$ 沿投影 $Y\to K\backslash Y$ 的拉回（回忆 $K\backslash Y$ 是光滑叠）。这是分次线丛，位于次数 $\dim K\backslash Y$。

若 $K_1,K_2\subset G$ 是如上两个合理的群子概形，则有典范同构
\begin{equation}\tag{378}
\omega_{(K_1\backslash Y)}=\lambda_{(\mathfrak{k}_1,\mathfrak{k}_2)}\otimes\omega_{(K_2\backslash Y)}
\end{equation}
它满足显然的传递性。事实上，为定义（378），只需考虑 $K_2\subset K_1$ 的情形。把光滑投影 $K_2\backslash Y\to K_1\backslash Y$ 的相对切丛拉回到 $Y$，得 $(\mathfrak{k}_1/\mathfrak{k}_2)\otimes\mathcal{O}_Y$，由此得到（378）。传递性显然。

今 $C_Y^{\cdot}\in\mathcal{C}(Y)$ 是 Clifford 模，并附有一族同构 $\lambda_{\mathfrak{k}}(C_Y^{\cdot})\simeq\omega_{(K\backslash Y)}$，对每个合理的子群 $K\subset G$ 给出，且与（377）、（378）相容。这样的 $C_Y^{\cdot}$ 存在，并在唯一同构的意义下唯一。

$G$ 在 $Y$ 上的作用典范地提升为 $C_Y^{\cdot}$ 上的 $G$-作用，且与 $G$ 在 Clifford 算子 $\mathfrak{g}\oplus\mathfrak{g}^{*}$ 上的伴随作用相容。实际上，$G(\mathbb{C})$ 作用在构造 $C_Y^{\cdot}$ 所涉及的全部对象上，故它作用在 $C_Y^{\cdot}$ 上。要定义 $A$-点 $G(A)$ 在 $C_Y^{\cdot}\otimes A$ 上的作用，须用 $G$ 的 $A$-族合理群子概形，把 $Y\times\operatorname{Spec}A$ 上 Clifford 模 $C_Y^{\cdot}\otimes A$ 的特征性质写清楚。我们把它留给读者。

\emph{注.} 任取 $y\in Y$。$C_y^{\cdot}$ 在 $y$ 处的纤维 $C_y^{\cdot}$ 是不可约分次 $\mathrm{Cl}^{\cdot}_{\mathfrak{g}}$-模，可如下描述。% [Tutor Note: 原书 p.356 此处印作「$C_y^{\cdot}$ 在 $y$ 处的纤维 $C_y^{\cdot}$」，第二个下标当作 $Y$，底本如此，保留原样。]
考虑「作用」映射 $\mathfrak{g}\to\Theta_y$。其核 $\mathfrak{g}_y$（$y$ 的稳定化子）是 $\mathfrak{g}$ 中的 d-格。余核 $T$ 是有限维向量空间。令 $C_{y\mathfrak{g}_y}^{\cdot}$ 为 $C_y^{\cdot}$ 中 $\mathfrak{g}_y$-余不变量（关于 $\mathfrak{g}_y$ 的 Clifford 作用）构成的分次向量空间。则有典范等同 $C_{y\mathfrak{g}_y}^{\dim T}\simeq\det(T^{*})$，而 $C_y^{\cdot}$ 由这一规范化唯一确定。

\bdp{7.14.4} 设 $\mathcal{L}=\mathcal{L}_Y$ 是 $Y$ 上的线丛，带有提升 $Y$ 上 $G$-作用的 $G'$-作用；我们假定 $\mathbb{G}_m\subset G$ 以与标准相反的特征标作用在 $\mathcal{L}$ 上。

\SourcePage{357}
取 $V\in\mathcal{M}(\mathfrak{g})'$，于是 $V$ 是离散 $\mathfrak{g}'$-模，其中 $\mathbb{C}\subset\mathfrak{g}'$ 按标准特征标作用。此时张量积 $\mathcal{L}\otimes V$ 是 $\mathfrak{g}$-模，$C_Y^{\cdot}\otimes\mathcal{L}\otimes V$ 亦然（即 $Y$ 上的 $\mathfrak{g}$-作用提升为这些 $\mathcal{O}$-模上的连续 $\mathfrak{g}$-作用）。我们把 $\alpha\in\mathfrak{g}$ 在 $C_Y^{\cdot}\otimes\mathcal{L}\otimes V$ 上的作用记为 $\operatorname{Lie}_{\alpha}$。注意 $C_Y^{\cdot}\otimes\mathcal{L}\otimes V$ 还是 Clifford 模，且上述 $\mathfrak{g}$-作用与 Clifford 算子相容。照例把 $\alpha\in\mathfrak{g}$ 的 Clifford 作用记为 $i_{\alpha}$。于是与 7.13.26 一样，$C_Y^{\cdot}\otimes\mathcal{L}\otimes V$ 是分次 $(\Omega_{\mathfrak{g}},\mathfrak{g}_{\Omega})$-模。

下述命题与 7.13.27 类似，其证明亦然，我们留给读者。

\bdp{7.14.5} \emph{命题.} 存在唯一层态射
\[
d:C_Y^{\cdot}\otimes\mathcal{L}\otimes V\to C_Y^{\cdot+1}\otimes\mathcal{L}\otimes V,
\]
使得对任意 $\alpha\in\mathfrak{g}$ 有 $\operatorname{Lie}_{\alpha}=di_{\alpha}+i_{\alpha}d$。这个 $d$ 是一阶微分算子，$d^2=0$，且 $C_{\mathcal{L}}(V):=(C_Y^{\cdot}\otimes\mathcal{L}\otimes V,d)$ 是 DG $(\Omega_{\mathfrak{g}},\mathfrak{g}_{\Omega})$-模。

\emph{注.} 也可由 7.13.27 直接推出 7.14.5。取任意如 7.14.3 的 $K$。则 $C_Y^{\cdot}\otimes\omega^{\ast}_{(K\backslash Y)}$ 是「常」Clifford 模：它典范同构于 $C^{\cdot}\otimes\mathcal{O}_Y$，其中 $C^{\cdot}$ 是某个不可约 Clifford 模。$C^{\cdot}$ 上的 $\mathfrak{g}^{\flat}$-作用与 $C_Y^{\cdot}$ 上的 $\mathfrak{g}^{\flat}$-作用给出 $\omega_{(K\backslash Y)}=\operatorname{Hom}(C_Y^{\cdot},C^{\cdot}\otimes\mathcal{O}_Y)$ 上的 $\mathfrak{g}^{\flat}$-作用，它提升 $Y$ 上的 $\mathfrak{g}$-作用。于是 $\mathfrak{g}^{\flat}$ 作用在 $\omega_{(K\backslash Y)}\otimes\mathcal{L}\otimes V$ 上，而 7.14.5 中的 $d$ 与 7.13.27 中关于 $C^{\cdot}\otimes(\omega_{(K\backslash Y)}\otimes\mathcal{L}\otimes V)$ 的 $d$ 一致。

\bdp{7.14.6} 这样我们定义了 $Y$ 上的 Ω-复形 $C_{\mathcal{L}}(V)$。当 $V$ 是 $\mathcal{M}(\mathfrak{g})'$ 中的复形时，用显然的方式延拓此定义。今设 $K$ 如 7.14.1。对 Harish-Chandra 复形 $V\in C(\mathfrak{k}_{\Omega}\times\mathfrak{g},K)'$，Ω-复形 $C_{\mathcal{L}}(V)$ 是 $K_{\Omega}$-等变的。实际上，$K$ 依其在 $C_Y^{\cdot}$、$\mathcal{L}$、$V$ 上的作用作用在 $C_{\mathcal{L}}(V)$ 上，而算子 $i_{\xi}$（$\xi\in\mathfrak{k}$）等于 $C_Y^{\cdot}$ 的相应 Clifford 算子与 $V$ 上 $\mathfrak{k}_{\Omega}^{-1}$-作用的算子之和。

\SourcePage{358}
置 $\Delta_{\Omega\mathcal{L}}(V):=C_{\mathcal{L}}(V)[\dim(K\backslash Y)]$。我们已定义 DG 函子
\begin{equation}\tag{379}
\Delta_{\Omega}=\Delta_{\Omega\mathcal{L}}:C(\mathfrak{k}_{\Omega}\times\mathfrak{g},K)'\longrightarrow C(K\backslash Y,\Omega)
\end{equation}

\bdp{7.14.7} \emph{注.} Ω-复形 $\Delta_{\Omega}(V)$ 带有典范滤过 $\Delta_{\Omega}(V)_{\cdot}$，其中 $\Delta_{\Omega}(V)_a$ 由这样的截面组成：$\mathfrak{k}^{\perp}\subset\mathfrak{g}^{*}$ 中任意 $a+1$ 个 Clifford 算子的乘积作用在它上面为零（见 7.13.29）。% [Tutor Note: 原书 p.358 此处印作「filtrartion」，应为「filtration」，底本如此，保留原样。]
由（376）有 $K_{\Omega}$-等变 Ω-复形的典范同构
\begin{equation}\tag{380}
\operatorname{gr}_a\Delta_{\Omega}(V)\xrightarrow{\ \sim\ }C_{\mathfrak{k}}(\omega_{(K\backslash Y)}\otimes\mathcal{L}\otimes V\otimes\Lambda^{a}(\mathfrak{g}/\mathfrak{k}))[a]
\end{equation}

\bdp{7.14.8} \emph{引理.} (i) 函子 $\Delta_{\Omega}$ 把拟同构送到 D-拟同构，故给出三角函子
\begin{equation}\tag{381}
L\Delta=L\Delta_{\mathcal{L}}:D(\mathfrak{g},K)'\to D(K\backslash Y)
\end{equation}
(ii) 函子 $L\Delta$ 右 t-正合，相应的右正合函子 $\Delta=\Delta_{\mathcal{L}}:\mathcal{M}((\mathfrak{g},K)'\to\mathcal{M}^{\ell}(K\backslash Y)$ 为
\begin{equation}\tag{382}
\Delta_{\mathcal{L}}(V)_Y=(\mathcal{D}_Y\otimes\mathcal{L})\widehat{\otimes}_{U(\mathfrak{g})}V=\mathcal{L}^{*}\otimes\mathcal{D}_{Y,\mathcal{L}}\widehat{\otimes}_{U(\mathfrak{g}')}V
\end{equation}
这里 $\mathcal{D}_Y$ 是 $Y$ 上微分算子的拓扑代数（见 1.2.6），$\mathcal{D}_{Y,\mathcal{L}}:=\mathcal{L}\otimes\mathcal{D}_Y\otimes\mathcal{L}^{*}$ 是相应的 $\mathcal{L}$-扭代数。% [Tutor Note: 原书 p.358 此处印作「$\mathcal{M}((\mathfrak{g},K)'$」，左括号多了一个，底本如此，保留原样。]

\emph{证明.} (i) 我们的命题是局部的，故必要时缩小 $K$，可设 $K$ 在 $Y$ 上的作用自由。把 $\Delta_{\Omega}(V)$ 视为 $K\backslash Y$ 上的滤过 Ω-复形。对 $K$-模 $P$，记 $P^{\sim}$ 为其 $Y$-扭，即 $K\backslash Y$ 上的 $\mathcal{O}$-模。投影 $C_{\mathfrak{k}}\to C_{\mathfrak{k}}/C_{\mathfrak{k}}^{\ge1}$ 按（380）给出典范同构
\begin{equation}\tag{383}
\operatorname{gr}_a\Delta_{\Omega}(V)_{K\backslash Y}=\omega_{K\backslash Y}\otimes\mathcal{L}_{K\backslash Y}\otimes V^{\sim}\otimes\Lambda^{a}(\mathfrak{g}/\mathfrak{k})^{\sim}[a]
\end{equation}
右端是 $\mathcal{O}$-复形，故 $V$ 之间的拟同构给出 $\Delta_{\Omega}(V)$ 之间的（滤过）D-拟同构。

(ii) 同上可设 $K$ 的作用自由。对 $V\in\mathcal{M}(\mathfrak{g},K)'$，可把（383）改写为同构 $\Delta_{\Omega}(V)^{a}_{K\backslash Y}=\omega_{K\backslash Y}\otimes\mathcal{L}_{K\backslash Y}\otimes V^{\sim}\otimes\Lambda^{-a}(\mathfrak{g}/\mathfrak{k})^{\sim}$。这说明 $\Delta_{\Omega}$ 右 t-正合。\SourcePage{359}%
$\Delta_{\Omega}(V)_{K\backslash Y}$ 中的微分如下。% [Tutor Note: 原书 p.359 此处印作「The differential in $\Delta_{\Omega}(V)_{K\backslash Y}$ as follows.」，缺谓语（疑脱 is），底本如此，译文按「如下」处理。]
$Y$ 上的 $\mathfrak{g}$-作用在 $(\mathfrak{g}/\mathfrak{k})^{\sim}$ 上定义了 $K\backslash Y$ 上 Lie 代数胚的结构。$\mathcal{L}_Y\otimes V$ 上的 $\mathfrak{g}$-作用在 $\mathcal{L}_{K\backslash Y}\otimes V^{\sim}$ 上定义了左 $(\mathfrak{g}/\mathfrak{k})^{\sim}$-模结构，故 $\omega_{K\backslash Y}\otimes\mathcal{L}_{K\backslash Y}\otimes V^{\sim}$ 是右 $(\mathfrak{g}/\mathfrak{k})^{\sim}$-模。于是 $\Delta_{\Omega}(V)_{K\backslash Y}$ 是 $(\mathfrak{g}/\mathfrak{k})^{\sim}$ 的、系数在 $\omega_{K\backslash Y}\otimes\mathcal{L}_{K\backslash Y}\otimes V^{\sim}$ 中的 Chevalley 同调复形。$K\backslash Y$ 上的右 D-模 $H^0_D(L\Delta(V))$ 为 $(\omega_{K\backslash Y}\otimes\mathcal{L}_{K\backslash Y}\otimes V^{\sim})\otimes_{(\mathfrak{g}/\mathfrak{k})^{\sim}}\mathcal{D}_{K\backslash Y}$；相应的左 D-模是 $\mathcal{D}_{K\backslash Y}\otimes_{(\mathfrak{g}/\mathfrak{k})^{\sim}}(\mathcal{L}_{K\backslash Y}\otimes V^{\sim})$。把这一同构提升到 $Y$ 上即得（382）。

\bdp{7.14.9} \emph{例.} 我们来计算 $L\Delta(Vac')$。嵌入 $\mathbb{C}\to Vac'$ 给出 $Y$ 上 Ω-复形的嵌入 $C_{\mathfrak{k}}(\omega_{(K\backslash Y)}\otimes\mathcal{L}_Y)\to\Delta_{\Omega\mathcal{L}}(Vac')_0$。相应的 $K_{\Omega}$-等变 Ω-复形态射 $C_{\mathfrak{k}}(\omega_{(K\backslash Y)}\otimes\mathcal{L}_Y)\to\Delta_{\Omega\mathcal{L}}(Vac')$ 是 D-拟同构，这一验证留给读者。今左端是 $K\backslash Y$ 上的 Ω-复形 $\Omega(\mathcal{D}_{K\backslash Y}\otimes\mathcal{L}_{K\backslash Y}$（见 7.3.3）。% [Tutor Note: 原书 p.359 此处括号不配对，右括号缺失，底本如此，保留原样。]
因此若 $K\backslash Y$ 是良叠，则
\[
L\Delta(Vac')=\Delta(Vac')=\mathcal{D}_{K\backslash Y}\otimes\mathcal{L}_{K\backslash Y}.
\]

\emph{注.} 由于 $\operatorname{End}Vac'$ 反同构于 1.2.5 中的代数 $D'_{(\mathfrak{g},K)}$（另见 1.2.2），我们有 $D'_{(\mathfrak{g},K)}$ 在 $\Delta(Vac')=\mathcal{D}_{K\backslash Y}\otimes\mathcal{L}_{K\backslash Y}$ 上的右作用，即从 $D'_{(\mathfrak{g},K)}$ 到扭微分算子环 $\Gamma(K\backslash Y,D'_{K\backslash Y})$ 的同态。这就是 1.2.5 中的同态 $h$（另见 1.2.3 与 1.2.4）。

\bdp{7.14.10} \emph{命题.} 函子 $L\Delta:D(\mathfrak{g},K)'\to D(K\backslash Y)$ 是 H-模的态射。

\emph{证明.} 7.8.8 的构造与论证以显然的方式移植到我们的无穷维情形。7.9 的无穷维版本是直接的。

\SourcePage{360}
\bdsec{7.15. 仿射旗空间是 \texorpdfstring{$\mathcal{D}$}{D}-仿射的}

本节说明：水平低于临界值的仿射李代数表示与仿射旗空间上的 $\mathcal{D}$-模之间的关系，与通常的有限维情形相同。

\bdp{7.15.1} 下文恒设 $K=\mathbb{C}((t))$，$O=\mathbb{C}[[t]]$。设 $\mathfrak{g}$ 是单（有限维）李代数\footnote{当 $\mathfrak{g}$ 是任意 reductive 李代数时，此推广是直接的。}，$G$ 是相应的单连通单群。我们有群归纳概形 $G(K)$ 及其群子概形 $G(O)$（见 7.11.2(iv)）。$G(K)$ 在 Tate 向量空间 $\operatorname{Lie}G(K)=\mathfrak{g}(K)$ 上的伴随作用给出 $G(K)$ 被 $\mathbb{G}_m$ 的中心扩张 $G(K)^{\flat}$（见 ??）。其李代数是 $\mathfrak{g}(K)$ 的中心扩张 $\mathfrak{g}(K)^{\flat}$，由余圈 $\phi,\psi\mapsto\operatorname{Res}(d\phi,\psi)$ 定义，其中 $(a,b):=\operatorname{Tr}(\operatorname{ad}_a\cdot\operatorname{ad}_b)$（见 ??）。令 $G(O)^{\flat}\subset G(K)^{\flat}$ 为 $G(O)$ 的原像。$G(O)$ 的伴随作用保持 c-格 $\mathfrak{g}(O)\subset\mathfrak{g}(K)$，故有典范等同 $s:G(O)^{\flat}\xrightarrow{\sim}G(O)\times\mathbb{G}_m$\footnote{由于 $G$ 单，分裂 $G(O)\to G(O)^{\flat}$ 唯一。}。% [Tutor Note: 原书 p.360 两处交叉引用均印作「??」（引用未解析），底本如此，保留原样。]

设 $N\subset B\subset G$，其中 $B$ 是 Borel 子群，$N$ 是其幂幺根，故 $H=B/N$ 是 $G$ 的嘉当群。令 $N^{+},B^{+}$ 为 $N,B$ 在显然投影 $G(O)\to G$ 下的原像，故 $B^{+}/N^{+}=H$，$G(O)/B^{+}=G/B$。令 $B^{\dagger}\subset G(K)^{\flat}$ 为 $B^{+}$ 的原像。存在唯一截面 $N^{+}\to G(K)$；置 $H^{\flat}:=B^{\dagger}/N^{+}$，$\mathfrak{h}^{\flat}=\operatorname{Lie}H^{\flat}$。截面 $s$ 给出同构 $B^{+}\times\mathbb{G}_m\xrightarrow{\sim}B^{\dagger}$，从而有同构 $H\times\mathbb{G}_m\xrightarrow{\sim}H^{\flat}$、$\mathfrak{h}\times\mathbb{C}\xrightarrow{\sim}\mathfrak{h}^{\flat}$。

置 $X:=G(K)/B^{+}=G(K)^{\flat}/B^{\dagger}$（层的商，取平坦或 Zariski 拓扑都一样，这一点由 4.5.1 推出）。$X$ 称为\emph{仿射旗空间}。它是既约、连通、归纳投射、形式光滑的归纳概形\footnote{$X$ 光滑地纤维化到仿射 Grassmann 流形 $G(K)/G(O)$ 上，见 4.5.1。}。置 $X^{\dagger}:=G(K)^{\flat}/N^{+}$：这是 $X$ 上的左 $H^{\flat}$-torsor（主齐性空间），作用为 $h^{\flat}\cdot x^{\dagger}=x^{\dagger}h^{\flat-1}$。它带有显然的 $G(K)^{\flat}$-作用。记投影 $X^{\dagger}\to X$ 为 $p$。

\SourcePage{361}
\bdp{7.15.2} 令 $\mathcal{M}^{\dagger}(X)$ 为 $X^{\dagger}$ 上弱 $H^{\flat}$-等变 $\mathcal{D}$-模的范畴（见 7.11.11）。这是 Abel 范畴。对 $M\in\mathcal{M}^{\dagger}(X)$ 置 $M_X:=(p_{\cdot}M)^{H^{\flat}}\in\mathcal{M}(X,\mathcal{O})$。函子 $\mathcal{M}^{\dagger}(X)\to\mathcal{M}(X,\mathcal{O})$，$M\mapsto M_X$，正合且忠实。

置 $\mathcal{D}^{\dagger}:=(p_{\cdot}\mathcal{D}^{\dagger}_X)^{H^{\flat}}$。这是 $X$ 上的 Diff-代数。映射
\begin{equation}\tag{384}
\mathfrak{h}^{\flat}\to\Gamma(X,\mathcal{D}^{\dagger})=\Gamma(X^{\dagger},\mathcal{D}^{\dagger}_{X^{\dagger}})^{H^{\flat}}
\end{equation}
等于沿 $p$ 的纤维的左作用的相反数，其像落在 $\mathcal{D}^{\dagger}$ 的中心内。实际上，$\mathcal{D}^{\dagger}$ 是 $\operatorname{Sym}(\mathfrak{h}^{\flat})$-族的 tdo（见 7.11.11(b)）。

注意 $\mathcal{D}^{\dagger}$ 以显然的方式（从右）作用在如上任何 $M_X$ 上，故有函子
\begin{equation}\tag{385}
\mathcal{M}^{\dagger}(X)\to\mathcal{M}(X,\mathcal{D}^{\dagger}).
\end{equation}
有（见 7.11.11 中的注 (ii)）：

\bdp{7.15.3} \emph{引理.} 函子（385）是范畴等价。\hfill$\Box$

\bdp{7.15.4} 对 $\chi=(\chi_0,c)\in\mathfrak{h}^{\flat*}=\mathfrak{h}^{*}\times\mathbb{C}$，记 $\mathcal{D}^{\chi}$ 为我们族 $\mathcal{D}^{\dagger}$ 中相应的 tdo。于是 $\mathcal{D}^{(0,0)}=\mathcal{D}_X$。置 $\mathcal{M}^{\chi}(X):=\mathcal{M}(X,\mathcal{D}^{\chi})\subset\mathcal{M}(X,\mathcal{D}^{\dagger})$。考虑拓扑代数 $\Gamma\mathcal{D}^{\chi}=\Gamma(X,\mathcal{D}^{\chi})$（见 7.11.9、7.11.10）。我们有函子
\begin{equation}\tag{386}
\Gamma:\mathcal{M}^{\chi}(X)\to\mathcal{M}^{r}(\Gamma\mathcal{D}^{\chi})
\end{equation}
其中 $\mathcal{M}^{r}(\Gamma\mathcal{D}^{\chi})$ 是离散右 $\Gamma\mathcal{D}^{\chi}$-模的范畴，$\Gamma M:=\Gamma(X,M)$。

$\mathfrak{g}(K)^{\flat}$ 在 $X^{\dagger}$ 上的作用给出连续态射 $\mathfrak{g}(K)^{\flat}\to\Gamma(X,\mathcal{D}^{\dagger})$。相应的态射 $\mathfrak{g}(K)^{\flat}\to\Gamma\mathcal{D}^{\chi}$ 把 $1^{\flat}\in\mathfrak{g}(K)^{\flat}$ 送到 $-c$。

\bdp{7.15.5} 若 Verma $\mathfrak{g}(K)^{\flat}$-模 $M(\chi)$ 不可约，则称 $\chi$ 是\emph{反占优}的。据 [KK] 3.1，这等价于以下三个条件：
\begin{enumerate}[label=(\roman*),leftmargin=2.2em]
\item 有 $c\ne-1/2$。
\SourcePage{362}%
\item 对 $\mathfrak{g}$ 的任意正余根 $h_{\alpha}\in\mathfrak{h}$ 有 $(\chi_0+\rho_0)(h_{\alpha})\ne1,2,\ldots$。
\item 对如上任意 $h_{\alpha}$ 及任意整数 $n>0$ 有
\[
\pm(\chi_0+\rho_0)(h_{\alpha})+2n\frac{c+1/2}{(\alpha,\alpha)}\ne1,2,\ldots
\]
\end{enumerate}

这里 $\rho_0\in\mathfrak{h}^{*}$ 是 $\mathfrak{g}$ 的正根之和的一半，$(\cdot,\cdot)$ 是 $\mathfrak{h}^{*}$ 上与 $\mathfrak{h}$ 上 $(\cdot,\cdot)$ 相对应的标量积（见 7.15.1）。

\emph{注.} 为从 [KK] 3.1 推出上述论断，只需注意 $\mathfrak{g}(K)^{\flat}$ 的「实」正余根是 $h_{\alpha}$ 与 $\pm h_{\alpha}+2n(\alpha,\alpha)^{-1}\cdot1^{\flat}$，其中 $h_{\alpha}$、$n$ 如上；而 [KK] 中的权 $\rho$ 由下式给出。

置 $\rho:=(\rho_0,1/2)\in\mathfrak{h}^{\flat*}$。若 $\chi+\rho$ 在仿射 Weyl 群 $W_{\mathrm{aff}}$ 中的稳定化子平凡，则称 $\chi$ 是\emph{正则}的\footnote{回忆 $W_{\mathrm{aff}}$ 在 $\mathfrak{h}^{\flat*}$ 上的作用来自 $G(K)$ 在 $\mathfrak{g}(K)^{\flat}$ 上的伴随作用。}。

\bdp{7.15.6} \emph{定理.} 设 $\chi$ 反占优且正则。则（386）是范畴等价。

我们在下文 7.15.8--?? 中证明 7.15.6。

\bdp{7.15.7} \emph{注.} (i) 令 $\mathcal{M}^{c}(\mathfrak{g}(K))$ 为离散 $\mathfrak{g}(K)^{\flat}$-模的范畴，其中 $1^{\flat}$ 以 $c$ 为乘子作用。令
\begin{equation}\tag{387}
\Gamma:\mathcal{M}^{\chi}(X)\to\mathcal{M}^{c}(\mathfrak{g}(K))
\end{equation}
为（386）与显然的「限制」函子 $\mathcal{M}^{r}(\Gamma\mathcal{D}^{\chi})\to\mathcal{M}^{c}(\mathfrak{g}(K))$ 的合成。据 7.15.6，此函子正合且忠实。

(ii) 人们可以期望 $\mathfrak{g}(K)^{\flat}$ 在 $\Gamma\mathcal{D}^{\chi}$ 中生成稠密子代数\footnote{这等价于：对 $M\in\mathcal{M}^{\chi}(X)$，$\Gamma M$ 的任何 $\mathfrak{g}(K)^{\flat}$-子模都来自 $M$ 的 $\mathcal{D}^{\chi}$-子模。}。换言之，$\Gamma\mathcal{D}^{\chi\circ}$ 是水平 $c$ 的包络代数 $\bar U^{c}=\bar U^{c}\mathfrak{g}(K)$ 按某种拓扑的完备化。% [Tutor Note: 原书 p.362 此处上标印作 $\chi\circ$（$\chi$ 后带一小圆圈），同段前文印作 $\Gamma\mathcal{D}^{\chi}$，底本如此，保留原样。]
能否显式确定这一拓扑？

注意，在有限维情形（见 [BB81] 或 [Kas]），人们通常由相应的「经典」论断推出\SourcePage{363}%
（利用 Kostant 的正规性定理）。这个「经典」论断（它说 $\mathfrak{g}(K)\hookrightarrow\Gamma(X,\Theta_X)$ 在 $\bigoplus_{n\ge0}\Gamma(X,\Theta_X^{\otimes n})$ 中生成稠密子代数）对仿射旗情形是\emph{假的}（例如，映射 $\mathfrak{g}(K)\to\Gamma(X,\Theta_X)$ 不满）。

如同 [BB81] 或 [Kas]，容易看出 7.15.6 由下述论断推出：

\bdp{7.15.8} \emph{定理.} (i) 若 $\chi$ 反占优，则对任意 $M\in\mathcal{M}^{\chi}(X)$ 及任意 $r>0$ 有 $H^{r}(X,M)=0$\footnote{这里 $H^{r}(X,M):=\varinjlim H^{r}(Y,M_{(Y)})$；我们用 7.11.4 的记号。}。

(ii) 若此外 $\chi$ 正则且 $M\ne0$，则 $\Gamma M\ne0$。

\emph{注.} 7.15.8(i) 的证明与相应的有限维论断的证明十分相似（见 [BB81] 或 [Kas]）。若能找到与 [BB81]（用平移函子）类似的方法来证明 7.15.8(ii)，那将很好，因为它有助于理解 7.15.7(ii)。

\bdp{7.15.9} 我们开始证明 7.15.8(i)。设 $\psi=(\psi_0,b)$ 是 $H^{\flat}$ 的特征标，$\mathcal{L}=\mathcal{L}^{\psi}$ 是相应的 $X$ 上 $G(K)^{\flat}$-等变线丛（由 $X^{\dagger}$ 定义）。假定 $\mathcal{L}$ 丰沛。这等价\footnote{见 7.15.5 中的注。}于 $\psi$ 的下述性质：对 $\mathfrak{g}$ 的任意正余根 $h_{\alpha}$ 有 $\frac{2b}{(\alpha,\alpha)}<\psi_0(h_{\alpha})<0$。

记 $V$ 为 pro-有限维向量空间 $\Gamma(X,\mathcal{L})$ 的对偶。% [Tutor Note: 原书 p.363 此处印作「Denote by $V$ be the dual…」，多出一 be，底本如此，保留原意。]
它以显然的方式成为 $G(K)^{\flat}$-模，从而是水平 $-b$ 的可积 $\mathfrak{g}(K)^{\flat}$-模\footnote{据 Borel--Weil 定理的一个变体（例如见 [?]），$V$ 是不可约 $\mathfrak{g}(K)^{\flat}$-模。}。考虑 $V\widehat{\otimes}\mathcal{L}$ 的典范截面；这是 $\mathcal{O}^{\flat}$-模的 $G(K)^{\flat}$-等变态射 $\mathcal{O}_X\to V\widehat{\otimes}\mathcal{L}$。把它与 $M$ 张量，得 $\mathcal{O}^{!}$-模的态射
\begin{equation}\tag{388}
i:M\to V\otimes\mathcal{L}\otimes M
\end{equation}
它与 $\mathfrak{g}(K)^{\flat}$ 的作用交换。% [Tutor Note: 原书 p.363 此两处 $\mathcal{O}$ 的上标在扫描件中字形不一（前一处似 $\flat$，后一处为细竖笔），据 7.15.10--7.15.11 的 $\mathcal{O}^{!}$-模与 !-层记法，统一作 $\mathcal{O}^{!}$。]

\SourcePage{364}
\bdp{7.15.10} 下面我们考虑 $X$ 上向量空间的 \emph{!-层}。这种对象 $F$ 是一条规则：对闭子概形 $Y\subset X$ 指定 $Y$ 的 Zariski 拓扑上的层 $F_{(Y)}$，并对 $Y\subset Y'$ 指定等同 $i^{!}_{YY'}F_{(Y')}=F_{(Y)}$\footnote{这里 $i^{!}_{YY'}F_{(Y')}:=$ 在 $Y$ 上（集合论意义下）支撑的截面构成的子层。}，它们满足显然的传递性（对照 7.11.4 的注 (i)）。注意 !-层构成 Abel 范畴。它把 $Y$ 上的层范畴作为对子商与扩张封闭的全子范畴包含在内。$X$ 上任何 $\mathcal{O}^{!}$-模 $M$ 给出 $X$ 上的 !-层 $\varinjlim M_{(Y)}$（于是 $Y$ 上相应的层是 $M_{(Y^{\wedge})}$\footnote{记号见 7.11.4。}）；为简便起见，我们仍把它记为 $M$。我们还将考虑 $\mathfrak{g}(K)^{\flat}$-模的 !-层，即配有 $\mathfrak{g}(K)^{\flat}$-作用的向量空间 !-层，使得在每个 $F_{(Y)}$ 上的作用以显然的意义离散。任何带有 $\mathfrak{g}(K)^{\flat}$-作用的 $\mathcal{O}^{!}$-模都可视为 $\mathfrak{g}(K)^{\flat}$-模的 !-层。

\bdp{7.15.11} \emph{命题.} 作为 $\mathfrak{g}(K)^{\flat}$-模的 !-层之间的态射，（388）是直和项嵌入。

\bdp{7.15.12} \emph{7.15.8(i) 的证明.} 任取 $\alpha\in H^{r}(X,M)=\varinjlim H^{r}(X_{(Y)},M_{(Y)})$。它来自某个闭子概形 $Y\subset X$ 及 $\mathcal{O}$-凝聚子模 $F\subset M_{(Y)}$。选取如上丰沛的 $\mathcal{L}$，使 $H^{r}(Y,\mathcal{L}\otimes F)=0$。由于 $i(\alpha)$ 属于 $H^{r}(Y,V\otimes\mathcal{L}\otimes F)$ 的像，它为零。由 7.15.11 得证。\hfill$\Box$

\bdp{7.15.13} \emph{7.15.11 的证明.} 我们将定义 $V\otimes\mathcal{L}\otimes M$ 的自同态 $A$，使得
\begin{equation}\tag{389}
\operatorname{Ker}A=M,\quad V\otimes\mathcal{L}\otimes M=\operatorname{Ker}A\oplus\operatorname{Im}A.
\end{equation}
这就证明了 7.15.11。

令 $\bar U:=\bar U\mathfrak{g}(K)^{\flat}$ 为 $\mathfrak{g}(K)^{\flat}$ 的通常完备化包络代数。考虑由公式（85）定义的 Sugawara 元 $\widetilde{\mathfrak{L}}_0\in\bar U$。对任意 $ft^{r}\in\mathfrak{g}((t))\subset\bar U$ 有 $[\widetilde{\mathfrak{L}}_0,ft^{r}]=(1^{\flat}+1/2)rft^{r}$（见（87））。对任意\SourcePage{365}%
$N\in\mathcal{M}^{c}(\mathfrak{g}(K))$（其中 $e\ne-1/2$），考虑作用在 $N$ 上的算子 $\Delta_N:=(e+1/2)^{-1}\widetilde{\mathfrak{L}}_0$。若又有 $e-b\ne-1/2$，则置
\begin{equation}\tag{390}
A_{V,N}:=\Delta_{V\otimes N}-\Delta_V\otimes\operatorname{id}_N-\operatorname{id}_V\otimes\Delta_N\in\operatorname{End}(V\otimes N).
\end{equation}
此算子与 $\mathfrak{g}(K)^{\flat}$ 的作用交换。

我们把这一构造用于 $\mathfrak{g}(K)^{\flat}$-模的 !-层 $N:=\mathcal{L}\otimes M$（于是 $e=b+c$，水平的条件满足）。置
\begin{equation}\tag{391}
A:=A_{V,\mathcal{L}\otimes M}\in\operatorname{End}(V\otimes\mathcal{L}\otimes M).
\end{equation}
我们来证明 $A$ 满足（389）。\hfill$\Box$

\bdp{7.15.14} 现在转向 7.15.8(ii)。它是下述命题的直接推论，该命题特别说明如何用 $\Gamma M$ 计算 $M$ 的纤维。先从记号说起。

考察 $X$ 按 $N^{+}$-轨道（Schubert 胞腔）的分层。胞腔由仿射 Weyl 群 $W_{\mathrm{aff}}$ 的元素标号。对 $w\in W_{\mathrm{aff}}$，相应的胞腔为 $i_w:Y_w\hookrightarrow X$，其维数为 $l(w)$。$H^{\flat}$-torsor（主齐性空间）$X^{\dagger}$ 到 $Y_w$ 的限制平凡\footnote{在 $X^{\dagger}$ 中、位于 $Y_w$ 上方的任何 $N^{+}$-轨道都给出一个截面。}。由于 $Y_w$ 上任何可逆函数都是常数，平凡化在相差常数平移的意义下唯一。因此 tdo $\mathcal{D}^{\chi}$ 拉回到 $Y_w$ 上是典范平凡化的。

令 $M$ 是导出范畴 $D(X,\mathcal{D}^{\chi})$ 的任意对象\footnote{其定义与 7.11.14 在非扭情形给出的定义类似。}。对任意 $w\in W_{\mathrm{aff}}$，我们有（非扭的，正如刚才所说明）$\mathcal{D}$-复形 $i_w^{!}M\in D(Y_w)$。

我们想计算李代数（连续）上同调 $H^{a}(\mathfrak{n}^{+},\Gamma M)$（注意，由于 7.15.8(i)，$\Gamma=\mathrm{R}\Gamma$）。由于 $\mathfrak{h}^{\flat}=\mathfrak{b}^{\dagger}/\mathfrak{n}^{+}$，这些是 $\mathfrak{h}^{\flat}$-模。我们假定 $\chi$ 正则。

\bdp{7.15.15} \emph{命题.} 存在典范同构
\[
H^{a}(\mathfrak{n}^{+},\Gamma M)\simeq\bigoplus_{w\in W_{\mathrm{aff}}}H^{a-l(w)}_{DR}(Y_w,i_w^{!}M),
\]
\SourcePage{366}%
使得 $\mathfrak{h}^{\flat}$ 在第 $w$ 个直和项上以 $w(\chi)$ 为乘子作用\footnote{回忆 $G(K)$ 在 $\mathfrak{g}(K)^{\flat}$ 上的伴随作用给出 $\mathfrak{h}^{\flat}$ 上的 $W_{\mathrm{aff}}$-作用。}。

\bdp{7.15.16} \emph{7.15.8(ii) 的证明.} 由于 $\Gamma$ 正合，可设 $M$ 紧支且有限生成。令 $Y\subset X$ 是 $M$ 的（约化）支撑中的光滑 Zariski 开子集。则 $M_{(Y)}$ 是光滑概形 $Y$ 上的凝聚 $\mathcal{D}$-模。故再缩小 $Y$，可设 $M_{(Y)}$ 是自由 $\mathcal{O}_Y$-模。今对任意 $x\in Y$ 有 $H^{\cdot}i_x^{!}M\ne0$。平移 $M$ 可设 $x=Y_1$。由 7.15.15 有 $H^{\cdot}(\mathfrak{n}^{+},\Gamma M)\ne0$，故 $\Gamma M\ne0$。\hfill$\Box$

\bdp{7.15.17} \emph{7.15.15 的证明.} 可设 $M=i_{w*}N$，其中 $N\in D(Y_w)$。确实，任何 $M\in D(X,\mathcal{D}^{\chi})$ 带有典范滤过，其 $\operatorname{gr}_iM=\bigoplus_{l(w)=i}i_{w*}i_w^{!}M$。今 7.15.15 对 $M$ 的同构来自对 $i_{w*}i_w^{!}M$ 的相应同构，连同 $\mathfrak{h}^{\flat}$ 作用的谱分解。这里我们用到 $\chi$ 正则的假设；论证的其余部分只需 $\chi$ 反占优。

先考虑 $M=\delta$ 的情形，此时 $\Gamma\delta$ 是 7.15.5 中的 Verma 模（见 7.15.7(iii)）。% [Tutor Note: 原书 p.366 引「7.15.7(iii)」，但 7.15.7 只有 (i)(ii)，底本如此，保留原样。]
这个 Verma 模是秩 1 的余自由 $N^{+}$-模（它由任何不消去真空向量的泛函 $\nu$ 余自由生成）\footnote{$\nu$ 的核不含非平凡 $\mathfrak{n}^{+}$-子模（否则，由于 $\mathfrak{n}^{+}$ 幂零，它必含 $\mathfrak{n}^{+}$-不变向量，这与 7.15.5(i) 矛盾）。故由 $\nu$ 定义的从 $\Gamma\delta$ 到余自由 $N^{+}$-模的态射是单射。再由维数理由，它是同构。}。于是 $H^{\cdot}(\mathfrak{n}^{+}_x,\Gamma\delta)=H^{0}(\mathfrak{n}^{+}_x,\Gamma\delta)^{\chi}=\mathbb{C}\cdot vac$。又 $H^{\cdot}i_x^{!}\delta=H^{0}i_x^{!}\delta=\mathbb{C}\cdot vac$，故得所要的同构。
%% end units/ch07_u13.tex
